\documentclass[11pt]{article}
\usepackage[T1]{fontenc}
\usepackage[utf8]{inputenc}
\usepackage{lmodern}
\usepackage[margin=1in]{geometry}
\usepackage{amsmath,amssymb,amsthm,mathtools}
\usepackage{aliascnt}
\usepackage{microtype}
\usepackage{booktabs,array,longtable,tabularx}
\usepackage{xcolor}
\usepackage{enumitem}
\usepackage{fancyhdr}
\usepackage{needspace}
\usepackage{hyperref}
\usepackage[nameinlink,noabbrev]{cleveref}
\hypersetup{colorlinks=true,linkcolor=blue!45!black,citecolor=blue!45!black,
 urlcolor=blue!45!black,
 pdftitle={Cofinality, Factorization and Scalar Extension for Entire Hahn Functions at Arbitrary Rank},
 pdfauthor={Merged research report prepared for Vladimir Reshetnikov}}
\setlength{\headheight}{14pt}
\pagestyle{fancy}\fancyhf{}
\fancyhead[L]{\small Entire Hahn functions at arbitrary rank}
\fancyhead[R]{\small Cofinality, B\'ezout failure, scalar extension}
\fancyfoot[C]{\thepage}
\setlength{\parskip}{3pt}
\setcounter{tocdepth}{1}
\widowpenalty=10000\clubpenalty=10000
\setlength{\emergencystretch}{2em}
\allowdisplaybreaks[2]
\newtheorem{theorem}{Theorem}[section]
\theoremstyle{plain}
\newaliascnt{lemma}{theorem}
\newtheorem{lemma}[lemma]{Lemma}
\aliascntresetthe{lemma}
\crefname{lemma}{lemma}{lemmas}
\Crefname{lemma}{Lemma}{Lemmas}
\theoremstyle{plain}
\newaliascnt{proposition}{theorem}
\newtheorem{proposition}[proposition]{Proposition}
\aliascntresetthe{proposition}
\crefname{proposition}{proposition}{propositions}
\Crefname{proposition}{Proposition}{Propositions}
\theoremstyle{plain}
\newaliascnt{corollary}{theorem}
\newtheorem{corollary}[corollary]{Corollary}
\aliascntresetthe{corollary}
\crefname{corollary}{corollary}{corollaries}
\Crefname{corollary}{Corollary}{Corollaries}
\theoremstyle{definition}
\newaliascnt{definition}{theorem}
\newtheorem{definition}[definition]{Definition}
\aliascntresetthe{definition}
\crefname{definition}{definition}{definitions}
\Crefname{definition}{Definition}{Definitions}
\theoremstyle{definition}
\newaliascnt{convention}{theorem}
\newtheorem{convention}[convention]{Convention}
\aliascntresetthe{convention}
\crefname{convention}{convention}{conventions}
\Crefname{convention}{Convention}{Conventions}
\theoremstyle{definition}
\newaliascnt{example}{theorem}
\newtheorem{example}[example]{Example}
\aliascntresetthe{example}
\crefname{example}{example}{examples}
\Crefname{example}{Example}{Examples}
\theoremstyle{definition}
\newaliascnt{question}{theorem}
\newtheorem{question}[question]{Question}
\aliascntresetthe{question}
\crefname{question}{question}{questions}
\Crefname{question}{Question}{Questions}
\theoremstyle{remark}
\newaliascnt{remark}{theorem}
\newtheorem{remark}[remark]{Remark}
\aliascntresetthe{remark}
\crefname{remark}{remark}{remarks}
\Crefname{remark}{Remark}{Remarks}
\newcommand{\C}{\mathbb C}\newcommand{\R}{\mathbb R}
\newcommand{\Q}{\mathbb Q}\newcommand{\Z}{\mathbb Z}\newcommand{\N}{\mathbb N}
\newcommand{\No}{\mathbf{No}}
\newcommand{\K}{K_\Gamma}\newcommand{\KD}{K_\Delta}
\newcommand{\E}{\mathcal E_\Gamma}
\newcommand{\A}{\mathcal A_\Gamma}\newcommand{\OO}{\mathcal O_\Gamma}
\newcommand{\TT}{\mathcal T_\Gamma}
\newcommand{\mm}{\mathfrak m_\Gamma}
\newcommand{\Dom}{\mathcal D_{\Gamma,\Delta}}
\DeclareMathOperator{\supp}{supp}\DeclareMathOperator{\lc}{lc}
\DeclareMathOperator{\ord}{ord}\DeclareMathOperator{\cf}{cf}
\DeclareMathOperator{\Div}{Div}\DeclareMathOperator{\inpol}{in}
\DeclareMathOperator{\conv}{conv}
\DeclareMathOperator{\remp}{rem}\DeclareMathOperator{\Res}{Res}
\DeclareMathOperator{\ev}{ev}
\DeclareMathOperator{\Spec}{Spec}\DeclareMathOperator{\red}{red}
\newcommand{\RD}{\mathcal R_D}\newcommand{\Ucal}{\mathcal U}
\newcommand{\ggA}{\mathrel{\gg_{\!A}}}
\newcommand{\abs}[1]{\lvert#1\rvert}
\newcommand{\norm}[1]{\lVert#1\rVert}
\newcommand{\res}{\operatorname{res}}
\newcommand{\kk}{\mathbf k}\newcommand{\PP}{\mathbb P}
\DeclareMathOperator{\Pdir}{Pdir}\DeclareMathOperator{\Adm}{Adm}
\newcommand{\Oz}{\mathbf{Oz}}\newcommand{\Ed}{\mathcal E_{\Gamma,d}}
\newcommand{\Kle}{K_\Gamma^{\le0}}\DeclareMathOperator{\Jac}{Jac}
\newcounter{entrcsaveeq}
\newcounter{entassaveeq}
\newcommand{\Zk}{\Z_{\kk}}\newcommand{\OzK}{\Oz_{\K}}
\newcommand{\PiK}{\Pi_{\K}}\newcommand{\pg}{\mathsf p}
\DeclareMathOperator{\Int}{Int}\DeclareMathOperator{\ct}{ct}
\DeclareMathOperator{\NF}{NF}\DeclareMathOperator{\Syz}{Syz}
\newcommand{\inner}[2]{\langle#1,#2\rangle}
\newcolumntype{Y}{>{\raggedright\arraybackslash}X}
\newcolumntype{P}[1]{>{\raggedright\arraybackslash}p{#1}}

\title{\textbf{Cofinality, Factorization and Scalar Extension\\
for Entire Hahn Functions at Arbitrary Rank}\\[5pt]
\large A classification of the entire ring, its ideal theory,\\
and its behaviour under enlargement of the value group}
\author{Merged research report prepared for Vladimir Reshetnikov}
\date{21 September 2026}

\begin{document}
\maketitle
\begin{abstract}
Let $\Gamma$ be a nonzero set-sized divisible ordered abelian group and
$\K=\C((t^\Gamma))$ the full Hahn field over the trivially valued complex
numbers. We study ordinary power series $\sum a_nZ^n$ whose evaluation
families are strongly Hahn-summable at every point of this one fixed
field. An exact coefficient criterion holds: this happens precisely when
$v(a_n)/n$ tends cofinally to $+\infty$ in $\Gamma$. Consequently
nonpolynomial entire series exist exactly when $\cf(\Gamma)=\aleph_0$.
Support-controlled preparation, proved by well-founded recursion over a
well-ordered monoid rather than by any contraction mapping, yields exact
root counts on every valuation ball, shell and residue direction, a
genus-zero canonical product for every ball-finite divisor with no
exponential correction factors, and complete factorization by zeros at
arbitrary rank.

Three sharp classifications follow. The first is ideal-theoretic. The entire
ring $\E$ is always a GCD domain with units $\K^\times$; it is the
polynomial ring, hence a PID, when $\cf(\Gamma)>\aleph_0$; it is a
non-Noetherian B\'ezout domain with unrestricted Hermite interpolation when
$\Gamma$ has an order unit; and otherwise it is a non-Noetherian GCD domain
that is \emph{not} B\'ezout, witnessed by two explicit entire functions with
disjoint simple zero sets generating a proper ideal. A separate order-unit
criterion decides invertibility in the valuation-restricted algebra
$\TT$, and shows that the positive case of the trichotomy is sharp rather
than merely sufficient.

The second classification is the anatomy of that third case. Grouping the
nodes of a radially finite divisor into valuation shells of radii
$\gamma_1<\gamma_2<\cdots$ and normalizing each shell to valuation zero, a
complete infinite Hermite prescription is realized by an entire function if
and only if, on all sufficiently distant shells, the coefficients of that
shell's finite Hermite interpolation polynomial lie in the valuation ring
obtained by killing the Archimedean scale of the shell radius. The quotient
of $\E$ by the divisor's canonical product is therefore an algebraic
restricted product of finite shell algebras. This gives a
necessary-and-sufficient B\'ezout criterion, a close-pair obstruction
invisible to the sizes of the individual prescribed values, and explicit
non-evaluation maximal ideals whose residue fields are ultraproducts of
lower-rank Hahn fields. For one node per shell, with arbitrary
multiplicities, these ideals and the evaluation ideals are all the maximal
ideals containing the canonical product; the maximal spectrum of the
quotient is the Stone--\v{C}ech space $\beta\N$ of the shell indices, and
its boundary residue fields are algebraically closed. A unit-ideal
criterion for finitely many generators and an exact separation threshold
for paired zeros complete the ideal-theoretic picture. Apart from the
statement that those residue fields are algebraically closed, these results
use no algebraic closedness, and all of them hold with real Hahn
coefficients, the boundary residue fields then being real closed. Beyond the
maximal ideals, every prime ideal containing a product with one \emph{simple}
node per shell is given by a nonprincipal ultrafilter together with a convex
subgroup of an ultraproduct of quotient value groups, or is an evaluation
ideal. That quotient is zero-dimensional exactly when $\Gamma$ has an order
unit, and otherwise every nonprincipal fibre of its spectrum contains a
strict chain of continuum many primes, although the quotient is a reduced
B\'ezout ring in which finite matrices have Smith forms. A cofinal
enlargement of the value group leaves the product and its zeros unchanged
and can nevertheless split one prime into continuum many; in the explicit
surreal example even an old maximal ideal acquires continuum many nonmaximal
extensions. The third
classification is analytic. For every
ordered-group extension $\Gamma\subseteq\Delta$, all nonpolynomial entire
series over $\K$ share one exact strong evaluation domain in $\KD$: the
valuation ring of $v_\Delta$ coarsened by the convex hull of $\Gamma$, whose
residue field is $\C((t^H))$ rather than $\C$. Entireness survives exactly
for cofinal extensions, no new zeros appear anywhere in the surviving
domain, and on a noncofinal extension no alternative entire series agrees
with the old one even on all ordinary complex constants. The scale boundary
is intrinsic: the value group, the coefficient slopes and the negated zero
valuations generate the same convex subgroup. In finitely many ordinary
variables the uniform domain has an exact replacement, proved without
divisibility and over any trivially valued coefficient field: the monomial
family of an entire series is strongly summable at a point of the larger
field exactly when the image of its ordinary support under the quotient
weights of the point is well ordered. For the separate convention of
homogeneous-block summation, the ordinary complex directions along which a
new dominating scale can be crossed are exactly the proper countable unions
of projective algebraic sets.

In $d\ge2$ variables and countable cofinality, every radially finite
configuration of points can be carried into a coordinate line by an entire
automorphism exactly when $\Gamma$ has an order unit. Without one,
configurations with polynomial-part coordinates, among them omnific ones,
still can be, but an escaping sequence of infinitesimal simplices cannot,
because every entire automorphism and its inverse reduce to mutually inverse
polynomial maps along an exhaustion by convex subgroups. On rectifiable
configurations the fat-point ideals have explicit generators,
$\binom{\bar m+d-1}{d-1}$ of them for maximal multiplicity $\bar m$, and are
not finitely generated when the multiplicities are unbounded.

At ordinary points of $\kk^d$, $\kk=\R$ or $\C$, and over an arbitrary,
possibly nondivisible, value group, the coefficients of an entire series at a
fixed Hahn exponent form an ordinary polynomial. A degree-controlled lifting
lemma then makes the entire ring faithfully flat over $\kk[X]$, with exact
coefficientwise membership and normal forms for every ideal defined over
$\kk$. Values in the polynomial-part ring on a Zariski-dense set of ordinary
points force a polynomial; the entire maps preserving the omnific integers of
a workspace are the integer-valued polynomials plus polynomials with purely
infinite coefficients; the ordinary exceptional locus is the zero set of a
finitely generated positive-tail ideal, every nonzero ideal occurs when
$\cf(\Gamma)=\aleph_0$, and the same finitely many scales control every jet
order. These finite certificates cannot be computed in general: an explicit
computable family with quadratic valuation growth has exceptional set $\{0\}$
or $\{0,1\}$ according as a machine halts.

This report is the union of eight manuscripts. Two were written
independently and proved the same engine; a third, written against their
merge, answered the question about the interpolation image that the merge
had left open; a fourth, written in parallel against a snapshot in which
that question was still open, proved the same image theorem independently by
a different construction and
added the ideal-theoretic refinements above; a fifth answered the
scalar-extension clause of the several-variable question; a sixth described
the prime ideals above a product with one simple node per shell; a seventh
proved the rectification dichotomy and the fat-point ideals in several
variables; an eighth proved the arithmetic-sampling, flatness,
exceptional-ideal and computability results at ordinary points. The
arbitrary-rank
classifications, the explicit non-B\'ezout construction, the exact shellwise
image and its ideal theory, the universal extension domain and its
several-variable replacement, the analytic prime fibres with their
splitting under cofinal extension, the rectification dichotomy with its
reduction obstruction and fat-point ideals, and the coefficientwise division,
faithful flatness, arithmetic sampling, exceptional ideals and halting barrier
at ordinary points are proposed original
contributions, not a
claimed solution of any named published conjecture. Classical Hahn-field
inputs and rank-one antecedents are identified and credited, and so is the
classical ultraproduct and valuation-domain mechanism behind the prime
fibres. Complete proofs
are supplied; they have not been independently refereed or verified in a
proof assistant.
\end{abstract}
\newpage
\tableofcontents
\newpage

\section{The question, the answer, and its scope}
\label{ent:sec:intro}

\subsection{Two questions that turn out to have one engine}
Over an algebraically closed field, finite polynomials with disjoint zero
sets generate the unit ideal. In familiar one-variable analytic theories,
interpolation extends this fact to entire functions, and it is tempting to
infer the same conclusion as soon as an infinite product theorem is
available. The inference is false in the higher-rank setting studied here.
We obtain a ring in which every admissible zero divisor is realized, every
function factors completely according to its zeros, and greatest common
divisors exist, yet two functions without a common zero need not satisfy a
B\'ezout identity.

The second question is what the word \emph{entire} names. A formal series
with coefficients in one Hahn field may be evaluable at every point of that
field and lose even strong summability at a single point of a larger one.
Both questions are decided by the same object: the ordered hierarchy of
scales above the coefficient slopes of the series. The first is decided by
whether those scales have a greatest Archimedean class; the second by
whether they remain cofinal after the enlargement.

The obstruction behind the first question is genuinely infinite. At
increasingly large arguments, any one entire function is restricted by a
common lower bound on the valuations of its coefficients. We place a second
sequence of zeros so close to the first that the reciprocal values required
by a B\'ezout identity escape those bounds. Every finite interpolation
problem remains solvable by a polynomial. It is the simultaneous infinite
problem that fails.

\subsection{The analytic category}
Throughout, except in \cref{ent:sec:multivariable}, which states its own
hypotheses (weaker ones in its first eight subsections), $\Gamma$ is a \emph{set-sized divisible ordered abelian
group}, written additively. Except where the degenerate case $\Gamma=0$ is
explicitly discussed, assume $\Gamma\ne0$. Put
\[
 \K=\C((t^\Gamma)),\qquad
 v\left(\sum_\eta c_\eta t^\eta\right)=\min\{\eta:c_\eta\ne0\}.
\]
A Hahn series has well-ordered support. We put $v(0)=+\infty$. This is the
\emph{full} Hahn field: all well-ordered supports are admitted, not a
support-restricted subfield. The complex coefficient field carries the
\emph{trivial} valuation, so every nonzero complex coefficient has valuation
zero and an infinite ordinary complex sum at one exponent is not an allowed
Hahn coefficient operation. The valuation ring and its maximal ideal are
\[
 \OO=\{x\in\K:v(x)\ge0\},\qquad
 \mm=\{x\in\K:v(x)>0\},
\]
with residue field exactly $\C$.

\begin{definition}\label{ent:def:entire}
The ring of \emph{strongly entire Hahn functions} is
\[
 \E=\left\{F=\sum_{n\ge0}a_nZ^n\in\K[[Z]]:
       (a_nz^n)_{n\ge0}\text{ is strongly Hahn-summable for every }z\in\K
       \right\}.
\]
Here strong summability means that the union of the term supports is well
ordered and each exponent belongs to only finitely many term supports. That
this set is a ring is proved below.
\end{definition}

We use $\N=\{0,1,2,\ldots\}$. The variable $Z$ is an ordinary single
indeterminate; its exponents are nonnegative \emph{integers}. The Hahn
exponents occur only in the coefficients, and the two infinities must not be
confused: a finite index $n$ in a power series is not a Hahn exponent, and
integer multiples of a positive Hahn exponent need not be cofinal. Nothing
is asserted about arbitrary generalized series in $Z$ with ordinal support
order types, about arbitrary local functions for the fine surreal topology,
or about the Berarducci--Mantova derivation. All derivatives in this report
hold $\K$ constant.

\begin{convention}[Ring discipline]\label{ent:conv:ring}
Every numbered statement below is asserted over one named coefficient ring,
and no statement is transferred to another. The rings used here are
\[
 \K[Z]\ \subseteq\ \E\ \subseteq\ \TT\ \subseteq\ \A\ \subseteq\ \K[[Z]],
\]
defined in \cref{ent:sec:support,ent:sec:units}, all of which have
\emph{Hahn series over $\C$} as coefficients. In particular no theorem here
is a theorem about $\mathcal O(U)((t^\Gamma))$ or $\C\{z\}((t^\Gamma))$,
whose coefficients are ordinary holomorphic functions on a common complex
domain or convergent germs. That distinction is not a formality: the
companion report \emph{Analytic Geometry over Surcomplex Hahn Fields}
adopts the same discipline by name for its own four coefficient rings, and
records there an explicit element of the radius-free algebra lying in no
common-domain algebra. We invoke that convention here so that no reader
transports a statement of this report into those categories. See
\cref{ent:sec:collection}.
\end{convention}

\subsection{The two classifications}
A subset of $\Gamma$ is cofinal if every element of $\Gamma$ is at most some
member of that subset; $\cf(\Gamma)$ is the least cardinality of a cofinal
subset, and it is at least $\aleph_0$ for $\Gamma\ne0$, since an ordered
group has no greatest element.

\begin{definition}[Order unit]\label{ent:def:order-unit}
An element $\delta>0$ of $\Gamma$ is an \emph{order unit} if for every
$\gamma\in\Gamma$ some positive integer $m$ satisfies $m\delta>\gamma$; that
is, $\{m\delta:m\ge1\}$ is cofinal in $\Gamma$.
\end{definition}

\begin{remark}[One name for one invariant]\label{ent:rem:one-name}
This report uses exactly two order-theoretic invariants of $\Gamma$, and
they must not be conflated.
\begin{itemize}[leftmargin=1.4em]
\item \emph{Cofinality} always means $\cf(\Gamma)$, the least cardinality of
a cofinal subset of the underlying ordered \emph{set} of the value group. It
is never the cofinality of an ordinal index, of a diagram, or of the image of
some displacement or spectral map; other reports in this collection use the
word for such other invariants, and nothing here transfers to them.
\item \emph{Order unit} is the condition of
\cref{ent:def:order-unit}. The first two source manuscripts merged here named it
twice: one called $\delta$ an order unit, the other called it a
\emph{cofinal-cyclic element}, meaning exactly that the cyclic subgroup it
generates is cofinal. These are the same condition, stated twice; we keep
the single name \emph{order unit} and record the synonym here once.
For a fixed $\delta>0$, it is equivalent to each of the following:
$\delta$ represents the greatest positive Archimedean class of $\Gamma$;
$m\delta\to+\infty$; the monomial
$t^\delta$ is topologically nilpotent in the intrinsic valuation topology of
$\K$. Thus an order unit exists exactly when the positive Archimedean
classes have a greatest member; a positive element in a smaller class is
not an order unit.
\end{itemize}
An order unit forces $\cf(\Gamma)=\aleph_0$, but not conversely: the group
$\Gamma_\infty$ of \cref{ent:sec:examples} has countable cofinality and no
order unit. That gap is exactly the third case of the trichotomy below, and
also the difference between the two existence questions of
\cref{ent:thm:main,ent:thm:units}. The report
\path{docs/surcomplex/expanding-polynomial-dynamics/} uses the same order unit
and renamed its delivered phrase ``cofinal scale'' accordingly
(\texttt{epd:rem:one-name}). The report on transcendence over bounded support uses $\cf(G)$ in
exactly this sense (\texttt{bst:subsec:terms}); its ordinal-index use of
``cofinal'' is confined to \texttt{bst:lem:coding}. The same cofinality
dichotomy governs completion of the algebraic scalar extension in Hahn vector
spaces (\path{docs/surcomplex/three-duals-of-hahn-vector-spaces/},
\texttt{duals:thm:cofinality-completion}, \texttt{duals:thm:basechange}).
\emph{Related (batch~34).} The omnific continued-fraction report
(\path{docs/surreal/omnific-continued-fractions/}) meets both invariants in
the full Hahn field $\R((\omega^G))$: some infinite continued-fraction code
has a unique realization there exactly when $G$ has countable cofinality
(\texttt{ocf:cor:cofinality}), and some eventually periodic code has one
exactly when $G$ has a positive order unit (\texttt{ocf:thm:orderunit}).
Neither report uses the other.
\end{remark}

\Needspace{18\baselineskip}
\begin{theorem}[Cofinality--rank trichotomy]\label{ent:thm:main}
Let $\Gamma\ne0$ be a divisible ordered abelian group. Then $\E$ is a GCD
domain, its units are exactly $\K^\times$, and exactly one of the following
cases occurs.
\begin{enumerate}[label=\textnormal{(\roman*)},leftmargin=2em]
\item If $\cf(\Gamma)>\aleph_0$, then $\E=\K[Z]$, a principal ideal domain.
\item If $\Gamma$ has an order unit, then $\E$ contains nonpolynomial
functions and is a non-Noetherian B\'ezout domain. Arbitrary finite-order
jets on every radially finite set can be interpolated by an entire function.
\item If $\cf(\Gamma)=\aleph_0$ and $\Gamma$ has no order unit, then $\E$
contains nonpolynomial functions and is a non-Noetherian GCD domain that is
not B\'ezout. More strongly, there are $f,g\in\E$ with disjoint simple zero
sets and $(f,g)\ne\E$.
\end{enumerate}
For $\Gamma=0$, the strong-summation convention gives $\mathcal E_0=\C[Z]$.
\end{theorem}

A GCD domain is an integral domain in which every pair has a greatest common
divisor, defined up to a unit. A B\'ezout domain is one in which every
finitely generated ideal is principal. Case (iii) therefore separates two
properties that coincide for finite polynomial algebra. This is not a
distinction about algebraic closedness: every $\K$ here is algebraically
closed.

The trichotomy says that case~(iii) occurs. The next theorem is the anatomy
of that third branch: it measures exactly how far interpolation fails, for
every radially finite divisor and every group, and it specializes to the other
two branches. It was proved after the first two source manuscripts, against
the gap their merged report itself recorded as open
(\cref{ent:q:image}), and is merged in here. A fourth manuscript, written in
parallel against a snapshot in which the gap was still open, proved it
independently by a different construction; its proof is kept as a marked
second route (\cref{ent:sm:subsec:cardinal}).

\Needspace{16\baselineskip}
\begin{theorem}[Exact image of the infinite jet map]\label{ent:thm:main-image}
Let $D$ be an infinite radially finite divisor. Group its nodes into
valuation shells of radii $0<\gamma_1<\gamma_2<\cdots$, cofinal in $\Gamma$,
normalize shell $n$ by $\tau_n=t^{-\gamma_n}$, and let $R_n$ be the unique
finite Hermite interpolation polynomial of that shell's prescribed jets in
the normalized coordinate. Let
\[
 V_n=\{x\in\K:v(x)\ge-m\gamma_n\text{ for some }m\in\N\}
\]
be the valuation ring obtained by killing the Archimedean scale of
$\gamma_n$. Then:
\begin{enumerate}[label=\textnormal{(\alph*)},leftmargin=2em]
\item an entire function realizing all the prescribed jets exists if and only
if $R_n\in V_n[T]$ for all sufficiently large $n$, with no condition on
finitely many exceptional shells and no uniform bound on the number of nodes
per shell, on the jet orders, or on the admitted scale exponents;
\item consequently $\E/(P_D)$ is the algebraic restricted product
$\Lambda_0\times\{(r_n):r_n\in\Lambda_n^\circ\text{ eventually}\}$ of the
finite shell algebras $\Lambda_n=\K[T]/(P_n)$ and
$\Lambda_n^\circ=V_n[T]/(P_n)$;
\item a B\'ezout identity $AP_D+BF=1$ with $F\in\E$ exists if and only if $F$
is nonzero at every node of $D$ \emph{and}, on all sufficiently large shells,
$v(F(\rho_{n,j}))\le c_n\gamma_n$ for some finite integer $c_n$;
\item for one simple node per shell, every nonprincipal ultrafilter on the
positive integers yields an explicit maximal ideal of $\E$ that is not an
evaluation ideal, with residue field the ultraproduct of the lower-rank Hahn
fields $\C((t^{H_n}))$, $H_n$ the convex subgroup generated by $\gamma_n$.
\end{enumerate}
These statements also hold over the real Hahn field $\R((t^\Gamma))$, with no
use of algebraic closedness.
\end{theorem}

Part (a) is \cref{ent:thm:image}, (b) is \cref{ent:thm:quotient}, (c) is
\cref{ent:thm:bezout-criterion}, (d) is \cref{ent:thm:maximal}, and the real
form is \cref{ent:prop:real-form}; all are proved in
\cref{ent:sec:obstruction}. The restricted product here is an algebraic
union, not a topological or adelic object; no norm, no local compactness and
no completion is asserted. Case (ii) of \cref{ent:thm:main} is the
degenerate case $V_n=\K$, and case (i) is the case where no infinite radially
finite divisor exists at all.

The fourth source adds to parts (c) and (d). It decides when finitely many
entire functions generate the unit ideal together with $P_D$
(\cref{ent:sm:thm:finite-generators}), gives the exact threshold at which a
pair of canonical products with paired simple zeros stops generating the
unit ideal (\cref{ent:sm:thm:paired}), and shows that for one node per shell,
with arbitrary multiplicities, the ideals of part (d) and the evaluation
ideals are \emph{all} the maximal ideals containing $P_D$, forming a copy of
the Stone--\v{C}ech space $\beta\N$
(\cref{ent:sm:thm:maximal-all,ent:sm:thm:beta}). None of these classifies
the whole maximal spectrum of $\E$; see \cref{ent:q:spectrum}.

The sixth source goes below the maximal ideals, for one \emph{simple} node
per shell (\cref{ent:ps:sec:spectrum}). Starting from the quotient of part
(b), which it credits, it shows that every prime ideal of $\E$ containing
such a product is an evaluation ideal or is given by a nonprincipal
ultrafilter $\Ucal$ and a convex subgroup of the ultraproduct
$\prod_\Ucal(\Gamma/H_n)$, the fibre over $\Ucal$ being the spectrum of an
explicit ultraproduct valuation domain (\cref{ent:ps:thm:spectrum}). The
quotient is zero-dimensional exactly when $\Gamma$ has an order unit, and
otherwise each nonprincipal fibre contains a chain of continuum many primes
(\cref{ent:ps:thm:dichotomy}); it is nevertheless a reduced semihereditary
B\'ezout ring with Smith forms (\cref{ent:ps:thm:bezout,ent:ps:thm:smith}).
Under a cofinal enlargement $\Gamma\subseteq\Delta$ the contraction of these
primes is intersection of convex subgroups, with all extensions of a prime
described by an interval of convex subgroups
(\cref{ent:ps:thm:contraction}); for $\Gamma_\infty$ inside its
rational-indexed enlargement one prime splits into continuum many while the
product and its zeros do not change
(\cref{ent:ps:thm:splitting,ent:ps:cor:maximal-splitting}). Bounded
multiplicities do not change the spectrum, and unbounded ones produce primes
that the reduced description misses (\cref{ent:ps:thm:multiplicity}). This
answers the ``prime ideals'' clause of \cref{ent:q:spectrum} for one simple
node per shell. The ultraproduct and valuation-domain mechanism is classical
\cite{FFW,StacksValuation}; what is proposed is its exact analytic
realization, the dichotomy and the splitting construction.

\Needspace{14\baselineskip}
\begin{theorem}[Universal scalar-extension package]\label{ent:thm:main-extension}
Let $\Gamma\subseteq\Delta$ be nonzero divisible ordered abelian groups and
let $H=\conv_\Delta(\Gamma)$ be the convex hull of $\Gamma$ in $\Delta$.
Write $\Dom$ for the valuation ring of $v_\Delta$ coarsened by $H$. Then for
every \emph{nonpolynomial} $f\in\E$:
\begin{enumerate}[label=\textnormal{(\alph*)},leftmargin=2em]
\item the family $(a_nz^n)_n$ is strongly summable in $\KD$ exactly for
$z\in\Dom$, so the domain does not depend on $f$, and $f(\Dom)\subseteq\Dom$;
\item $\Dom=\KD$, equivalently $f\in\mathcal E_\Delta$, exactly when
$\Gamma$ is cofinal in $\Delta$;
\item every zero of the extended function on $\Dom$ already lies in $\K$
with unchanged multiplicity, even when $f$ has infinitely many zeros;
\item if $\Gamma$ is not cofinal in $\Delta$, there is no
$g\in\mathcal E_\Delta$ agreeing with $f$ on $\C$, let alone on $\K$;
\item $\Gamma$, the slope set $\{v(a_n)/n:n\ge1,\ a_n\ne0\}$ and the
negated zero valuations $\{-v(r):r\ne0,\ f(r)=0\}$ generate the same convex
subgroup of any larger ordered group,
so $\Dom$ is recoverable from the coefficients alone or from the zeros alone.
\end{enumerate}
The residue field of $\Dom$ is canonically $\C((t^H))$, not $\C$.
\end{theorem}

\begin{remark}[The two extension statements are one statement]
\label{ent:rem:conflict}
Part (b) is often phrased negatively --- in the noncofinal case the series
is simply \emph{not entire over} $\KD$, and indeed
$\mathcal E_\Delta\cap\K[[Z]]=\K[Z]$ there
(\cref{ent:thm:extension-criterion}) --- while parts (a) and (c) say that
the very same series remains strongly summable on all of $\Dom$ and keeps
every one of its zeros there. Read side by side these look like a
disagreement about whether the function survives an enlargement. They are
not in conflict: they are the coarse and the exact form of one fact. ``Not
entire over $\KD$'' means precisely that the strong evaluation domain is a
\emph{proper} subring of $\KD$, and \cref{ent:thm:extension} computes
exactly which subring it is, namely $\Dom$, on which the function is
unchanged. The negative statement locates the boundary; the positive
statement describes everything inside it. Throughout this report ``entire
over $K$'' is reserved for strong summability at \emph{every} point of $K$,
and any weaker survival is described by naming its domain.
\end{remark}

The proof of \cref{ent:thm:main} is distributed through the article. The
coefficient criterion and cofinality dichotomy are in
\cref{ent:sec:growth}; the order-unit criterion for the restricted algebra,
which is what makes case (ii) sharp rather than merely sufficient, is in
\cref{ent:sec:units}; preparation and root counting are in
\cref{ent:sec:prep,ent:sec:zeros}; factorization and GCDs are in
\cref{ent:sec:products}; the negative and positive B\'ezout results are in
\cref{ent:sec:obstruction,ent:sec:positive}.
\Cref{ent:sec:classification} assembles them.
\Cref{ent:thm:main-extension} is proved in \cref{ent:sec:extension}.
\Cref{ent:thm:main-image} is proved in \cref{ent:sec:obstruction}, whose
shell notation is fixed in \cref{ent:subsec:notation,ent:subsec:shells}.

\paragraph{The same alternative for finite extensions (batch 32).}
The Galois part of the report on transcendence over bounded support
(\path{docs/surreal/transcendence-over-bounded-support/}) proves an algebraic
counterpart of parts (b) and (d) of \cref{ent:thm:main-extension}
(\texttt{bst:gr:thm:transition}). For an inclusion of nonzero set-sized
ordered groups, the smaller one divisible, consider the complex
bounded-support fraction fields at the two scales. If the inclusion is
cofinal, every finite extension of the smaller field stays a field of the
same degree over the larger one; if it is not, every finite extension
splits completely there. The objects differ (finite extensions of
bounded-support fraction fields there, entire series over full Hahn fields
here), and neither theorem is used for the other. What they share is that
cofinality of the exponent inclusion alone decides the outcome.

\paragraph{Several ordinary variables.}
Part (a) of \cref{ent:thm:main-extension} has no uniform analogue in
several variables: a series independent of one coordinate imposes no
constraint in that coordinate. \Cref{ent:sec:multivariable}, the fifth
source's contribution, gives the exact replacement. For a whole-family
entire series $F=\sum_\alpha a_\alpha X^\alpha$ in finitely many variables
and a point $z$ of the larger field with nonzero coordinates, the family
$(a_\alpha z^\alpha)_\alpha$ is strongly summable if and only if
\[
 \Bigl\{\textstyle\sum_j\alpha_j\bigl(v(z_j)+H\bigr):a_\alpha\ne0\Bigr\}
 \ \text{is well ordered in }\Delta/H
\]
(\cref{ent:mv:thm:domain}). That section works under weaker hypotheses than
the rest of the report --- no divisibility, any trivially valued coefficient
field --- and keeps two summation conventions apart; only the whole-family
one is called entire. It answers the ``In particular'' clause of
\cref{ent:q:several} and no other part of that question.

\paragraph{Rectification and fat points.}
Subsections \crefrange{ent:rc:subsec:setting}{ent:rc:subsec:limits} of
\cref{ent:sec:multivariable}, the seventh source's, answer part of two further clauses of \cref{ent:q:several}, for
configurations of points. For $d\ge2$ and $\cf(\Gamma)=\aleph_0$, every
radially finite subset of $\K^d$ is carried into a coordinate line by some
entire automorphism if and only if $\Gamma$ has an order unit
(\cref{ent:rc:thm:main}). Without an order unit, configurations with
polynomial-part coordinates are still rectifiable, which covers omnific
configurations (\cref{ent:rc:thm:omnific-main}), while an escaping sequence
of infinitesimal simplices is not (\cref{ent:rc:thm:obstruction}): every
entire automorphism and its inverse eventually have coefficients integral for
the convex subgroups of an exhaustion, and reduce there to mutually inverse
polynomial maps (\cref{ent:rc:thm:reduction,ent:rc:cor:automorphism-reduction}).
On a rectifiable configuration with multiplicities, the fat-point ideal has
explicit generators and exactly $\binom{\bar m+d-1}{d-1}$ minimal generators
when the multiplicities are bounded with maximum $\bar m$, and it is not
finitely generated when they are unbounded (\cref{ent:rc:thm:fat-main}); the
values taken on such a configuration are described exactly by the
one-variable image theorem (\cref{ent:rc:cor:evaluation}). Preparation, GCD
and B\'ezout theory, positive-dimensional divisors, and ideals of
configurations that are not line-rectifiable remain open in several
variables.

\paragraph{Arithmetic sampling at ordinary points.}
The last subsections of \cref{ent:sec:multivariable},
\crefrange{ent:as:subsec:setting}{ent:as:subsec:limits}, are source~11's (the
eighth manuscript, batch~32). They work over $\kk\in\{\R,\C\}$ with an
arbitrary, possibly nondivisible, value group and ask what arithmetic values
at ordinary points force. At a fixed Hahn exponent the coefficients of an
entire series form an ordinary polynomial (\cref{ent:as:lem:scale}); with a
degree-controlled lifting lemma this gives exact coefficientwise membership in
every ideal extended from $\kk[X]$, and $\Ed$ is faithfully flat over $\kk[X]$
(\cref{ent:as:thm:matrix,ent:as:thm:flat}). An entire function whose values
at a Zariski-dense set of ordinary points lie in the polynomial-part ring
$\Kle$ is a polynomial (\cref{ent:as:thm:sampling}); the entire maps preserving
the omnific integers $\OzK=\Zk+\PiK$ of the workspace are exactly
$\Int(\Zk^d,\Zk)\oplus\PiK[X]$ (\cref{ent:as:thm:integer}); the ordinary
exceptional locus is the zero set of a finitely generated positive-tail ideal,
every nonzero ideal occurs when $\cf(\Gamma)=\aleph_0$, and the same finitely
many scales control every jet order
(\cref{ent:as:thm:tail,ent:as:thm:realization,ent:as:thm:jets}). The finite
certificate cannot be computed in general, by an explicit halting family with
quadratic growth (\cref{ent:as:thm:undecidable}). This partly answers the
ideal-theory clause of \cref{ent:q:several}: for ideals extended from
$\kk[X]$, including the vanishing ideals of ordinary point sets and, for
$\kk=\C$, of algebraic sets defined over $\kk$ (\cref{ent:as:cor:vanishing},
the merge's). Ideals not of that form, preparation, and GCD or B\'ezout theory
stay open in several variables.

\subsection{Relationship to the repository and the literature}
\label{ent:subsec:repo}
A note on how this report grew. Its first two source manuscripts were
independently written against pinned snapshots of this repository, proved the
same engine, and carried complementary payloads; their merge is described in
\cref{ent:app:sources}. A third manuscript was written later, against the
merged report itself at a pinned snapshot, and answers the question that the
merged report had recorded as open in \cref{ent:q:image}. Its contribution is
\cref{ent:thm:main-image} and the material of
\cref{ent:subsec:shells,ent:subsec:jet-image,ent:subsec:image-proof,ent:subsec:no-bezout,ent:subsec:invisible};
it credits, rather than claims, the coefficient criterion
(\cref{ent:thm:criterion}), the canonical product
(\cref{ent:thm:product}), the order-unit interpolation theorem
(\cref{ent:thm:interpolation}) and the explicit non-B\'ezout pair
(\cref{ent:thm:no-bezout}) as already present here. Where it re-proved one of
those for dependency control, both proofs are printed and marked; see
\cref{ent:prop:kernel-direct} and the two proofs of
\cref{ent:thm:interpolation,ent:thm:bezout}. Its symbols $\rho_n$ and
$a_{n,j}$ were the exact inverses of this report's, and the single
renormalizing pass that fixed that is recorded in
\cref{ent:subsec:notation}.

Three further manuscripts were merged later. The \emph{fourth} was written
against a snapshot of the merged report taken after the third source's pin
but before the third source was merged in, and it too set out to identify
the missing image of the infinite jet map. It did so independently and at the same
generality, with a different sufficiency construction, so the image,
quotient and B\'ezout theorems of \cref{ent:sec:obstruction} now have two
independent sources; they are printed once, and its construction is kept as
a marked second route (\cref{ent:sm:subsec:cardinal}). Its framing as
\emph{the} resolution of that gap is accurate for its snapshot and stale for
this tree, where the third source answered the question first; the ledger in
\cref{ent:app:sources} records this. Its genuinely additional results are
the explicit truncation bound of the second route, the finite-generator and
membership criteria, the exact threshold for paired zeros, and the
classification of all maximal ideals above a one-node-per-shell product with
their topology, residue-field type and boundary coordinate
(\cref{ent:sm:subsec:all-maximal}). Its letters collide with this report's
more seriously than the third source's did --- two of them are swapped with
each other --- and \cref{ent:subsec:notation} records the translation. The
\emph{fifth} was written against a later snapshot, after the third source
had been merged, and answers the scalar-extension clause of
\cref{ent:q:several}; it is \cref{ent:sec:multivariable}, which absorbs the
finite-variable criteria that earlier versions kept in an appendix. The
\emph{sixth} was written against a still later snapshot, after the fifth
source had been merged; this report's directory is identical at that pin and
at the base of the sixth merge. It takes
\cref{ent:thm:quotient,ent:cor:single-node} and the maximal-ideal results of
the fourth source as its starting point, credits them, and describes all
prime ideals above a product with one simple node per shell, their
behaviour under cofinal enlargement of the value group, and the
multiplicity threshold; it is \cref{ent:ps:sec:spectrum}. It reproves the
simple-node interpolation theorem for dependency control, by the argument of
the first route (\cref{ent:ps:rem:route}), so that proof is printed once.
Its sequence ring, convex subgroups and several other letters collide with
this report's and are renamed (\cref{ent:subsec:notation}).

A \emph{seventh} manuscript was written against a snapshot taken after the
sixth source had been merged; this report's directory is identical at that
pin and at the base of the seventh merge. It proves the rectification
dichotomy, the reduction obstruction and the fat-point ideal theory of
\crefrange{ent:rc:subsec:setting}{ent:rc:subsec:limits}. It credits the
coefficient criterion, canonical products and division, order-unit
interpolation and the image theorem to this report, and reproves the part it
needs for dependency control; that part is printed once, by citation. It
described its problem as different from the several-variable question here;
it is a partial answer to the ideal-theory and interpolation clauses of
\cref{ent:q:several} and is placed accordingly. It did not cite the
one-variable polynomial-reduction lemma of the holonomic-rigidity report, or
this report's coarsened residue map and restricted-unit argument, which are
antecedents of its reduction theorem; the merge credits them
(\cref{ent:rc:subsec:reduction}). Its letters are renamed where they collide
(\cref{ent:subsec:notation}).

An \emph{eighth} manuscript, source~11 (batch~32, manuscript~09; file prefix
\texttt{11-arithmetic-sampling}), was written against a snapshot pinned at
\texttt{343dc2c}, after the seventh source had been merged; this report's
directory is identical at that pin and at the base of the eighth merge. It
proves the arithmetic-sampling, flatness, exceptional-ideal and computability
results of \crefrange{ent:as:subsec:setting}{ent:as:subsec:limits}. It
credits the coefficient criterion and the cofinality principle to this report
and reproves them for dependency control; they are printed once, by citation,
and its statement that its proof allows arbitrary rank without divisibility
adds nothing, since \cref{ent:thm:multi-criterion} already does. Its
whole-class rigidity corollary is \cref{ent:cor:all-no,ent:mv:prop:allscale}.
It did not cite the exceptional-direction theorems of the fifth source, which
run parallel to its exceptional ideals (\cref{ent:as:subsec:tail}), or the
numerical-polynomial theorems of the set-sized-quotients report, which are the
polynomial half of its classification of omnific-preserving entire maps; the
merge records both. It presented its question as different from those studied
here (``what happens when arithmetic values are prescribed at ordinary
coefficient-field points''); it is a partial answer to the ideal-theory clause
of \cref{ent:q:several} and is placed accordingly. Its letters are renamed
where they collide (\cref{ent:subsec:notation}).

The repository supplied with the question already contains a rank-one report
over $\C((t^\R))$, including preparation, zero counting, zero-free rigidity
and a partial-theta canonical product \cite{RepoRankOne}. None of the source
manuscripts claims to originate those points, and neither do we; the present
report is their arbitrary-rank generalization, and the rank-one report
itself proves the obstruction that makes rank one non-generic. The
repository also contains a global-divisor report in a different ring,
$\mathcal O(U)((t^\Gamma))$, with ordinary holomorphic coefficient functions
on a common complex domain \cite{RepoDivisors}; its global support
obstructions belong to a different analytic category and our argument is not
a restatement of its Picard-group or moving-divisor results.
\Cref{ent:sec:collection} states the cross-references in full.

Classical non-Archimedean factorization is an essential antecedent. Lazard's
work is a foundational reference for one-variable zeros \cite{Lazard};
Cherry records the one-variable product factorization and develops GCD and
factorization results under a complete real-valued non-Archimedean absolute
value \cite{Cherry,CherryLectures}, and himself calls those factorization
results well known. Hahn-field algebraic closedness is standard, used here
in the form of Poonen's Corollary~4 \cite{Poonen}; the support calculus is
Neumann's \cite{Neumann}; support-based operator methods are developed
extensively by van der Hoeven \cite{vdH}. The distinction between positive
valuation and topological nilpotence in rank two is established, and we
follow Conrad's discussion of it \cite{Conrad}. We claim none of these facts
as new.

\paragraph{Novelty statement.}
The proposed contributions are the arbitrary-rank coefficient, unit and
workspace criteria, the explicit non-B\'ezout construction, the exact
shellwise interpolation image with its restricted-product quotient, sharp
B\'ezout criterion and explicit residue fields, the universal
scalar-extension domain with zero conservation, and their combination into
\cref{ent:thm:main,ent:thm:main-image,ent:thm:main-extension}. The literature
comparison for the image theorem was targeted in the same bounded sense as
the others and located no matching statement of the shellwise
restricted-product theorem; that is a reason to present it as proposed new
work, not a proof that no such statement occurs anywhere. The ultrafilter
construction of \cref{ent:thm:maximal} is standard algebra; what is proposed
there is the analytic surjection onto the restricted product and the
consequent identification of the residue fields. The support-controlled
preparation and product arguments provide a self-contained route to those
results; their rank-one specializations are not new. The literature checks
performed for the first two source manuscripts did not locate the stated
arbitrary-rank classifications or the explicit obstruction in the primary
sources consulted. This is a bounded literature statement, not a
certification of priority. The repository review was targeted at the
catalogue, relevant READMEs and stated scope, not a line-by-line audit of
every manuscript. No named published open problem is claimed solved.

The fourth source proposes as new the exact block restricted-product
quotient, its cardinal construction with cofinal truncation bounds
independent of block dimension and multiplicity, the unit-ideal criterion and
paired-root threshold, and the scale-dependent boundary residue fields with
the full ultrafilter locus of the paired example. The first and third of
these were already in this tree from the third source when it was merged, so
here they count as an independent second derivation. It credits as
antecedents Henriksen's free maximal ideals and transcendental boundary
coordinate for the classical entire ring \cite{Henriksen}, Bruno's
non-Archimedean maximal-ideal classifications under a real value group
\cite{Bruno}, Helmer's divisibility theory \cite{Helmer}, and Cherry; it does
not present the existence of ultrafilter maximal ideals, or of a
transcendental boundary coordinate, as new. The fifth source proposes
\cref{ent:mv:thm:domain} with its repeated-weight proof, and the combination
of the extension geometry with the exceptional-direction classification and
realization. It credits the Ma--Neelon realization of countable algebraic
unions as convergence sets of divergent power series \cite{MaNeelon}, and
does not reclaim the finite-variable growth criterion, the one-variable
extension domain or the whole-class rigidity. The sixth source proposes the
description of all prime fibres above a simple one-node-per-shell product,
the order-unit dichotomy for that reduced divisor and the cofinal splitting
construction, stating that these were not found in the sources it compared
and that this is not proof of their absence. It credits the ultraproduct and
valuation-domain mechanism for primes of product rings to Finocchiaro,
Frisch and Windisch \cite{FFW} and the valuation-domain prime correspondence
to the standard literature \cite{StacksValuation}, and it does not reclaim
the interpolation quotient, the maximal ideals or the maximal spectrum.
The seventh source proposes the order-unit rectification dichotomy, the
eventual polynomial reductions of an entire automorphism and its inverse, the
escaping-simplex obstruction, polynomial-part and omnific rectification
without an order unit, and the fat-point generator formula with its sharp
count. It states that the projection-and-shear strategy is classical and the
jet counting elementary, that its novelty assessment is a bounded search and
not a priority claim, and it cites Rosay and Rudin \cite{RosayRudin},
Winkelmann \cite{Winkelmann}, Cherry \cite{CherryLectures},
Mantova and Matusinski \cite{MantovaMatusinski} and L'Innocente and Mantova
\cite{LInnocenteMantova} for the classical and omnific background.
Source~11 proposes the bounded-degree transfer and the exact
coefficientwise image over $\kk[X]$ with faithful flatness and normal forms,
dense and affine arithmetic sampling with the restriction theorem, the
analytic collapse behind the classification of omnific-preserving entire
maps, the finite positive-tail ideal with its sharp one-variable bound and jet
certificates, universal ideal realization, and the halting family. It states
that the growth criterion, Hahn summation, the Hilbert basis theorem,
Gr\"obner division, the flatness criteria, finite differences, integer-valued
polynomials, the constant-term pullback and the halting obstruction are
established; it cites L'Innocente and Mantova \cite{LInnocenteMantova},
Kuhlmann and Serra \cite{KuhlmannSerra}, the Stacks Project
\cite{Stacks,StacksFlat} and the abstract of Garcia-Fritz and Pasten
\cite{GarciaFritzPasten}, and calls its search limited, a failed search being
no evidence of originality.
\Cref{ent:sec:scope,ent:app:sources} state the limits in full, with the
seventh source's also in \cref{ent:rc:subsec:limits} and source~11's in
\cref{ent:as:subsec:limits}; every limitation recorded by any of the eight
source manuscripts is preserved there.

\section{Hahn supports, the surreal interpretation, and the algebra chain}
\label{ent:sec:support}

\subsection{Foundational setting}
The proofs can be carried out in ZFC using only the set-sized field $\K$. To
interpret them inside the surcomplex numbers, take $\Gamma$ to be a divisible
ordered additive subgroup of $\No$ and send
\begin{equation}\label{ent:eq:embedding}
 \sum_\gamma c_\gamma t^\gamma
 \longmapsto \sum_\gamma c_\gamma\omega^{-\gamma}.
\end{equation}
The convention $t^\gamma:=\omega^{-\gamma}$ is Conway's monomial map; it is
explicitly \emph{not} a surreal exponential. The real and imaginary parts of
the image are Conway normal forms, and the standard normal-form arithmetic
makes \eqref{ent:eq:embedding} an embedding of fields: the real coefficient
subfield $\R((t^\Gamma))$ lies in $\No$, and its coefficientwise
complexification lies in $\No[i]$ \cite{Conway,Gonshor}. Conjugation acts on
coefficients only. Our explicit examples use displayed subgroups of $\No$, so
they do not require an unspecified universal embedding theorem. An element
$x\in\OO$ has a unique expression $x=c+\varepsilon$ with $c\in\C$ and
$\varepsilon\in\mm$; the residue map $\res(x)=c$ is defined on $\OO$, not on
all of $\K$.

\subsection{The support facts actually used}
We use the following classical Hahn--Neumann support facts
\cite{Neumann}\cite[Section~3]{Poonen}\cite{vdH}. They are stated explicitly
because every infinite construction below depends on them, and they are
imported, not proved here.

\begin{lemma}[Hahn--Neumann support calculus]\label{ent:lem:neumann}
Let $S,T$ be well-ordered subsets of an ordered abelian group.
\begin{enumerate}[label=\textnormal{(\alph*)},leftmargin=2em]
\item $S\cup T$ and $S+T$ are well ordered. For every $\eta$, there are only
finitely many pairs $(s,t)\in S\times T$ with $s+t=\eta$.
\item If $S\subseteq\Gamma_{>0}$, its additive monoid
$S^*=\{s_1+\cdots+s_r:r\ge0,\ s_i\in S\}$ is well ordered. For each $\eta$,
there are only finitely many finite ordered words in $S$ whose sum is $\eta$.
\item Consequently, if $u$ is a Hahn series over a commutative coefficient
ring and $v(u)>0$, the geometric series $\sum_{r\ge0}u^r$ is strongly
summable and is the inverse of $1-u$.
\end{enumerate}
\end{lemma}

Part (b) is a support theorem, not an assertion that $nv(u)$ tends cofinally
to infinity. The latter can fail. In $\Q e_1\oplus\Q e_2$ with $e_2>ne_1$ for
every $n$, the Hahn identity $\sum_{n\ge0}t^{ne_1}=(1-t^{e_1})^{-1}$ is valid
but its partial sums do not converge in the intrinsic valuation topology: no
tail has valuation greater than $e_2$. \Cref{ent:ex:ranktwo-unit} turns this
observation into a theorem about units.

\begin{lemma}[Cofinal tails]\label{ent:lem:cofinal-tails}
Suppose $u_n\in\K$ and $v(u_n)\to+\infty$ cofinally in $\Gamma$; that is, for
every $B\in\Gamma$, eventually $v(u_n)>B$. Then $(u_n)$ is strongly summable.
Its sum has valuation at least $\min_n v(u_n)$ when at least one term is
nonzero.
\end{lemma}
\begin{proof}
For a fixed $B$, only finitely many nonzero term supports meet
$(-\infty,B]$. A descending sequence in their union, beginning at $B$, would
therefore lie in a finite union of well-ordered supports, which is
impossible. Likewise, any exponent occurs in only finitely many term
supports. The valuation estimate follows from coefficientwise addition. A
nonempty set of term valuations has a minimum because it is a subset of the
well-ordered support union.
\end{proof}

We use the usual ZFC equivalence between well-ordering of a linear order and
the absence of an infinite strictly descending sequence. All supports and
index sets are sets; no proper-class support or proper-class summation is
involved.

\subsection{Three meanings of an infinite expression}
\label{ent:subsec:three-meanings}
It is essential to keep separate
\begin{enumerate}[label=(\arabic*),leftmargin=2em]
\item a strong Hahn sum, which requires only a well-ordered support union
and pointwise finiteness;
\item convergence in the intrinsic valuation topology of the fixed field
$\K$, which requires cofinal progress through every threshold of $\Gamma$;
\item convergence in the topology inherited from all surreal scales.
\end{enumerate}
These need not agree. \Cref{ent:def:entire} starts with (1). The
\emph{all-point} condition will force (2) for evaluation series
(\cref{ent:thm:criterion}), even though (1) alone does not, and on a single
disk the two genuinely differ. We make no assertion of convergence in (3),
and no density of $\C$ in any fine surreal topology is asserted or used.
The repository's foundational reports emphasize this distinction
\cite{RepoCatalogue,RepoRankOne}.

\subsection{The chain of coefficient algebras}
\label{ent:subsec:chain}
Four algebras of series in $Z$ over $\K$ are kept distinct throughout.
Define the \emph{polynomial-coefficient Hahn algebra}
\begin{equation}\label{ent:eq:Halg}
 \A=\C[X]((t^\Gamma)),
\end{equation}
whose elements are $A=\sum_{\eta\in S}P_\eta(X)t^\eta$ with $S$ well ordered
and every $P_\eta$ a finite polynomial, with no common bound on
$\deg P_\eta$; and the \emph{valuation-restricted algebra}
\begin{equation}\label{ent:eq:Talg}
 \TT=\left\{\sum_{n\ge0}a_nZ^n\in\K[[Z]]:
  v(a_n)\to+\infty\text{ cofinally in }\Gamma\right\}.
\end{equation}
Taking the coefficient of each $X^n$ gives an injective ring map
$\A\hookrightarrow\K[[X]]$, and under it $\A$ is exactly the set of series
whose coefficient family $(a_n)_n$ is itself strongly summable: finite degree
at each Hahn exponent is precisely the pointwise-finiteness half of strong
summability. With $X$ and $Z$ identified under that map,
\begin{equation}\label{ent:eq:chain}
 \K[Z]\ \subseteq\ \E\ \subseteq\ \TT\ \subseteq\ \A\ \subseteq\ \K[[Z]].
\end{equation}
The first two inclusions need not be strict; the last two record genuinely
different summation requirements. The first two source manuscripts named the ring
\eqref{ent:eq:Halg} differently --- one as the auxiliary algebra $\A$ in the
variable $X$, the other as the strong disk algebra $\mathcal H_\Gamma$ in the
variable $Z$ --- but it is one ring, and we use the single symbol $\A$. That
the inclusions in \eqref{ent:eq:chain} hold, and that $\TT$ and $\A$ are
rings, follows from \cref{ent:lem:neumann} together with
\cref{ent:prop:operations} below.

For $A\ne0$ in $\A$, set $w(A)=\min\{\eta:P_\eta\ne0\}$; the polynomial
$P_{w(A)}$ is its \emph{initial polynomial}. Since $\C[X]$ is a domain,
\begin{equation}\label{ent:eq:gauss-mult}
 w(AB)=w(A)+w(B),
\end{equation}
and initial polynomials multiply. Equivalently, writing
$\inpol_0(A)=\overline{t^{-w(A)}A}\in\C[X]\setminus\{0\}$ for the reduction,
$\inpol_0(AB)=\inpol_0(A)\inpol_0(B)$. This Gauss multiplicativity is the
engine of \cref{ent:sec:units}.

\subsection{Notation: radii live in \texorpdfstring{$\Gamma$}{Gamma}, nodes
live in \texorpdfstring{$\K$}{K}}
\label{ent:subsec:notation}
Eight manuscripts are merged here, and six of them arrived with letters
that collide with this report's. The third arrived with the opposite
convention for the two letters that matter most. Its $\rho_n$ was a positive
element of the \emph{value group} and its $a_{n,j}$ was a point of the
\emph{field}. This report keeps the convention of the first two throughout,
in which $\rho$ always names a point of $\K$ and $\gamma$ always names an
element of $\Gamma$. Every statement taken from the third, fourth, fifth,
sixth, seventh or eighth manuscript has been rewritten in the convention below; the tables are
recorded so that a reader comparing the texts is not misled by a symbol that
typesets identically and means something else.

\begin{center}\small
\begin{tabularx}{\textwidth}{@{}P{0.155\textwidth}P{0.255\textwidth}YY@{}}
\toprule
Symbol here & Meaning & In the third manuscript & Note\\\midrule
$\gamma_n$ & shell radius, an element of $\Gamma$; $0<\gamma_1<\gamma_2
<\cdots$, cofinal & $\rho_n$ & \emph{inverted}: its $\rho_n$ is our
$\gamma_n$\\[0.35em]
$\rho_{n,j}$ & a node, a point of $\K$ with $v(\rho_{n,j})=-\gamma_n$;
written $\rho_n$ when the shell carries one node
& $a_{n,j}$ & \emph{inverted}: its $a_{n,j}$ is our $\rho_{n,j}$\\[0.35em]
$\tau_n=t^{-\gamma_n}$ & the shell normalizer, $v(\tau_n)=-\gamma_n$
& $s_n$ & \\[0.35em]
$\lambda_n$ & the valuation $v(g(\rho_n))$ of \eqref{ent:eq:lambda}
& --- & its $\lambda_n$ is a shell \emph{parameter}, here always written
$\gamma_{n+1}$\\[0.35em]
$m_{n,j}$, $M_n$ & node multiplicity and shell total $\sum_jm_{n,j}$
& $r_{n,j}$, $M_n$ & $r$ is reserved here for roots and for the jet
polynomials of \eqref{ent:eq:interpolation-r}\\[0.35em]
$m_\gamma(F)$ & weighted Gauss valuation \eqref{ent:eq:weighted},
$\min_n\{v(a_n)+n\gamma\}$ & $w_\beta(F)=\min_k\{v(f_k)-k\beta\}$
& \emph{sign-inverted}: $w_\beta=m_{-\gamma}$ with $\gamma=-\beta$\\
\bottomrule
\end{tabularx}
\end{center}

The fourth manuscript, like the third, writes $a_{n,j}$ for the nodes, and
it writes $\lambda_n$ for the radii. More seriously, two of its letters are
\emph{swapped} with two of ours: its $C_n$ is our convex subgroup $H_n$ and
its $H_n$ is a correction, our $C_n$. Several other letters shift. Every statement taken from it has been
rewritten as follows, and ``shell'' is used for its ``block''.

\begin{center}\small
\begin{tabularx}{\textwidth}{@{}P{0.2\textwidth}P{0.3\textwidth}P{0.17\textwidth}Y@{}}
\toprule
Symbol here & Meaning & In the fourth manuscript & Note\\\midrule
$H_n$ & convex subgroup generated by $\gamma_n$ & $C_n$ &
\emph{swapped} with the next row\\[0.3em]
$C_{n,d}$, $\widehat C_{n,d}$ & sequential correction \eqref{ent:eq:correction};
cardinal correction of \cref{ent:sm:subsec:cardinal} & ---, $H_{n,N}$ &
\emph{swapped}: its $H_n$ is our cardinal $\widehat C_n$\\[0.3em]
$\gamma_n$ & shell radius & $\lambda_n$ & our $\lambda_n$ is the valuation
\eqref{ent:eq:lambda}\\[0.3em]
$\rho_{n,j}$, $\xi_{n,j}$, $\tau_n$ & node, normalized node, normalizer &
$a_{n,j}$, $u_{n,j}$, $s_n$ & our $\xi_{n,j}$ never denotes a boundary
coordinate\\[0.3em]
$q_n$, $M_n$ & number of nodes and total degree of shell $n$ & $r_n$, $d_n$
& its $M_n$ is our $c_n$\\[0.3em]
$c_n$ & integer with $v\ge-c_n\gamma_n$, as in \eqref{ent:eq:maincondition}
& $M_n$ & its ``scale-moderate'' means eventually in $V_n$\\[0.3em]
$d$ & damping or suppression exponent & $N$ & \\[0.3em]
$P_n$, $\Pi_n$, $P_D$, $\widehat P_n$ & monic shell polynomial, shell factor,
canonical product, cofactor $P_D/\Pi_n$ & $W_n$, $p_n$, $P$, $Q_n$ & our
$Q_n$ is the product \eqref{ent:eq:Qn} over already treated nodes\\[0.3em]
$\Lambda_n$, $\Lambda_n^\circ$, $\RD$ & shell algebras, restricted product
& $A_n$, $B_n$, $\prod_n'(A_n,B_n)$ & our $\RD$ carries the finite factor
$\Lambda_0$\\[0.3em]
$R_n$, $\remp_n$, $\ev_D$ & normalized shell Hermite polynomial, shell class,
jet map & $J_n$, $\mathcal J_n$, $\mathcal J_D$ & our $J_{n,j}$ are the
prescribed jets\\[0.3em]
$\nu$, $\varpi_n$, $\varkappa_n$, $\Xi_{n,d}$ & least coefficient valuation,
cofactor class, inverse-shell gain, cardinal remainder & $\nu_n$, $q_n$,
$\beta_n$, $R_{n,N}$ & our $q_n$ counts nodes; our $\beta_i$ are reciprocal
node valuations\\[0.3em]
$m_\gamma(F)$ & weighted Gauss valuation & $w_r(F)$ & \emph{sign-inverted}:
$w_r=m_{-r}$; its $r\le\lambda_n/2$ is our $\gamma\ge-\gamma_n/2$\\[0.3em]
$\vartheta_n$ & displacement $\sigma_n-\rho_n=t^{\vartheta_n}$ & $\theta_n$
& our $\theta_n$ is a real radius in the second proof of
\cref{ent:thm:interpolation}\\[0.3em]
$\mathfrak q_{\gamma_n}$, $\kappa_n$ & coarsened maximal ideal and residue
field & $\mathfrak p_n$, $\kappa_n$ & \\[0.3em]
$\Phi_\Ucal$, $\iota_\Ucal$, $L_\Ucal$ & ultrafilter character, its
restriction to $\K$, its target field & $\chi_\Ucal$, $j_\Ucal$,
$\prod_\Ucal\kappa_n$ & its boundary coordinate $\xi_\Ucal$ is our
$\Phi_\Ucal(Z)$\\[0.3em]
nonprincipal; non-evaluation maximal ideal & & free; boundary ideal & not
the ``free ideals'' or disk boundaries of the analytic-geometry
report\\[0.3em]
$\Gamma_\infty=\bigoplus_{j\ge0}\Q e_j$ & $H_n$ has rank $n+1$ &
$\Gamma_*=\bigoplus_{j\ge1}\Q e_j$ & index shift: its group omits $e_0$, and
its $C_n$ has rank $n$\\
\bottomrule
\end{tabularx}
\end{center}

The fifth manuscript's letters are used only in \cref{ent:sec:multivariable}.
Its $H_n$, a homogeneous block, would have been a third meaning of that
letter, after our convex subgroup and the fourth source's correction; it is
renamed.

\begin{center}\small
\begin{tabularx}{\textwidth}{@{}P{0.2\textwidth}P{0.3\textwidth}P{0.17\textwidth}Y@{}}
\toprule
Symbol here & Meaning & In the fifth manuscript & Note\\\midrule
$F_{(n)}$ & homogeneous component of degree $n$ & $H_n$ & $H_n$ is the
convex subgroup here\\[0.3em]
$\kk$ & trivially valued coefficient field & $k$ & our $k$ is an
index\\[0.3em]
$\Sigma(F)$, $\Sigma_J$ & ordinary monomial support, its faces & $E(F)$,
$E_J$ & our $\E$ is the entire ring\\[0.3em]
$\Delta/H$, $\pi_H$, $\bar\delta$ & ordered quotient, quotient map,
quotient weights $\pi_H(v(z_j))$ & $Q$, $\pi$, $q$ & our $Q_n$, the $\pi$ of
\cref{ent:lem:coarsening}, and our $q_n$ are unrelated\\[0.3em]
$\inner\Sigma{\bar\delta}$, $\Adm_\Sigma(G)$ & weight set, admissible weight
set & $W_E(q)$, $\mathcal A_E(Q)$ & our $\A$ is the Hahn algebra\\[0.3em]
$\Dom$ & coarsened valuation ring & $\mathcal O_H$ & the same
object\\[0.3em]
$e_*$ & a new scale above $H$ & $\tau$ & our $\tau_n$ are
normalizers\\[0.3em]
$\Pdir(F)$ & polynomial-direction locus & $\operatorname{Pol}(F)$ & avoids
the nonabelian-support report's $\operatorname{Pol}$\\[0.3em]
$\mathcal Z_N$, $Y_m$, $\mathcal L$, $\varphi_n$ & vanishing loci,
prescribed algebraic sets, relation lattice, damping sequence & $Z_N$, $V_m$,
$L$, $\beta_n$ & our $Z$ is the variable, $V_n$ the coarsened rings\\[0.3em]
$\mathcal D^{\mathrm{hom}}(F)$, $\omega^x$ & homogeneous-block domain,
Conway omega-map & $\operatorname{Dom}_{\mathrm{hom}}(F)$, $\Omega(x)$ & \\
\bottomrule
\end{tabularx}
\end{center}

The sixth manuscript's letters are used only in \cref{ent:ps:sec:spectrum}.
Its nodes, radii, $H_n$, $V_n$, $\mathfrak q_n$ (our $\mathfrak
q_{\gamma_n}$), ultrafilters $\Ucal$, residue fields $L_\Ucal$, order units
and cofinality agree with ours; the following are renamed. The word
\emph{spectrum} means prime spectrum in that section, and \emph{free}
ultrafilters and fibres are called nonprincipal, as elsewhere here.

\begin{center}\small
\begin{tabularx}{\textwidth}{@{}P{0.2\textwidth}P{0.3\textwidth}P{0.17\textwidth}Y@{}}
\toprule
Symbol here & Meaning & In the sixth manuscript & Note\\\midrule
$\RD$, $\RD^\Gamma$ & sequence ring $\cong\mathcal E_\Gamma/(P_D)$ &
$\mathcal A$, $\mathcal A_\Gamma$ & our $\A$ is the Hahn algebra\\[0.3em]
$P_D$, $P_{D_{\mathbf m}}$ & simple and multiple one-node products & $P$,
$Q$ & our $P$ in \eqref{ent:eq:infinite-product} runs over $j\ge0$; our $Q_n$
is \eqref{ent:eq:Qn}\\[0.3em]
$\kk$ & $\R$ or $\C$ & $k$ & our $k$ is an index\\[0.3em]
$\mathfrak e_n$ & evaluation prime of $\RD$ & $\mathfrak p_n$ & the fourth
source's $\mathfrak p_n$ is our $\mathfrak q_{\gamma_n}$\\[0.3em]
$\mathfrak n_\Ucal$, $\overline{\mathfrak M}_\Ucal$ & minimal and maximal
prime of a fibre & $\mathfrak n_\Ucal$, $\mathfrak m_\Ucal$ & the inverse
image of $\overline{\mathfrak M}_\Ucal$ in $\E$ is our $\mathfrak M_\Ucal$;
our $\mathfrak m_\Gamma$ is the maximal ideal of $\OO$\\[0.3em]
$\mathfrak p_{\Ucal,\mathcal C}$, $\mathcal C$, $\mathcal C'$ & fibre prime,
convex subgroups & $\mathfrak p_{\Ucal,C}$, $C$, $D$ & our $C_n$ are
corrections, our $D$ a divisor\\[0.3em]
$\mathcal C_{\min}$, $\mathcal C_{\max}$, $\mathcal C_{\mathrm{int}}$,
$\mathcal C'_\theta$, $\mathcal C_\theta$ & convex extensions, splitting and
gap subgroups & $C_{\min}$, $C_{\max}$, $C$, $D_r$, $C_r$ & real parameter
$\theta\in(0,1)$ for its $r$\\[0.3em]
$\mathcal V_\Ucal$, $\mathcal K_\Ucal$, $\Gamma_\Ucal$, $w_\Ucal$ &
ultraproduct valuation ring, its fraction field, value group, valuation &
$W_\Ucal$, $F_\Ucal$, $G_\Ucal$, $w_\Ucal$ & our $W$ and $F$ are entire
functions\\[0.3em]
$\red_\Ucal$ & quotient map to $\mathcal V_\Ucal$ & $q_\Ucal$ & our $q_n$
counts nodes; our $\Phi_\Ucal$ is $\red_\Ucal$ followed by residues\\[0.3em]
$\bar\eta_n$, $\zeta_\theta$ & gap elements & $d_n$, $d_r$ & our $d$ is a
damping exponent\\[0.3em]
$m_{-r}(F)$ & weighted Gauss valuation & $\mu_r(F)$ & \emph{sign-inverted},
as in the first table\\[0.3em]
$\Delta_\Q=\bigoplus_{q\in\Q_{\ge0}}\Q e_q$ & rational-indexed
enlargement of $\Gamma_\infty$ & $\bigoplus\Q f_q$, $f_j=e_j$ & our $f$ is
\eqref{ent:eq:f-explicit}\\[0.3em]
$\No[i]$ & surcomplex numbers & $\No(i)$ & collection notation\\
\bottomrule
\end{tabularx}
\end{center}

The seventh manuscript's letters are used only in
\crefrange{ent:rc:subsec:setting}{ent:rc:subsec:limits}. Its order unit,
cofinality, strong summability, whole-family entireness and $\Gamma_\infty$
agree with ours, and its weighted Gauss value has the same sign as ours; the
following are renamed.

\begin{center}\small
\begin{tabularx}{\textwidth}{@{}P{0.2\textwidth}P{0.3\textwidth}P{0.17\textwidth}Y@{}}
\toprule
Symbol here & Meaning & In the seventh manuscript & Note\\\midrule
$\Ed$, $\E$, $\K$, $\kk$ & whole-family entire ring, one-variable ring,
$\kk((t^\Gamma))$, $\R$ or $\C$ & $\mathcal E_d$, $\mathcal E_1$, $K$, $k$ &
our $k$ is an index\\[0.3em]
$m_\gamma(F)$ & weighted Gauss valuation & $w_\gamma(F)$ & same sign, unlike
the third and fourth manuscripts\\[0.3em]
$v(p)$, $B^d_\gamma$ & $\min_jv(p_j)$, valuation polydisk & $r(p)$, polydisk
& $r$ is reserved here for roots and jet polynomials\\[0.3em]
$\mathcal X$, $p_i$ & configuration, its points & $D$, $a_i$ & our $D$ is a
divisor and $a_n$ are coefficients\\[0.3em]
$\rho_i$, $\beta_i$ & nodes, their radii $-v(\rho_i)$ & $z_i$, $\eta_i$ &
$\beta_i$ as in \cref{ent:prop:kernel-direct}\\[0.3em]
$\widehat P_i$, $\varpi_i$ & node cofactor, its value at $\rho_i$ & $P_i$,
$d_i$ & $d$ is the number of variables\\[0.3em]
$c_i$ & integer bounding a loss & $C_i$ & our $C_n$ are corrections\\[0.3em]
$\delta$ & shift scale of \cref{ent:rc:lem:shift} & $\beta$ & \\[0.3em]
$\Kle$ & polynomial-part ring $\kk((t^{\Gamma_{\le0}}))$ & $\mathcal A_\Gamma$
& our $\A$ is $\C[X]((t^\Gamma))$\\[0.3em]
$g_c$, $\Psi$ & constant linear map, triangular shear & $L$, $S$ & \\[0.3em]
$V_H$, $\mathfrak q_H$, $\kappa_H$, $\red_H$ & coarsened ring, its maximal
ideal, residue field, residue map & $\mathcal O_H$, $\mathfrak m_H$, $K_H$,
$\red_n$ & our $\mm$ is the maximal ideal of $\OO$; $V_{H_\gamma}=V_\gamma$\\[0.3em]
$\gamma_n$, $H_n$ & separated scales, $H_n=H_{\gamma_n}$ & $\eta_n$, $H_n$ &
the same convex subgroups\\[0.3em]
$\mathbf u_j$, $\mathbf w$ & coordinate vectors, a vector in
\cref{ent:rc:lem:taylor-remainder} & $e_j$, $u$ & our $e_j$ span
$\Gamma_\infty$; our $u_j$ are units\\[0.3em]
$\vartheta_n$ & perturbation exponent & $\theta_n$ & our $\theta_n$ is a real
radius\\[0.3em]
$\Jac\Phi$ & Jacobian matrix & $J_\Phi$ & our $J_{n,j}$ are prescribed
jets\\[0.3em]
$X'$ & transverse variables $(X_2,\ldots,X_d)$ & $Y$ & our $Y_m$ are
prescribed algebraic sets\\[0.3em]
$P_{\langle r\rangle}$, $D_{\langle r\rangle}$ & product and divisor of the
reduced multiplicities $(m_i-r)_+$ & $P_r$ & our $P_n$ are shell
polynomials\\[0.3em]
$\bar m$, $\mathsf h_i$ & maximal multiplicity, leading Taylor part & $M$,
$\tau_i$ & our $M_n$ are shell totals, our $\tau_n$ normalizers\\[0.3em]
$N_1$, $N_3$, $P_{N_1}$, $P_{N_3}$ & node sets and their products & $L$,
$H$, $Q_L$, $Q_H$ & our $H$ is a convex subgroup, our $Q_n$ is
\eqref{ent:eq:Qn}\\[0.3em]
$\Gamma\subseteq\No$ with \eqref{ent:eq:embedding} & surreal reading &
$\iota:\Gamma\hookrightarrow\No$ & \\[0.3em]
$\No[i]$ & surcomplex numbers & $\No(i)$ & collection notation\\
\bottomrule
\end{tabularx}
\end{center}

The eighth manuscript, source~11, is used only in
\crefrange{ent:as:subsec:setting}{ent:as:subsec:limits}. Its strong
summability, whole-family entireness, cofinality and Hasse derivatives agree
with ours, and its valuation has the same sign; the following are renamed.
The tempting false readings are recorded in the last column: its $A$ is not
our Hahn algebra $\A$, its $D$ is not a divisor, and its $H_F$ is not a convex
subgroup.

\begin{center}\small
\begin{tabularx}{\textwidth}{@{}P{0.2\textwidth}P{0.3\textwidth}P{0.17\textwidth}Y@{}}
\toprule
Symbol here & Meaning & In source~11 & Note\\\midrule
$\kk$, $\K$, $\Ed$, $\Ed^{+}$ & $\R$ or $\C$; $\kk((t^\Gamma))$, any
$\Gamma$; whole-family entire ring; its positive-scale part & $k$, $K$,
$\mathcal E_d(K)$, $\mathcal E_d^+(K)$ & our $k$ is an index\\[0.3em]
$\PiK$ & series supported in $\Gamma_{<0}$ & $\Pi$ & our $\Pi$ is the class of
purely infinite surreals and $\Pi_n$ a shell factor; the
set-sized-quotients report's $J_\Gamma$, sign reversed\\[0.3em]
$\Kle$ & $\kk+\PiK$, support in $\Gamma_{\le0}$ & $B$ & the seventh source's
polynomial-part ring; our $B$ is a bound, $B_\gamma$ a ball\\[0.3em]
$\Zk$ & $\Z$ for $\R$, $\Z[i]$ for $\C$ & $D$ & our $D$ is a divisor\\[0.3em]
$\OzK$ & $\Zk+\PiK$, the omnific integers of the workspace & $A$ & our $\A$ is
$\C[X]((t^\Gamma))$ and $A$ a B\'ezout cofactor; the set-sized-quotients
report's $A_\Gamma$\\[0.3em]
$\Int(\Zk^d,\Zk)$ & integer-valued polynomials & $\Int(D^d,D)$ & the
set-sized-quotients report's $\Int(\Z^r)$, $\Int(\Z[i]^r)$\\[0.3em]
$\eta$, $\gamma$, $b$ & Hahn scale, weight, bound & $\gamma$, $\rho$, $\beta$ &
our $\rho$ is a node; as in \cref{ent:thm:multi-criterion} and
\eqref{ent:eq:Halg}\\[0.3em]
$\alpha$, $\alpha'$; $\beta$, $\beta'$ & monomial indices; Hasse orders &
$\nu$, $\mu$; $\alpha$, $\beta$ (bold) & as in \cref{ent:mv:def:entire} and
\cref{ent:rc:lem:sum}\\[0.3em]
$\pg_\eta(F)$, $\ct$ & scale polynomial $[t^\eta]F$; coefficient at $0$ on
$\Kle$ & $p_\gamma(F)$, $\operatorname{ct}_t$ & our $p_i$ are points;
$\pg_\eta$ is the $P_\eta$ of \eqref{ent:eq:Halg}\\[0.3em]
$F_{\le0}$, $F_+$ & nonpositive truncation, positive part & $P_B(F)$, $F_+$ &
our $P_D$ is a canonical product\\[0.3em]
$\kk[X]$ & coefficient polynomial ring & $R_0$ & our $R_n$ are Hermite
polynomials\\[0.3em]
$\mathbf M$, $\mathbf N$, $\Syz$ & matrices over $\kk[X]$, syzygy generators,
syzygy module & $M$, $H$, $N$ & our $H$ is a convex subgroup; the $M$ of
\cref{ent:mv:prop:linear} is over $\K$\\[0.3em]
$\mathsf q_\eta$, $c$, $c_{\mathbf M}$ & polynomial lifts, degree-loss
constants & $q_\gamma$, $C$, $C_M$ & our $q_n$ count nodes, $C_n$ are
corrections\\[0.3em]
$\mathbf a$, $\mathbf L$, $Y$ & affine translation, linear part, new variables
& $\boldsymbol a$, $H$, $\boldsymbol Z$ & our $Z$ is the one-variable
coordinate, and $W$ a prepared polynomial\\[0.3em]
$G_j$ & positive-scale cofactors in \cref{ent:as:thm:restriction} & $H_j$ &
\\[0.3em]
$\mathcal S$, $\mathbf s$, $I_\kk(\mathcal S)$, $\mathbb V_\kk(I)$ & sample
set, ordinary point, vanishing ideal in $\kk[X]$, zero set in $\kk^d$ & $S$,
$\boldsymbol s$, $I(S)$, $\operatorname V_k(I)$ & our $S$ are well-ordered
sets and $s$ a scale; the seventh source's $I(\mathcal X)$ is an ideal of
$\Ed$; our $V_n$ are coarsened rings\\[0.3em]
\bottomrule
\end{tabularx}
\end{center}

\begin{center}\small
\begin{tabularx}{\textwidth}{@{}P{0.2\textwidth}P{0.3\textwidth}P{0.17\textwidth}Y@{}}
\toprule
Symbol here & Meaning & In source~11 & Note\\\midrule
$\gcd_+(F)$, $n_0$ & monic generator of $I_+(F)$, degree of a positive scale &
$H_F$, $D_0$ & our $H$ is a convex subgroup, $D_0$ a finite part\\[0.3em]
$f_j$ & polynomials of the sharpness example & $q_j$ & \\[0.3em]
$I_{+,m}(F)$, $m$ & jet ideal, jet order & $J_r(F)$, $r$ & our $J_{n,j}$ are
prescribed jets; $r$ is reserved for roots\\[0.3em]
$\nabla_j$, $\nabla^\beta$, $\mathbf u_j$, $b_\beta$ & forward differences,
coordinate vectors, binomial coefficients & $\Delta_j$, $\Delta^\alpha$,
$e_j$, $b_\alpha$ & our $\Delta$ is a larger value group, $e_j$ span
$\Gamma_\infty$; $\nabla$ is not a gradient\\[0.3em]
$u_n$, $nu_n$ & cofinal sequence, realization scales & $\beta_n$,
$\lambda_n$ & our $\lambda_n$ is the valuation \eqref{ent:eq:lambda}\\[0.3em]
$\mathtt M$, $h_{\mathtt M}$, $G_{\mathrm{sq}}$, $F_{\mathtt M}$ & machine,
halting indicator, base series, machine family & $M$, $h_M$, $G$, $F_M$ & $M$
is a matrix above\\[0.3em]
$F_{\mathrm{loc}}$, $P_{\mathrm{esc}}$, $\Theta$ & local counterexample,
escaping product, finite sample & $L$, $P$, $T$ & our $L_\Ucal$ are fields,
$P$ products, $T$ a variable, $\Theta_n$ a shell remainder\\
\bottomrule
\end{tabularx}
\end{center}

Two consequences of the last row of the first table are used repeatedly and
are recorded here once. First, $m_\gamma(F)$ is nondecreasing in $\gamma$, because each term
$v(a_n)+n\gamma$ is. Second, on polynomials and on entire series
\[
 m_\gamma(F+G)\ge\min\{m_\gamma(F),m_\gamma(G)\},\qquad
 m_\gamma(FG)\ge m_\gamma(F)+m_\gamma(G),
\]
with $m_\gamma(0)=+\infty$; the second is the Cauchy-product estimate, and
equality holds by \eqref{ent:eq:gauss-mult}, applied to the rescalings
$F(t^\gamma X),G(t^\gamma X)\in\A$ of \cref{ent:thm:criterion}(2). The first
proof of the image theorem never needs the equality; the cardinal second
route of \cref{ent:sm:subsec:cardinal} does, in the one form
$m_\gamma(\widehat P_n)=m_\gamma(P_D)$. Only these relations are used, never
a maximum-modulus principle.

Finally, we write $F^{[k]}(b)=[Z^k]F(Z+b)$ for the $k$th Hasse jet
coefficient at $b$, which is $F^{(k)}(b)/k!$ by \eqref{ent:eq:taylor}. For
$v(b)\ge\gamma$ the same expansion gives the bound
\begin{equation}\label{ent:eq:jet-bound}
 v\bigl(F^{[k]}(b)\bigr)\ \ge\ m_\gamma(F)-k\gamma ,
\end{equation}
since every summand $\binom nk a_nb^{n-k}$ has valuation at least
$v(a_n)+(n-k)\gamma=(v(a_n)+n\gamma)-k\gamma$.

\section{An exact coefficient criterion and the cofinality dichotomy}
\label{ent:sec:growth}

For $n\ge1$, division of $v(a_n)$ by $n$ is legitimate because $\Gamma$ is
divisible and torsion-free. We interpret $(+\infty)/n=+\infty$, so that zero
coefficients impose no condition.

\begin{theorem}[All-point strong summability]\label{ent:thm:criterion}
Let $\Gamma\ne0$ and $F=\sum_{n\ge0}a_nZ^n\in\K[[Z]]$. The following
conditions are equivalent.
\begin{enumerate}[label=\textnormal{(\arabic*)},leftmargin=2em]
\item $F\in\E$.
\item For every $\gamma\in\Gamma$, the family $(a_nt^{n\gamma})_{n\ge0}$ is
strongly summable; equivalently $F(t^\gamma X)\in\A$.
\item For every $h\in\Gamma$, eventually $v(a_n)>nh$.
\item For every $\gamma,B\in\Gamma$, eventually $v(a_n)+n\gamma>B$.
\end{enumerate}
In particular, strong entire evaluation is equivalent to the cofinal
condition $v(a_n)/n\to+\infty$, and it is enough to test the monomial points
$z=t^\gamma$.
\end{theorem}
\begin{proof}
The implication (1)$\Rightarrow$(2) is immediate by testing $z=t^\gamma$; the
reformulation of (2) inside $\A$ holds because the common Hahn support is the
union of the supports of the scaled coefficients, and the coefficient at a
fixed Hahn exponent is a polynomial in $X$ exactly when that exponent occurs
at only finitely many indices $n$. The zero series satisfies all four
conditions, so assume $F\ne0$.

Suppose (2) holds. Applying it at $\gamma=0$ shows that the union of the
supports of the $a_n$ is well ordered. Thus
\[
 m=\min\{v(a_n):a_n\ne0\}\in\Gamma
\]
exists. Suppose (3) fails. There is $h\in\Gamma$ for which infinitely many
positive integers $n$ satisfy $a_n\ne0$ and $v(a_n)\le nh$. Choose $q>0$ with
$q>|h|$ and $q>|m|$, where $|\cdot|$ is the ordered group absolute value.
By (2), the set of leading exponents
\[
 \{v(a_n)-2nq:a_n\ne0\}
\]
is well ordered. Its minimum is attained at some \emph{finite} index $k$;
write that minimum as $\mu$. Then
\[
 \mu=v(a_k)-2kq\ge m-2kq>-(2k+1)q.
\]
Choose a violating index $n>2k+1$. It gives
\[
 v(a_n)-2nq\le nh-2nq<-nq<-(2k+1)q<\mu,
\]
a contradiction. This proves (2)$\Rightarrow$(3). Notice that the argument
used a minimum attained at a finite index; it did not assume that the
multiples of $q$ are cofinal in $\Gamma$. See \cref{ent:app:audit}.

Suppose (3) holds and fix $\gamma,B$. Choose $r>0$ and set
$h=-\gamma+|B|+r$. Eventually
\[
 v(a_n)+n\gamma>n(|B|+r)>B
\]
for $n\ge1$. Thus (4) holds.

Finally, fix $z\ne0$ and put $\gamma=v(z)$. Condition (4) implies
$v(a_nz^n)=v(a_n)+n\gamma\to+\infty$. Apply
\cref{ent:lem:cofinal-tails}. At $z=0$ there is only the constant term.
This proves (4)$\Rightarrow$(1).
\end{proof}

\begin{remark}\label{ent:rem:criterion-strength}
Condition (3) is strictly stronger than $v(a_n)\to+\infty$, that is, than
membership in $\TT$. Over $\Gamma=\Q$ the coefficients $a_n=t^n$ tend to zero
in valuation, but $\sum t^nZ^n$ is not entire: at $Z=t^{-1}$ its terms all
have exponent zero. Conversely $a_n=t^{n^2}$ satisfies (3) over $\Q$ and over
$\R$, but not over a group containing an element larger than every real
number. That last example is exactly the one recorded in the foundations
report of this collection; see \cref{ent:sec:collection}.
\end{remark}

\begin{corollary}[Cofinality dichotomy]\label{ent:cor:cofinality}
For $\Gamma\ne0$, the following are equivalent: $\E\ne\K[Z]$; $\E$ contains a
nonpolynomial series; $\Gamma$ has a countable cofinal subset;
$\cf(\Gamma)=\aleph_0$. If $\cf(\Gamma)>\aleph_0$, every strongly entire
ordinary power series is a polynomial.
\end{corollary}
\begin{proof}
If infinitely many $a_n$ are nonzero, their valuations divided by their
positive indices give, by \cref{ent:thm:criterion}, a countable cofinal
subset of $\Gamma$. Conversely, choose a positive increasing cofinal sequence
$(q_n)_{n\ge1}$. Then
\[
 1+\sum_{n\ge1}t^{nq_n}Z^n
\]
satisfies the coefficient criterion and is nonpolynomial. A nonzero ordered
group has no finite cofinal subset, because it has no greatest element.
\end{proof}

\begin{example}[Two different sources of polynomial rigidity]
\label{ent:ex:two-rigidities}
If $\Gamma=0$, every nonzero coefficient has support $\{0\}$, so strong
summability at $1$ forces a polynomial. If $\Gamma$ has uncountable
cofinality, there are many infinitesimal coefficients, but countably many
coefficient slopes cannot reach all of its scales. The conclusion is again
polynomiality, for a different reason. The same countability argument shows
that in the uncountable-cofinality case even the larger algebra $\TT$
collapses to $\K[Z]$; see \eqref{ent:eq:uncountable-rigidity}.
\end{example}

\begin{proposition}[Entire operations]\label{ent:prop:operations}
$\E$ is a $\K$-subalgebra of $\K[[Z]]$, and so are $\TT$ and $\A$.
Evaluation at each $z\in\K$ is a homomorphism on $\E$. Translation
$F(Z)\mapsto F(Z+b)$, differentiation, and the primitive with zero constant
coefficient preserve $\E$. The translation coefficients are
\begin{equation}\label{ent:eq:taylor}
 [Z^k]F(Z+b)=\sum_{n\ge k}\binom nk a_n b^{n-k}
 =\frac{F^{(k)}(b)}{k!}.
\end{equation}
Entire composition $F\circ G$ is also defined and belongs to $\E$. Moreover,
at each fixed scale $s\in\Gamma$ the partial sums of $F\in\E$ converge to $F$
in the weighted Gauss valuation $m_s$ of \eqref{ent:eq:weighted}.
\end{proposition}
\begin{proof}
The ring assertion for $\E$ follows from \cref{ent:thm:criterion}(2):
addition and multiplication preserve every rescaled copy of $\A$. Explicitly,
at a fixed $s$ write $A_i=v(a_i)+is$ and $B_j=v(b_j)+js$ for two members;
both tend cofinally to infinity and have lower bounds $A,B$. Given $\beta$,
choose $N$ with $A_i>\beta-B$ and $B_i>\beta-A$ for $i>N$; if $n>2N$ then
each pair $i+j=n$ has $i>N$ or $j>N$, so the convolution coefficient has
weighted valuation greater than $\beta$. Taking $s=0$ gives the ring
assertion for $\TT$; the same tail estimate gives the stated weighted
convergence of partial sums. That $\A$ is a ring is
\cref{ent:lem:neumann}(a).

Evaluation can be performed by choosing $\gamma=v(z)$, writing
$z=t^\gamma x$ with $x\in\OO$, and applying \cref{ent:lem:A-units}.

For translation and a fixed testing scale $\gamma$, expand
$F(b+t^\gamma X)$ into its finite binomial blocks. Put
$\beta=\min\{v(b),\gamma\}$, with $v(0)=+\infty$. Every coefficient in the
$n$th block has valuation at least $v(a_n)+n\beta$, which tends cofinally to
infinity, and each block has finitely many $X$-coefficients. The
support-union and local-finiteness proof of \cref{ent:lem:cofinal-tails}
therefore applies to the blocks with polynomial coefficients. This gives an
element of $\A$ at every scale and proves \eqref{ent:eq:taylor}. Regrouping
the same jointly summable family proves the evaluation identity and the
inverse translation by $-b$.

For differentiation, $v(na_n)=v(a_n)$ for positive $n$, since $n$ is a
nonzero complex coefficient of valuation zero. Shifting the index by one
preserves the criterion in \cref{ent:thm:criterion}: for a fixed $h$, choose
$q>0$ with $q>|h|$; the bound $v(a_n)>n(h+q)$ for large $n$ implies the
required bounds with $n-1$ or $n+1$ in place of $n$. No hidden Archimedean
estimate is used. Division of $a_n$ by $n+1$ also preserves valuation, so
termwise integration preserves the criterion as well.

For composition, fix $\gamma$ and let $G_\gamma=G(t^\gamma X)\in\A$. If
$G_\gamma\ne0$, put $\mu=w(G_\gamma)$. Then
$w(a_nG_\gamma^n)=v(a_n)+n\mu\to+\infty$, so the sum is strongly summable in
$\A$. Its coefficients define the same formal composite at every scale, since
rescaling commutes with each coefficient's strong sum. The case $G=0$ is
constant. Evaluation and associativity follow from the same
support-controlled regrouping.
\end{proof}

\emph{Related (batch~34).} For a polynomial inner function the
holonomic-rigidity report proves the composition clause again
(\texttt{hol:pc:prop:composition} \cite{RepoHolonomic}) and computes the
composite exactly at large scales (\texttt{hol:pc:prop:pullback}). If
$P=cZ^d+\cdots$ with $d\ge1$, then at every sufficiently large scale the
Gauss value, the active exponents and the residue polynomial of $F\circ P$
are those of $F$ at a rescaled scale, with the active exponents multiplied by
$d$ and the residue polynomial composed with $\bar cX^d$. That report writes
a scale with the sign opposite to the one used here ($z=t^{-\delta}X$), and it
uses this profile for its rigidity theorems on equations with polynomial
pullbacks. Nothing here depends on it.

\section{Restricted series: the sharp order-unit criterion}
\label{ent:sec:units}

The trichotomy of \cref{ent:thm:main} splits its two nonpolynomial cases by
the presence of an order unit. This section proves that the order unit is not
an artefact of the proof technique of \cref{ent:sec:positive}: it is exactly
the condition for a nonconstant element of the restricted algebra $\TT$ of
\eqref{ent:eq:Talg} to be invertible there. That is what makes case (ii) of
the trichotomy sharp rather than merely sufficient, and it is a statement
about $\TT$, not about $\E$; the ring discipline of \cref{ent:conv:ring}
applies.

\begin{theorem}[Restricted-unit criterion]\label{ent:thm:units}
Let $f=\sum a_nZ^n\in\TT$ be nonconstant. Then $f$ is a unit of $\TT$ if and
only if $a_0\ne0$ and
\begin{equation}\label{ent:eq:unit-gap}
 \delta(f)=\min_{n\ge1,\ a_n\ne0}v(a_n/a_0)
\end{equation}
is a positive order unit of $\Gamma$. Nonzero constant series are always
units, and the minimum in \eqref{ent:eq:unit-gap} exists whenever $a_0\ne0$.
\end{theorem}
\begin{proof}
The minimum exists because the coefficient valuations tend cofinally to
infinity: below the valuation of any chosen nonzero coefficient there are
only finitely many nonzero coefficients.

Suppose $fg=1$ with $g\in\TT$. The constant coefficient forces $a_0\ne0$.
Divide $f$ by $a_0$ and multiply $g$ by $a_0$, so that both constant
coefficients are $1$. Then $w(f),w(g)\le0$, while \eqref{ent:eq:gauss-mult}
gives their sum as zero. Hence both are zero. Their reductions in $\C[X]$
multiply to $1$, so both reductions are constant. Every nonconstant
coefficient of $f$ and of $g$ consequently has positive valuation; in
particular $\delta=\delta(f)>0$.

Assume for contradiction that $\delta$ is not an order unit. The convex
subgroup it generates,
\begin{equation}\label{ent:eq:Hdelta}
 H_\delta=\{\gamma\in\Gamma:|\gamma|\le m\delta
                       \text{ for some }m\in\N\},
\end{equation}
is then proper. Coarsen the valuation to the ordered quotient
$\Gamma/H_\delta$. Every coefficient of $f$ and $g$ is integral for this
coarsening. Their coarsened reductions are \emph{polynomials} over the
coarsened residue field: choose $\beta>H_\delta$, meaning $\beta>h$ for every
$h\in H_\delta$; since coefficient valuations tend cofinally to infinity in
$\Gamma$, only finitely many coefficients can have value at most $\beta$, and
all the remaining ones reduce to zero.

The reduced polynomial of $f$ is nonconstant, because a nonconstant
coefficient of valuation $\delta\in H_\delta$ has nonzero coarsened residue.
On the other hand the reduced polynomials of $f$ and $g$ multiply to $1$ in a
polynomial ring over a field, and a nonconstant polynomial cannot be a unit
there. This contradiction proves that $\delta$ is an order unit.

Conversely, suppose $a_0\ne0$ and $\delta(f)$ is a positive order unit.
Write $f=a_0(1+h)$, where $h$ has zero constant coefficient and
$w(h)=\delta$. The formal inverse is
\begin{equation}\label{ent:eq:restricted-inverse}
 f^{-1}=a_0^{-1}\sum_{m\ge0}(-h)^m.
\end{equation}
Its summands have Gauss valuation $m\delta\to+\infty$, so the sum exists
strongly in $\A$. To prove that it belongs to $\TT$, fix a valuation
threshold $\beta$. All sufficiently large $m$ contribute only coefficients of
valuation greater than $\beta$; the remaining finitely many powers of $h$
belong to $\TT$, so their coefficients also eventually exceed $\beta$. Thus
the coefficients of \eqref{ent:eq:restricted-inverse} tend cofinally to
infinity. The finite geometric identity, followed by the justified strong
sum, proves $ff^{-1}=1$.
\end{proof}

The coarsening argument is what makes this a \emph{necessity} theorem rather
than the observation that one particular geometric-series construction fails.
It shows that no cancellation in a hypothetical inverse can evade the
order-unit obstruction.

\begin{corollary}[Unit group of the restricted algebra]
\label{ent:cor:unit-structure}
The algebra $\TT$ has a nonconstant unit if and only if $\Gamma$ has an order
unit; otherwise $\TT^\times=\K^\times$. The rule ``the constant coefficient
strictly dominates all the others'' is a sufficient unit criterion for every
restricted series if and only if $\Gamma$ is Archimedean.
\end{corollary}
\begin{proof}
An order unit $\delta$ makes $1-t^\delta Z$ a nonconstant unit by the
theorem, and the converse follows by applying the theorem to any nonconstant
unit. An ordered group is Archimedean exactly when every positive element has
cofinal positive integer multiples; apply that characterization to
$1-t^\delta Z$ for the last assertion.
\end{proof}

\begin{example}[Zero-free does not mean restricted-invertible]
\label{ent:ex:ranktwo-unit}
Let $\Gamma=\Q\epsilon\oplus\Q\eta$ with $0<\epsilon\ll\eta$, the $\eta$
coordinate being compared first. The polynomial
\[
 f(Z)=1-t^\epsilon Z
\]
has residue $1$ and no zero in $\OO$: its unique zero $t^{-\epsilon}$ has
negative valuation. Yet its only formal inverse,
$\sum_{n\ge0}t^{n\epsilon}Z^n$, is not restricted, since $n\epsilon<\eta$ for
every $n$. It does belong to $\A$ and can be strongly evaluated on the
integral disk. By contrast $1-t^\eta Z$ has a restricted inverse, because
$n\eta$ is cofinal.
\end{example}

\begin{remark}[Why a rank-one contraction proof may not be imported]
\label{ent:rem:no-contraction}
\Cref{ent:ex:ranktwo-unit} is the reason the preparation theorem of
\cref{ent:sec:prep} is proved by well-founded recursion over a well-ordered
monoid and not by a convergent iteration. A rank-one Banach argument
containing the phrase ``the valuation increases by a positive amount at each
step'' is not valid here: the standard iteration may improve the error by
$\epsilon$ repeatedly while remaining forever below $\eta$, and its purported
convergence is then simply unproved. For a valid valuation-convergence proof
the positive amount must be an order unit. The preparation proof below uses
positive \emph{supports} and finite-word control instead, and works with no
order unit at all. The same care applies to \cref{ent:sec:positive}, whose
real-valued coarsening is constructed only after an order unit has been
assumed.
\end{remark}

\begin{remark}[Two existence questions, genuinely different]
\label{ent:rem:two-existence}
Nonpolynomial entire series need only $\cf(\Gamma)=\aleph_0$
(\cref{ent:cor:cofinality}); a nonconstant restricted unit needs an order
unit (\cref{ent:cor:unit-structure}). The group $\Gamma_\infty$ of
\cref{ent:sec:examples} has the first and not the second, so the two
questions separate. In topological language: the valuation topology of $\K$
has a countable neighbourhood base at zero exactly when
$\cf(\Gamma)=\aleph_0$, whereas a nonzero topologically nilpotent element
exists exactly when $\Gamma$ has an order unit, since $v(x^n)=nv(x)$. Neither
of these underlying topological facts is claimed as new; the higher-rank
framework in which they sit is discussed by Conrad \cite{Conrad}.
\end{remark}

\section{Support-controlled preparation}
\label{ent:sec:prep}

\subsection{The algebra that retains support certificates}
Recall $\A=\C[X]((t^\Gamma))$ from \eqref{ent:eq:Halg}. The support facts in
\cref{ent:lem:neumann} justify multiplication in $\A$ and show that the
coefficient map $\A\hookrightarrow\K[[X]]$ respects it.

\begin{lemma}[Units and integral substitution]\label{ent:lem:A-units}
A nonzero element of $\A$ is a unit if and only if its initial polynomial is a
nonzero constant. Every $A\in\A$ can be evaluated at any $x\in\OO$, and this
evaluation is a ring homomorphism; if $A\ne0$ the value has valuation at
least $w(A)$. For $x\in\OO$, substitution $X\mapsto X+x$ is an automorphism
of $\A$.
\end{lemma}
\begin{proof}
If the initial polynomial is $c\ne0$ at exponent $\eta$, write
$A=ct^\eta(1+u)$ with $w(u)>0$ and invert $1+u$ using
\cref{ent:lem:neumann}(c). Conversely, if $AB=1$, the product of the two
initial polynomials is $1$, so both are nonzero constants.

For substitution, put $S=\supp_t A$ and write $x=c+\varepsilon$ with $c\in\C$ and
$v(\varepsilon)>0$, treating $\varepsilon=0$ separately. For each $\eta$,
expand the finite polynomial $P_\eta(X+c+\varepsilon)$ in powers of
$\varepsilon$. All supports lie in $S+(\supp\varepsilon)^*$, which is well
ordered. At a fixed exponent there are only finitely many decompositions from
$S+(\supp\varepsilon)^*$, finitely many words in $\supp\varepsilon$ for each
such decomposition, and finitely many terms from each $P_\eta$. Thus
substitution is a legitimate element of $\A$. Finite polynomial identities
and strong regrouping prove multiplicativity, and its inverse is substitution
by $-x$. Evaluation follows by setting $X=0$ after this substitution. If
$A\ne0$, every contribution has exponent at least $\min S=w(A)$, because
$v(\varepsilon)>0$; cancellation can only raise the valuation. This gives
the claimed bound without any assertion of cofinal tail convergence.
\end{proof}

\subsection{Preparation without an Archimedean contraction}
\begin{theorem}[Support-controlled preparation]\label{ent:thm:prep}
Let $A\in\A$ satisfy $w(A)=0$, and suppose its initial polynomial
$P\in\C[X]$ is monic of degree $d$. There are unique $W,U$ such that
\[
 A=WU,\qquad
 W=P+R,\quad R\in\mm[X],\quad\deg R<d,\qquad
 U=1+Q,\quad w(Q)>0.
\]
Here $W$ is a monic polynomial of degree $d$ and $U$ is a unit of $\A$. If
$S=\supp_t(A-P)\subseteq\Gamma_{>0}$, both corrections are supported in
$S^*\setminus\{0\}$.
\end{theorem}
\begin{proof}
Let $M=S^*$, a well-ordered monoid. Write $A=P+\sum_{\eta>0}A_\eta t^\eta$,
with $A_\eta=0$ off $S$. Recursively for $\eta\in M\setminus\{0\}$, define
polynomials $R_\eta,Q_\eta$ by ordinary division by $P$:
\begin{align}
 H_\eta&=A_\eta-
   \sum_{\substack{\alpha+\beta=\eta\\\alpha,\beta>0}}
       R_\alpha Q_\beta,\label{ent:eq:prep-rec1}\\
 H_\eta&=P Q_\eta+R_\eta,\qquad\deg R_\eta<d.
 \label{ent:eq:prep-rec2}
\end{align}
The sum is finite by \cref{ent:lem:neumann}, and every index on its right is
smaller than $\eta$. Thus this is a well-founded recursion on a set. Set
$R=\sum R_\eta t^\eta$ and $Q=\sum Q_\eta t^\eta$. Because the remainders
have a uniform degree bound, $R$ is a finite polynomial over $\mm$, and
$Q\in\A$. Equations
\eqref{ent:eq:prep-rec1}--\eqref{ent:eq:prep-rec2} give $A=(P+R)(1+Q)$
coefficient by coefficient, and the second factor is a unit by
\cref{ent:lem:A-units}.

For uniqueness, suppose two pairs differ. The union of their positive
supports is well ordered, so there is a least exponent $\eta$ at which either
pair differs. At that exponent, equality of the products gives
\[
 P(Q_\eta-Q'_\eta)+(R_\eta-R'_\eta)=0.
\]
The second summand has degree less than $d$. Uniqueness of polynomial
division forces both differences to vanish, a contradiction. For $d=0$, the
normalization gives $P=1$, $W=1$ and $U=A$.
\end{proof}

This proof is deliberately not a claim that a Banach iteration converges; see
\cref{ent:rem:no-contraction}. It remains valid when no positive element of
$\Gamma$ has cofinal integer multiples. The uniformly bounded remainder
degree is the mechanism that extracts a finite zero scheme from an infinite
Hahn series.

\begin{corollary}[The certificate survives enlargement of the value group]
\label{ent:cor:prep-extension}
Let $\Gamma\subseteq\Delta$ be an ordered-group inclusion and regard
$A\in\A$ as an element of $\mathcal A_\Delta$ under the canonical inclusion.
The factorization $A=WU$ of \cref{ent:thm:prep} computed over $\Gamma$ is
also the factorization computed over $\Delta$: the same $W$ and $U$, with the
same support certificate $S^*\setminus\{0\}$.
\end{corollary}
\begin{proof}
The recursion \eqref{ent:eq:prep-rec1}--\eqref{ent:eq:prep-rec2} is indexed
by $S^*$, computed inside $\C[X]$, and refers to no element of the value
group beyond the exponents already occurring in $S$. It therefore produces
the same data over $\Delta$, and uniqueness over $\Delta$ identifies the two
factorizations.
\end{proof}

\Cref{ent:cor:prep-extension} is not a formality: it is exactly what lets the
same prepared polynomial be used on both sides of a scalar extension in
\cref{ent:thm:zero-conservation}, and it is available only because the
preparation was proved by a support recursion rather than by an iteration
whose convergence would have to be re-established in the larger field.

\begin{example}[A concrete preparation]\label{ent:ex:cubic}
Take $P=X^2-X$, $d=2$, and
\begin{equation}\label{ent:eq:cubic}
 A=X^2-X+t(X^3+1)\in\A,
\end{equation}
so that $S=\{1\}$ and $S^*=\N$. The recursion gives $Q_1=X+1$, $Q_2=-1$,
$Q_3=3$ and the first three remainders
\[
 R_1=X+1,\qquad R_2=-3X-1,\qquad R_3=8X+2,
\]
each of degree less than $2$ as required. For instance
$X^3+1=(X^2-X)(X+1)+(X+1)$ gives the first step, and
$H_2=-R_1Q_1=-(X+1)^2=(X^2-X)(-1)+(-3X-1)$ the second. The recursion is run
to perturbation order eight by the accompanying finite check suite
(\cref{ent:app:checks}).
\end{example}

\section{Exact root counting at every scale}
\label{ent:sec:zeros}

\subsection{Initial polynomials on a ball}
For $\gamma\in\Gamma$, define the closed valuation ball
\[
 B_\gamma=\{z\in\K:v(z)\ge\gamma\}.
\]
A larger ball corresponds to a smaller, typically negative, $\gamma$. For
$F\in\E\setminus\{0\}$, put
\begin{align}
 m_\gamma(F)&=\min_{n:a_n\ne0}\{v(a_n)+n\gamma\},\label{ent:eq:weighted}\\
 I_\gamma(F)(X)&=
 \sum_{v(a_n)+n\gamma=m_\gamma(F)}\lc(a_n)X^n.
 \label{ent:eq:initial}
\end{align}
The minimum is attained, and attained at only finitely many indices, and the
sum is a finite nonzero polynomial, precisely because the weighted
coefficient valuations diverge (\cref{ent:thm:criterion}). Equivalently
$I_\gamma(F)$ is the reduction of $t^{-m_\gamma(F)}F(t^\gamma X)$. Write
$d_\gamma(F)=\deg I_\gamma(F)$ and $e_\gamma(F)=\ord_{X=0}I_\gamma(F)$.

A zero's multiplicity is the order of the translated formal power series
$F(z+Z)\in\K[[Z]]$; the translation identity \eqref{ent:eq:taylor} makes this
intrinsic.

\begin{corollary}[Finite zero geometry on the integral disk]
\label{ent:cor:integral-zeros}
Let $A\in\A$ be nonzero with $w(A)=0$ and initial polynomial monic of degree
$d$ after normalization, and let $W$ be its prepared polynomial. The zeros of
$A$ in $\OO$ are exactly the $d$ roots of $W$ with multiplicity, all of them
integral, and for each $c\in\C$ the number of them with residue $c$ is the
multiplicity of $c$ as a root of the initial polynomial. In particular a
nonzero element of $\A$ has only finitely many zeros on the integral disk.
No complex compactness is used.
\end{corollary}
\begin{proof}
This is the content of the first two paragraphs of the proof of
\cref{ent:thm:roots} below, applied at $\gamma=0$.
\end{proof}

\begin{theorem}[Ball, shell, and residue-direction counts]
\label{ent:thm:roots}
For $F\in\E\setminus\{0\}$ and $\gamma\in\Gamma$:
\begin{enumerate}[label=\textnormal{(\alph*)},leftmargin=2em]
\item $B_\gamma$ contains exactly $d_\gamma(F)$ zeros of $F$, counted with
multiplicity.
\item The open ball $\{z:v(z)>\gamma\}$ contains exactly $e_\gamma(F)$ zeros,
and the shell $\{z:v(z)=\gamma\}$ contains $d_\gamma(F)-e_\gamma(F)$ zeros.
\item For each $c\in\C$, the zeros for which $v(z)\ge\gamma$ and
$\res(t^{-\gamma}z)=c$ have total multiplicity $\ord_{X=c}I_\gamma(F)$.
\end{enumerate}
\end{theorem}
\begin{proof}
Let $c_0$ be the leading coefficient of $I_\gamma(F)$ and apply
\cref{ent:thm:prep} to
\[
 A=\frac{F(t^\gamma X)}{c_0t^{m_\gamma(F)}}.
\]
It gives $A=WU$ with $W\in\OO[X]$ monic of degree $d_\gamma(F)$,
$\res(W)=I_\gamma(F)/c_0$, and $U$ of initial polynomial $1$. The field $\K$
is algebraically closed by the standard Hahn-field theorem
\cite[Corollary~4]{Poonen}, so $W$ has exactly its degree many roots in $\K$,
counted with multiplicity.

Every root $x$ of $W$ is integral: if $v(x)<0$, the leading term $x^d$ has
valuation $dv(x)$, strictly smaller than that of every lower-degree term
whose coefficient is integral, and such terms cannot sum to zero. On $\OO$,
evaluation of $U$ has residue $1$ and is nonzero. After translation by an
integral root, $U$ has nonzero constant term, so its formal local power
series is a unit. Hence zeros and multiplicities of $A$ in $\OO$ coincide
exactly with those of $W$.

Factor $W=\prod_j(X-x_j)$ and reduce modulo $\mm$. Then
$\res(W)=\prod_j(X-\res(x_j))$, which proves the residue-direction count.
The direction $c=0$ is exactly the open ball, and the nonzero directions make
up the shell. Rescale back by $Z=t^\gamma X$.
\end{proof}

\begin{corollary}[Units, faithfulness, and values]\label{ent:cor:units}
A nonzero entire function without zeros is a nonzero constant. Distinct
elements of $\E$ define distinct functions on $\K$. Every nonconstant element
of $\E$ takes every value in $\K$.
\end{corollary}
\begin{proof}
If a zero-free $F$ has $a_0=0$ it vanishes at zero, a contradiction. If
$a_n\ne0$ for some $n>0$, choose $\delta>0$ and
\[
 \gamma=\frac{v(a_0)-v(a_n)}{n}-\delta.
\]
Then $v(a_n)+n\gamma<v(a_0)$, so the constant coefficient is not an active
term in \eqref{ent:eq:initial} and the initial polynomial has positive
degree; \cref{ent:thm:roots} supplies a zero. Therefore $F$ is constant.

For faithfulness, a nonzero $F$ has only finitely many zeros in any
$B_\gamma$, whereas that ball contains infinitely many elements $t^\gamma c$
with $c\in\C$; apply this to the difference of two series. Finally, for any
$b\in\K$ a nonconstant $F-b$ cannot be zero-free, so $F$ takes the value $b$.
\end{proof}

\begin{lemma}[Division by a linear zero factor]\label{ent:lem:linear-division}
If $F\in\E$ and $F(r)=0$, then $F=(Z-r)G$ with $G\in\E$.
\end{lemma}
\begin{proof}
Set $H(Z)=F(r+Z)$, which is entire by \cref{ent:prop:operations}, and note
that $H(0)=0$. For each $\gamma$, write
$H(t^\gamma X)=\sum_\eta P_\eta(X)t^\eta\in\A$.
Its constant coefficient is zero, so every $P_\eta$ is divisible by $X$.
Dividing these finite polynomials by $X$ preserves the well-ordered Hahn
support and gives an element of $\A$. Consequently
\[
 (H/Z)(t^\gamma X)=t^{-\gamma}\,H(t^\gamma X)/X\in\A.
\]
By \cref{ent:thm:criterion}(2), $H/Z$ is entire. Translating back gives
$G(Z)=(H/Z)(Z-r)\in\E$ and $F=(Z-r)G$.
\end{proof}

\subsection{Stability of the counts}
\begin{corollary}[Initial-form stability]\label{ent:cor:stability}
If $F,G\in\E$ satisfy $w((F-G)(t^\gamma X))>m_\gamma(F)$, then
$m_\gamma(G)=m_\gamma(F)$ and $I_\gamma(G)=I_\gamma(F)$. All three counts in
\cref{ent:thm:roots} are therefore unchanged.
\end{corollary}
\begin{proof}
A perturbation supported strictly above the leading exponent changes no
coefficient at that exponent; apply \cref{ent:thm:roots}.
\end{proof}

\begin{proposition}[Locally finite Newton profile]\label{ent:prop:newton}
Let $F\in\E\setminus\{0\}$ and let $[\gamma_0,\gamma_1]$ be an interval of
$\Gamma$. On that interval, $\gamma\mapsto m_\gamma(F)$ is the minimum of
\emph{finitely many} affine functions $\gamma\mapsto v(a_n)+n\gamma$. The set
of active indices changes only at the slopes
\[
 \frac{v(a_i)-v(a_j)}{j-i}\in\Gamma\qquad(i<j),
\]
and between two consecutive such values the three counts of
\cref{ent:thm:roots} are constant.
\end{proposition}
\begin{proof}
By \cref{ent:thm:criterion}(4) applied at $\gamma_0$ and at $\gamma_1$, only
finitely many indices $n$ can attain the minimum \eqref{ent:eq:weighted} for
some $\gamma$ in the interval: all but finitely many satisfy
$v(a_n)+n\gamma>m_{\gamma_0}(F)+|m_{\gamma_1}(F)-m_{\gamma_0}(F)|$ throughout
it. A minimum of finitely many affine functions is piecewise affine, and two
of them exchange dominance exactly at the displayed slope, which lies in
$\Gamma$ by divisibility. \Cref{ent:cor:stability} gives the constancy.
\end{proof}

This is a valuation statement, a non-Archimedean analogue of a Rouch\'e
stability argument; it is not a theorem about Euclidean circles in the fine
surcomplex topology. Its purpose here is to retain exact multiplicities in
the factorization argument of \cref{ent:sec:products}.

\section{Radial divisors, canonical products and factorization}
\label{ent:sec:products}

\subsection{Which zero sets are admissible?}
\begin{definition}\label{ent:def:radial}
An effective divisor on $\K$ is a function $D:\K\to\N$ with set-sized
support. It is \emph{radially finite}, equivalently \emph{ball-finite}, if
\[
 \sum_{z\in B_\gamma}D(z)<\infty\qquad(\gamma\in\Gamma).
\]
The sum means a finite total multiplicity, not a Hahn sum. Let
$\Div^+_{\mathrm{rad}}(\K)$ denote the monoid of these divisors.
\end{definition}

The first two source manuscripts used the two displayed adjectives for this one
condition; we keep \emph{radially finite} and record \emph{ball-finite} as
the synonym. This is a restriction on valuation balls, not on discrete
subsets for the fine surreal topology, and the distinction is necessary:
fine discreteness alone would impose no useful control on an infinite
collection inside a fixed bounded valuation ball. Radial finiteness is
exactly what the root count of \cref{ent:thm:roots} supplies.

\begin{lemma}[Countability and escape]\label{ent:lem:divisor-countable}
An infinite radially finite divisor has countably many points counted with
multiplicity. Its nonzero points can be enumerated as $(\rho_j)_{j\ge1}$,
with repetitions for multiplicities, and
\[
 v(\rho_j)\longrightarrow-\infty \quad\hbox{coinitially in }\Gamma.
\]
In particular an infinite radially finite divisor forces
$\cf(\Gamma)=\aleph_0$; and when $\cf(\Gamma)>\aleph_0$, every radially
finite divisor is finite.
\end{lemma}
\begin{proof}
Choose countably many distinct occurrences from the infinite divisor. Their
valuations cannot have a common lower bound, for otherwise one ball would
contain infinitely many occurrences. They form a countable coinitial subset
of $\Gamma$, whose negatives are a countable cofinal subset. Balls at those
scales exhaust $\K$ and each contains only finitely many divisor
occurrences, so their union contains the whole divisor and is countable. In
any one-to-one enumeration of its occurrences, only finitely many indices
have valuation at least any fixed $\gamma$, which is the asserted escape. The
occurrence at zero, if present, has finite multiplicity and is removed before
enumerating the nonzero points.
\end{proof}

\subsection{Construction and exact factorization}
\begin{theorem}[Canonical product for a radial divisor]
\label{ent:thm:product}
Let $D\in\Div^+_{\mathrm{rad}}(\K)$, let $m=D(0)$, and enumerate its nonzero
occurrences as $\rho_j$. The product
\begin{equation}\label{ent:eq:product}
 P_D(Z)=Z^m\prod_j\left(1-\frac Z{\rho_j}\right)
\end{equation}
is a well-defined element of $\E$, independent of the enumeration. Its
divisor is exactly $D$, and the coefficient of $Z^m$ is $1$. The product is
finite when the nonzero divisor is finite. Its partial products converge to
it in the weighted Gauss valuation at every scale. \emph{No exponential
correction factors are required}: the theory is genus zero at every rank.
\end{theorem}
\begin{proof}
Only the infinite case needs a support argument. Put $b_j=\rho_j^{-1}$. Fix
$\gamma\in\Gamma$ and substitute $Z=t^\gamma X$. By
\cref{ent:lem:divisor-countable}, $v(b_j)+\gamma\to+\infty$. There are only
finitely many factors for which $v(b_j)+\gamma\le0$; set them aside. For the
rest, the supports of $b_jt^\gamma$ have a well-ordered positive union and
are locally finite, by \cref{ent:lem:cofinal-tails}.

Expand the tail product over finite subsets of indices. Every exponent
belongs to the monoid generated by that positive union. For a fixed exponent,
\cref{ent:lem:neumann} gives finitely many words in the support, and local
finiteness gives finitely many choices of factor labels carrying those
exponents. Thus only finitely many product terms contribute, and the
coefficient at that exponent is a polynomial in $X$. The tail product is an
element of $\A$ with initial polynomial $1$, hence a unit. Multiplying by the
finitely many set-aside factors and by $t^{m\gamma}X^m$ gives an element of
$\A$; the same estimate shows that the finite partial products approach it in
$m_\gamma$.

For each fixed ordinary degree, this construction sums the same finite-subset
coefficient family multiplied by $t^{k\gamma}$. Consequently the elements
obtained for different $\gamma$ are rescalings of one series in $\K[[Z]]$,
which belongs to $\E$ by \cref{ent:thm:criterion}(2). Strong regrouping also
proves independence of the enumeration.

On $B_\gamma$, the set-aside factors give precisely the roots in that ball,
while the remaining unit is nowhere zero on it. Local translation shows that
multiplicities are the displayed factor multiplicities. Since the balls cover
$\K$, the divisor is exactly $D$. Every factor of the nonzero-root product
has constant term $1$, so the coefficient of $Z^m$ is $1$.
\end{proof}

\begin{theorem}[Factorization by zeros]\label{ent:thm:factorization}
Let $F=\sum a_nZ^n\in\E\setminus\{0\}$ and let $m=\ord_{Z=0}F$. Its zero
divisor is radially finite. If $\rho_j$ are all its nonzero zeros, repeated
with their multiplicities, then
\begin{equation}\label{ent:eq:factorization}
 F(Z)=a_mZ^m\prod_j\left(1-\frac Z{\rho_j}\right).
\end{equation}
Thus the zero divisor determines $F$ up to a nonzero scalar.
\end{theorem}
\begin{proof}
Radial finiteness is \cref{ent:thm:roots}. Let $G=P_{\Div(F)}$ be the
canonical product from \cref{ent:thm:product}. Since $F$ and $G$ have the
same order $m$ at zero, the formal quotient
\[
 Q=\frac{F/Z^m}{G/Z^m}\in\K[[Z]]
\]
exists, because $G/Z^m$ has constant coefficient $1$.

We claim $Q\in\E$. Fix $\gamma$. Preparation of $F(t^\gamma X)$ and
$G(t^\gamma X)$ writes them as
\[
 F(t^\gamma X)=c_F W_F U_F,\qquad
 G(t^\gamma X)=c_G W_G U_G,
\]
with $c_F,c_G\in\K^\times$, units $U_F,U_G$ of $\A$, and monic $W_F,W_G$ all
of whose roots lie in $\OO$. Those roots are exactly the normalized roots of
$F$ and of $G$ in $B_\gamma$, with multiplicity, so $W_F=W_G$. Cancellation
in $\K[[X]]$, including the common factor $X^m$, gives
\[
 Q(t^\gamma X)=\frac{c_F}{c_G}\frac{U_F}{U_G}\in\A.
\]
This holds at every scale, so \cref{ent:thm:criterion}(2) proves the claim.

The identity $F=GQ$ is now an identity of entire functions. Orders at any
translated centre add in the formal power-series ring. Since $F$ and $G$ have
identical zero multiplicities, $Q$ has no zeros, so $Q$ is constant by
\cref{ent:cor:units}. Comparing the coefficient of $Z^m$ gives $Q=a_m$.
\end{proof}

The quotient step is important and is the reason the preparation theorem was
proved at \emph{every} scale. Equal zero sets alone do not justify dividing
two infinite series in the entire category, and no compactness or transfinite
convergence assertion replaces that argument. Note also the order of the
argument: the quotient is shown strongly evaluable at all scales
\emph{before} \cref{ent:thm:criterion} upgrades it to an entire,
valuation-convergent series. In particular the proof never assumes that a
zero-free restricted series is restricted-invertible --- by
\cref{ent:ex:ranktwo-unit} that is exactly what can fail in higher rank.

\begin{corollary}[Divisibility, GCDs, and least common multiples]
\label{ent:cor:gcd}
The map $F\mapsto\Div(F)$ induces an isomorphism of commutative monoids
\[
 (\E\setminus\{0\})/\K^\times
 \ \cong\ \Div^+_{\mathrm{rad}}(\K).
\]
For nonzero $F,G\in\E$, $G$ divides $F$ in $\E$ if and only if
$\Div(G)\le\Div(F)$ pointwise. Greatest common divisors and least common
multiples are given by the pointwise minimum and maximum of divisors. In
particular $\E$ is a GCD domain.
\end{corollary}
\begin{proof}
Surjectivity is \cref{ent:thm:product}; injectivity modulo constants is
\cref{ent:thm:factorization}. Divisors add under products by Taylor
expansion. If the divisor of $G$ is contained in that of $F$, their
difference is radially finite, so its canonical product supplies the
quotient, up to the scalar fixed by \cref{ent:thm:factorization}. Pointwise
minima and maxima of two radial divisors remain radially finite; their
canonical products therefore have the defining universal divisibility
properties of a GCD and an LCM.
\end{proof}

\begin{corollary}[Finite jets and evaluation ideals]\label{ent:cor:evaluation}
For $x\in\K$ and $r\ge1$, the kernel of the Taylor jet map
\[
 \E\longrightarrow\K[U]/(U^r),\qquad
 F\longmapsto F(x+U)\bmod U^r
\]
is $(Z-x)^r\E$, and the map is surjective. In particular evaluation at $x$
has kernel $(Z-x)\E$ and residue field $\K$.
\end{corollary}
\begin{proof}
Vanishing of the jet is exactly a zero of order at least $r$; apply
\cref{ent:cor:gcd}. Every jet is realized by a finite polynomial in $Z-x$.
\end{proof}

\begin{corollary}[Rigidity and value distribution]
\label{ent:cor:value-distribution}
Let $F\in\E$ be nonzero.
\begin{enumerate}[label=\textnormal{(\alph*)},leftmargin=2em]
\item $F$ has finitely many zeros if and only if $F$ is a polynomial; a
zero-free $F$ is a nonzero constant.
\item If $F$ is nonconstant it is surjective onto $\K$; if $F$ is
nonpolynomial, every fibre $F^{-1}(b)$ is infinite.
\item If $F$ is injective as a function on $\K$, then $F$ is affine of
degree one.
\end{enumerate}
\end{corollary}
\begin{proof}
(a) A polynomial has finitely many zeros, and conversely
\eqref{ent:eq:factorization} with a finite product is a polynomial.
Zero-freeness is \cref{ent:cor:units}. (b) Surjectivity is
\cref{ent:cor:units}; if $F$ is nonpolynomial then $F-b$ is nonpolynomial,
so by (a) it has infinitely many zeros. (c) An injective $F$ is nonconstant,
hence a polynomial by (b), say of degree $d\ge1$ with leading coefficient
$a_d$. For each $b$, injectivity makes the divisor of $F-b$ a single point
$r_b$ with multiplicity $d$, so $F(Z)-b=a_d(Z-r_b)^d$ by
\cref{ent:thm:factorization}. Comparing two values $b\ne b'$ gives
$a_d\bigl[(Z-r_b)^d-(Z-r_{b'})^d\bigr]=b'-b$, a nonzero constant. For
$d\ge2$ the left side has a nonzero coefficient of $Z^{d-1}$ unless
$r_b=r_{b'}$, and $r_b=r_{b'}$ makes it zero; either way we contradict
$b\ne b'$. Hence $d=1$.
\end{proof}

Part (b) and part (c) are familiar non-Archimedean value distribution, stated
here for completeness at arbitrary rank and not offered as independent
priority claims.

\begin{corollary}[Conjugation and real-coefficient functions]
\label{ent:cor:conjugation}
A nonzero entire series has all coefficients in $\R((t^\Gamma))$ if and only
if its zero divisor is invariant under coefficientwise conjugation and its
first nonzero coefficient lies in $\R((t^\Gamma))$.
\end{corollary}
\begin{proof}
Conjugating coefficients conjugates each value and each zero with its
multiplicity. Conversely, under the stated conditions the conjugate series
has the same divisor and the same first nonzero coefficient, so uniqueness in
\cref{ent:thm:factorization} makes the two series equal.
\end{proof}

\subsection{What ``factorization'' does not imply}
\begin{proposition}\label{ent:prop:nonnoetherian}
If $\E$ contains a nonpolynomial function, it is not Noetherian and is not
atomic. Its irreducible elements are precisely the scalar multiples of linear
polynomials.
\end{proposition}
\begin{proof}
By \cref{ent:thm:factorization}, a nonpolynomial entire function has
infinitely many zeros with multiplicity. Let $\rho_1,\rho_2,\ldots$ be an
escaping sequence of nonzero roots and put
$T_n(Z)=\prod_{j\ge n}(1-Z/\rho_j)$. Then $T_n=(1-Z/\rho_n)T_{n+1}$, so
\[
 (T_1)\subsetneq(T_2)\subsetneq(T_3)\subsetneq\cdots,
\]
strictness following either by evaluating multiplicities at $\rho_n$ or from
\cref{ent:cor:gcd}. Thus the ring is not Noetherian.

A nonunit has a zero. If its divisor has at least two occurrences, separating
one occurrence from the remaining divisor gives a product of two nonunits. A
singleton divisor gives a scalar multiple of one linear factor and is
irreducible. Consequently a function with an infinite divisor cannot be a
\emph{finite} product of irreducibles, which is the failure of atomicity.
\end{proof}

The canonical product theorem is therefore not a unique-factorization-domain
theorem. Its product is infinite in a precisely controlled Hahn sense; an
algebraic UFD requires finite irreducible factorizations.

\section{The higher-rank obstruction and the exact image of the jet map}
\label{ent:sec:obstruction}

\subsection{A sequence of genuinely new Archimedean scales}
For positive $\alpha,\beta\in\Gamma$, write
\[
 \beta\ggA\alpha
 \quad\Longleftrightarrow\quad
 \beta>k\alpha\text{ for every integer }k\ge1.
\]
This is Archimedean domination: a genuinely new Archimedean class, not merely
a large finite ratio.

\begin{lemma}\label{ent:lem:escaping-classes}
Suppose $\cf(\Gamma)=\aleph_0$ and $\Gamma$ has no order unit. There is a
positive cofinal sequence $(\gamma_n)_{n\ge1}$ with
$\gamma_{n+1}\ggA\gamma_n$ for every $n$.
\end{lemma}
\begin{proof}
Choose a positive cofinal sequence $(p_n)$. After selecting $\gamma_n$, the
set $\{k\gamma_n:k\ge1\}$ is not cofinal, since $\gamma_n$ is not an order
unit; hence it has a strict upper bound $b_n\in\Gamma$. Choose
$\gamma_{n+1}>\max\{b_n,p_{n+1},\gamma_n\}$ and start with
$\gamma_1\ge p_1>0$. The resulting sequence has the required domination and
is cofinal because it dominates the $p_n$.
\end{proof}

\subsection{Shells, normalized nodes, and the coarsened valuation rings}
\label{ent:subsec:shells}
The rest of this section identifies, exactly, the image of the full jet map
of a radially finite divisor. The objects it needs are finite algebras
attached to one valuation shell at a time, together with one coarsened
valuation ring per shell. They are introduced here. Throughout this
subsection $D$ is a radially finite effective divisor in the sense of
\cref{ent:def:radial}, and $D$ is infinite unless the contrary is said;
$\Gamma$ is arbitrary, and in particular the radii below are \emph{not}
assumed to satisfy the Archimedean domination of
\cref{ent:lem:escaping-classes}.

\begin{lemma}[Shell decomposition]\label{ent:lem:shells}
Let $D$ be an infinite radially finite divisor. There is a unique
decomposition
\[
 D=D_0+\sum_{n\ge1}D_n,\qquad
 D_n=\sum_{j=1}^{q_n}m_{n,j}\,[\rho_{n,j}],\qquad
 v(\rho_{n,j})=-\gamma_n,
\]
in which $D_0$ is a finite divisor supported in $\OO$, the $q_n$ and
$m_{n,j}$ are positive integers, the points $\rho_{n,1},\ldots,\rho_{n,q_n}$
are distinct, and
\[
 0<\gamma_1<\gamma_2<\cdots,\qquad(\gamma_n)\text{ cofinal in }\Gamma.
\]
Each shell total $M_n=\sum_{j\le q_n}m_{n,j}$ is finite. In particular an
infinite radially finite divisor forces $\cf(\Gamma)=\aleph_0$.
\end{lemma}
\begin{proof}
Radial finiteness at $\gamma=0$ makes $D_0$, the restriction of $D$ to
$\OO$, a finite divisor; it contains the occurrence at $0$, whose valuation
is $+\infty$. The remaining support has negative valuations, and their
negatives form the set of positive radii. Only finitely many radii lie below
any fixed bound, by radial finiteness, so every nonempty set of radii has a
least member: choose one member and minimize over the finitely many members
at most it. Listing the radii by repeatedly taking the next least member
exhausts them, because each has only finitely many predecessors. The list is
infinite and cofinal, since otherwise one ball would carry the whole infinite
divisor. Each shell lies in a ball and therefore has finite total
multiplicity. Uniqueness is the uniqueness of the valuation of a point.
Cofinality of $(\gamma_n)$ is the last assertion; it also follows from
\cref{ent:lem:divisor-countable}, which this lemma refines by grouping the
enumeration there into shells.
\end{proof}

Fix the normalizing monomials and the normalized shell data
\begin{equation}\label{ent:eq:shellpoly}
 \tau_n=t^{-\gamma_n},\qquad
 \xi_{n,j}=\rho_{n,j}/\tau_n,\qquad
 P_n(T)=\prod_{j=1}^{q_n}(T-\xi_{n,j})^{m_{n,j}} .
\end{equation}
Each $\xi_{n,j}$ is a unit of $\OO$, so $P_n$ is monic of degree $M_n$ with
coefficients in $\OO$ and $P_n(0)$ a unit. Distinct normalized nodes need not
have distinct residues in $\C$, and their differences may have arbitrarily
large positive valuation; nothing below assumes otherwise.

A prescribed jet at $\rho_{n,j}$ is a polynomial
$J_{n,j}(U)=\sum_{\ell<m_{n,j}}b_{n,j,\ell}U^\ell$ over $\K$, the intended
meaning being $F(\rho_{n,j}+U)\equiv J_{n,j}(U)\pmod{U^{m_{n,j}}}$. Finite
Hermite interpolation over the field $\K$ supplies a unique
\emph{normalized shell Hermite polynomial} $R_n$ with $\deg R_n<M_n$ and
\begin{equation}\label{ent:eq:normalizedjets}
 R_n(\xi_{n,j}+U)\equiv J_{n,j}(\tau_nU)\pmod{U^{m_{n,j}}}
 \qquad(1\le j\le q_n).
\end{equation}
It exists for arbitrary data, with no integrality assumption: the question
answered below is whether its coefficients lie in a particular subring. If
$F\in\E$ realizes the jets then $R_n=\remp\bigl(F(\tau_nT),P_n\bigr)$, since
$F(\tau_n(\xi_{n,j}+U))=F(\rho_{n,j}+\tau_nU)$.

\begin{definition}[Killing one Archimedean scale]\label{ent:def:coarse}
For $\gamma>0$ in $\Gamma$ put
\begin{align}
 H_\gamma&=\{\eta\in\Gamma:\abs{\eta}\le m\gamma\text{ for some }m\in\N\},
 \label{ent:eq:Hgamma}\\
 V_\gamma&=\{x\in\K:v(x)\ge-m\gamma\text{ for some }m\in\N\}.
 \label{ent:eq:Vgamma}
\end{align}
Here $\N$ contains $0$, and $v(0)=+\infty$ puts $0$ in $V_\gamma$. Write
$H_n=H_{\gamma_n}$ and $V_n=V_{\gamma_n}$ for the radii of a fixed divisor.
\end{definition}

$H_\gamma$ is the convex subgroup of $\Gamma$ generated by $\gamma$, and it is
divisible because dividing by a positive integer does not increase absolute
value. Both families increase with $n$, and
$\bigcup_nH_n=\Gamma$, $\bigcup_nV_n=\K$, because $(\gamma_n)$ is cofinal.
Membership in $V_n$ is a one-sided bound; arbitrarily small elements are
admitted. The unit condition is two-sided:
\begin{equation}\label{ent:eq:Vunits}
 x\in V_n^\times
 \quad\Longleftrightarrow\quad
 x\ne0\text{ and }\abs{v(x)}\le m\gamma_n\text{ for some }m\in\N .
\end{equation}

\begin{lemma}[Coarsened residues]\label{ent:lem:coarse-residue}
$V_\gamma$ is the valuation ring of the coarsening of $v$ by $H_\gamma$, with
value group $\Gamma/H_\gamma$. Its maximal ideal and residue field are
\[
 \mathfrak q_\gamma=\{x:v(x)>m\gamma\text{ for every }m\in\N\},
 \qquad
 \kappa_\gamma=V_\gamma/\mathfrak q_\gamma\cong\C\bigl((t^{H_\gamma})\bigr),
\]
the residue map being restriction of the Hahn support to $H_\gamma$.
\end{lemma}
\begin{proof}
An element of $V_\gamma$ has nonnegative class modulo $H_\gamma$, so all of
its support exponents lie in nonnegative cosets. Restricting the support to
the zero coset is additive, is surjective onto the full Hahn field on
$H_\gamma$, and is multiplicative because two nonnegative cosets sum to zero
only when both are zero. Its kernel consists of the series all of whose
exponents exceed $H_\gamma$, which is the displayed $\mathfrak q_\gamma$.
This also gives the valuation-ring and unit descriptions.
\end{proof}

This is the same computation as \cref{ent:lem:coarsened-ring}, performed
inside $\Gamma$ by a convex subgroup rather than across an ordered-group
extension; the two coarsenings are analogous, not the same object, and no
statement is transferred between them. Both produce a residue field of the
form $\C((t^H))$ rather than $\C$, and in both cases that is the point.

Define the two finite algebras of shell $n$:
\begin{equation}\label{ent:eq:shell-algebras}
 \Lambda_n=\K[T]/(P_n),
 \qquad
 \Lambda_n^\circ=V_n[T]/(P_n)\ \subseteq\ \Lambda_n .
\end{equation}
The inclusion is injective because division by a monic polynomial gives a
unique representative of degree below $M_n$ over either coefficient ring.
A class in $\Lambda_n$ therefore lies in $\Lambda_n^\circ$ exactly when every
coefficient of its degree-bounded representative lies in $V_n$. Over $\K$,
finite Chinese remainders identify $\Lambda_n$ with the product of the local
jet rings of the shell, the identification with the original jets including
the rescaling in \eqref{ent:eq:normalizedjets}. Over $V_n$ the analogous
decomposition may fail, because two normalized nodes can differ by a nonunit
of $V_n$. That failure is the source of the close-node obstruction of
\cref{ent:thm:close-pair}, not a technical inconvenience to be discarded.

For the finite part set $\Lambda_0=\prod_{x\in\supp D_0}\K[U]/(U^{D(x)})$,
omitting the factor when $D_0$ is empty; the same convention is used in every
formula below that names $\Lambda_0$.

\begin{definition}[The shellwise restricted product]\label{ent:def:restricted}
\begin{equation}\label{ent:eq:restricted}
 \RD=\Lambda_0\times
 \left\{(r_n)\in\prod_{n\ge1}\Lambda_n:
    r_n\in\Lambda_n^\circ\text{ for all sufficiently large }n\right\},
\end{equation}
equivalently the union over $N$ of
$\Lambda_0\times\prod_{n<N}\Lambda_n\times\prod_{n\ge N}\Lambda_n^\circ$.
\end{definition}

This is a ring under componentwise operations, and a $\K$-algebra because a
fixed scalar lies in $V_n$ for all large $n$. The phrase ``restricted
product'' here asserts \emph{no} locally compact topology, no adelic
structure and no Banach norm; it names the algebraic union just displayed and
nothing else.

\subsection{The exact image of the infinite jet map}
\label{ent:subsec:jet-image}
Every finite projection of the jet map of a radially finite divisor is
surjective, by ordinary finite Hermite interpolation over the field $\K$.
Neither a finite truncation nor a finite root calculation can therefore
detect a global failure of surjectivity; see \cref{ent:app:checks}. What the
finite projections cannot see is a condition on the \emph{tail} of the family
of shells, and the following theorem states it.

\begin{theorem}[Shellwise Hermite interpolation; the exact image]
\label{ent:thm:image}
Let $D$ be an infinite radially finite divisor, with the shell data of
\cref{ent:lem:shells}, and prescribe arbitrary finite jets $J_{n,j}$ at all of
its nodes together with an arbitrary element of $\Lambda_0$ at its finite
part. Let $R_n$ be the normalized shell Hermite polynomials
\eqref{ent:eq:normalizedjets}. There is an $F\in\E$ realizing every
prescribed jet if and only if
\begin{equation}\label{ent:eq:maincondition}
 R_n\in V_n[T]\qquad\text{for all sufficiently large }n,
\end{equation}
that is, if and only if for all sufficiently large $n$ there is an integer
$c_n\ge0$ with
\[
 v\bigl([T^j]R_n\bigr)\ge-c_n\gamma_n\qquad(0\le j<M_n).
\]
No condition is imposed on $D_0$ or on finitely many exceptional shells, and
no uniform bound on $q_n$, on $m_{n,j}$, on $M_n$ or on $c_n$ is required.
Consequently the image of the full jet map
\[
 \ev_D:\E\longrightarrow
 \Lambda_0\times\prod_{n\ge1}\Lambda_n
\]
is exactly $\RD$.
\end{theorem}

The condition does not depend on the choice of normalizing monomial.
Replacing $\tau_n$ by any $\tau_n'$ of the same valuation replaces $T$ by
$uT$ for a unit $u$ of $\OO$, which multiplies each coefficient by a power of
a valuation-zero unit and leaves membership in $V_n[T]$ unchanged.

\begin{theorem}[The entire quotient is the restricted product]
\label{ent:thm:quotient}
For a radially finite divisor $D$ with canonical product $P_D$ of
\cref{ent:thm:product}, the shellwise jet map induces an isomorphism of
$\K$-algebras
\[
 \E/(P_D)\ \cong\ \RD,
\]
so that
$0\to P_D\E\to\E\xrightarrow{\ \ev_D\ }\RD\to0$
is exact. Every finite projection of $\ev_D$ is surjective; the full map into
$\Lambda_0\times\prod_n\Lambda_n$ need not be.
\end{theorem}
\begin{proof}
The image is \cref{ent:thm:image}. The kernel is $P_D\E$ by
\cref{ent:cor:gcd,ent:cor:evaluation}: vanishing of all prescribed jets is
exactly divisibility of the divisor. That $\ev_D$ is an algebra homomorphism
is \cref{ent:prop:operations} together with the multiplicativity of the
truncated jet map. Surjectivity of a finite projection is finite Hermite
interpolation.
\end{proof}

The kernel statement invoked in that proof is the report's own
\cref{ent:cor:gcd}, which is downstream of support-controlled preparation
(\cref{ent:thm:prep}) and of algebraic closedness of $\K$
(\cref{ent:thm:roots}). A second and genuinely different proof of the same
statement, restricted to radially finite canonical products, is recorded next.
Both are kept. The second uses neither preparation nor algebraic closedness,
which is what makes \cref{ent:prop:real-form} available; the first proves far
more, namely that \emph{every} nonzero entire function is such a product
(\cref{ent:thm:factorization}), and that stronger statement does use
algebraic closedness.

\begin{proposition}[Principal vanishing ideal, second proof]
\label{ent:prop:kernel-direct}
Let $D$ be radially finite with canonical product $P_D$ of
\cref{ent:thm:product}, and let $F\in\E$. If every jet prescribed by $D$
vanishes on $F$, then $F\in P_D\E$. The converse is the jet homomorphism. Let
$P$ denote the subproduct of $P_D$ over the nodes of \emph{negative}
valuation, listed with multiplicity, and let $0<\beta_1\le\beta_2\le\cdots$ be
the valuations $v(1/x)$ of the reciprocals of those nodes. The proof uses only
the coefficient bound
\begin{equation}\label{ent:eq:productcoef}
 v\bigl([Z^k]P\bigr)\ \ge\ \beta_1+\cdots+\beta_k ,
\end{equation}
which also re-proves the entireness half of \cref{ent:thm:product} for $P$
without preparation, since for $k\ge2$ the right-hand side divided by $k$ is
at least $\beta_{\lfloor k/2\rfloor}/2$, and that tends cofinally to $+\infty$.
\end{proposition}
\begin{proof}
Divide first by $Z^{D(0)}$ and by the finitely many factors at nonzero nodes
of nonnegative valuation, using \cref{ent:lem:linear-division}; the quotient
is entire and retains the required zeros at every remaining node. We may
therefore replace $P_D$ by $P$, which has $P(0)=1$ and only negative-valuation
nodes; note $\beta_i\to+\infty$ cofinally by
\cref{ent:lem:divisor-countable}.

Expanding $P$ over finite subsets of the reciprocals $1/x$ gives
\eqref{ent:eq:productcoef}: every degree-$k$ term is a product of $k$ of the
listed reciprocals with distinct labels, so its valuation is at least the sum
of the $k$ smallest $\beta_i$, and the valuation of a Hahn sum is at least the
minimum of its terms. Among $\beta_1,\ldots,\beta_k$ more than half are at
least $\beta_{\lfloor k/2\rfloor}$ and the rest are positive, so the average
is at least $\beta_{\lfloor k/2\rfloor}/2$;
\cref{ent:thm:criterion} then gives entireness of $P$ directly.

Since $P(0)=1$ there is a unique formal quotient $G=F/P\in\K[[Z]]$; we show
$G\in\E$. Let $P_{\le N}$ be the finite product over the first $N$ shells and
$F_N=F/P_{\le N}$, entire by finitely many applications of
\cref{ent:lem:linear-division}. Every coefficient of the formal inverse
$P_{\le N}^{-1}$ is integral, because each factor $(1-Z/x)^{-1}$ expands with
$v(1/x)>0$; so if $A\ge0$ bounds $-v$ on the coefficients of $F$, then every
coefficient of every $F_N$ has valuation at least $-A$, \emph{independently of
$N$}. Write $T_N$ for the formal inverse of the tail $P_{>N}$. By
\eqref{ent:eq:productcoef} applied to that tail, $v([Z^k]P_{>N})\ge
k\gamma_{N+1}$, and the reciprocal recursion gives the same bound for
$[Z^k]T_N$. In $\K[[Z]]$ we have $G=F_NT_N$.

Fix a target $\gamma\in\Gamma$ and choose $\theta>\max\{\gamma,0\}$; choose
$N$ with $\gamma_{N+1}>4\theta+2A$. As $F_N$ is entire, for all large $j$ its
coefficient has valuation at least $4j\theta$. Split the convolution for
$[Z^k]G$ at $j=k/2$: a summand with $j\ge k/2$ has valuation at least
$4j\theta\ge2k\theta$, and a summand with $j<k/2$ has valuation at least
\[
 -A+(k-j)\gamma_{N+1}>-A+k\gamma_{N+1}/2
 \ge k(\gamma_{N+1}/2-A)>2k\theta .
\]
Hence $v([Z^k]G)>k\gamma$ for all large $k$; cancellation only improves the
bound. As $\gamma$ was arbitrary, \cref{ent:thm:criterion} gives $G\in\E$.
\end{proof}

The theorem identifies the missing object rather than a necessary inequality
inside it. It also explains why solving each finite Chinese remainder problem
and passing to an unspecified limit is not enough: the relevant restricted
product is in general a proper subring of the full product. The element
$g\bmod f$ of \cref{ent:thm:no-bezout} maps to an everywhere nonzero
sequence, hence to a unit of the full product ring, but it is not a unit of
$\E/(f)$: an inverse there would give representatives satisfying a B\'ezout
identity. Thus the elementary implication ``invertible at each node implies
invertible in the global quotient'' is exactly where unrestricted infinite
Chinese remainder reasoning fails, and \cref{ent:thm:bezout-criterion} below
says precisely how much more is needed.

\subsection{Proof of the exact image theorem}
\label{ent:subsec:image-proof}
Two finite-algebra facts and one assembly lemma carry the proof.

\begin{lemma}[Integral polynomial operations]\label{ent:lem:integral}
Let $P\in\OO[T]$ be monic of degree $M$.
\begin{enumerate}[label=\textnormal{(\alph*)},leftmargin=2em]
\item If every coefficient of $G\in\K[T]$ has valuation at least $B$, then so
does every coefficient of $\remp(G,P)$. In particular $\remp$ maps $\OO[T]$
into $\OO[T]$.
\item Suppose in addition $P=\prod_j(T-\xi_j)^{m_j}$ with $v(\xi_j)=0$, and
let $Q(T)=\prod_h(T-\eta_h)^{m_h'}$ be a finite product with $v(\eta_h)>0$,
where $\eta_h=0$ is permitted. Then $T$ and $Q$ are units of $\OO[T]/(P)$,
and for every integer $d\ge0$ the inverse of $T^dQ$ has a representative in
$\OO[T]$ of degree below $M$.
\end{enumerate}
\end{lemma}
\begin{proof}
(a) Monic division subtracts only integral multiples of $P$, so the bound $B$
survives every step.

(b) Consider multiplication on the free $\OO$-module $\OO[T]/(P)$ with basis
$1,T,\ldots,T^{M-1}$; its matrices have entries in $\OO$. After extending
scalars to $\K$, the determinant of multiplication by $L\in\K[T]$ is
\begin{equation}\label{ent:eq:detnorm}
 \det(m_L)=\prod_jL(\xi_j)^{m_j}=\Res(P,L),
\end{equation}
because in $\K[U]/(U^{m_j})$ multiplication by $L(\xi_j+U)$ is triangular
with constant diagonal $L(\xi_j)$, and finite Chinese remainders over $\K$
decompose $\Lambda_n$ accordingly. For $L=T$ each $\xi_j$ is a unit of $\OO$.
For $L=Q$ each $\xi_j-\eta_h$ has valuation zero, since $v(\xi_j)=0<v(\eta_h)$.
The determinants are therefore units of $\OO$, and the adjugate supplies
inverses over $\OO$. Multiplying the two inverses gives the assertion for
$T^dQ$.
\end{proof}

No lower bound on $v(\xi_j-\xi_k)$ for two nodes of the \emph{current} shell
was used; they may be separated only at an arbitrarily higher scale. Only the
already-treated nodes, of strictly smaller radius, enter $Q$, and they are
uniformly separated from the current shell after normalization. This is the
one place where the same finite algebra that creates the close-node
obstruction also makes the construction possible.

\begin{lemma}[Assembly]\label{ent:lem:assembly}
Let $C_n\in\K[Z]$ be polynomials with $Z^n\mid C_n$, and suppose that for
every $\gamma\in\Gamma$ the weighted Gauss valuations $m_\gamma(C_n)$ tend
cofinally to $+\infty$. Then for any polynomial $S$ the series
$F=S+\sum_nC_n$ is a well-defined element of $\E$, every coefficient of $F$
being a finite sum, and evaluation and each fixed finite jet may be
interchanged with the sum over $n$.
\end{lemma}
\begin{proof}
Only $C_n$ with $n\le k$ contributes to the coefficient of $Z^k$, so each
coefficient is a finite sum. Fix $\gamma$ and $B$. All but finitely many
$C_n$ have $m_\gamma(C_n)>B$, and the finitely many exceptions have finite
degree; beyond those degrees every weighted coefficient of $F$ exceeds $B$.
This is \cref{ent:thm:criterion}(4), so $F\in\E$. For evaluation at $b$ and
for a fixed jet order, choose $\gamma\le v(b)$; estimate
\eqref{ent:eq:jet-bound} makes the family of jet coefficients cofinally
summable, and the full family of evaluated monomials is strongly summable
because below a fixed exponent only finitely many polynomials contribute,
each with finitely many terms. \Cref{ent:lem:cofinal-tails} and regrouping
give the interchange.
\end{proof}

\begin{proof}[Proof of necessity in \cref{ent:thm:image}]
Let $F=\sum_ka_kZ^k\in\E$ be nonzero and fix $n$. Put
$E_{n,k}(T)=\remp(T^k,P_n)$, whose coefficients lie in $\OO$ by
\cref{ent:lem:integral}(a). By \cref{ent:thm:criterion} the valuations
$v(a_k\tau_n^k)=v(a_k)-k\gamma_n$ tend cofinally to $+\infty$, so
\begin{equation}\label{ent:eq:analyticremainder}
 \remp_n(F)=\sum_{k\ge0}a_k\tau_n^kE_{n,k}(T)
\end{equation}
is a polynomial of degree below $M_n$ whose coefficients are strong Hahn sums
(\cref{ent:lem:cofinal-tails}). At each normalized node its finite jets agree
with those of $F(\tau_nT)$: the finite remainder identity holds term by term,
and jet evaluation at the integral nodes preserves the cofinal estimates by
\eqref{ent:eq:jet-bound}. Hence $\remp_n(F)$ is the normalized shell Hermite
polynomial $R_n$ of the jets of $F$.

By strong summability at $Z=1$ there is a fixed $A\ge0$ with $v(a_k)\ge-A$
for every $k$. Every coefficient of \eqref{ent:eq:analyticremainder} has
valuation at least
\[
 m_{-\gamma_n}(F)=\min_k\{v(a_k)-k\gamma_n\},
\]
a minimum attained at a finite index $k_n$, which may depend on $n$. For all
$n$ with $\gamma_n>A$ this gives
\[
 m_{-\gamma_n}(F)\ \ge\ -A-k_n\gamma_n\ \ge\ -(k_n+1)\gamma_n ,
\]
so every coefficient of $R_n$ lies in $V_n$. The zero series is immediate,
and $D_0$ imposes nothing.
\end{proof}

The attained minimum is load-bearing. An infimum of infinitely many estimates
lying strictly above a cut need not itself lie above that cut in a
non-Archimedean ordered group; here there is an actual minimizing index, not
an appeal to an unproved compactness principle. This is the same point as in
\cref{ent:lem:growth-barrier} and is recorded again in
\cref{ent:app:audit}.

For the converse, suppose a polynomial $S$ already fits every node of radius
smaller than $\gamma_n$, together with $D_0$. Put
\begin{align}
 Q_n(T)&=\prod_{a\text{ already treated}}\bigl(T-a/\tau_n\bigr)^{D(a)},
 \label{ent:eq:Qn}\\
 \Theta_n(T)&=R_n(T)-\remp\bigl(S(\tau_nT),P_n\bigr).
 \label{ent:eq:theta}
\end{align}
Every already-treated node has strictly smaller radius, so each $a/\tau_n$
has positive valuation or is zero, and $Q_n$ satisfies the hypothesis of
\cref{ent:lem:integral}(b). For an integer $d\ge0$ set
\begin{align}
 \Psi_{n,d}(T)&=\remp\Bigl(\Theta_n(T)\bigl(T^dQ_n(T)\bigr)^{-1},P_n\Bigr),
 \label{ent:eq:psi}\\
 C_{n,d}(Z)&=(Z/\tau_n)^d\,Q_n(Z/\tau_n)\,\Psi_{n,d}(Z/\tau_n),
 \label{ent:eq:correction}
\end{align}
the inverse in \eqref{ent:eq:psi} being taken in $\Lambda_n$ with the
integral representative supplied by \cref{ent:lem:integral}(b). Then every
earlier prescribed jet of $C_{n,d}$ vanishes, the current-shell remainder of
$C_{n,d}$ is exactly $\Theta_n$, so $S+C_{n,d}$ fits shell $n$ as well, and
$Z^d$ divides $C_{n,d}$. The parameter $d$ changes none of these identities;
it is free, and is used only for size.

\begin{lemma}[Uniform damping estimate]\label{ent:lem:damping-shell}
Suppose every coefficient of $\Theta_n$ has valuation at least $B$. Then the
same holds for $\Psi_{n,d}$, independently of $d$, and
\begin{equation}\label{ent:eq:damping-shell}
 m_{-\gamma_n/2}(C_{n,d})\ \ge\ B+d\gamma_n/2 .
\end{equation}
If $\Theta_n\in V_n[T]$, so that $B=-c\gamma_n$ may be taken with an integer
$c\ge0$, then every $d\ge\max\{n,2c+2\}$ gives
\begin{equation}\label{ent:eq:damping-target}
 m_{-\gamma_n/2}(C_{n,d})\ \ge\ \gamma_n .
\end{equation}
\end{lemma}
\begin{proof}
Multiplication by an integral inverse and integral remainder reduction
preserve the lower bound $B$, by \cref{ent:lem:integral}. Under $T=Z/\tau_n
=t^{\gamma_n}Z$ the monomial $T$ has weighted value
$\gamma_n-\gamma_n/2=\gamma_n/2>0$ at $\gamma=-\gamma_n/2$; the substituted
$Q_n$ has weighted value at least $0$ because its coefficients are integral;
and the substituted $\Psi_{n,d}$ has weighted value at least $B$. Adding
these, by the product estimate of \cref{ent:subsec:notation}, gives
\eqref{ent:eq:damping-shell}. For the last assertion,
$-c\gamma_n+d\gamma_n/2\ge\gamma_n$ holds as soon as $d\ge2c+2$.
\end{proof}

The estimate never uses $d\gamma_n\to+\infty$ in $\Gamma$ for fixed $n$; that
assertion is false when $\gamma_n$ is not an order unit, and it is exactly
the false step that \cref{ent:lem:escaping-classes} forbids. The finite
damping needed at stage $n$ lives inside the convex subgroup $H_n$ generated
by $\gamma_n$. It is the stages, not the powers at one stage, that traverse
the cofinal hierarchy.

\begin{proof}[Proof of sufficiency in \cref{ent:thm:image}]
Suppose $R_n\in V_n[T]$ for every $n>N$. Fit $D_0$ and the first $N$ shells by
a finite Hermite polynomial $S_N$ over $\K$. Since $S_N$ has finitely many
coefficients and $\bigcup_nV_n=\K$, there is an $L>N$ with all coefficients of
$S_N$ in $V_L$.

Treat the finitely many shells $N<n\le L$ by
\eqref{ent:eq:Qn}--\eqref{ent:eq:correction} with $d=0$. Throughout these
steps the coefficients stay in $V_L$: for such $n$ we have
$R_n\in V_n[T]\subseteq V_L[T]$, and both $\tau_n$ and $\tau_n^{-1}$ are units
of $V_L$ because $\gamma_n\in H_L$, so normalizing, taking integral
remainders and multiplying by the integral inverse of
\cref{ent:lem:integral}(b) all preserve $V_L$-coefficients. This produces a
polynomial $S_L$ fitting $D_0$ and the first $L$ shells, with coefficients in
$V_L$. Choose $A\ge0$ in $H_L$ --- a sufficiently large finite multiple of
$\gamma_L$ will do --- so that every coefficient of $S_L$ has valuation at
least $-A$. This one bound is maintained for the rest of the construction.

Let $n>L$ and suppose $S_{n-1}$ fits everything of radius below $\gamma_n$
and has all coefficient valuations at least $-A$. Since $A\in H_L\subseteq
H_n$, all coefficients of $S_{n-1}$ lie in $V_n$, and $\tau_n$ is a unit of
$V_n$; hence $\Theta_n$ of \eqref{ent:eq:theta} lies in $V_n[T]$. Choose
$d_n\ge n$ by \cref{ent:lem:damping-shell} so that $C_n=C_{n,d_n}$ satisfies
$m_{-\gamma_n/2}(C_n)\ge\gamma_n$, and set $S_n=S_{n-1}+C_n$. Then every
coefficient $c_{n,k}$ of $C_n$ obeys
\begin{equation}\label{ent:eq:strongcoefbound}
 v(c_{n,k})\ \ge\ \gamma_n+k\gamma_n/2\ \ge\ \gamma_n\ >\ 0,
\end{equation}
so the bound $-A$ is preserved and the induction continues. A zero residual
is handled by $C_n=0$.

Put $F=S_L+\sum_{n>L}C_n$. Since $Z^n$ divides $C_n$, each coefficient is a
finite sum. For fixed $\gamma$ we eventually have $\gamma\ge-\gamma_n/2$, and
$m_\gamma$ is nondecreasing, so
\[
 m_\gamma(C_n)\ \ge\ m_{-\gamma_n/2}(C_n)\ \ge\ \gamma_n
 \longrightarrow+\infty\quad\text{cofinally}.
\]
\Cref{ent:lem:assembly} makes $F$ entire and lets every fixed finite jet pass
through the sum. At any prescribed node the required jet is attained at a
finite stage, and every later correction has zero jet there. Hence $F$
realizes all the prescribed data.
\end{proof}

\begin{remark}[Why the finite initialization is not a formality]
An exceptional early jet may carry a coefficient at a scale larger than many
later radii. ``Absorb finitely many exceptions into a polynomial'' does not by
itself justify the later damping estimates, which need one permanent lower
bound. The intermediate finite run over $V_L$, followed by the single fixed
bound $-A\in H_L$, is what prevents such a coefficient from propagating into
an uncontrolled obstruction. This is the only initialization issue, and no
uniform bound on the later jet orders is required.
\end{remark}

\begin{corollary}[Exact one-node interpolation]\label{ent:cor:single-node}
Suppose each shell carries one node $\rho_n$, of multiplicity $m_n$, and
prescribe jets $\sum_{\ell<m_n}b_{n,\ell}(Z-\rho_n)^\ell$. These are realized
by an entire function if and only if, for all sufficiently large $n$,
\[
 b_{n,\ell}\in V_n\qquad\text{for every }\ell<m_n .
\]
For values alone ($m_n=1$ for all $n$) this reads: the prescription
$F(\rho_n)=b_n$ is realizable if and only if $b_n\in V_n$ eventually, that is,
if and only if for all large $n$ there is an integer $c_n\ge0$ with
$v(b_n)\ge-c_n\gamma_n$.
\end{corollary}
\begin{proof}
The normalized Hermite polynomial is
$R_n(T)=\sum_{\ell<m_n}b_{n,\ell}\tau_n^\ell(T-\xi_n)^\ell$. Both $\tau_n$ and
$\xi_n$ are units of $V_n$, and the change between powers of $T$ and powers of
$T-\xi_n$ is given by mutually inverse triangular matrices over $V_n$. Hence
$R_n\in V_n[T]$ if and only if every $b_{n,\ell}$ lies in $V_n$. Apply
\cref{ent:thm:image}.
\end{proof}

\begin{corollary}[Coarsely separated shells]\label{ent:cor:coarse-separated}
Suppose that on every sufficiently large shell the differences
$\xi_{n,j}-\xi_{n,k}$, $j\ne k$, are units of $V_n$. Then arbitrary
prescribed jets are realizable exactly when all of their coefficients, in the
original coordinate, lie in $V_n$ eventually.
\end{corollary}
\begin{proof}
Under that hypothesis the ideals $(T-\xi_{n,j})^{m_{n,j}}$ are already
pairwise comaximal in $V_n[T]$: two linear ideals are comaximal because the
difference of their generators is a unit, and their powers are comaximal by
expanding a sufficiently high power of an identity $u+w=1$. The finite
Chinese remainder isomorphism therefore holds over $V_n$, and rescaling jets
by powers of the $V_n$-unit $\tau_n$ does not change membership. Apply
\cref{ent:thm:image}.
\end{proof}

The hypothesis is \emph{coarsened} separation, not separation in $v$.
Normalized differences may have arbitrarily large positive fine valuation,
provided that valuation still lies in $H_n$. Without it the conclusion is
false, by \cref{ent:thm:close-pair}.

\begin{corollary}[Order units and the disappearance of the restriction]
\label{ent:cor:order-unit-image}
If $\Gamma$ has an order unit $h$, then $H_\gamma=\Gamma$ and $V_\gamma=\K$
for every $\gamma\ge h$, so $V_n=\K$ and $\Lambda_n^\circ=\Lambda_n$ for all
large $n$. Condition \eqref{ent:eq:maincondition} is then vacuous, $\RD$ is
the full product, and every radially finite Hermite prescription is realized.
\end{corollary}
\begin{proof}
An order unit makes every element of $\Gamma$ bounded in absolute value by a
finite multiple of $h$, hence of any $\gamma\ge h$; a cofinal $(\gamma_n)$
eventually exceeds $h$. Apply \cref{ent:thm:image,ent:thm:quotient}.
\end{proof}

This recovers \cref{ent:thm:interpolation} and is discussed again in
\cref{ent:sec:positive}, where the independent Gauss-norm proof is retained.
At the other extreme, if $\cf(\Gamma)>\aleph_0$ then \cref{ent:lem:shells}
leaves no infinite radially finite divisor at all and every interpolation
problem here is finite. The new content lies in the intervening groups.

\begin{corollary}[Without an order unit the restriction never disappears]
\label{ent:sm:cor:no-order-unit}
Conversely, if $\Gamma$ has no order unit, then every $H_n$ is a proper
subgroup of $\Gamma$. In particular, for a divisor with one simple node
$\rho_n$ on each shell and empty finite part, the image of the value map
$F\mapsto(F(\rho_n))_n$ is a proper subring of $\K^{\N_{>0}}$.
\end{corollary}
\begin{proof}
If $H_n=\Gamma$, every element of $\Gamma$ is bounded in absolute value by a
finite multiple of $\gamma_n$, so $\gamma_n$ is an order unit. If every $H_n$
is proper, choose $\eta_n>0$ outside $H_n$. It exceeds every finite multiple
of $\gamma_n$, so $t^{-\eta_n}\notin V_n$ for every $n$, and
\cref{ent:cor:single-node} shows that the values $F(\rho_n)=t^{-\eta_n}$ are
realized by no entire $F$.
\end{proof}

This converse is the fourth source's; \cref{ent:cor:no-interpolation}(a) is
its instance $\eta_n=\gamma_{n+1}$ for the escaping sequence of
\cref{ent:lem:escaping-classes}.

\subsection{A second proof of sufficiency: cardinal corrections}
\label{ent:sm:subsec:cardinal}
The fourth source proved \cref{ent:thm:image} independently. Its necessity
argument is the one above, an attained Gauss minimum. Its sufficiency
construction is different, and is recorded here as a second route. The
first route corrects shell by shell with \emph{polynomials} that vanish at
every earlier node, and needs the finite initialization over $V_L$. The
second builds, for each shell separately, an \emph{entire} correction that
vanishes to the prescribed orders on every other shell: the canonical product
with that shell removed, times a finite polynomial. One integral estimate for
the inverse of that cofactor on its own shell replaces the sequential
bookkeeping, and the construction comes with an explicit truncation bound.
Neither route is offered as a correction of the other, and like the first,
the second uses neither preparation nor algebraic closedness.

Throughout this subsection $D$ is an infinite radially finite divisor with
the shell data of \cref{ent:lem:shells}, and $D_0=0$ until the last step. Put
\begin{equation}\label{ent:sm:eq:cofactor}
 \Pi_n(Z)=\frac{P_n(Z/\tau_n)}{P_n(0)}
 =\prod_{j=1}^{q_n}\Bigl(1-\frac Z{\rho_{n,j}}\Bigr)^{m_{n,j}},
 \qquad
 P_D=\prod_{n\ge1}\Pi_n,
 \qquad
 \widehat P_n=\prod_{h\ne n}\Pi_h .
\end{equation}
The cofactor $\widehat P_n$ is the canonical product of the radially finite
divisor $D-D_n$, entire by \cref{ent:thm:product}; it is obtained by
removing finitely many factors, not by an entire division theorem. For a
class $r\in\Lambda_n$ write $\nu(r)$ for the least valuation of the
coefficients of its representative of degree below $M_n$, with
$\nu(0)=+\infty$. Monic division never lowers that minimum
(\cref{ent:lem:integral}(a)), so $\nu(rs)\ge\nu(r)+\nu(s)$. For $F\in\E$ let
$\remp_n(F)\in\Lambda_n$ be the shell class \eqref{ent:eq:analyticremainder}.
It is the $n$th component of $\ev_D$, hence a $\K$-algebra homomorphism, and
the necessity proof above shows
\begin{equation}\label{ent:sm:eq:class-bound}
 \nu\bigl(\remp_n(F)\bigr)\ \ge\ m_{-\gamma_n}(F).
\end{equation}
Finally put
\begin{equation}\label{ent:sm:eq:gain}
 \varpi_n=\remp_n(\widehat P_n)\in\Lambda_n,
 \qquad
 \varkappa_n=\sum_{h<n}M_h(\gamma_n-\gamma_h)\ \ge\ 0 .
\end{equation}

\begin{lemma}[Summing entire corrections]\label{ent:sm:lem:entire-sum}
Let $F_n\in\E$ and suppose that $m_\gamma(F_n)\to+\infty$ cofinally for every
$\gamma\in\Gamma$. Then the coefficientwise sum $F=\sum_nF_n$ is defined and
belongs to $\E$, $m_\gamma\bigl(F-\sum_{n\le N}F_n\bigr)\to+\infty$ for every
$\gamma$, and $\remp_m(F)=\sum_n\remp_m(F_n)$ for every shell $m$, the right
side being a coefficientwise strong sum.
\end{lemma}
\begin{proof}
At $\gamma=0$ the hypothesis makes the valuations in each coefficient family
$([Z^k]F_n)_n$ tend cofinally to $+\infty$, so \cref{ent:lem:cofinal-tails}
defines its sum. Fix $\gamma$ and $B$. All sufficiently late $F_n$ have every
weighted coefficient above $B$, hence so does the sum of those late terms,
and the finitely many earlier $F_n$ are entire, so they have this property in
all sufficiently large degrees. Thus the weighted coefficients of $F$ tend
cofinally to $+\infty$, which is \cref{ent:thm:criterion}(4); the same
argument applied to the late terms alone gives the tail assertion. For the
last assertion, \eqref{ent:sm:eq:class-bound} makes the coefficient families
of $(\remp_m(F_n))_n$ cofinally small, and
$\nu\bigl(\remp_m(F)-\sum_{n\le N}\remp_m(F_n)\bigr)\ge
m_{-\gamma_m}\bigl(F-\sum_{n\le N}F_n\bigr)\to+\infty$; a class whose $\nu$
exceeds every element of $\Gamma$ is zero.
\end{proof}

Compare \cref{ent:lem:assembly}: there the corrections are polynomials
divisible by $Z^n$; here they are entire functions, and only their weighted
Gauss valuations are controlled.

\begin{lemma}[Integral inverse of the cofactor]\label{ent:sm:lem:inverse-shell}
The class $\varpi_n$ is a unit of $\Lambda_n$, and
\begin{equation}\label{ent:sm:eq:inverse-bound}
 \nu\bigl(\varpi_n^{-1}\bigr)\ \ge\ \varkappa_n .
\end{equation}
\end{lemma}
\begin{proof}
Write $A_n^{\OO}=\OO[T]/(P_n)$, a finite free $\OO$-algebra inside
$\Lambda_n$ in which $T$ is a unit, by \cref{ent:lem:integral}(b). Every
element of $1+\mm A_n^{\OO}$ is a unit of $A_n^{\OO}$: the determinant of
its multiplication matrix in the basis $1,T,\ldots,T^{M_n-1}$ is $1$ modulo
$\mm$, hence a unit of $\OO$, and the adjugate gives an integral inverse. No
geometric series is used.

Split $\widehat P_n=P_{<n}P_{>n}$ into the finite product over earlier shells
and the canonical product of the later ones. For a node $\rho=\rho_{h,j}$ of
an earlier shell, the class of its factor is
\[
 1-\frac{\tau_n}{\rho}T
 =-\frac{\tau_n}{\rho}\,T\,\Bigl(1-\frac{\rho}{\tau_n}T^{-1}\Bigr),
 \qquad v(\rho/\tau_n)=\gamma_n-\gamma_h>0,
\]
so the last factor lies in $1+\mm A_n^{\OO}$. For a node $\rho$ of a later
shell $h$ the class $1-(\tau_n/\rho)T$ itself lies in $1+\mm A_n^{\OO}$,
since $v(\tau_n/\rho)=\gamma_h-\gamma_n>0$; so do the classes $u_N$ of the
finite partial products of $P_{>n}$. Those partial products converge to
$P_{>n}$ in every $m_\gamma$ (\cref{ent:thm:product}), so by
\eqref{ent:sm:eq:class-bound} some $u_N$ satisfies
$\nu\bigl(\remp_n(P_{>n})-u_N\bigr)>0$, and then $\remp_n(P_{>n})$ lies in
$1+\mm A_n^{\OO}$ as well. Hence, with
\[
 c_n=\prod_{h<n}\prod_{j=1}^{q_h}\bigl(-\tau_n/\rho_{h,j}\bigr)^{m_{h,j}},
 \qquad v(c_n)=-\varkappa_n,
\]
the class $\varpi_n/c_n$ is a power of the integral unit $T$ times an element
of $1+\mm A_n^{\OO}$. Its inverse is integral, and multiplying by
$c_n^{-1}$ gives \eqref{ent:sm:eq:inverse-bound}.
\end{proof}

The argument never inverts a difference of two nodes of the \emph{same}
shell. Those may be arbitrarily close, and whatever they cost is already
measured by $R_n$. Only the strict radial separation of distinct shells
enters, exactly as in \cref{ent:lem:integral}(b).

\begin{lemma}[Cardinal corrections]\label{ent:sm:lem:cardinal}
Let $r\in\Lambda_n$ with representative $R$ of degree below $M_n$, let
$d\ge0$ be an integer, and put
\begin{align}
 \Xi_{n,d}(T)&=\remp\bigl(T^{-d}\varpi_n^{-1}R,\ P_n\bigr),
 \label{ent:sm:eq:Xi}\\
 \widehat C_{n,d}(Z)&=\widehat P_n(Z)\,(Z/\tau_n)^d\,
   \Xi_{n,d}(Z/\tau_n),\label{ent:sm:eq:cardinal}
\end{align}
the inverses being taken in $\Lambda_n$. Then $\widehat C_{n,d}\in\E$,
$\remp_n(\widehat C_{n,d})=r$, every prescribed jet of $D$ at a node outside
shell $n$ vanishes on $\widehat C_{n,d}$, and for every $\gamma\ge-\gamma_n/2$
\begin{equation}\label{ent:sm:eq:suppression}
 m_\gamma\bigl(\widehat C_{n,d}\bigr)\ \ge\
 m_\gamma(P_D)+\nu(r)+\varkappa_n+d(\gamma_n+\gamma).
\end{equation}
\end{lemma}
\begin{proof}
It is an entire function times a polynomial. Its class on shell $n$ is
$\varpi_n\cdot T^d\cdot\Xi_{n,d}=r$, and on any other shell $h$ it vanishes
to the prescribed orders because $\widehat P_n$ contains the factor $\Pi_h$.
By \cref{ent:sm:lem:inverse-shell} and integral multiplication,
$\nu(\Xi_{n,d})\ge\nu(r)+\varkappa_n$. Under $T=Z/\tau_n=t^{\gamma_n}Z$ the
coefficient of $Z^k$ acquires the factor $t^{k\gamma_n}$, so its weighted
value at $\gamma$ exceeds the valuation of the original coefficient by
$k(\gamma_n+\gamma)\ge0$. Hence
$m_\gamma(\Xi_{n,d}(Z/\tau_n))\ge\nu(\Xi_{n,d})$, and
$m_\gamma\bigl((Z/\tau_n)^d\bigr)=d(\gamma_n+\gamma)$. Next,
$m_\gamma(\Pi_n)=0$ for $\gamma>-\gamma_n$: the constant term is $1$, and the
coefficient of $Z^k$, $k\ge1$, is a sum of products of $k$ reciprocals of
nodes of valuation $\gamma_n$, so its weighted value is at least
$k(\gamma_n+\gamma)>0$. The equality case of Gauss multiplicativity
(\cref{ent:subsec:notation}) applied to $P_D=\Pi_n\widehat P_n$ gives
$m_\gamma(\widehat P_n)=m_\gamma(P_D)$. The product estimate now gives
\eqref{ent:sm:eq:suppression}.
\end{proof}

\begin{proof}[Second proof of sufficiency in \cref{ent:thm:image}, for $D_0=0$]
Suppose $R_n\in V_n[T]$ for all $n>N_0$. For such $n$ with $R_n\ne0$ choose an
integer $c_n\ge0$ with $\nu(R_n)\ge-c_n\gamma_n$ and put
\begin{equation}\label{ent:sm:eq:d-choice}
 d_n=2c_n+2;
\end{equation}
for $n\le N_0$ put $d_n=0$. Let $\widehat C_n=\widehat C_{n,d_n}$ be the
cardinal correction of the class of $R_n$, and $\widehat C_n=0$ if $R_n=0$.
For $n>N_0$ and $\gamma\ge-\gamma_n/2$, \cref{ent:sm:lem:cardinal} and
$\varkappa_n\ge0$ give
\[
 m_\gamma(\widehat C_n)\ \ge\ m_\gamma(P_D)-c_n\gamma_n
   +(2c_n+2)\frac{\gamma_n}{2}
 \ =\ m_\gamma(P_D)+\gamma_n .
\]
For fixed $\gamma$ the condition $\gamma\ge-\gamma_n/2$ holds for all large
$n$, and $m_\gamma(P_D)+\gamma_n\to+\infty$ cofinally. By
\cref{ent:sm:lem:entire-sum}, $F=\sum_n\widehat C_n$ lies in $\E$ and
$\remp_m(F)=\sum_n\remp_m(\widehat C_n)=R_m$ for every $m$, since exactly one
summand has a nonzero class on shell $m$. By the identification in the
necessity proof, $F$ realizes every prescribed jet.
\end{proof}

\begin{corollary}[An explicit truncation certificate]\label{ent:sm:cor:error}
In that construction, if $L\ge N_0$ and $\gamma\ge-\gamma_{L+1}/2$, then
\[
 m_\gamma\Bigl(F-\sum_{n\le L}\widehat C_n\Bigr)\ \ge\
 m_\gamma(P_D)+\gamma_{L+1}.
\]
Arbitrarily large shell degrees $M_n$ and integers $c_n$ do not weaken it.
\end{corollary}
\begin{proof}
Every remaining summand satisfies the bound of the preceding proof with
$\gamma_n\ge\gamma_{L+1}$, and a strong sum has weighted valuation at least
the least weighted valuation of its terms.
\end{proof}

\begin{proof}[Completion of the second route: a nonzero finite part]
Let $P_0=Z^{D(0)}\prod_x(Z-x)^{D(x)}$, the product over the nonzero points of
$\supp D_0$, and let $F_0$ be a polynomial with the prescribed jets on $D_0$.
A solution has the form $F=F_0+P_0G$, where $G$ must carry the shell data
$\bigl(R_n-\remp_n(F_0)\bigr)\remp_n(P_0)^{-1}$. Each class $\remp_n(P_0)$ is
a unit of $\Lambda_n^\circ$: it is $\tau_n^{\deg P_0}$ times $T^{\deg P_0}$
times the classes of the factors $1-(x/\tau_n)T^{-1}$, where
$v(x/\tau_n)=v(x)+\gamma_n>0$ because $v(x)\ge0$, and $\tau_n$ is a unit of
$V_n$. Since $\remp_n(F_0)\in\Lambda_n^\circ$ for large $n$ by the necessity
half, the new data lie in $\RD$ exactly when the original data do. Apply the
case $D_0=0$ to $D-D_0$ to obtain $G$. The kernel half of
\cref{ent:thm:quotient} is \cref{ent:prop:kernel-direct}, which the fourth
source proved in the same way, by coefficient division without preparation.
\end{proof}

\begin{remark}[One simple node per shell]\label{ent:sm:rem:simple-cardinal}
If shell $n$ carries one simple node $\rho_n$ and the data are values $b_n$,
then $\Lambda_n=\K$, $\remp_n(G)=G(\rho_n)$, and the cardinal correction is
\[
 \widehat C_{n,d}(Z)=b_n\,\frac{\widehat P_n(Z)}{\widehat P_n(\rho_n)}
 \Bigl(\frac Z{\rho_n}\Bigr)^{d}.
\]
For instance, in the group $\Gamma_\infty$ of \cref{ent:sec:examples}, with
$\rho_n=t^{-e_n}$, the admissible values $b_n=t^{-n!e_n}$ of
\cref{ent:cor:no-interpolation}(b) are interpolated explicitly by
\[
 F(Z)=\sum_{n\ge1}t^{-n!e_n}\,
 \frac{\widehat P_n(Z)}{\widehat P_n(t^{-e_n})}\,
 \bigl(t^{e_n}Z\bigr)^{2\,n!+2},
\]
whereas the values $t^{-e_{n+1}}$ are admissible at no shell. For a single
shell, $d_n$ is an ordinary finite exponent: it absorbs a negative coefficient
valuation bounded by a finite multiple of $\gamma_n$, but not one below every
such multiple. That is why coarsening by $H_n$ is both necessary and
sufficient, rather than a convenient growth condition.
\end{remark}

\subsection{A universal bound on values at escaping monomials}
The next lemma was the original core obstruction, and its necessity half is
still what the negative results use. It is short, but it uses both halves of
\cref{ent:def:entire}: a common lower bound on the coefficient supports, and
legitimate evaluation at the large argument. \Cref{ent:thm:image} now
supplies its converse, so the bound can be stated as an exact criterion.

\begin{lemma}[Archimedean growth barrier, exact form]
\label{ent:lem:growth-barrier}
Let $(\gamma_n)_{n\ge1}$ be positive and strictly increasing with cofinal
range, so that in particular $\gamma_n\to+\infty$ cofinally.
\begin{enumerate}[label=\textnormal{(\alph*)},leftmargin=2em]
\item \emph{(Necessity.)} For every $F\in\E$ there is an $n_0$ such that
$F(t^{-\gamma_n})\in V_n$ for all $n\ge n_0$; that is, for each such $n$ some
integer $c_n\ge0$ satisfies $v(F(t^{-\gamma_n}))\ge-c_n\gamma_n$. In
particular, if $\lambda_n>0$ satisfies $\lambda_n\ggA\gamma_n$ for every
sufficiently large $n$, then eventually
\begin{equation}\label{ent:eq:growth-barrier}
 v\bigl(F(t^{-\gamma_n})\bigr)>-\lambda_n .
\end{equation}
The value $+\infty$ is allowed when $F$ vanishes at the argument.
\item \emph{(Sufficiency.)} Conversely, assume $(\gamma_n)$ is strictly
increasing, and let $b_n\in\K$ satisfy $b_n\in V_n$ for all sufficiently
large $n$. Then there is an $F\in\E$ with $F(t^{-\gamma_n})=b_n$ for every
$n\ge1$, with no condition at all on the finitely many exceptional indices.
\end{enumerate}
Hence for a strictly increasing positive cofinal $(\gamma_n)$ the prescription
$F(t^{-\gamma_n})=b_n$ is solvable in $\E$ if and only if $b_n\in V_n$
eventually.
\end{lemma}
\begin{proof}
(a) The assertion is immediate for $F=0$. Write $F=\sum_{k\ge0}a_kZ^k$ and let
$m=\min_{a_k\ne0}v(a_k)$, which exists by strong summability at $Z=1$. The
terms' valuations at $Z=t^{-\gamma_n}$ tend cofinally to infinity by
\cref{ent:thm:criterion}, so their nonempty set has a minimum, attained at a
finite index $k_n$, and the valuation of their Hahn sum is at least that
minimum. For large $n$ we have $\gamma_n>\abs m$ and therefore
\[
 v\bigl(F(t^{-\gamma_n})\bigr)\ \ge\ m-k_n\gamma_n\ >\ -(k_n+1)\gamma_n ,
\]
which is the first assertion with $c_n=k_n+1$. If moreover
$\lambda_n\ggA\gamma_n$, then $-(k_n+1)\gamma_n>-\lambda_n$, which is
\eqref{ent:eq:growth-barrier}.

(b) The points $t^{-\gamma_n}$ are distinct and their divisor is radially
finite with one simple node per shell, so
\cref{ent:cor:single-node} applies and supplies the interpolant.
\end{proof}

Part~(a) is the form used in \cref{ent:thm:no-bezout} and in
\cref{ent:cor:no-interpolation}; it is elementary and independent of the
constructive half of \cref{ent:thm:image}. Part~(b) is not elementary, and
the dependency is recorded in \cref{ent:app:audit}. The monotonicity
hypothesis is stated because the proof of (a) needs $\gamma_n$ to exceed a
fixed bound \emph{eventually}, which a merely cofinal set of values does not
supply; every application in this report --- \cref{ent:lem:escaping-classes}
and \cref{ent:lem:shells} --- produces a strictly increasing sequence. No
conclusion has been weakened by making the hypothesis explicit.

The eventual condition on $\gamma_n$ cannot be replaced by cofinality of
its range. In the reverse-lexicographic group
$\bigoplus_{j\ge0}\Q e_j$, set
$\gamma_{2n}=e_{n+2}$, $\gamma_{2n+1}=e_0$ and
$\lambda_{2n}=e_{n+3}$, $\lambda_{2n+1}=e_1$ for $n\ge0$.
The range of $(\gamma_n)$ is cofinal and $\lambda_n\ggA\gamma_n$ always,
but the constant entire function $H=t^{-e_2}$ violates the conclusion at
every odd index. The strictly increasing sequences used here avoid this failure.


\begin{corollary}[Explicit failure of value interpolation, and its exact
boundary]
\label{ent:cor:no-interpolation}
Assume $\cf(\Gamma)=\aleph_0$ and $\Gamma$ has no order unit, and take the
sequence $(\gamma_n)$ of \cref{ent:lem:escaping-classes}.
\begin{enumerate}[label=\textnormal{(\alph*)},leftmargin=2em]
\item \emph{(The negative half.)} The prescription
\[
 F(t^{-\gamma_n})=t^{-\gamma_{n+1}}\qquad(n\ge1)
\]
has no solution $F\in\E$, although each finite subcollection of these
prescriptions has a polynomial solution.
\item \emph{(The positive half.)} For \emph{every} sequence $(d_n)$ of
nonnegative integers the prescription $F(t^{-\gamma_n})=t^{-d_n\gamma_n}$
does have an entire solution. No rate restriction on $(d_n)$ is imposed. More
generally one may prescribe, at each node, finitely many derivatives whose
number grows arbitrarily with $n$, provided each shell's finitely many
coefficients satisfy the same scale condition.
\end{enumerate}
Together, (a) and (b) are the two sides of the criterion of
\cref{ent:lem:growth-barrier}: what an entire function may not do is jump
repeatedly to a strictly larger Archimedean class, while arbitrarily high
\emph{finite} powers of the current scale are always available.
\end{corollary}
\begin{proof}
(a) The prescribed values have valuation $-\gamma_{n+1}$, and
$\gamma_{n+1}\ggA\gamma_n$, so they lie in no $V_n$; this contradicts
\cref{ent:lem:growth-barrier}(a). The nodes are distinct, so ordinary finite
Lagrange interpolation over the field $\K$ realizes any finite subcollection.
(b) Each $t^{-d_n\gamma_n}$ lies in $V_n$ by definition, so
\cref{ent:lem:growth-barrier}(b) applies; the jet version is
\cref{ent:cor:single-node}.
\end{proof}

This failure is not caused by nodes accumulating in a bounded ball: their
divisor is radially finite and is therefore the divisor of an entire function
by \cref{ent:thm:product}. Realizability of zeros is strictly weaker than
realizability of values, and \cref{ent:thm:image} measures the difference
exactly.

\subsection{Two functions with no common zero and no B\'ezout identity}
\label{ent:subsec:no-bezout}
The quotient description now decides the B\'ezout question for a canonical
product against an arbitrary entire second factor. Two finite unit
computations are needed first.

\begin{lemma}[Units of a shell algebra]\label{ent:lem:shellunits}
For $R\in V_n[T]$ of degree below $M_n$ the following are equivalent:
its class is a unit of $\Lambda_n^\circ$; $\Res(P_n,R)$ is a unit of $V_n$;
every value $R(\xi_{n,j})$ is a unit of $V_n$.
\end{lemma}
\begin{proof}
Multiplication by $R$ on the finite free $V_n$-module $\Lambda_n^\circ$ is
invertible if and only if its determinant is a unit, by determinants in one
direction and the adjugate in the other. Formula \eqref{ent:eq:detnorm} gives
$\Res(P_n,R)=\prod_jR(\xi_{n,j})^{m_{n,j}}$. All the $\xi_{n,j}$ and all the
coefficients lie in $V_n$, so every factor lies in that local ring, and a
finite product in a local ring is a unit exactly when each factor is. The
case of a vanishing factor is included.
\end{proof}

The statement keeps repeated roots and does not require the normalized nodes
to have distinct residues. The determinant appears here as an exact
finite-algebra unit test, not as an approximate stability threshold.

\begin{lemma}[Units of the restricted product]\label{ent:lem:restrictedunits}
An element $(r_0,r_1,r_2,\ldots)$ of $\RD$ is a unit if and only if $r_0$ is a
unit of $\Lambda_0$, every $r_n$ is a unit of $\Lambda_n$, and
$r_n\in(\Lambda_n^\circ)^\times$ for all sufficiently large $n$.
\end{lemma}
\begin{proof}
A ring inverse must be the componentwise inverse. If both the element and its
inverse lie in $\RD$, then both components lie in $\Lambda_n^\circ$ for all
large $n$, making those components units there. Conversely those conditions
put the componentwise inverse in $\RD$.
\end{proof}

\begin{theorem}[Sharp B\'ezout criterion]\label{ent:thm:bezout-criterion}
Let $D$ be radially finite with canonical product $P_D$, and let $F\in\E$.
There are $A,B\in\E$ with
\begin{equation}\label{ent:eq:bezout-identity}
 AP_D+BF=1
\end{equation}
if and only if both of the following hold:
\begin{enumerate}[label=\textnormal{(\arabic*)},leftmargin=2em]
\item $F(x)\ne0$ at every node $x$ of $D$;
\item for all sufficiently large occupied shells $n$ there is a finite
integer $c_n\ge0$ with
\begin{equation}\label{ent:eq:valueupper}
 v\bigl(F(\rho_{n,j})\bigr)\le c_n\gamma_n\qquad(1\le j\le q_n).
\end{equation}
\end{enumerate}
For a finite divisor the second condition is vacuous, and for a $\Gamma$ with
an order unit it is automatic. Moreover the criterion depends only on the
support of $D$ and on the values of $F$ there: it is unchanged by any
reassignment of positive finite multiplicities, even unbounded ones.
\end{theorem}
\begin{proof}
For a finite divisor, finite Hermite interpolation supplies the inverse jets
when all values are nonzero, and \cref{ent:cor:gcd} converts them into
\eqref{ent:eq:bezout-identity}; necessity is evaluation. Assume $D$ infinite.
An identity \eqref{ent:eq:bezout-identity} exists exactly when $F\bmod P_D$ is
a unit of $\E/(P_D)$, so \cref{ent:thm:quotient,ent:lem:restrictedunits}
apply. The unit condition in $\Lambda_0$ and in each $\Lambda_n$ is
nonvanishing of the constant terms of all local jets, which is (1).

For the tail condition, necessity in \cref{ent:thm:image} puts the normalized
shell remainder $\remp_n(F)$ in $V_n[T]$ for all large $n$, and its values at
the normalized nodes are $F(\rho_{n,j})$. By \cref{ent:lem:shellunits} its
inverse lies in $\Lambda_n^\circ$ exactly when those values are units of
$V_n$. The lower halves of the unit condition \eqref{ent:eq:Vunits} are
automatic from membership in $V_n$; what remains is exactly the upper bound
\eqref{ent:eq:valueupper}, and a maximum over the finitely many nodes of a
shell produces one $c_n$ per shell.

Conversely, under (1) and (2) the local inverse jets of $F$ have shell
representatives in $V_n[T]$ for all large $n$, so \cref{ent:thm:image} yields
a $B\in\E$ carrying them. Then $1-BF$ vanishes on $D$ to the required orders,
and \cref{ent:cor:gcd} supplies $A$.

Finally, both conditions refer only to $\supp D$ and to the values of $F$
there, and any assignment of positive finite multiplicities to a radially
finite set of nodes is again radially finite.
\end{proof}

The simplification to a pointwise test is special to the \emph{inverse jets of
an already entire function}. Arbitrary Hermite data require the full
coefficient criterion of \cref{ent:thm:image}; replacing it by a pointwise
condition would lose \cref{ent:thm:close-pair}.

The fourth source proved \cref{ent:thm:bezout-criterion} independently, as a
consequence of its own proof of \cref{ent:thm:quotient}, and extended it to
finitely many generators and to membership of arbitrary elements.

\begin{theorem}[Finitely many generators]\label{ent:sm:thm:finite-generators}
Let $D$ be an infinite radially finite divisor with $D_0=0$ and canonical
product $P_D$, and let $G_1,\ldots,G_s\in\E$, $s\ge1$. Then
$(P_D,G_1,\ldots,G_s)=\E$ if and only if
\begin{enumerate}[label=\textnormal{(\arabic*)},leftmargin=2em]
\item at every node of $D$ at least one $G_\ell$ is nonzero, and
\item for every node $\rho_{n,j}$ of every sufficiently large shell, at least
one $G_\ell$ satisfies $G_\ell(\rho_{n,j})\ne0$ and
$v\bigl(G_\ell(\rho_{n,j})\bigr)\le c\gamma_n$ for some integer $c\ge0$;
equivalently, $v\bigl(G_\ell(\rho_{n,j})\bigr)\in H_n$.
\end{enumerate}
On a sufficiently large shell, the condition in (2) is equivalent to
\[
 \gcd\bigl(\overline{P_n},\overline{\remp_n(G_1)},\ldots,
   \overline{\remp_n(G_s)}\bigr)=1\quad\text{in }\kappa_n[T],
\]
the bars denoting reduction modulo $\mathfrak q_{\gamma_n}$; this holds even
when distinct normalized nodes have the same reduction.
\end{theorem}
\begin{proof}
By \cref{ent:thm:quotient} the ideal is the unit ideal exactly when the images
$g_\ell=\ev_D(G_\ell)$ generate the unit ideal of $\RD$, that is, when there
are $a_{\ell,n}\in\Lambda_n$ with $\sum_\ell a_{\ell,n}g_{\ell,n}=1$ for every
$n$ and $a_{\ell,n}\in\Lambda_n^\circ$ for all large $n$. In $\Lambda_n$,
which is a finite product of local jet rings over $\K$, the $g_{\ell,n}$
generate the unit ideal exactly when at each node one of them has a nonzero
constant term; this is (1).

For large $n$ every $g_{\ell,n}$ lies in $\Lambda_n^\circ$, by the necessity
half of \cref{ent:thm:image}. The reduction of $\Lambda_n^\circ$ modulo
$\mathfrak q_{\gamma_n}$ is $\kappa_n[T]/(\overline{P_n})$, and
$\overline{P_n}$ splits at the reductions of the $\xi_{n,j}$. The reduced
classes generate the unit ideal exactly when they have no common root among
those reductions, which is the gcd condition. The value of $g_{\ell,n}$ at
$\xi_{n,j}$ is $G_\ell(\rho_{n,j})\in V_n$, whose coarsened residue is nonzero
exactly when its valuation lies in $H_n$; so the gcd condition is (2) at
shell $n$, repeated residues included. A residue identity lifts
coefficientwise to $\sum_\ell a_\ell g_{\ell,n}=1+e$ with
$e\in\mathfrak q_{\gamma_n}\Lambda_n^\circ$. The multiplication determinant of
$1+e$ is $1$ modulo $\mathfrak q_{\gamma_n}$, so $1+e$ is a unit of
$\Lambda_n^\circ$ and the identity can be divided by it. Hence the
$g_{\ell,n}$ generate the unit ideal of $\Lambda_n^\circ$ exactly when (2)
holds at shell $n$.

If (1) and (2) hold, choose the $a_{\ell,n}$ in $\Lambda_n$ for the finitely
many exceptional shells and in $\Lambda_n^\circ$ afterwards, and realize each
tuple $(a_{\ell,n})_n$ by some $B_\ell\in\E$ (\cref{ent:thm:image}). Then
$1-\sum_\ell B_\ell G_\ell$ has zero image, so it lies in $P_D\E$.
Conversely, a global identity reduces to identities in every $\Lambda_n$ and,
for large $n$, in $\Lambda_n^\circ$.
\end{proof}

For $s=1$ this is \cref{ent:thm:bezout-criterion} in the case $D_0=0$. As
there, finite nilpotent directions do not enter: a unit of a local jet ring
is detected by its constant term, so only values matter, even for unbounded
multiplicities.

\begin{corollary}[Membership, one simple node per shell]
\label{ent:sm:cor:membership}
Suppose $D$ has one simple node $\rho_n$ on each shell and $D_0=0$, and let
$F,G\in\E$. Then $F\in(P_D,G)$ if and only if $F(\rho_n)=0$ whenever
$G(\rho_n)=0$, and the ratios $F(\rho_n)/G(\rho_n)$ at the remaining nodes lie
in $V_n$ for all sufficiently large $n$.
\end{corollary}
\begin{proof}
If $F=AP_D+BG$, then $F(\rho_n)=B(\rho_n)G(\rho_n)$, and $B(\rho_n)\in V_n$
for all large $n$ by the necessity half of \cref{ent:cor:single-node}.
Conversely, prescribe $B(\rho_n)$ to be the ratio where $G(\rho_n)\ne0$ and
$0$ elsewhere; \cref{ent:cor:single-node} realizes $B$, and $F-BG$ vanishes
at every node, so it lies in $P_D\E$ by \cref{ent:thm:quotient}.
\end{proof}

\begin{theorem}[An explicit pointwise B\'ezout failure]
\label{ent:thm:no-bezout}
Assume the hypotheses of \cref{ent:lem:escaping-classes} and choose its
sequence $(\gamma_n)$. Define
\begin{align}
 \rho_n&=t^{-\gamma_n},&
 \sigma_n&=t^{-\gamma_n}+t^{\gamma_{n+1}},\label{ent:eq:paired-roots}\\
 f(Z)&=\prod_{n\ge1}(1-Z/\rho_n),&
 g(Z)&=\prod_{n\ge1}(1-Z/\sigma_n).\label{ent:eq:paired-products}
\end{align}
Then $f,g\in\E$, both have constant coefficient $1$, their roots are simple
and disjoint, and
\[
 \gcd(f,g)=1,\qquad (f,g)\ne\E.
\]
In particular there are no $A,B\in\E$ with $Af+Bg=1$. The failure is exactly
a failure of condition~(2) of \cref{ent:thm:bezout-criterion}: condition~(1)
holds, while
\begin{equation}\label{ent:eq:lambda}
 \lambda_n:=v(g(\rho_n))
   =\gamma_{n+1}+(2-n)\gamma_n+\sum_{j<n}\gamma_j
\end{equation}
exceeds every finite multiple of $\gamma_n$.
\end{theorem}
\begin{proof}
We have $v(\rho_n)=v(\sigma_n)=-\gamma_n$, so both root sequences escape
every valuation ball and \cref{ent:thm:product} defines the two entire
functions. Different indices have different valuations; at the same index
$\sigma_n-\rho_n=t^{\gamma_{n+1}}\ne0$. Their zero sets are consequently
disjoint and simple, and \cref{ent:cor:gcd} gives their GCD.

We compute $v(g(\rho_n))$ exactly. For $j<n$, $\rho_n/\sigma_j$ has negative
valuation, so
\[
 v(1-\rho_n/\sigma_j)=\gamma_j-\gamma_n.
\]
For $j=n$,
\[
 1-\rho_n/\sigma_n=\frac{\sigma_n-\rho_n}{\sigma_n},\qquad
 v(1-\rho_n/\sigma_n)=\gamma_{n+1}+\gamma_n.
\]
For $j>n$ the ratio has positive valuation, and the infinite tail product has
residue $1$ and valuation zero by the positive-support product argument used
in \cref{ent:thm:product}. This gives \eqref{ent:eq:lambda}.
The sum there is finite. Since $0<\gamma_j<\gamma_n$ for $j<n$, the part of
\eqref{ent:eq:lambda} other than $\gamma_{n+1}$ is bounded in absolute value
by a finite multiple of $\gamma_n$. Hence $\lambda_n\ggA\gamma_n$.

Suppose $Af+Bg=1$. Evaluate at $\rho_n$. The values are well defined and
evaluation is multiplicative, so
\[
 B(\rho_n)=\frac1{g(\rho_n)},\qquad
 v(B(\rho_n))=-\lambda_n.
\]
For all sufficiently large $n$ this contradicts
\cref{ent:lem:growth-barrier}(a). Thus the ideal is proper.

This argument is self-contained and uses only part~(a) of
\cref{ent:lem:growth-barrier}. Alternatively, and equivalently,
$D=\sum_n[\rho_n]$ has one simple node per shell, $f=P_D$ exactly, since
$D(0)=0$ in \eqref{ent:eq:product}; condition~(1) of
\cref{ent:thm:bezout-criterion} holds because $g$ has no zero in common with
$f$, and \eqref{ent:eq:lambda} with $\lambda_n\ggA\gamma_n$ is exactly the
failure of \eqref{ent:eq:valueupper}. The criterion therefore both recovers
this example and says what would have had to be true instead.
\end{proof}

The cancellation-sensitive step is \eqref{ent:eq:lambda}. No unproved
assertion that nearby zeros force a small value is used: the contribution of
each earlier root is retained, and the infinite tail is proved to be a unit.
The earlier-root terms can lower the valuation by finite multiples of
$\gamma_n$, but cannot offset the new Archimedean scale $\gamma_{n+1}$.

The fourth source turned this single instance into an exact threshold.

\begin{theorem}[The exact threshold for paired zeros]\label{ent:sm:thm:paired}
Let $0<\gamma_1<\gamma_2<\cdots$ be cofinal in $\Gamma$, with no domination
assumed, let $\vartheta_n>0$ in $\Gamma$, and put
\[
 \rho_n=t^{-\gamma_n},\qquad
 \sigma_n=\rho_n+t^{\vartheta_n},\qquad
 f=\prod_{n\ge1}(1-Z/\rho_n),\qquad
 g_\vartheta=\prod_{n\ge1}(1-Z/\sigma_n).
\]
Then $f,g_\vartheta\in\E$ have simple and disjoint zero sets,
\begin{equation}\label{ent:sm:eq:paired-value}
 v\bigl(g_\vartheta(\rho_n)\bigr)
 =\vartheta_n+(2-n)\gamma_n+\sum_{h<n}\gamma_h,
\end{equation}
and $(f,g_\vartheta)=\E$ if and only if $\vartheta_n\in H_n$ for all
sufficiently large $n$.
\end{theorem}
\begin{proof}
Both root sequences have valuation $-\gamma_n$ at index $n$, distinct indices
give distinct valuations, and $\sigma_n-\rho_n\ne0$; \cref{ent:thm:product}
gives the zero assertions. The computation of \eqref{ent:eq:lambda} applies
verbatim with $\gamma_{n+1}$ replaced by $\vartheta_n$: the factor of index
$h<n$ has valuation $\gamma_h-\gamma_n$, the factor of index $n$ has
valuation $\vartheta_n+\gamma_n$, and the tail is a unit of valuation zero.
The terms of \eqref{ent:sm:eq:paired-value} other than $\vartheta_n$ lie in
$H_n$, so $v(g_\vartheta(\rho_n))\in H_n$ exactly when $\vartheta_n\in H_n$.
Apply \cref{ent:thm:bezout-criterion} to $D=\sum_n[\rho_n]$, for which
$f=P_D$ and condition~(1) holds.
\end{proof}

Under the hypotheses of \cref{ent:lem:escaping-classes} the displacement
$\vartheta_n=\gamma_{n+1}$ lies in no $H_n$, and this is
\cref{ent:thm:no-bezout}. At the other extreme a displacement
$\vartheta_n=n!\,\gamma_n$ causes no obstruction. The dividing line is not
the size of an ordinary numerical coefficient but whether the displacement
leaves the convex subgroup of the current radius, and infinitely many such
shells suffice to make the ideal proper.

\begin{theorem}[A close pair with individually harmless values]
\label{ent:thm:close-pair}
Keep $\rho_n,\sigma_n$ of \eqref{ent:eq:paired-roots} and regard
$\{\rho_n,\sigma_n\}$ as one shell of two simple nodes, with
$\tau_n=t^{-\gamma_n}$, normalized nodes $1$ and $1+\varepsilon_n$, and
\[
 \varepsilon_n=t^{\gamma_{n+1}+\gamma_n}.
\]
The prescription $F(\rho_n)=0$, $F(\sigma_n)=c_n$ is realized by some
$F\in\E$ if and only if
\begin{equation}\label{ent:eq:closeexact}
 c_n/\varepsilon_n\in V_n\qquad\text{for all sufficiently large }n .
\end{equation}
In particular there is \emph{no} $F\in\E$ with $F(\rho_n)=0$ and
$F(\sigma_n)=1$ for all $n$, even though every prescribed value lies in every
$V_n$ and every finite subproblem has a polynomial solution. By contrast
$c_n=\varepsilon_n$ and $c_n=t^{\gamma_{n+1}}$ are both realizable.
\end{theorem}
\begin{proof}
The unique shell interpolation polynomial of degree below two with
$R_n(1)=0$ and $R_n(1+\varepsilon_n)=c_n$ is
$R_n(T)=(c_n/\varepsilon_n)(T-1)$, and both of its coefficients are
$\pm c_n/\varepsilon_n$. \Cref{ent:thm:image} gives
\eqref{ent:eq:closeexact}. For $c_n=1$ both nonzero coefficients have
valuation $-\gamma_{n+1}-\gamma_n$, which lies strictly below $-m\gamma_n$
for every $m\in\N$ because $\gamma_{n+1}\ggA\gamma_n$; the criterion fails at
every shell. Finite subproblems are ordinary Lagrange interpolation over
$\K$. For $c_n=\varepsilon_n$ one gets $R_n=T-1$, and for
$c_n=t^{\gamma_{n+1}}$ one gets $R_n=t^{-\gamma_n}(T-1)$; both lie in
$V_n[T]$.
\end{proof}

This obstruction is genuinely joint. Neither the list of node radii nor the
separate valuations of the prescribed values detects it: the divided
difference carries an additional negative Archimedean scale. It is also the
exact reason why \cref{ent:cor:coarse-separated} needs its hypothesis, and
why applying \cref{ent:cor:single-node} node by node on a multi-node shell is
wrong.

\subsection{Maximal ideals that are not evaluations}
\label{ent:subsec:invisible}

\begin{corollary}[A finitely generated ideal invisible to point tests]
\label{ent:cor:invisible-ideal}
In the situation of \cref{ent:thm:no-bezout} there is a proper two-generated
ideal contained in no evaluation maximal ideal $(Z-x)\E$. Every maximal ideal
containing it is therefore not the kernel of evaluation at any $x\in\K$.
\end{corollary}
\begin{proof}
The ideal $(f,g)$ is proper but $f$ and $g$ do not vanish simultaneously. By
\cref{ent:cor:evaluation}, containment in $(Z-x)\E$ would mean that both
vanish at $x$. A maximal ideal containing a proper ideal exists by Zorn's
lemma; none of these maximal ideals is an evaluation ideal.
\end{proof}

This is a failure of the assertion that \emph{every finitely generated proper
ideal has a common zero}. We do not call it a metric corona counterexample,
because no norm lower bound and no bounded-function algebra has been
specified. Nor do we claim that non-evaluation maximal ideals are absent in
the rank-one entire ring; the statement concerns a \emph{two-generated} ideal
already lacking a common zero.

\Cref{ent:cor:invisible-ideal} produces its maximal ideals by Zorn's lemma
and names none of them. It is kept because it is weaker and cheaper: it needs
only \cref{ent:cor:evaluation} and \cref{ent:thm:no-bezout}, and in
particular none of the image theorem. The next theorem constructs such
maximal ideals explicitly, with their residue fields, at the cost of using
\cref{ent:thm:image} and a nonprincipal ultrafilter.

For the rest of this subsection let $D$ have one simple node $\rho_n$ per
occupied shell and empty finite part, so that
\[
 P_D=\prod_{n\ge1}(1-Z/\rho_n),
 \qquad
 \RD=\{(b_n)\in\K^{\N_{>0}}:b_n\in V_n\text{ eventually}\}
\]
by \cref{ent:thm:quotient,ent:cor:single-node}. Adjoining a finite part
changes none of the constructions below. Fix a nonprincipal ultrafilter
$\Ucal$ on the positive integers --- this is an ordinary ZFC choice, not an
algorithmically computed object --- and set
\begin{equation}\label{ent:eq:ultraproduct}
 L_\Ucal=\Bigl(\prod_{n\ge1}\kappa_n\Bigr)\big/\Ucal,
 \qquad
 \kappa_n=\C\bigl((t^{H_n})\bigr),
\end{equation}
two sequences being identified when they agree on a member of $\Ucal$. This
quotient is a field: a nonzero class has nonzero components on a member of
$\Ucal$, and the reciprocals there give its inverse.

\begin{theorem}[Explicit non-evaluation maximal ideals]\label{ent:thm:maximal}
The formula
\begin{equation}\label{ent:eq:Phi}
 \Phi_\Ucal(F)=\bigl[\res_n\bigl(F(\rho_n)\bigr)\bigr]_\Ucal ,
\end{equation}
where $\res_n:V_n\to\kappa_n$ is the coarsened residue map of
\cref{ent:lem:coarse-residue} and the residue is taken at all sufficiently
large $n$, defines a surjective ring homomorphism $\E\to L_\Ucal$. Arbitrary
values at the finitely many remaining indices do not affect the class. Its
kernel is the maximal ideal
\begin{equation}\label{ent:eq:Mexplicit}
 \mathfrak M_\Ucal=
 \bigl\{F\in\E:\{n:v(F(\rho_n))>m\gamma_n\text{ for every }m\in\N\}
 \in\Ucal\bigr\},
\end{equation}
which contains $P_D$ and is not the kernel of evaluation at any point of
$\K$. Distinct nonprincipal ultrafilters give distinct maximal ideals, and
$\E/\mathfrak M_\Ucal\cong L_\Ucal$.
\end{theorem}
\begin{proof}
Necessity in \cref{ent:thm:image}, in the form
\cref{ent:lem:growth-barrier}(a), gives the eventual membership
$F(\rho_n)\in V_n$ that \eqref{ent:eq:Phi} needs. Addition and multiplication
are respected on a common cofinite set of indices, so $\Phi_\Ucal$ is a ring
homomorphism, and changing finitely many components does not change the class
at a nonprincipal ultrafilter. For surjectivity, given $c_n\in\kappa_n$ use
the canonical inclusion $\kappa_n=\C((t^{H_n}))\subseteq\K$ to lift to
$b_n\in V_n$ and apply \cref{ent:cor:single-node}.
\Cref{ent:lem:coarse-residue} identifies the kernel with
\eqref{ent:eq:Mexplicit}, and since $L_\Ucal$ is a field that kernel is
maximal. It contains $P_D$, which vanishes at every node.

Suppose the kernel were evaluation at a point $x$. Then $P_D(x)=0$, so
$x=\rho_r$ for some $r$ by \cref{ent:thm:product}. The bounded prescription
$b_r=1$ and $b_n=0$ for $n\ne r$ has an entire interpolant $F_r$ by
\cref{ent:cor:single-node}; its residue class is zero at a nonprincipal
ultrafilter, so $F_r\in\mathfrak M_\Ucal$, while $F_r(\rho_r)=1$. This
excludes every evaluation point.

For distinct nonprincipal ultrafilters $\Ucal$ and $\mathcal W$, choose
$S\in\Ucal$ with $S\notin\mathcal W$ and interpolate the values $1$ on $S$
and $0$ off $S$; the image is $1$ under $\Phi_\Ucal$ and $0$ under
$\Phi_\mathcal W$, so the kernels differ. The residue-field identification is
the first isomorphism theorem.
\end{proof}

\begin{remark}[How the coefficient field acts]\label{ent:rem:coeff-action}
The copy of $\K$ inside $L_\Ucal$ is not a constant-coordinate embedding into
each $\kappa_n$. A scalar $c\in\K$ maps to the class of its eventual support
restrictions $[c|_{H_n}]_\Ucal$. For $c\ne0$ the leading exponent of $c$ lies
in $H_n$ for all large $n$, so those residues are nonzero; the map is an
embedding and supplies the $\K$-algebra structure on the target. It remains
valid when $\supp c$ is contained in no single $H_n$.
\end{remark}

\begin{corollary}[Finite-rank residue fields in an infinite-rank example]
\label{ent:cor:residue-fields}
Take $\Gamma_\infty=\bigoplus_{j\ge0}\Q e_j$ of
\cref{ent:subsec:gamma-infinity} with $\gamma_n=e_n$ for $n\ge1$. Then
\[
 H_n=\bigoplus_{j=0}^{n}\Q e_j,
 \qquad
 \E/\mathfrak M_\Ucal\ \cong\
 \Bigl(\prod_{n\ge1}
  \C\bigl(\bigl(t^{\oplus_{j=0}^{n}\Q e_j}\bigr)\bigr)\Bigr)\big/\Ucal .
\]
Every nonprincipal $\mathfrak M_\Ucal$ contains both of the functions $f,g$ of
\eqref{ent:eq:paired-products}, although they have no common zero.
\end{corollary}
\begin{proof}
In the largest-index order, every element supported in $e_0,\ldots,e_n$ is
bounded in absolute value by a finite multiple of $e_n$, while an element
with a nonzero coordinate of index above $n$ is not; this is $H_n$. Note that
the index $0$ is included: $H_n$ has rank $n+1$, not $n$.
\Cref{ent:thm:maximal} gives the residue field. Now $f(\rho_n)=0$ for every
$n$, so $f\in\mathfrak M_\Ucal$; and \eqref{ent:eq:lambda} with
$\lambda_n\ggA\gamma_n$ puts $g(\rho_n)$ in $\mathfrak q_{\gamma_n}$ for
every $n$, so $g\in\mathfrak M_\Ucal$ as well.
\end{proof}

Each of these residue fields is a full Hahn field on a convex subgroup of
$\Gamma$, exactly the shape produced by the coarsening of
\cref{ent:lem:coarsened-ring} in an entirely different situation. The
coincidence of shape is worth noticing and is not an identification: one
coarsening happens inside $\Gamma$ at a shell radius, the other across an
ordered-group extension.

This construction gives a specified family of maximal ideals together with a
field presentation of each. By itself it does \emph{not} classify all maximal
ideals of $\E$, and it does not assert that every non-evaluation maximal
ideal arises this way. The fourth source's
\cref{ent:sm:thm:maximal-all} below shows that nothing else occurs
\emph{above} $P_D$; the maximal ideals of $\E$ that contain no such product
remain undescribed. The ultrafilter mechanism itself is standard algebra;
what is new is the analytic surjection onto $\RD$ that feeds it. No
model-theoretic transfer principle is used.

\subsection{All maximal ideals above a one-node-per-shell product}
\label{ent:sm:subsec:all-maximal}
This subsection is the fourth source's. It keeps the one-node-per-shell
hypothesis of \cref{ent:subsec:invisible} and the empty finite part, but
allows each node $\rho_n$ an arbitrary multiplicity $m_n\ge1$, so that
$P_D=\prod_n(1-Z/\rho_n)^{m_n}$. In the coordinate
$U=(Z-\rho_n)/\rho_n=(T-\xi_n)/\xi_n$, where $\xi_n$ is a unit of $V_n$,
triangular changes of basis over $V_n$, as in the proof of
\cref{ent:cor:single-node}, identify $\Lambda_n$ with $\K[U]/(U^{m_n})$ and
$\Lambda_n^\circ$ with $V_n[U]/(U^{m_n})$. By \cref{ent:thm:quotient},
\begin{equation}\label{ent:sm:eq:local-product}
 \E/(P_D)\ \cong\ \RD=\Bigl\{(r_n)\in\prod_n\K[U]/(U^{m_n}):
   r_n\in V_n[U]/(U^{m_n})\text{ eventually}\Bigr\},
\end{equation}
and each $V_n[U]/(U^{m_n})$ is a local ring with maximal ideal
$(\mathfrak q_{\gamma_n},U)$ and residue field $\kappa_n$. Formula
\eqref{ent:eq:Phi} still defines $\Phi_\Ucal$, since $F(\rho_n)\in V_n$
eventually by \cref{ent:cor:single-node}.

\begin{theorem}[Every maximal ideal above $P_D$]\label{ent:sm:thm:maximal-all}
The maximal ideals of $\E$ containing $P_D$ are exactly
\begin{enumerate}[label=\textnormal{(\alph*)},leftmargin=2em]
\item the evaluation ideals $(Z-\rho_n)\E$, $n\ge1$, with residue field
$\K$; and
\item the kernels $\mathfrak M_\Ucal$ of $\Phi_\Ucal$, one for each
nonprincipal ultrafilter $\Ucal$ on $\N_{>0}$, with
$\E/\mathfrak M_\Ucal\cong L_\Ucal=\bigl(\prod_n\C((t^{H_n}))\bigr)/\Ucal$.
\end{enumerate}
Distinct ultrafilters give distinct ideals, and inserting arbitrary elements
of $\kappa_n$ at finitely many coordinates of \eqref{ent:eq:Phi} does not
change $\Phi_\Ucal$.
\end{theorem}
\begin{proof}
For simple nodes, (b) is \cref{ent:thm:maximal}, and its proof works verbatim
with multiplicities: \cref{ent:cor:single-node} realizes arbitrary values in
$V_n$ with zero higher jets, which gives surjectivity, and the kernel and
distinctness arguments use values only. The evaluation ideals contain $P_D$
and are maximal by \cref{ent:cor:evaluation}.

For exhaustiveness it suffices to find the maximal ideals $M$ of $\RD$. For
$S\subseteq\N_{>0}$ let $e_S\in\RD$ be the idempotent equal to $1$ on $S$ and
$0$ elsewhere; its coordinates lie in every $V_n[U]/(U^{m_n})$. In the field
$\RD/M$ each $e_S$ maps to $0$ or $1$, and $e_S+e_{\N_{>0}\setminus S}=1$,
$e_Se_T=e_{S\cap T}$; hence $\Ucal_M=\{S:e_S\notin M\}$ is an ultrafilter. If
$\Ucal_M$ is principal at $n$, then $1-e_{\{n\}}\in M$, the quotient factors
through $\K[U]/(U^{m_n})$, whose unique maximal ideal is $(U)$, and $M$
corresponds to $(Z-\rho_n)\E$.

Suppose $\Ucal_M$ is nonprincipal. If $re_S=0$ with $S\in\Ucal_M$, the image
of $e_S$ is $1$ and so the image of $r$ is $0$; thus $M$ contains every tuple
vanishing on a member of $\Ucal_M$. Modulo those tuples, $\RD$ becomes the
ultraproduct $\bigl(\prod_nV_n[U]/(U^{m_n})\bigr)/\Ucal_M$, because changing
finitely many coordinates makes any element of $\RD$ integral without
changing its class. An ultraproduct of local rings is local: a tuple is a
unit exactly when its coordinates are units on a member of $\Ucal_M$, the
nonunits are the tuples whose constant coefficients lie in
$\mathfrak q_{\gamma_n}$ on a member of $\Ucal_M$, and they form an ideal
with quotient $\bigl(\prod_n\kappa_n\bigr)/\Ucal_M$. So $M$ is the kernel of
$\Phi_{\Ucal_M}$. The idempotents recover $\Ucal_M$ from $M$, which gives
distinctness again.
\end{proof}

Write $\beta\N_{>0}$ for the set of all ultrafilters on $\N_{>0}$ with the
basic clopen sets $\widehat S=\{\Ucal:S\in\Ucal\}$: the Stone--\v{C}ech space
of the discrete set $\N_{>0}$, whose principal ultrafilters form a copy of
$\N_{>0}$.

\begin{theorem}[The maximal spectrum of the quotient]\label{ent:sm:thm:beta}
With the Zariski topology on the left, the bijection of
\cref{ent:sm:thm:maximal-all} is a homeomorphism
\[
 \operatorname{Max}\bigl(\E/(P_D)\bigr)\ \cong\ \beta\N_{>0},
\]
sending $(Z-\rho_n)\E$ to the principal ultrafilter at $n$. The evaluation
ideals form a discrete dense subset.
\end{theorem}
\begin{proof}
Work in \eqref{ent:sm:eq:local-product}. For $r=(r_n)$ let $c_n$ be the
constant coefficient of $r_n$ and put
$Z_r=\{n:c_n=0\}\subseteq T_r=\{n:c_n\in\mathfrak q_{\gamma_n}\}$. At the
principal ultrafilter at $n$, $r$ lies outside the maximal ideal exactly when
$n\notin Z_r$; at a nonprincipal $\Ucal$, exactly when
$\N_{>0}\setminus T_r\in\Ucal$, finitely many nonintegral coordinates being
irrelevant. So the basic open set $D(r)$ corresponds to
\[
 \widehat{\N_{>0}\setminus T_r}\ \cup\
 \{\text{principal ultrafilter at }n:\ n\in T_r\setminus Z_r\},
\]
which is open in $\beta\N_{>0}$ because principal points are open.
Conversely $D(e_S)$ corresponds exactly to $\widehat S$. Every nonempty
$\widehat S$ contains the principal ultrafilter at some $n\in S$, which gives
density.
\end{proof}

The ultrafilter topology is familiar. The new information is the varying
residue fields, which a topological description of the spectrum does not
see.

\begin{corollary}[The whole maximal-ideal locus of a paired example]
\label{ent:sm:cor:boundary-pair}
In \cref{ent:sm:thm:paired} let $S=\{n:\vartheta_n\notin H_n\}$. The maximal
ideals of $\E$ containing both $f$ and $g_\vartheta$ are exactly the
$\mathfrak M_\Ucal$ with $\Ucal$ nonprincipal and $S\in\Ucal$. For finite $S$
there are none and $(f,g_\vartheta)=\E$; for infinite $S$ they exist. In the
situation of \cref{ent:thm:no-bezout}, $S=\N_{>0}$, and the locus is all of
$\beta\N_{>0}\setminus\N_{>0}$.
\end{corollary}
\begin{proof}
No evaluation ideal contains both functions, since they have no common zero.
By \eqref{ent:sm:eq:paired-value}, $v(g_\vartheta(\rho_n))$ is $\vartheta_n$
plus an element of $H_n$. Since $\vartheta_n>0$, its class modulo $H_n$ is
positive exactly when $n\in S$, so $\Phi_\Ucal(g_\vartheta)=0$ exactly when
$S\in\Ucal$. Apply \cref{ent:sm:thm:maximal-all}. A nonprincipal ultrafilter
contains no finite set, and the ultrafilter lemma extends an infinite $S$ to
a nonprincipal ultrafilter.
\end{proof}

In particular every maximal ideal that \cref{ent:cor:invisible-ideal}
obtains from Zorn's lemma is one of the explicit ideals of
\cref{ent:thm:maximal}, and \cref{ent:cor:residue-fields} lists all of them.
The description still uses ultrafilters and choice; it computes no
particular nonprincipal ultrafilter.

\begin{corollary}[Algebraic type of the residue fields]
\label{ent:sm:cor:boundary-fields}
Every field $L_\Ucal$ is algebraically closed. Over the real Hahn field
(\cref{ent:sm:rem:real}) the corresponding field
$\bigl(\prod_n\R((t^{H_n}))\bigr)/\Ucal$ is real closed for its ultraproduct
order.
\end{corollary}
\begin{proof}
Each $H_n$ is divisible, so $\C((t^{H_n}))$ is algebraically closed by the
Hahn-field theorem \cite[Corollary~4]{Poonen}; $\R((t^{H_n}))$ is ordered and
its complexification is that field, so it is real closed. A monic polynomial
of fixed degree over an ultraproduct has monic coordinates of that degree on
a member of the ultrafilter, and roots can be chosen coordinatewise; the same
argument with square roots of positive elements and roots of odd-degree
polynomials gives real closedness.
\end{proof}

\begin{proposition}[A transcendental boundary coordinate]
\label{ent:sm:prop:transcendental}
Let $\Ucal$ be nonprincipal and let $\iota_\Ucal:\K\to L_\Ucal$ be the
restriction of $\Phi_\Ucal$, injective by \cref{ent:rem:coeff-action}. Then
$\Phi_\Ucal(Z)$ is transcendental over $\iota_\Ucal(\K)$. Over the real Hahn
field with the monomial nodes $\rho_n=t^{-\gamma_n}$, $\Phi_\Ucal(Z)$ exceeds
every element of $\iota_\Ucal\bigl(\R((t^\Gamma))\bigr)$ in the ultraproduct
order.
\end{proposition}
\begin{proof}
Let $h=\sum_{j\le e}c_jZ^j\in\K[Z]$ with $c_e\ne0$. Once $\gamma_n$ exceeds
the finitely many differences $v(c_j)-v(c_e)$, the term $c_e\rho_n^e$ is the
unique one of least valuation, so $v(h(\rho_n))=v(c_e)-e\gamma_n$, which lies
in $H_n$ for large $n$. Its coarsened residue is then nonzero, so
$\Phi_\Ucal(h)\ne0$. Since $\Phi_\Ucal$ restricted to $\K[Z]$ is the
$\K$-algebra map sending $Z$ to $\Phi_\Ucal(Z)$, no nonzero polynomial over
$\iota_\Ucal(\K)$ vanishes there. In the real monomial case, for fixed $c$
and large $n$ the element $\rho_n-c$ has leading exponent $-\gamma_n\in H_n$
with coefficient $1$, so its coarsened residue is positive, and the
ultraproduct inequality gives $\Phi_\Ucal(Z)>\iota_\Ucal(c)$.
\end{proof}

The transcendence has a classical analogue in Henriksen's Theorem~5 for the
ring of complex entire functions \cite{Henriksen}, and it is not claimed as a
new phenomenon. What is specific here is that the coefficient embedding and
the residue fields are computed through the varying coarsenings.

\begin{remark}[Not an illicit infinite substitution]
\label{ent:sm:rem:no-substitution}
The equation $\Phi_\Ucal(P_D)=0$ is consistent with the transcendence of
$\Phi_\Ucal(Z)$. The character is a ring homomorphism and respects finite
algebra; it has not been shown to preserve strong Hahn sums, and
$\Phi_\Ucal(P_D)$ must not be computed by substituting $\Phi_\Ucal(Z)$
termwise into the infinite product or its power series. These maximal ideals
form an algebraic spectrum. No additional zero reachable by strong
evaluation in some field extension is claimed, and the characters are not
asserted to be evaluation maps on any larger surreal domain.
\end{remark}

\paragraph{What this subsection does not do.}
The classification is stated for one node per shell. With several nodes on a
shell the finite residue algebra $\kappa_n[T]/(\overline{P_n})$ can have
several maximal ideals, and a description of the spectrum must track that
varying residue-direction structure together with the ultrafilter; this is
not done. Prime ideals below these maximal ideals are not classified,
especially for unbounded multiplicities. Maximal ideals of $\E$ that contain
no product of this form are not described. The unit-ideal criteria of
\cref{ent:subsec:no-bezout} are not a corona theorem: no bounded function
algebra and no norm lower bound is involved. See \cref{ent:q:spectrum}.

\emph{Re-scoped by the sixth source.} The fourth source's non-claim about
prime ideals is kept above as it was written. For one \emph{simple} node per
shell with empty finite part, the prime ideals below these maximal ideals are
now classified in \cref{ent:ps:sec:spectrum}
(\cref{ent:ps:thm:spectrum}); the classification transfers to uniformly
bounded multiplicities as a statement about prime spaces and dimension, and
for unbounded multiplicities \cref{ent:ps:thm:multiplicity} shows only that
primes not containing the reduced product exist. Shells with several nodes,
unbounded multiplicities, and maximal or prime ideals of $\E$ containing no
product of this form remain unclassified.

\subsection{Real Hahn coefficients}
\begin{proposition}[Real form]\label{ent:prop:real-form}
\Cref{ent:thm:image,ent:thm:quotient,ent:thm:bezout-criterion} and the
interpolation corollaries
\cref{ent:cor:single-node,ent:cor:coarse-separated,ent:cor:order-unit-image},
together with \cref{ent:thm:close-pair}, remain true with $\K$ replaced by the
real Hahn field $\R((t^\Gamma))$ and $\E$ by its ring of strongly entire
series, provided all interpolation nodes and all prescribed data lie in that
field.
\end{proposition}
\begin{proof}
Every algebraic inverse used above concerns polynomials that are split by
hypothesis at the given nodes, not arbitrary polynomials; finite Hermite
interpolation uses only field arithmetic; and the coefficient, support and
coarsening arguments are unchanged. For the kernel statement use the second
proof recorded in \cref{ent:prop:kernel-direct}, which does not invoke
\cref{ent:thm:prep} or \cref{ent:thm:roots}.
\end{proof}

No algebraic closedness of the coefficient field is used anywhere in this
subsection's list. This is a genuine difference from the factorization theory
of \cref{ent:sec:products}: \cref{ent:thm:roots,ent:thm:factorization} do use
algebraic closedness essentially, as recorded in
\cref{ent:subsec:scope-limits}, and \cref{ent:prop:real-form} makes no claim
about them. Under the embedding \eqref{ent:eq:embedding}, the close pair of
\cref{ent:thm:close-pair} in $\Gamma_\infty$ consists of the real surreal
numbers $\rho_n=\omega^{\omega^n}$ and
$\sigma_n=\omega^{\omega^n}+\omega^{-\omega^{n+1}}$ of
\eqref{ent:eq:surreal-roots}. So: no entire power series in that fixed
\emph{real} Hahn workspace takes the value $0$ at every $\rho_n$ and $1$ at
every $\sigma_n$, while changing the second values to $\omega^{-\omega^{n+1}}$
makes the problem solvable. These remain statements about a fixed set-sized
workspace embedded in $\No$ and $\No[i]$; nothing here asserts a nonpolynomial
series summable at every point of the whole proper class, and the
distinction is the one already drawn in \cref{ent:subsec:three-meanings}.

\begin{remark}[The fourth source over the real field]\label{ent:sm:rem:real}
The fourth source proved all of its results for the real and the complex
Hahn field simultaneously. Accordingly the second route of
\cref{ent:sm:subsec:cardinal}, \cref{ent:sm:cor:no-order-unit},
\cref{ent:sm:thm:finite-generators,ent:sm:cor:membership,ent:sm:thm:paired},
\cref{ent:thm:maximal} and all of \cref{ent:sm:subsec:all-maximal} hold over
$\R((t^\Gamma))$, with nodes and data in that field and residue fields
$\R((t^{H_n}))$ and their ultraproducts, which are real closed by
\cref{ent:sm:cor:boundary-fields}. Their proofs use only the ingredients
listed in the proof of \cref{ent:prop:real-form}, Gauss multiplicativity over
$\R[X]$, and, for the type of the residue fields alone, algebraic closedness
of the complexified Hahn fields $\C((t^{H_n}))$. The fourth source also
proves, without root counts, that distinct strongly entire series define
distinct functions, so this part of \cref{ent:cor:units} holds over
$\R((t^\Gamma))$ as well: for $F\ne0$ and a scale $\gamma$, choose a real
$c\ne0$ at which the initial polynomial $I_\gamma(F)$ of
\eqref{ent:eq:initial} does not vanish; then $F(t^\gamma c)$ has valuation
$m_\gamma(F)$ and is nonzero.
\end{remark}


\section{The positive case: an order unit restores interpolation}
\label{ent:sec:positive}

\subsection{A rank-one coarsening with the same entire functions}
It would be incorrect to infer that the preceding obstruction occurs for
every higher-rank group. A group may have several convex levels and still
have a greatest Archimedean class. In that case a single cofinal scale
permits ordinary non-Archimedean analytic methods.

\begin{lemma}[Real-valued coarsening]\label{ent:lem:coarsening}
Suppose $h>0$ is an order unit of $\Gamma$. There is an order-preserving
homomorphism $\pi:\Gamma\to\R$ with $\pi(h)=1$, whose kernel is the convex
subgroup
\[
 H=\{\gamma\in\Gamma:n|\gamma|<h\text{ for all }n\ge1\}.
\]
The formula
\[
 |a|_\pi=\exp(-\pi(v(a)))\quad(a\ne0),\qquad |0|_\pi=0
\]
defines a nontrivial real-valued non-Archimedean absolute value on $\K$. Its
topology agrees with the intrinsic valuation topology, and $\K$ is complete
for it. Furthermore
\begin{equation}\label{ent:eq:entire-coarsening}
 \E=\left\{\sum_{n\ge0}a_nZ^n:
           |a_n|_\pi R^n\longrightarrow0\text{ for every real }R>0
     \right\}.
\end{equation}
\end{lemma}
\begin{proof}
Being an order unit bounds every $\gamma$ above and below by integer
multiples of $h$. Define
\[
 \pi(\gamma)=\sup\{q\in\Q:qh\le\gamma\}\in\R.
\]
Rational upper and lower approximations show
$\pi(\alpha+\beta)=\pi(\alpha)+\pi(\beta)$ and $\pi(h)=1$; monotonicity is
immediate. The kernel consists exactly of the elements infinitesimal relative
to $h$, giving the displayed $H$. In particular $\pi(\gamma)>n$ implies
$\gamma>nh$, and $\gamma>(n+1)h$ implies $\pi(\gamma)>n$. The balls specified
by integer multiples of $h$ are cofinal among all valuation neighbourhoods of
zero, and these inequalities show that the real-valued and ordered-group
topologies agree.

For completeness, take a Cauchy sequence for $|\cdot|_\pi$ and choose a
subsequence $(x_{j_n})$ with $v(x_{j_{n+1}}-x_{j_n})>nh$. The differences
have valuations tending cofinally to infinity, so their Hahn sum exists by
\cref{ent:lem:cofinal-tails}; adding it to $x_{j_1}$ gives a limit of the
subsequence, and the Cauchy property gives the same limit for the whole
sequence. Algebraic closedness of $\K$ was already supplied by the Hahn-field
theorem.

By \cref{ent:thm:criterion}, the left side of
\eqref{ent:eq:entire-coarsening} is characterized by $v(a_n)/n\to+\infty$ in
$\Gamma$. This is equivalent to $\pi(v(a_n))/n\to+\infty$ in $\R$: integer
multiples of $h$ test the forward implication, and for the reverse one
chooses a real bound strictly exceeding $\pi(\gamma)$ for a prescribed
$\gamma\in\Gamma$. The usual elementary coefficient test identifies this last
condition with the right side of \eqref{ent:eq:entire-coarsening}:
explicitly, it implies $\pi(v(a_n))-n\log R\to+\infty$ for every $R$;
conversely, failure of the ratio condition gives infinitely many ratios at
most some $C$, and then $R=e^{C+1}$ violates the coefficient test.
\end{proof}

This is a \emph{coarsening}, not a valuation-preserving embedding of $\Gamma$
in $\R$. The kernel $H$ may be nonzero, and residue-direction data is not
identified across the two valuations: the coarsening changes residue
information even while preserving the topology and the entire ring. It is the
cofinality of the topology and the all-radius coefficient test that permit
its use here.

\subsection{A complete Hermite interpolation argument}
This theorem is the order-unit special case of \cref{ent:thm:image}: an order
unit makes every coarsened ring $V_n$ equal to $\K$ on all large shells, so
the shellwise condition \eqref{ent:eq:maincondition} is vacuous. That is the
first proof below, and it is short. The second proof, which is the one the
source manuscripts wrote, is kept in full: it is a genuinely different
argument, by a real-valued coarsening and Gauss norms rather than by
valuation estimates inside $\Gamma$, it records the damping mechanism that
fails in \cref{ent:subsec:no-bezout} in its own idiom, and it is
self-contained --- in particular it does not depend on \cref{ent:thm:image},
so the two routes to case~(ii) of \cref{ent:thm:main} remain independent.
Neither proof is offered as a correction of the other.

\begin{theorem}[Arbitrary radial Hermite interpolation]
\label{ent:thm:interpolation}
Assume $\Gamma$ has an order unit. Let $(x_n)$ be distinct points of a
radially finite set, and assign a positive integer $m_n$ to each point so
that the divisor $\sum_nm_n[x_n]$ is radially finite. For arbitrary jets
$J_n\in\K[U]/(U^{m_n})$ there exists $F\in\E$ with
\[
 F(x_n+U)\equiv J_n(U)\pmod{U^{m_n}}
 \qquad\text{for every }n.
\]
\end{theorem}
\begin{proof}[First proof, from the exact image theorem]
If the divisor is finite this is ordinary Hermite interpolation over the
field $\K$ by a polynomial. Otherwise \cref{ent:lem:shells} supplies its
shell decomposition and \cref{ent:cor:order-unit-image} says that
\eqref{ent:eq:maincondition} holds automatically, because $V_n=\K$ for all
$n$ with $\gamma_n\ge h$. Apply \cref{ent:thm:image}. Note the direction of
the dependence: this uses the constructive half of \cref{ent:thm:image},
whose damping estimate \eqref{ent:eq:damping-shell} is carried out inside the
convex subgroup $H_n$ and never assumes that the multiples of one $\gamma_n$
are cofinal --- the order unit is used only to make the shellwise condition
empty, not to run the construction.
\end{proof}
\begin{proof}[Second proof, by a real-valued coarsening and Gauss norms]
Use $|\cdot|=|\cdot|_\pi$ from \cref{ent:lem:coarsening}. A radially finite
infinite set can be enumerated with $|x_n|\to\infty$: a ball of any fixed
real absolute-value radius is contained in some original valuation ball, and
conversely an original ball has bounded $|\cdot|_\pi$ radius. Handle the
finitely many nodes with $|x_n|\le1$, including zero, first by a finite
Hermite polynomial, and include their factors in every later vanishing
product. Reindex the remaining nodes by $n\ge1$, so that $|x_n|>1$ and
$|x_n|\to\infty$.

Suppose a polynomial $S_{n-1}$ has the required jets at earlier nodes and put
\[
 Q_{n-1}(Z)=\prod_{j<n}(Z-x_j)^{m_j},
\]
including the finitely many initially handled nodes. It is a unit in the
local jet ring at $x_n$, so there is a unique jet
\[
 B_n(U)=\frac{J_n(U)-S_{n-1}(x_n+U)}{Q_{n-1}(x_n+U)}
              \pmod{U^{m_n}}.
\]
For a positive integer $N$ define the polynomial of degree less than $m_n$
\begin{equation}\label{ent:eq:interpolation-r}
 r_{n,N}(U)=B_n(U)(1+U/x_n)^{-N}\pmod{U^{m_n}}
\end{equation}
and the polynomial correction
\begin{equation}\label{ent:eq:interpolation-p}
 p_{n,N}(Z)=Q_{n-1}(Z)(Z/x_n)^N r_{n,N}(Z-x_n).
\end{equation}
It vanishes to the required orders at all earlier nodes, and by
\eqref{ent:eq:interpolation-r} its jet at $x_n$ is the required residual. The
parameter $N$ can therefore be used purely to control its size.

Set $\theta_n=|x_n|^{1/2}$, so $1<\theta_n<|x_n|$ and $\theta_n\to\infty$.
(The radius is written $\theta_n$, not $R_n$, because $R_n$ denotes the
normalized shell Hermite polynomial of \eqref{ent:eq:normalizedjets}
throughout this report.) For a polynomial
$P=\sum p_kZ^k$ write $\|P\|_R=\max_k|p_k|R^k$. If
$B_n(U)=\sum_{j<m_n}b_jU^j$, the coefficient of $U^\ell$ in
\eqref{ent:eq:interpolation-r} has absolute value at most
$\max_{j\le\ell}|b_j|\,|x_n|^{j-\ell}$. This bound is independent of $N$,
because the binomial coefficients $\binom{-N}{k}$ are integers and so have
absolute value at most $1$. Since $\theta_n<|x_n|$, substitution $U=Z-x_n$
gives
\[
 \|r_{n,N}(Z-x_n)\|_{\theta_n}
 \le\max_{j<m_n}|b_j|\,|x_n|^j=:C'_n,
\]
and consequently
\begin{equation}\label{ent:eq:damping}
 \|p_{n,N}\|_{\theta_n}
 \le \|Q_{n-1}\|_{\theta_n} C'_n
          \left(\frac{\theta_n}{|x_n|}\right)^N.
\end{equation}
The right side tends to zero as the ordinary integer $N$ tends to infinity.
Choose $N_n$ so that this norm is at most $2^{-n}$ and put
$S_n=S_{n-1}+p_{n,N_n}$.

For every fixed real $R>0$, eventually $R\le\theta_n$, so
$\|p_{n,N_n}\|_R\le2^{-n}$. Hence $(S_n)$ is Cauchy in the Gauss norm at
every radius. The space of series satisfying $|a_k|R^k\to0$ is complete for
that norm: coefficientwise limits exist by completeness of $\K$, and uniform
approximation by a finite polynomial bounds the tail of the limit. Thus the
limits for all radii are one common coefficient series $F$ satisfying
\eqref{ent:eq:entire-coarsening}, so $F\in\E$.

Finally, every specified finite jet is preserved. For a fixed node $x$ and
order $m$, choose $R>|x|$. Coefficients of the translated polynomial up to
degree $m-1$ are continuous for $\|\cdot\|_R$, by the finite binomial formula
and its convergent tail. All sufficiently late $S_n$ have the prescribed jet,
so their limit does too.
\end{proof}

\begin{theorem}[B\'ezout in the order-unit case]\label{ent:thm:bezout}
If $\Gamma$ has an order unit then $\E$ is a B\'ezout domain. In particular,
if $f,g\in\E$ have no common zero there are $A,B\in\E$ with $Af+Bg=1$.
\end{theorem}

This is the order-unit special case of \cref{ent:thm:bezout-criterion}: with
an order unit, condition~(2) of that criterion is automatic, so only the
common-zero condition~(1) remains. Both proofs are kept, for the reason given
before \cref{ent:thm:interpolation}.

\begin{proof}[First proof, from the sharp criterion]
Let $f,g$ be nonzero with no common zero and let $D=\Div(f)$, which is
radially finite by \cref{ent:thm:roots} and has $f=a_mP_D$ for a scalar
$a_m\in\K^\times$ by \cref{ent:thm:factorization}. Condition~(1) of
\cref{ent:thm:bezout-criterion} is the absence of a common zero. For
condition~(2), an order unit $h$ makes $H_{\gamma_n}=\Gamma$ for every
$\gamma_n\ge h$, so every nonzero $g(\rho_{n,j})$ has
$\abs{v(g(\rho_{n,j}))}\le c_n\gamma_n$ for some integer $c_n\ge0$; a maximum
over the finitely many nodes of a shell gives one $c_n$. The criterion
supplies $A',B$ with $A'P_D+Bg=1$, hence $(A'/a_m)f+Bg=1$. The reduction of
an arbitrary pair to this case, and the induction on the number of
generators, are as in the second proof below.
\end{proof}
\begin{proof}[Second proof, from the order-unit interpolation theorem]
First suppose $f,g$ are nonzero with no common zero. At each zero $x$ of $f$
let $m_x$ be its multiplicity. Since $g(x)\ne0$, the formal inverse of
$g(x+U)$ exists modulo $U^{m_x}$. By \cref{ent:thm:interpolation} choose
$B\in\E$ whose jet is that inverse at every such $x$. Then $1-Bg$ has a zero
of order at least $m_x$ at each zero of $f$, so it is divisible by $f$ in
$\E$ by \cref{ent:cor:gcd}; write $1-Bg=Af$.

For arbitrary nonzero $f,g$, divide both by their GCD $d$, which exists by
\cref{ent:cor:gcd}. The quotients have no common zero, so their ideal is the
unit ideal; multiplying the resulting identity by $d$ shows $(f,g)=(d)$.
Cases involving zero are immediate. Induct on the finite number of generators
to obtain the B\'ezout property.
\end{proof}

For a radially finite divisor $D$ the full jet map is therefore surjective:
\[
 \E\big/\bigl(P_D\bigr)
 \ \cong\ \prod_{x\in\supp D}\K[U]/(U^{D(x)}).
\]
The kernel is again given by divisibility, and surjectivity is
\cref{ent:thm:interpolation}. Equivalently, in the language of
\cref{ent:subsec:jet-image}, an order unit collapses the restricted product
$\RD$ onto the full product, since $\Lambda_n^\circ=\Lambda_n$ for all large
$n$ (\cref{ent:cor:order-unit-image}). This exact unrestricted Chinese
remainder statement is the boundary case of \cref{ent:thm:quotient}, and it
is exactly what fails, in a now completely measured way, when $\Gamma$ has
countable cofinality and no order unit.

\section{Change of workspace: the universal scalar-extension domain}
\label{ent:sec:extension}

The adjective ``entire'' must name its field of arguments. A formal series
with coefficients in one Hahn field may be entire there and fail even strong
summability at a single point of a larger field. This section gives the exact
boundary, and then the exact description of everything inside it.

Throughout, $\Gamma\subseteq\Delta$ are nonzero divisible ordered abelian
groups with the canonical inclusion $\K\subseteq\KD$, both full Hahn fields
over the \emph{same} coefficient field $\C$.

\subsection{The sharp criterion for remaining entire}
\begin{theorem}[Sharp change-of-workspace criterion]
\label{ent:thm:extension-criterion}
With the above notation,
\begin{equation}\label{ent:eq:extension}
 \mathcal E_\Delta\cap\K[[Z]]=
 \begin{cases}
   \E,&\Gamma\text{ is cofinal in }\Delta,\\
   \K[Z],&\Gamma\text{ is not cofinal in }\Delta.
 \end{cases}
\end{equation}
In the second case a single monomial argument of $\KD$ makes \emph{every}
nonpolynomial series with coefficients in $\K$ fail strong summability.
\end{theorem}
\begin{proof}
If $\Gamma$ is cofinal in $\Delta$, a sequence of elements of
$\Gamma\cup\{+\infty\}$ tends cofinally to infinity in $\Gamma$ if and only
if it does so in $\Delta$. Apply \cref{ent:thm:criterion} to the coefficient
slopes.

If $\Gamma$ is not cofinal, choose $\delta\in\Delta$ strictly larger than
every element of $\Gamma$. For a nonpolynomial $F=\sum a_nZ^n\in\K[[Z]]$
consider evaluation at $Z=t^{-\delta}$. For any two nonzero coefficients with
$m>n$,
\[
 (v(a_m)-m\delta)-(v(a_n)-n\delta)
       =v(a_m)-v(a_n)-(m-n)\delta<0,
\]
since $v(a_m)-v(a_n)\in\Gamma$ and $(m-n)\delta\ge\delta$. The leading
exponents of the nonzero terms therefore form an infinite strictly descending
sequence, so their support union is not well ordered. No nonpolynomial series
over $\K$ can be strongly entire over $\KD$. Polynomials are entire over
either field, proving \eqref{ent:eq:extension}.
\end{proof}

\subsection{The convex-hull valuation ring}
\begin{definition}\label{ent:def:hull}
The convex hull of $\Gamma$ in $\Delta$ is
\begin{equation}\label{ent:eq:convex-hull}
 H=\conv_\Delta(\Gamma)
 =\{s\in\Delta:|s|\le\gamma\text{ for some }\gamma\in\Gamma_{>0}\},
\end{equation}
a divisible convex subgroup of $\Delta$. Coarsening $v_\Delta$ by $H$ gives a
valuation with value group $\Delta/H$.
\end{definition}

\begin{lemma}\label{ent:lem:coarsened-ring}
The valuation ring of that coarsening is
\begin{equation}\label{ent:eq:D}
 \Dom
 =\{0\}\cup\{z\ne0:v_\Delta(z)\ge\gamma
                       \text{ for some }\gamma\in\Gamma\}.
\end{equation}
Its maximal ideal consists of zero together with the elements satisfying
$v_\Delta(z)>\gamma$ for every $\gamma\in\Gamma$. Its residue field is
canonically $\C((t^H))$. Moreover $\Dom=\KD$ exactly when $\Gamma$ is cofinal
in $\Delta$.
\end{lemma}
\begin{proof}
A value $s$ has nonnegative class modulo $H$ exactly when it is bounded below
by some member of $H$, and every member of $H$ is itself bounded below by a
member of $\Gamma$; this proves \eqref{ent:eq:D}. The class is strictly
positive exactly when $s>H$, equivalently $s>\Gamma$.

For an element of this valuation ring, discard all support exponents outside
$H$. No negative coset can occur, and the discarded positive cosets form the
coarsened maximal ideal. The remaining Hahn series has support in $H$. This
is a surjective residue homomorphism onto $\C((t^H))$ with the stated kernel.
Finally the ring is the whole field exactly when $H=\Delta$, which by
\eqref{ent:eq:convex-hull} is equivalent to cofinality of $\Gamma$ in
$\Delta$.
\end{proof}

The residue field here is generally \emph{not} $\C$. It records all scales in
the convex hull $H$, whereas the original fine valuation's residue field
records only the exponent zero.

\subsection{Exact strong evaluation, not merely a sufficient domain}
\begin{theorem}[Universal scalar-extension domain]\label{ent:thm:extension}
Let $f=\sum a_nZ^n\in\E$ be nonpolynomial. In $\KD$ the family $(a_nz^n)_n$
is strongly summable if and only if $z\in\Dom$. The domain is therefore
independent of the choice of nonpolynomial $f$, and the values on it again
lie in $\Dom$.
\end{theorem}
\begin{proof}
Suppose $z\ne0$ lies in the stated ring. Choose $\gamma\in\Gamma$ with
$v_\Delta(z)\ge\gamma$ and put $u=t^{-\gamma}z\in\mathcal O_\Delta$. Since
$f$ is entire over $\K$, \cref{ent:thm:criterion}(2) gives
$f(t^\gamma X)\in\A\subseteq\mathcal A_\Delta$, and
\cref{ent:lem:A-units} gives strong summability of its evaluation family at
the integral point $u$. That family is the original family $(a_nz^n)_n$. With
$\mu=m_\gamma(f)\in\Gamma$, the same lemma gives $v_\Delta(f(z))\ge\mu$
unless the value is zero, so the value lies in $\Dom$. At $z=0$ the assertion
is immediate.

Conversely suppose $s=v_\Delta(z)<\gamma$ for every $\gamma\in\Gamma$. List
the infinitely many nonzero coefficients at increasing indices
$n_1<n_2<\cdots$ and write $\alpha_j=v(a_{n_j})\in\Gamma$. Since $s<0$,
\begin{align*}
 (\alpha_{j+1}+n_{j+1}s)-(\alpha_j+n_js)
 &=\alpha_{j+1}-\alpha_j+(n_{j+1}-n_j)s\\
 &\le\alpha_{j+1}-\alpha_j+s<0,
\end{align*}
the final inequality because $s$ is below the element
$\alpha_j-\alpha_{j+1}$ of $\Gamma$. These are actual leading exponents of
the nonzero summands $a_{n_j}z^{n_j}$, regardless of the higher terms of $z$.
Their strict descent prevents the support union from being well ordered, so
the family is not strongly summable.
\end{proof}

\begin{corollary}[Cofinal extensions are exactly the entire extensions]
\label{ent:cor:entire-extension}
For nonpolynomial $f\in\E$, the same coefficient series belongs to
$\mathcal E_\Delta$ if and only if $\Gamma$ is cofinal in $\Delta$.
Polynomials always survive.
\end{corollary}
\begin{proof}
Combine \cref{ent:thm:extension} with \cref{ent:lem:coarsened-ring};
polynomial evaluations are finite sums.
\end{proof}

\begin{remark}[The negative and the positive statement are compatible]
\label{ent:rem:reconcile}
\Cref{ent:thm:extension-criterion} and \cref{ent:thm:extension} must be read
together, and \cref{ent:rem:conflict} already stated how. The first says that
in the noncofinal case a nonpolynomial $f$ is \emph{not entire over} $\KD$,
so that $\mathcal E_\Delta\cap\K[[Z]]$ collapses to $\K[Z]$; the second says
that the very same series remains strongly summable on the proper subring
$\Dom$, with values in $\Dom$, and \cref{ent:thm:zero-conservation} below
says it keeps every one of its zeros there. There is no disagreement about
whether the function survives: it survives on $\Dom$ and nowhere else, and
``not entire over $\KD$'' is precisely the statement $\Dom\ne\KD$. Indeed the
noncofinal half of \cref{ent:thm:extension-criterion} is the special case
$z=t^{-\delta}$ of the necessity half of \cref{ent:thm:extension}, and its
independent descending-exponent proof is retained because it exhibits the
single witnessing point explicitly.
\end{remark}

A cofinal extension need not preserve rank, and a noncofinal one need not
introduce any algebraic element: the obstruction lies entirely in the ordered
hierarchy of scales. Nor does a strongly defined value in the larger field
have to be a limit there. If $\Delta$ contains a positive value above all of
$\Gamma$, the original cofinal coefficient valuations in $\Gamma$ no longer
tend to infinity in $\Delta$.

\subsection{No new zeros, even for an infinite divisor}
\begin{theorem}[Zero conservation under arbitrary enlargement]
\label{ent:thm:zero-conservation}
Let $0\ne f\in\E$. Every zero of its strongly defined extension on $\Dom$
belongs to $\K$, with unchanged multiplicity. This remains true when $f$ has
infinitely many zeros.
\end{theorem}
\begin{proof}
Suppose $z\in\Dom$ and $f(z)=0$. Choose $\gamma\in\Gamma$ with
$v_\Delta(z)\ge\gamma$. Preparation of $f(t^\gamma X)$ over $\K$
(\cref{ent:thm:prep}) gives
\[
 f(t^\gamma X)=c\,t^{\mu}W_\gamma(X)U_\gamma(X),
\]
with $W_\gamma$ monic over $\OO$ and $U_\gamma$ a unit of $\A$ with initial
polynomial $1$. By \cref{ent:cor:prep-extension} this is also the preparation
over $\Delta$, and by \cref{ent:lem:A-units} the unit factor is nonzero at
every point of $\mathcal O_\Delta$. Therefore $W_\gamma(t^{-\gamma}z)=0$. The
polynomial $W_\gamma$ already splits into linear factors over the
algebraically closed field $\K$, so $t^{-\gamma}z$, and hence $z$, lies in
$\K$.

The same prepared polynomial has the same linear-factor multiplicities in
both fields, and the evaluated unit is nonzero in both, so local orders of
vanishing are unchanged. The argument uses only the finitely many zeros in
one containing ball and applies separately to each zero; it does not require
the global divisor to be finite.
\end{proof}

This is the one statement of the section that imports algebraic closedness.
\Cref{ent:app:audit} records that dependency.

\subsection{No alternative entire continuation}
\begin{proposition}[Identity on constants]\label{ent:prop:identity}
Two formal series over $\KD$ lying in $\mathcal A_\Delta$ and agreeing at
infinitely many distinct points of $\mathcal O_\Delta$ are equal as formal
series. In particular agreement at all ordinary complex constants suffices.
\end{proposition}
\begin{proof}
Their difference $A$ lies in $\mathcal A_\Delta$. If it were nonzero,
multiplication by a nonzero scalar would normalize its Gauss value to zero
and its initial polynomial to monic. By \cref{ent:thm:prep}, the normalized
series is $WU$ with $W$ a nonzero finite-degree polynomial and $U$ a unit.
Evaluation on $\mathcal O_\Delta$ is a ring homomorphism by
\cref{ent:lem:A-units}, so $U$ never vanishes there. Thus $A$ has at most
$\deg W$ distinct integral zeros, a contradiction. This finite upper bound
uses no splitting of $W$ and hence no algebraic-closedness hypothesis.
\end{proof}

\begin{theorem}[Obstruction to any entire-series continuation]
\label{ent:thm:no-continuation}
Suppose $\Gamma$ is not cofinal in $\Delta$ and $f\in\E$ is nonpolynomial.
There is no $g\in\mathcal E_\Delta$ agreeing with $f$ at all points of $\K$.
In fact there is no such $g$ agreeing with $f$ merely at all points of $\C$.
\end{theorem}
\begin{proof}
The coefficient series of $f$ belongs to $\mathcal A_\Delta$, since its
strong support certificate at scale zero over $\Gamma$ remains valid under an
ordered embedding; the same is true of $g$. Agreement on $\C$ would force the
two coefficient series to be identical by \cref{ent:prop:identity}. But
\cref{ent:cor:entire-extension} says that this coefficient series is not
entire over $\KD$, a contradiction.
\end{proof}

This rules out an alternative entire \emph{power-series} continuation. It
does not rule out arbitrary functions extending the pointwise values, or
continuations in a different analytic category, and it asserts no density of
$\C$ in any fine surreal topology.

\subsection{The boundary is intrinsic to the function}
\begin{corollary}[Intrinsic description of the extension boundary]
\label{ent:cor:intrinsic-boundary}
Let $f\in\E$ be nonpolynomial. In any larger ordered group $\Delta$ the
following three subsets generate the same convex subgroup:
\[
 \Gamma,\qquad
 \{v(a_n)/n:n\ge1,\ a_n\ne0\},\qquad
 \{-v(r):r\ne0,\ f(r)=0\}.
\]
Consequently the universal domain of \cref{ent:thm:extension} can be
recovered either from the coefficient slopes alone or from the zero divisor
alone.
\end{corollary}
\begin{proof}
All the displayed values lie in $\Gamma$. The slopes are cofinal in $\Gamma$
by \cref{ent:thm:criterion}. There are infinitely many zeros by
\cref{ent:cor:value-distribution}(a), and their negated valuations are
cofinal by \cref{ent:lem:divisor-countable}. A cofinal subset of an ordered
group generates its full convex hull in any containing ordered group.
\end{proof}

This makes the extension obstruction intrinsic. It is not an artefact of
having embedded the coefficients in an unnecessarily large workspace whose
extra scales the function never detects.

\subsection{Consequences for the B\'ezout counterexample and for \texorpdfstring{$\No[i]$}{No[i]}}
\begin{corollary}[No repair of the counterexample by a Hahn extension]
\label{ent:cor:no-repair}
Let $f,g$ be the pair of \cref{ent:thm:no-bezout}. If $\Gamma$ is cofinal in
$\Delta$ then $f,g\in\mathcal E_\Delta$ and still no
$A,B\in\mathcal E_\Delta$ satisfy $Af+Bg=1$. If $\Gamma$ is not cofinal in
$\Delta$ then neither $f$ nor $g$ is entire on $\KD$, although both remain
strongly defined on $\mathcal D_{\Gamma,\Delta}$ with all their zeros
(\cref{ent:thm:extension,ent:thm:zero-conservation}).
\end{corollary}
\begin{proof}
In the cofinal case $(\gamma_n)$ remains cofinal in $\Delta$ and
$\gamma_{n+1}\ggA\gamma_n$ is unchanged, so the evaluation formula
\eqref{ent:eq:lambda} is unchanged. The proof of
\cref{ent:lem:growth-barrier} now applies to any candidate
$B\in\mathcal E_\Delta$, even with coefficients outside $\K$, giving the same
contradiction. In the noncofinal case apply
\cref{ent:thm:extension-criterion} to the two nonpolynomial series, and
\cref{ent:thm:extension,ent:thm:zero-conservation} for the positive half.
\end{proof}

\begin{corollary}[Whole-class rigidity on the surcomplex numbers]
\label{ent:cor:all-no}
Let $F(Z)=\sum_{n\ge0}a_nZ^n$ be an ordinary power series with coefficients
in $\No[i]$. If the family $(a_nz^n)_n$ is strongly normal-form-summable for
every $z\in\No[i]$, then only finitely many $a_n$ are nonzero.
\end{corollary}
\begin{proof}
The union of the normal-form exponent supports of the real and imaginary
parts of the countably many coefficients is a \emph{set}. Its divisible
additive span together with $1$, with signs chosen to match
$t^\gamma=\omega^{-\gamma}$, is a nonzero set-sized ordered subgroup
$\Gamma\subseteq\No$, and all $a_n$ lie in $\K$ under
\eqref{ent:eq:embedding}; the hypothesis makes $F$ entire over $\K$. Every
set of surreal numbers has a surreal strict upper bound: the set-sized cut
$\delta=\{\,\Gamma\cup\{0\}\mid\varnothing\,\}$ supplies a positive
$\delta\in\No$ larger than every member of $\Gamma$. Enlarge $\Gamma$ to its
divisible span with $\delta$. If $F$ were nonpolynomial, the
descending-leading-exponent argument of
\cref{ent:thm:extension-criterion} --- equivalently
\cref{ent:thm:extension} --- would prohibit strong evaluation at
$z=\omega^{\delta}=t^{-\delta}$, contradicting the hypothesis.
\end{proof}

\begin{remark}[This rigidity is a re-derivation, not a new theorem]
\label{ent:rem:rigidity-credit}
\Cref{ent:cor:all-no} is presented as a consequence of the extension theorem,
not as an independent contribution. The same conclusion is already in this
collection twice over: the \emph{Analysis} report proves all-scale polynomial
rigidity directly, and the \emph{Foundations} report records the exact
boundary of the phenomenon with a pair of series over $\C((t^\Q))$ differing
only in the sign of an exponent. The first two source manuscripts merged
here credit the repository's pre-existing all-scale theme rather than
claiming the corollary, and the fifth does the same for the whole-family half
of its several-variable version, \cref{ent:mv:prop:allscale}. The step that makes the reduction legitimate is the localization
principle recorded in the foundations report --- every \emph{set} of
surcomplex numbers lies in one divisible set-sized workspace --- and the
reduction consumes exactly that principle plus the cut $\delta$.
What is proposed as new here is only the exact cofinal/noncofinal criterion
for a specified pair of set-sized workspaces
(\cref{ent:thm:extension-criterion,ent:thm:extension}) and the persistence of
the B\'ezout obstruction under every cofinal extension
(\cref{ent:cor:no-repair}). See \cref{ent:sec:collection}.
\end{remark}

\section{Several ordinary variables: extension domains, exceptional
directions, rectification and fat points}
\label{ent:sec:multivariable}

This section is the fifth source's contribution, together with the
finite-variable criteria that earlier versions of this report kept in an
appendix, now \cref{ent:app:multivariate}. It answers the scalar-extension
clause of \cref{ent:q:several}, the clause beginning ``In particular'', and
nothing else of that question: preparation, interpolation, divisor and ideal
theory in several variables remain open (\cref{ent:mv:subsec:limits}).
Three warnings govern the whole section.

\begin{enumerate}[label=(\arabic*),leftmargin=2em]
\item \emph{Weaker hypotheses.} The fifth source assumes neither
divisibility of the value group nor anything about the coefficient field
beyond its trivial valuation. Nothing in this section uses divisibility,
algebraic closedness, preparation or root counts, except where a statement
says otherwise. Conversely, nothing proved elsewhere in the report under its
standing hypotheses is transferred here. In particular zero conservation
(\cref{ent:thm:zero-conservation}), which needs preparation and algebraic
closedness, and the inclusion $f(\Dom)\subseteq\Dom$ of
\cref{ent:thm:extension} are not asserted in this setting.
\item \emph{Two summation conventions.} ``Entire'' means summability of the
\emph{whole} monomial family, as in \cref{ent:def:entire}.
\Cref{ent:mv:subsec:blocks} studies a second, explicitly different
convention, homogeneous-block evaluation. It is never called entire, and a
statement about one convention is not a statement about the other.
\item \emph{Notation.} The fifth source's letters are translated by the
fifth manuscript's table in \cref{ent:subsec:notation}. In particular its homogeneous block
$H_n$ is written $F_{(n)}$ here, because $H_n$ is the convex subgroup of
\cref{ent:def:coarse}.
\end{enumerate}

The seven subsections after the first eight,
\crefrange{ent:rc:subsec:setting}{ent:rc:subsec:limits}, are the seventh
source's contribution and state their own hypotheses
(\cref{ent:rc:conv:standing}): divisible $\Gamma$ and $\kk\in\{\R,\C\}$.
Warning (1) does not apply to them, and the three warnings above were written
for the first eight subsections. They re-scope the sentence above: for
configurations of points that an entire automorphism carries into a
coordinate line, the ideal theory and the interpolation of values are settled
there, while preparation, positive-dimensional divisors and GCD or B\'ezout
theory remain open (\cref{ent:mv:subsec:limits}).

The final eleven subsections,
\crefrange{ent:as:subsec:setting}{ent:as:subsec:limits}, are source~11's
(batch~32) and state their own hypotheses (\cref{ent:as:conv:standing}):
$\kk\in\{\R,\C\}$, and $\Gamma$ arbitrary, not divisible, possibly zero.
Warning (1) applies to them with that restriction on $\kk$, and warnings (2)
and (3) apply unchanged. They re-scope the first paragraph of this section
once more: for ideals of $\Ed$ extended from $\kk[X]$, the ideal theory is
settled there by coefficientwise division and faithful flatness
(\cref{ent:as:thm:matrix,ent:as:cor:ideal,ent:as:thm:flat}), including the
vanishing ideals of ordinary point sets and, for $\kk=\C$, of algebraic sets
defined over $\kk$ (\cref{ent:as:cor:vanishing}, the merge's); preparation,
GCD or B\'ezout theory and ideals not extended from $\kk[X]$ remain open.

\subsection{Setting}
\label{ent:mv:subsec:setting}
In this section only, $\kk$ is an arbitrary field with the trivial valuation,
$\Gamma$ is a nonzero set-sized ordered abelian group, not assumed divisible,
of finite rank, or to have an order unit, and $\K=\kk((t^\Gamma))$; for
$\kk=\C$ and divisible $\Gamma$ this is the field of the rest of the report.
The number $d\ge1$ of ordinary variables is finite; for $\alpha\in\N^d$ write
$|\alpha|=\alpha_1+\cdots+\alpha_d$ and $X^\alpha=\prod_jX_j^{\alpha_j}$. For
$F=\sum_\alpha a_\alpha X^\alpha\in\K[[X_1,\ldots,X_d]]$ the \emph{ordinary
monomial support} is $\Sigma(F)=\{\alpha:a_\alpha\ne0\}$; each nonzero
$a_\alpha$ also has its own Hahn exponent support in $\Gamma$, and the two
kinds of support play different roles.

\begin{definition}[Whole-family entireness]\label{ent:mv:def:entire}
$\mathcal E_{\Gamma,d}$ is the set of $F\in\K[[X_1,\ldots,X_d]]$ for which the
\emph{whole} family $(a_\alpha x^\alpha)_{\alpha\in\N^d}$ is strongly
summable for every $x\in\K^d$. Summing first along one coordinate and then
another, or first within each total degree, is not the definition. For
$x_j=0$ the convention $x_j^0=1$ applies, so only the face $\alpha_j=0$ of
the support survives.
\end{definition}

Fix an ordered-group inclusion $\Gamma\subseteq\Delta$ and the canonical
inclusion $\K\subseteq\KD=\kk((t^\Delta))$. The convex hull
$H=\conv_\Delta(\Gamma)$ is defined by \eqref{ent:eq:convex-hull}; it is a
convex subgroup, with no divisibility needed, and $\Gamma$ is cofinal in it.
Write $\pi_H:\Delta\to\Delta/H$ for the quotient map of ordered groups and
\[
 \Dom=\{0\}\cup\{z\in\KD^\times:\pi_H(v(z))\ge0\},
\]
the coarsened valuation ring of \eqref{ent:eq:D}; it equals $\KD$ exactly
when $\Gamma$ is cofinal in $\Delta$. No splitting of $H\subseteq\Delta$ is
assumed. A convenient model of a new scale is $\Gamma=\Q$ and
$\Delta=\R e_*\oplus\Q$, ordered lexicographically with the $e_*$ coordinate
dominant: then $H=\Q$, $\Delta/H\cong\R$, and $t^{-e_*}$ is larger in
magnitude than every old monomial $t^{-\gamma}$.

The support calculus of \cref{ent:lem:neumann} is stated for arbitrary
ordered abelian groups and is used unchanged. Two further elementary facts
are used, for families indexed by arbitrary sets.

\begin{lemma}[Tail certificate and regrouping]\label{ent:mv:lem:tails}
Let $(f_i)_{i\in I}$ be a family of Hahn series over an ordered abelian
group $G$.
\begin{enumerate}[label=\textnormal{(\alph*)},leftmargin=2em]
\item If for each $b\in G$ only finitely many nonzero $f_i$ satisfy
$v(f_i)\le b$, the family is strongly summable.
\item A strongly summable family may be partitioned into arbitrary
subfamilies, summed within each part and then summed over the parts, with the
original sum as result. The product of two strongly summable families is
jointly strongly summable on the product index set, and its sum is the
product of the two sums.
\end{enumerate}
\end{lemma}
\begin{proof}
(a) Only the finitely many terms of valuation at most $b$ can contain the
exponent $b$, which gives finite incidence. A descending sequence
$b_1>b_2>\cdots$ in the support union would lie in the finitely many supports
of terms of valuation at most $b_1$, and a finite union of well-ordered sets
is well ordered by \cref{ent:lem:neumann}(a). (b) Regrouped supports lie in
the original well-ordered union, and at each exponent only finitely many
original indices, hence finitely many parts, contribute. Product supports lie
in the sumset of the two support unions, which is well ordered with finitely
many decompositions of each exponent by \cref{ent:lem:neumann}(a), and each
decomposition has finitely many contributing indices.
\end{proof}

The converse of (b) is false, and \cref{ent:mv:ex:cancel} exhibits the
failure that matters here.

\subsection{Finite-variable coefficient and unit criteria}
\label{ent:app:multivariate}
The two numbered statements of this subsection were Appendix~B of earlier
versions of this report, proved there under the report's standing
hypotheses, with the coefficient criterion in ratio form; their labels are
kept. The fifth source reproved the coefficient criterion in a
division-free form under the weaker hypotheses of this section, and that
form is printed here as the statement. The proof of the earlier appendix was
the same descending-leading-exponent argument and is not printed twice.

\begin{theorem}[Finite-variable coefficient criterion]
\label{ent:thm:multi-criterion}
For $F=\sum_\alpha a_\alpha X^\alpha\in\K[[X_1,\ldots,X_d]]$ the following are
equivalent:
\begin{enumerate}[label=\textnormal{(\roman*)},leftmargin=2em]
\item $F\in\mathcal E_{\Gamma,d}$;
\item the family is strongly summable at every diagonal monomial point
$(t^\gamma,\ldots,t^\gamma)$, $\gamma\in\Gamma$;
\item for all $\gamma,b\in\Gamma$ the set
\begin{equation}\label{ent:mv:eq:growth}
 \{\alpha\in\Sigma(F):v(a_\alpha)+|\alpha|\gamma\le b\}
 \quad\text{is finite.}
\end{equation}
\end{enumerate}
If $\Gamma$ is divisible, they are also equivalent to
$v(a_\alpha)/|\alpha|\to+\infty$ cofinally as $|\alpha|\to\infty$ through
nonzero coefficients of positive degree. Nonpolynomial members of
$\mathcal E_{\Gamma,d}$ exist exactly when $\cf(\Gamma)=\aleph_0$.
\end{theorem}
\begin{proof}
(i)$\Rightarrow$(ii) is immediate. Assume (ii), fix $\gamma$, and put
$b_\alpha=v(a_\alpha)+|\alpha|\gamma$. Strong summability at the diagonal
point gives a common lower bound $L\in\Gamma$ for these leading exponents.
Suppose infinitely many $\alpha\in\Sigma(F)$ satisfy $b_\alpha\le b$. Only
finitely many multi-indices have bounded total degree, so there are such
$\alpha_1,\alpha_2,\ldots$ with strictly increasing degrees $n_r=|\alpha_r|$.
Choose $\delta>0$ in $\Gamma$ with $\delta>b-L$. Then
\[
 (b_{\alpha_{r+1}}-n_{r+1}\delta)-(b_{\alpha_r}-n_r\delta)\ \le\ b-L-\delta<0,
\]
and these are actual leading exponents at the diagonal point of valuation
$\gamma-\delta$. Their strict descent contradicts (ii); so (iii) holds.

For (iii)$\Rightarrow$(i), take a point with nonzero coordinates and let
$\gamma=\min_jv(x_j)$. Then $v(a_\alpha x^\alpha)\ge v(a_\alpha)+|\alpha|\gamma$,
and \eqref{ent:mv:eq:growth} with \cref{ent:mv:lem:tails}(a) gives strong
summability. Zero coordinates discard terms and cause no further difficulty.

If $\Gamma$ is divisible, the ratio condition implies \eqref{ent:mv:eq:growth}
by testing a ratio bound above $-\gamma$ and a positive element dominating
$|b|$; conversely, for $c\in\Gamma$, \eqref{ent:mv:eq:growth} with
$\gamma=-c$ and $b=0$ gives $v(a_\alpha)>|\alpha|c$ for all but finitely many
$\alpha$. Finally, a nonpolynomial $F$ has infinitely many nonzero
coefficients, and \eqref{ent:mv:eq:growth} at $\gamma=0$ makes their
valuations a countable cofinal subset of $\Gamma$; a nonzero ordered group
has no finite cofinal subset. The converse is
\cref{ent:mv:prop:realize-support}.
\end{proof}

\begin{proposition}[Every ordinary support occurs]
\label{ent:mv:prop:realize-support}
Suppose $\cf(\Gamma)=\aleph_0$. There is a sequence $(\varphi_n)_{n\ge0}$ in
$\Gamma$ such that for every $\Sigma\subseteq\N^d$ and every choice of
$c_\alpha\in\kk^\times$, $\alpha\in\Sigma$, the series
\[
 F_\Sigma(X)=\sum_{\alpha\in\Sigma}c_\alpha t^{\varphi_{|\alpha|}}X^\alpha
\]
is entire and has ordinary support exactly $\Sigma$.
\end{proposition}
\begin{proof}
Choose a nondecreasing positive cofinal sequence $u_n\in\Gamma$ and put
$\varphi_0=0$, $\varphi_n=n^2u_n$ for $n\ge1$. For fixed $\gamma,b$ we
eventually have $u_n>|\gamma|+|b|$, and then for $n\ge2$
\[
 \varphi_n+n\gamma\ \ge\ n^2u_n-n|\gamma|\ >\ n(n-1)u_n\ >\ b .
\]
Nonzero elements of $\kk$ have valuation zero, whatever their size in any
other sense, so \eqref{ent:mv:eq:growth} holds.
\end{proof}

Over $\Q$ the choice $\varphi_n=n^2$ works. The arbitrary-rank construction
matters because $n^2\eta$ need not be cofinal when $\eta$ is not an order
unit. For $\Gamma=0$ strong summability at $(1,\ldots,1)$ already forces a
finite ordinary support, and every extension question below is trivial.

For the unit criterion put
\[
 \mathcal T_{\Gamma,d}
 =\Bigl\{\sum_\alpha a_\alpha X^\alpha:
              v(a_\alpha)\to+\infty\text{ as }|\alpha|\to\infty\Bigr\}.
\]
The following theorem is stated, as in the earlier appendix, under the
report's standing hypotheses, $\Gamma$ divisible and $\kk=\C$; the fifth
source does not treat it.

\begin{theorem}\label{ent:thm:multi-unit}
A nonconstant $F\in\mathcal T_{\Gamma,d}$ is a unit of that algebra if and
only if $a_0\ne0$ and
$\min_{\alpha\ne0,\ a_\alpha\ne0}v(a_\alpha/a_0)$ is a positive order unit of
$\Gamma$.
\end{theorem}
\begin{proof}
Regrouping into $\C[X_1,\ldots,X_d]((t^\Gamma))$ gives multiplicativity of
the Gauss valuation. If $FG=1$, normalize both constant coefficients to $1$;
their Gauss values are nonpositive and sum to zero, so their initial
polynomials multiply to $1$ and are constant, and all nonconstant coefficient
values are positive. If their minimum $\delta$ is not an order unit, coarsen
by the proper convex subgroup \eqref{ent:eq:Hdelta} it generates:
restrictedness reduces both series to finite polynomials over the coarsened
residue field, and at least one nonconstant coefficient of $F$ survives,
contradicting invertibility of its reduction. Conversely an order-unit gap
makes the geometric inverse restricted, with precisely the tail estimates of
\cref{ent:thm:units}; finite-variable convolution has finitely many terms at
each multi-index, so those estimates are unchanged.
\end{proof}

Finite-dimensionality is used in the finiteness of bounded-degree index sets.
These statements assert nothing for infinitely many variables, and they do
not extend the one-variable divisor classification to higher-dimensional zero
sets. The single-variable uniform extension domain of \cref{ent:thm:extension}
is \emph{false} in several variables, since a series independent of one
coordinate imposes no constraint in that coordinate; its exact replacement is
\cref{ent:mv:thm:domain}.

\subsection{The exact whole-family extension domain}
\label{ent:mv:subsec:domain}

\begin{lemma}[Uniform unit twisting]\label{ent:mv:lem:twist}
Let $(b_i)_{i\in I}$ be a family in $\KD$, let $\alpha(i)\in\N^d$, and let
$u_1,\ldots,u_d\in1+\mathfrak m_\Delta$, where
$\mathfrak m_\Delta=\{x\in\KD:v(x)>0\}$. Then $(b_i)_i$ is strongly summable
if and only if $\bigl(b_i\prod_ju_j^{\alpha(i)_j}\bigr)_i$ is. Multiplying
each term by a nonzero element of $\kk$ also preserves and reflects strong
summability.
\end{lemma}
\begin{proof}
Let $B$ be the support union of the first family and put
$S=\bigcup_j\supp(u_j-1)\subseteq\Delta_{>0}$, a finite union of
well-ordered sets. Every finite product $\prod_ju_j^{\alpha(i)_j}$ has
support in the one monoid $S^*$, regardless of $i$, so the twisted support
union lies in the well-ordered set $B+S^*$ (\cref{ent:lem:neumann}). A fixed
exponent has finitely many decompositions $b+s$ with $b\in B$, $s\in S^*$,
and each $b$ lies in finitely many of the original supports; this gives
finite incidence. For the converse apply the same argument with
$u_j^{-1}\in1+\mathfrak m_\Delta$. Nonzero scalars from $\kk$ do not change
supports.
\end{proof}

Every nonzero $z_j\in\KD$ has a unique leading-term decomposition
$z_j=c_jt^{\delta_j}u_j$ with $c_j\in\kk^\times$, $\delta_j=v(z_j)$ and
$u_j\in1+\mathfrak m_\Delta$. Consequently $(a_\alpha z^\alpha)_\alpha$ is
strongly summable if and only if $(a_\alpha t^{\inner\alpha\delta})_\alpha$
is, for arbitrary coefficient families. This removes the tails of the
arguments, not those of the coefficients.

\begin{lemma}[The finite lattice estimate]\label{ent:mv:lem:lattice}
Let $\delta_1,\ldots,\delta_d\in\Delta$, put $\bar\delta_j=\pi_H(\delta_j)$,
and let $\mathcal L=\{m\in\Z^d:\sum_jm_j\bar\delta_j=0\}$. There are a
positive integer $c_{\mathcal L}$ and $\eta\in\Gamma_{>0}$ with
\begin{equation}\label{ent:mv:eq:lattice}
 \Bigl|\sum_jm_j\delta_j\Bigr|\ \le\ c_{\mathcal L}\,\|m\|_1\,\eta
 \qquad(m\in\mathcal L),\qquad \|m\|_1=\sum_j|m_j|.
\end{equation}
\end{lemma}
\begin{proof}
A subgroup of $\Z^d$ is free of finite rank, by induction on $d$ through the
projection to the last coordinate. Choose a basis $\ell_1,\ldots,\ell_r$ of
$\mathcal L$. Each $h_s=\sum_j(\ell_s)_j\delta_j$ lies in $H$, so one
$\eta\in\Gamma_{>0}$ bounds all $|h_s|$. The matrix with columns $\ell_s$ has
a rational left inverse, and clearing its denominators shows that the
integer coordinates of $m=\sum_sk_s\ell_s$ satisfy
$\sum_s|k_s|\le c_{\mathcal L}\|m\|_1$ for a positive integer
$c_{\mathcal L}$ independent of $m$. Only this numerical estimate uses
rational linear algebra; no division in $\Gamma$ or $\Delta$ occurs. Hence
$|\sum_jm_j\delta_j|=|\sum_sk_sh_s|\le c_{\mathcal L}\|m\|_1\eta$. If
$\mathcal L=0$, take $c_{\mathcal L}=1$ and any $\eta$.
\end{proof}

\begin{theorem}[Exact multivariate extension domain]\label{ent:mv:thm:domain}
Let $F=\sum_\alpha a_\alpha X^\alpha\in\mathcal E_{\Gamma,d}$ and
$z\in(\KD^\times)^d$. Put $\bar\delta_j=\pi_H(v(z_j))$ and
\[
 \inner{\Sigma(F)}{\bar\delta}=\{\inner\alpha{\bar\delta}:\alpha\in\Sigma(F)\}
 \subseteq\Delta/H .
\]
Then
\begin{equation}\label{ent:mv:eq:main}
 (a_\alpha z^\alpha)_\alpha\text{ is strongly summable}
 \quad\Longleftrightarrow\quad
 \inner{\Sigma(F)}{\bar\delta}\text{ is well ordered.}
\end{equation}
Here $\inner{\Sigma(F)}{\bar\delta}$ is a set of distinct weights, not an
indexed family: infinitely many multi-indices may have the same weight. If
some coordinates of $z$ vanish, restrict to the face of $\Sigma(F)$ on which
those coordinates have exponent zero, delete them, and apply the same
criterion; if no variable remains, evaluation is the constant coefficient.
\end{theorem}
\begin{proof}
Write $\delta_j=v(z_j)$. By \cref{ent:mv:lem:twist} we may replace $z_j$ by
$t^{\delta_j}$, so the $\alpha$th term becomes
$b_\alpha=a_\alpha t^{\inner\alpha\delta}$, whose exponents all lie in
$\inner\alpha\delta+\Gamma$ and map under $\pi_H$ to $\inner\alpha{\bar\delta}$.

\emph{Necessity.} A strongly summable family has a well-ordered support
union, and its image under the order-preserving map $\pi_H$ is well ordered,
because a strictly descending sequence of images lifts to a strictly
descending sequence of chosen preimages. That image is
$\inner{\Sigma(F)}{\bar\delta}$, since every term with $\alpha\in\Sigma(F)$ is
nonzero.

\emph{Sufficiency within one fibre.} Assume the weight set is well ordered,
fix a weight $w$ in it and some $\alpha^0\in\Sigma(F)$ of weight $w$. For every
$\alpha$ of the same weight, $m=\alpha-\alpha^0$ lies in the lattice
$\mathcal L$ of \cref{ent:mv:lem:lattice}. After multiplying the whole fibre
by $t^{-\inner{\alpha^0}\delta}$, its terms have all their exponents in $H$
and valuations at least
\begin{equation}\label{ent:mv:eq:fiberbound}
 v(a_\alpha)+\inner{\alpha-\alpha^0}\delta
 \ \ge\ \bigl(v(a_\alpha)-c_{\mathcal L}|\alpha|\eta\bigr)
        -c_{\mathcal L}|\alpha^0|\eta .
\end{equation}
Condition \eqref{ent:mv:eq:growth} with $\gamma=-c_{\mathcal L}\eta$ says
that the bracket exceeds any given bound in $\Gamma$ except at finitely many
indices, and since $\Gamma$ is cofinal in $H$ the same holds for bounds in
$H$: for $b\in H$ choose $B\in\Gamma$ with $B\ge b+c_{\mathcal L}|\alpha^0|\eta$.
\Cref{ent:mv:lem:tails}(a), applied in the ordered group $H$, makes the fibre
strongly summable.

\emph{Assembly of the fibres.} Different fibres lie in different cosets of
$H$. A nonempty subset of the full support union has a least coset because
the weight set is well ordered, and within that coset a least exponent by the
fibre result; that exponent is least in the whole subset. Each exponent
belongs to one fibre, where finite incidence is already known. Finally, a
zero coordinate removes every term with a positive exponent in it, the
surviving coefficients still satisfy \eqref{ent:mv:eq:growth}, and the same
proof applies to the remaining variables.
\end{proof}

Finiteness of $d$ is used twice, for the finiteness of bounded-degree index
sets and for the finite basis of $\mathcal L$. Entireness is used in
\eqref{ent:mv:eq:fiberbound}, not merely to obtain a lower bound at one
point. Full Hahn fields license the support-controlled unit inverses and
sums, and convexity licenses the ordered quotient. No root extraction,
algebraic closedness, spherical completeness, real-valued norm or
divisibility enters.

\begin{corollary}[Independence of tails within the entire category]
\label{ent:mv:cor:tails}
Two members of $\mathcal E_{\Gamma,d}$ with the same ordinary support have the
same whole-family strong evaluation domain in every $\KD^d$, and for either
of them the domain depends only on the coarsened coordinate valuations, with
zero coordinates treated by faces as above.
\end{corollary}
\begin{proof}
The right side of \eqref{ent:mv:eq:main} and its face variants contain no
other data.
\end{proof}

This corrects the guess recorded in \cref{ent:q:several}. Within the entire
category, higher Hahn tails do \emph{not} create an additional boundary
obstruction; without entireness they can, by \cref{ent:mv:ex:hidden} below,
which is a rank-one analogue of the pair of \cref{ent:subsec:hidden-tails}.

\subsection{Consequences: one variable, cones and semilinear supports}
\label{ent:mv:subsec:cones}

\begin{corollary}[The one-variable domain, recovered]\label{ent:mv:cor:onevar}
For a nonpolynomial $f\in\mathcal E_{\Gamma,1}$ the strong evaluation domain in
$\KD$ is exactly $\Dom$; a polynomial is strongly evaluable everywhere. Hence
a cofinal enlargement preserves every entire series in any finite number of
variables.
\end{corollary}
\begin{proof}
For an infinite $\Sigma\subseteq\N$ the set $\{n\bar\delta:n\in\Sigma\}$ is a
singleton for $\bar\delta=0$, lies in a well-ordered positive cyclic monoid
for $\bar\delta>0$, and has a strictly descending sequence for $\bar\delta<0$.
Apply \cref{ent:mv:thm:domain}. A cofinal enlargement has $H=\Delta$, so every
weight set is empty or a singleton.
\end{proof}

This is the domain half of \cref{ent:thm:extension}, now without
divisibility and over any trivially valued coefficient field. The value
inclusion of \cref{ent:thm:extension}, zero conservation and the continuation
theorems of \cref{ent:sec:extension} are not part of it.

\begin{definition}[Admissible weights]\label{ent:mv:def:adm}
For $\Sigma\subseteq\N^d$ and an ordered abelian group $G$ put
$\Adm_\Sigma(G)=\{\bar\delta\in G^d:\inner\Sigma{\bar\delta}\text{ is well
ordered}\}$. It depends on the whole support, not on an enumeration.
\end{definition}

\begin{proposition}[Cone and monotonicity properties]\label{ent:mv:prop:cone}
$\Adm_\Sigma(G)$ contains $G_{\ge0}^d$, is closed under addition and under
multiplication by nonnegative integers, and is upward closed in the
coordinatewise order. If $G$ is divisible it is closed under nonnegative
rational scaling, and for $G=\R$ it is a convex cone, not necessarily closed.
If $\Sigma_1\subseteq\Sigma_2$ then $\Adm_{\Sigma_2}(G)\subseteq
\Adm_{\Sigma_1}(G)$, and a finite change of $\Sigma$ does not change
$\Adm_\Sigma(G)$.
\end{proposition}
\begin{proof}
For nonnegative weights every element of $\inner\Sigma{\bar\delta}$ lies in the
monoid generated by the positive $\bar\delta_j$, which is well ordered by
\cref{ent:lem:neumann}(b). Addition uses
$\inner\Sigma{\bar\delta+\bar\epsilon}\subseteq\inner\Sigma{\bar\delta}+
\inner\Sigma{\bar\epsilon}$ and \cref{ent:lem:neumann}(a). Positive integer,
rational or real scaling is an order isomorphism of the weight set where it
is defined, and zero scaling gives at most a singleton. Upward closure is
addition of a nonnegative vector. The remaining assertions are immediate,
since adding finitely many weights preserves well-ordering.
\end{proof}

On the torus $(\KD^\times)^d$ the strong domain of a fixed entire series is
therefore a multiplicative monoid, stable under coordinatewise
multiplication by nonzero elements of $\Dom$. Its values need not lie in
$\Dom$: a finite monomial prefactor can have negative quotient weight.

\begin{proposition}[Evaluation algebras at a fixed point]
\label{ent:mv:prop:algebra}
Fix $z\in\KD^d$. The $F\in\mathcal E_{\Gamma,d}$ whose whole-family evaluation
at $z$ exists form a $\K$-subalgebra, and evaluation is a $\K$-algebra
homomorphism from it to $\KD$.
\end{proposition}
\begin{proof}
By \cref{ent:mv:lem:tails}(b), the sum and the double product of two
evaluation families at $z$ are strongly summable. Regroup the product by the
multi-index sum; each formal coefficient has finitely many decompositions in
$\N^d$.
\end{proof}

\begin{definition}[Semilinear support]\label{ent:mv:def:semilinear}
$\Sigma\subseteq\N^d$ is semilinear if
$\Sigma=\bigcup_{\ell=1}^s\bigl(b_\ell+\N p_{\ell1}+\cdots+\N p_{\ell
r_\ell}\bigr)$ with $b_\ell,p_{\ell j}\in\N^d$; empty unions and components
without periods are allowed.
\end{definition}

\begin{theorem}[The period-inequality criterion]\label{ent:mv:thm:semilinear}
For a support presented as in \cref{ent:mv:def:semilinear},
$\bar\delta\in\Adm_\Sigma(G)$ if and only if $\inner{p_{\ell j}}{\bar\delta}
\ge0$ for every period vector. The offsets impose no inequalities, so the
extension domain of an entire series with semilinear support is decided by
finitely many ordered-group inequalities.
\end{theorem}
\begin{proof}
A period of negative weight gives the strictly descending weights of
$b_\ell+np_{\ell j}$. If every period weight is nonnegative, each component
has as its weight set a translate of a monoid generated by finitely many
nonnegative elements, well ordered by \cref{ent:lem:neumann}(b), and a finite
union of translates is well ordered.
\end{proof}

The inequalities are weak. An infinite support ray of quotient weight zero is
not an obstruction, because old coefficient growth controls the whole
collapsed fibre; replacing $\ge0$ by $>0$ would give a false boundary.

\begin{example}[A coupled domain]\label{ent:mv:ex:product}
Over $\C((t^\Q))$ the series $F_\times(X,Y)=\sum_{n\ge0}t^{n^2}(XY)^n$ is
entire, has the single period $(1,1)$, and at a point with nonzero
coordinates in any extension is strongly evaluable exactly when
$\bar\delta_X+\bar\delta_Y\ge0$. In particular $X=t^{-e_*}$, $Y=t^{e_*}$ is
admissible although $X\notin\Dom$: every quotient weight is zero and the
family is $(t^{n^2})_n$. The domain is not a coordinatewise copy of the
one-variable ring. By contrast
$G_+(X,Y)=1+\sum_{n\ge1}t^{n^2}X^n+\sum_{n\ge1}t^{n^2}Y^n$ has torus domain
$\bar\delta_X\ge0$, $\bar\delta_Y\ge0$, so two nonpolynomial series involving
both coordinates can have different domain geometries.
\end{example}

\begin{theorem}[A nonclosed domain in two variables]\label{ent:mv:thm:nonclosed}
Let $\Gamma=\Q$, $\Delta=\R e_*\oplus\Q$ as above, and
$F_{\mathrm{par}}(X,Y)=\sum_{n\ge1}t^{n^4}X^nY^{n^2}$. This series is entire
over $\C((t^\Q))$, and with $\Delta/H$ identified with $\R$,
\[
 \Adm_{\Sigma(F_{\mathrm{par}})}(\R)=\{(u,w):w>0\}\cup\{(u,0):u\ge0\}.
\]
This set is convex but not closed, and it is not a finite intersection of
weak linear half-spaces.
\end{theorem}
\begin{proof}
The total degree is $n+n^2$, and $n^4+(n+n^2)\gamma\to+\infty$ for every
rational $\gamma$. The weight set is $\{nu+n^2w:n\ge1\}$. For $w>0$ the
successive differences $u+(2n+1)w$ are eventually positive, so the set is
well ordered; for $w<0$ they are eventually negative; for $w=0$ the set is
well ordered exactly when $u\ge0$. Apply \cref{ent:mv:thm:domain}. The
admissible points $(-1,1/m)$ tend in the real plane to $(-1,0)$, which is not
admissible, and a finite intersection of weak half-spaces is closed.
\end{proof}

This concerns the Euclidean topology of the real weight parameter, not any
surreal fine topology.

\begin{example}[Bounded positive weights are not enough]\label{ent:mv:ex:dense}
Fix an irrational real $\theta>0$, put $\ell_n=\lceil n\theta\rceil$, and let
$F_\theta(X,Y)=\sum_{n\ge1}t^{(n+\ell_n)^2}X^nY^{\ell_n}$ over $\C((t^\Q))$,
entire by the degree-squared damping. At the quotient weights $(-\theta,1)$
its weights are $\ell_n-n\theta\in(0,1)$, never $0$, with infimum $0$: for any
$N$, two of the fractional parts of $0,\theta,\ldots,N\theta$ lie within
$1/N$, so some positive integer $q$ has $q\theta$ within $1/N$ of an integer;
if its fractional part exceeds $1-1/N$ this is a small positive gap, and if
it is $\epsilon\in(0,1/N)$ then, with $k=\lfloor1/\epsilon\rfloor$,
irrationality gives $0<1-k\epsilon<\epsilon$ and $kq$ has a gap below $1/N$.
So the weights have no least element and the point is not in the strong
domain, although every weight is positive. Well-ordering, not positivity or
boundedness, is the criterion.
\end{example}

\begin{example}[What fails without entireness]\label{ent:mv:ex:hidden}
In $\C((t^\Q))[[X]]$ let $\tilde a_n=t^{1-1/(n+1)}$ and
$\tilde b_n=t^{1-1/(n+1)}+t^{2+1/(n+1)}$. The series $A_\flat=\sum\tilde
a_nX^n$ and $B_\flat=\sum\tilde b_nX^n$ have the same ordinary support and the
same coefficient valuations. At $X=1$ the family of $A_\flat$ is strongly
summable, its exponents strictly increasing and each occurring once, while
that of $B_\flat$ is not, its second exponents strictly decreasing. Neither
series is entire: at $X=t^{-1}$ even the leading exponents strictly
decrease. So \cref{ent:mv:cor:tails} is not a general principle that leading
data decide Hahn summability; entireness is doing real work.
\end{example}

\subsection{Homogeneous blocks: a second, explicitly different convention}
\label{ent:mv:subsec:blocks}
For a formal series $F$ write $F_{(n)}(X)=\sum_{|\alpha|=n}a_\alpha X^\alpha$,
a finite homogeneous polynomial over $\K$, so that $F=\sum_nF_{(n)}$.

\begin{definition}[Homogeneous-block evaluation]\label{ent:mv:def:block}
At $z\in\KD^d$ the homogeneous-block evaluation of $F$ exists if
$(F_{(n)}(z))_{n\ge0}$ is strongly summable; its value is written
$F^{\mathrm{hom}}(z)$ and its domain $\mathcal D^{\mathrm{hom}}(F)$.
\end{definition}

Whole-family evaluation implies homogeneous-block evaluation with the same
value, by \cref{ent:mv:lem:tails}(b). The reverse implication is not part of
the definition and is false. Block evaluation is never called entire in this
report.

\begin{proposition}[Linear covariance]\label{ent:mv:prop:linear}
For a matrix $M$ over $\K$ defining the substitution $X=MY$, an entire $F$
gives an entire formal series $F\circ M$ whose degree-$n$ block is
$F_{(n)}\circ M$. Consequently $y\in\mathcal D^{\mathrm{hom}}(F\circ M)$ if and
only if $My\in\mathcal D^{\mathrm{hom}}(F)$, with equal values; this is a
covariance law under $\mathrm{GL}_d(\K)$, and the same argument applies to
non-square $M$, in particular to line restrictions.
\end{proposition}
\begin{proof}
For $M\ne0$ choose $\mu\in\Gamma$ not exceeding $0$ or the valuation of any
nonzero entry. Every new degree-$n$ coefficient is a finite sum of old
degree-$n$ coefficients times products of $n$ entries and integers, so its
valuation is at least $b_n+n\mu$, where $b_n$ is the least valuation of the
nonzero coefficients of $F_{(n)}$; cancellation only raises it. Condition
\eqref{ent:mv:eq:growth} at $\gamma+\mu$ gives the same condition for the new
coefficients at $\gamma$. By homogeneity the degree-$n$ block of $F\circ M$
at $y$ is exactly $F_{(n)}(My)$, so the two block families coincide.
\end{proof}

\begin{example}[Cancellation changes the domain]\label{ent:mv:ex:cancel}
Over $\C((t^\Q))$ let $F_{\mathrm{can}}(X,Y)=\sum_{n\ge0}t^{n^2}(X-Y)^n$. At
$(t^{-e_*},t^{-e_*})$, with $e_*$ above all old scales, every block of
positive degree is zero, so the block evaluation exists and equals $1$. The
whole monomial family is not strongly summable: for every $n$ it contains the
term $t^{n^2}X^n$ of quotient weight $-n\pi_H(e_*)$, and these weights
strictly decrease. In the coordinates $U=X-Y$, $V=Y$ the same series is
$\sum t^{n^2}U^n$, whose whole family is summable at $(0,t^{-e_*})$. So
whole-family domains depend on the chosen coordinate expansion, while block
domains obey \cref{ent:mv:prop:linear}; the value $1$ is a value of the block
convention, not a sum assigned to the nonsummable expanded family.
\end{example}

\begin{theorem}[Polynomial rays and new scales]\label{ent:mv:thm:rays}
Let $F\in\mathcal E_{\Gamma,d}$ and $c\in\K^d$, and let
$F_c(T)=F(Tc)=\sum_nF_{(n)}(c)T^n$, an entire one-variable series over $\K$.
For every ordered-group enlargement,
\[
 \{\lambda\in\KD:\lambda c\in\mathcal D^{\mathrm{hom}}(F)\}
 =\begin{cases}\KD,&F_c\text{ is a polynomial},\\
               \Dom,&F_c\text{ is not a polynomial}.\end{cases}
\]
In particular, if $e_*>H$ then $t^{-e_*}c\in\mathcal D^{\mathrm{hom}}(F)$ if
and only if $F_{(n)}(c)=0$ for all sufficiently large $n$.
\end{theorem}
\begin{proof}
\Cref{ent:mv:prop:linear} gives entireness of $F_c$. The block family at
$\lambda c$ is exactly $(F_{(n)}(c)\lambda^n)_n$; apply
\cref{ent:mv:cor:onevar}. Directly: if infinitely many $F_{(n)}(c)$ are
nonzero, the corresponding terms at $t^{-e_*}c$ have quotient valuations
$-n\pi_H(e_*)$, strictly descending, and otherwise the evaluation is a finite
sum; this needs no factorization or zero theorem.
\end{proof}

The one-variable theorem is an input here, not a new priority claim. Directions
not defined over $\K$ need a different analysis, and
\cref{ent:mv:thm:rays} does not classify block values at all points of
$\KD^d$.

For nonempty $J\subseteq\{1,\ldots,d\}$ put
$\Sigma_J=\{\alpha\in\Sigma(F):\alpha_j=0\text{ for }j\notin J\}$ and
$\PP_J=\{[c]\in\PP^{d-1}(\C):c_j=0\text{ for }j\notin J\}$.

\begin{proposition}[Directions for the whole family]\label{ent:mv:prop:coordinate}
Let $\kk=\C$, $F\in\mathcal E_{\Gamma,d}$ and $e_*>H$. The ordinary complex
directions $[c]$ at which the whole monomial family is strongly summable at
$t^{-e_*}c$ are exactly $\bigcup\{\PP_J:J\ne\varnothing,\ \Sigma_J\text{
finite}\}$, a finite union of coordinate projective subspaces. For
nonpolynomial $F$ it is a proper algebraic subset, and it is contained in the
block-exceptional locus $\Pdir(F)$ of \cref{ent:mv:subsec:exceptional}.
\end{proposition}
\begin{proof}
Let $J$ be the set of nonzero coordinates of $c$. Every surviving term has
quotient weight $-|\alpha|\pi_H(e_*)$, and these are well ordered exactly
when the degrees in $\Sigma_J$ are bounded, that is, when $\Sigma_J$ is finite.
The family of such $J$ is downward closed, so the admissible directions are
the union of the $\PP_J$. For nonpolynomial $F$ the full coordinate set is not
admissible, so a direction with all coordinates nonzero lies outside.
Inclusion in the block locus is regrouping.
\end{proof}

\begin{proposition}[Realizing the coordinate pattern]
\label{ent:mv:prop:coordinate-realize}
Suppose $\cf(\Gamma)=\aleph_0$. Every downward-closed family $\mathcal J$ of
subsets of $\{1,\ldots,d\}$ containing $\varnothing$ is the family of $J$
with $\Sigma_J$ finite, for some entire series; if the full set is not in
$\mathcal J$, the series may be chosen nonpolynomial.
\end{proposition}
\begin{proof}
If every subset lies in $\mathcal J$, take a polynomial. Otherwise let $p_B$
be the indicator vector of each inclusion-minimal $B\notin\mathcal J$, and put
$\Sigma=\bigcup_B\{np_B:n\ge1\}$. Its face $\Sigma_J$ is infinite exactly when
some such $B$ is contained in $J$, that is, exactly when $J\notin\mathcal J$,
and otherwise it is empty. \Cref{ent:mv:prop:realize-support} realizes
$\Sigma$.
\end{proof}

This finite coordinate classification contrasts with the much larger family
of block-exceptional loci below, and it detects when a claimed multivariate
extension theorem has silently changed its summation convention.

\subsection{Exceptional projective directions}
\label{ent:mv:subsec:exceptional}
In this subsection $\kk=\C$, and projective space means the ordinary set
$\PP^{d-1}(\C)$.

\begin{definition}[Polynomial-direction locus]\label{ent:mv:def:pdir}
For $F=\sum_nF_{(n)}$ over $\K$ put
$\Pdir(F)=\{[c]\in\PP^{d-1}(\C):F(Tc)\in\K[T]\}$, which does not depend on
the representative of $[c]$. For entire $F$, \cref{ent:mv:thm:rays}
identifies it with the ordinary directions admitting homogeneous-block
evaluation at a new dominating scale.
\end{definition}

The algebraic statements below hold for arbitrary formal $F$ over $\K$;
entireness is needed only for the radial domain formula and the analytic
interpretation.

\begin{theorem}[Necessary form and generic obstruction]
\label{ent:mv:thm:exceptional}
For every $F\in\K[[X_1,\ldots,X_d]]$ there are ascending projective algebraic
sets $\mathcal Z_N\subseteq\PP^{d-1}(\C)$ with
\[
 \Pdir(F)=\bigcup_{N\ge0}\mathcal Z_N,\qquad
 \mathcal Z_N=\{[c]:F_{(n)}(c)=0\text{ for all }n>N\}.
\]
If $F$ is not a polynomial, every $\mathcal Z_N$ is proper, $\Pdir(F)$ is not
all of projective space, and its complement is dense in the ordinary complex
topology. For $d=2$ the set $\Pdir(F)$ is at most countable.
\end{theorem}
\begin{proof}
Coefficientwise, $F_{(n)}(X)=\sum_{\eta\in S_n}t^\eta P_{n,\eta}(X)$ with
ordinary complex homogeneous polynomials $P_{n,\eta}$ of degree $n$, $S_n$
being the finite union of the Hahn supports of the finitely many coefficients
of $F_{(n)}$. For a complex vector $c$, $F_{(n)}(c)=0$ exactly when every
$P_{n,\eta}(c)=0$. So $\mathcal Z_N$ is the common zero set of a family of
homogeneous complex polynomials; the ideal they generate is finitely
generated by the Hilbert basis theorem \cite{Stacks}, so $\mathcal Z_N$ is
projective algebraic. The sets ascend, and their union is eventual vanishing.

If $F$ is nonpolynomial, each $N$ is followed by a nonzero block containing
some nonzero $P_{n,\eta}$, which vanishes on $\mathcal Z_N$ and not
identically; so $\mathcal Z_N$ is proper. Choose one nonzero homogeneous
$p_N$ vanishing on each $\mathcal Z_N$. Their coefficients generate a
countable subfield $L\subseteq\C$. In the chart $c_1=1$ the polynomials
$p_N(1,Y_2,\ldots,Y_d)$ are nonzero; choosing $d-1$ complex numbers
algebraically independent over $L$, successively outside the algebraic
closure of a countable field, gives a point outside every $\mathcal Z_N$. The
same choices can be made inside any product of nonempty open disks, since only
countably many values are forbidden at each step, which gives density. For
$d=1$ the locus of a nonpolynomial series is empty. Finally, in $\PP^1(\C)$ a
proper algebraic set is finite, and a countable union of finite sets is at
most countable.
\end{proof}

The covering argument is the countable-field device of Lemma
\texttt{hol:lem:generic} in the report
\texttt{docs/surcomplex/holonomic-rigidity-for-entire-hahn-functions}
\cite{RepoHolonomic}, which chooses $c\in\C^d$ avoiding countably many
nonzero polynomials over $\K$; here it is applied to complex polynomials and
refined to disks for density. The device is elementary, and neither report
claims it as new. No countability of Hahn supports is hidden: an individual
coefficient may have uncountable support, the proof never enumerates
exponents, and Noetherianity reduces each $\mathcal Z_N$ to finitely many
equations while there are only countably many $N$.

\begin{theorem}[Exact realization of the exceptional locus]
\label{ent:mv:thm:realization}
Suppose $\cf(\Gamma)=\aleph_0$. For every countable union
$A=\bigcup_{m\ge1}Y_m\subseteq\PP^{d-1}(\C)$ of projective algebraic sets
there is $F\in\mathcal E_{\Gamma,d}$ with $\Pdir(F)=A$. If $A$ is proper, $F$
can be chosen nonpolynomial; if $A$ is all of projective space, no
nonpolynomial formal series has that locus, and a polynomial realizes it.
\end{theorem}
\begin{proof}
The last case is $F=0$ together with \cref{ent:mv:thm:exceptional}. Let $A$ be
proper and replace $Y_m$ by $Y_1\cup\cdots\cup Y_m$, so that the sets ascend
and are proper. For each $m$ choose nonzero homogeneous polynomials
$p_{m,1},\ldots,p_{m,r_m}$ with common projective zero set $Y_m$ (the single
polynomial $1$ if $Y_m$ is empty). To each triple $(m,k,j)$, $1\le k\le r_m$,
$1\le j\le d$, assign a distinct integer $D_{m,k,j}>\deg p_{m,k}$, in
increasing order block by block in $m$, and put
$Q_{D_{m,k,j}}=X_j^{D_{m,k,j}-\deg p_{m,k}}p_{m,k}$, with $Q_n=0$ at unassigned
degrees. The block of index $m$ has common projective zero set exactly $Y_m$:
off $Y_m$ some $p_{m,k}$ and some coordinate $c_j$ are nonzero, and using every
coordinate prevents a spurious exceptional hyperplane. Set
$F=\sum_nt^{\varphi_n}Q_n$ with the damping of
\cref{ent:mv:prop:realize-support}; each nonzero monomial coefficient of
degree $n$ has valuation $\varphi_n$, so $F$ is entire, and it is
nonpolynomial. If $[c]\in A$ it lies in some $Y_M$ and all later blocks, so
$F(Tc)$ is a polynomial. If $[c]\notin A$, every block has a polynomial not
vanishing at $c$, and the assigned degrees are unbounded.
\end{proof}

\begin{corollary}[Sharp two-variable classification]\label{ent:mv:cor:countable}
If $\cf(\Gamma)=\aleph_0$, the ordinary projective directions admitting
new-scale homogeneous-block evaluation of a nonpolynomial entire series in two
variables are exactly the at most countable subsets of $\PP^1(\C)$.
\end{corollary}
\begin{proof}
Necessity is \cref{ent:mv:thm:exceptional}; a countable set is a countable
union of points, so \cref{ent:mv:thm:realization} gives sufficiency; the
identification with new-scale evaluation is \cref{ent:mv:thm:rays}.
\end{proof}

For pairwise distinct $c_1,c_2,\ldots\in\C$ the series
\[
 F_{\mathrm{dir}}(X,Y)=1+\sum_{n\ge1}t^{n^2}\prod_{j=1}^n(Y-c_jX)
\]
over $\C((t^\Q))$ is entire, and its restriction to $[1:c_m]$ is the
polynomial $1+\sum_{n<m}t^{n^2}\prod_{j\le n}(c_m-c_j)T^n$ of degree $m-1$,
while every other slope and the direction $[0:1]$ give nonpolynomial
restrictions. So $\Pdir(F_{\mathrm{dir}})=\{[1:c_m]:m\ge1\}$; a projective
linear change and \cref{ent:mv:prop:linear} place a prescribed point at
infinity. Enumerating $\Q+i\Q$ gives a countable dense exceptional locus:
polynomial behaviour on a dense set of complex lines does not force a Hahn
series to be polynomial, while polynomial behaviour on \emph{every} complex
line does, by \cref{ent:mv:thm:exceptional}.

The exceptional set is neither an ordinary convergence set nor a zero set of
$F$: it asks for eventually vanishing coefficients on a line, which the
enlargement turns into a summability condition because a negative new
quotient valuation overwhelms every old coefficient. Ma and Neelon already
realize countable unions of projective algebraic sets as convergence sets of
divergent complex power series \cite[Corollary~4.19]{MaNeelon}. That
realization idea is not claimed as new, their stronger analytic theorems are
not reproved here, and convergence sets must not be conflated with
polynomial-direction loci. What is proposed is the exact Hahn extension
interpretation with the complete classification and an entire realization
over every countably cofinal value group.

\subsection{Transfer to surreal and surcomplex numbers}
\label{ent:mv:subsec:surreal}
Conway's omega-map $x\mapsto\omega^x$ and his normal-form theorem identify
surreal numbers with real generalized series with set-sized reverse
well-ordered supports \cite[Theorem~10.2]{Berarducci}; this is the map
\eqref{ent:eq:embedding}, $t^\gamma\mapsto\omega^{-\gamma}$, and it must not
be confused with an exponential defined from Gonshor's $\exp$. For a
set-sized ordered subgroup $\Gamma\subseteq\No$ it embeds $\R((t^\Gamma))$ in
$\No$ and $\C((t^\Gamma))$ in $\No[i]$, preserving and reflecting the support
conditions that define strong sums.

\begin{lemma}[Localization of the data]\label{ent:mv:lem:localize}
Every set of surreal or surcomplex numbers lies in such an embedded workspace
for a nonzero divisible set-sized ordered subgroup $\Gamma\subseteq\No$, and
adjoining a finite tuple of arguments needs at most a set-sized enlargement
$\Gamma\subseteq\Delta$.
\end{lemma}
\begin{proof}
Take the union of the exponent supports of the real and imaginary parts of
all given numbers, adjoin $1$, and take the rational linear span in $\No$;
this is a divisible set, unchanged by negating exponents. Repeat after
adjoining the arguments.
\end{proof}

This is the localization principle of the foundations report \cite{RepoFoundations},
already consumed by \cref{ent:cor:all-no}; the fifth source gives the short
proof above. Localization says nothing about entireness, which remains a
hypothesis on a specified old workspace.

\begin{corollary}[Surcomplex support-domain theorem]\label{ent:mv:cor:surcomplex}
Let $\Gamma\subseteq\No$ be a set-sized ordered subgroup and $F$ whole-family
entire over the embedded field $\C((t^\Gamma))$. For $z\in\No[i]^d$ choose a
set-sized $\Delta$ containing $\Gamma$ and the exponent supports of $z$, and
put $H=\conv_\Delta\Gamma$. After discarding the terms killed by zero
coordinates, the family defining $F(z)$ is strongly summable exactly when
$\{\sum_j\alpha_j(v(z_j)+H):a_\alpha\ne0\}$ is well ordered in $\Delta/H$.
The answer does not depend on the choice of $\Delta$.
\end{corollary}
\begin{proof}
Transport \cref{ent:mv:thm:domain}. For $\Delta\subseteq\Delta'$ the hull
$H'$ satisfies $H'\cap\Delta=H$, so $\Delta/H$ order-embeds in $\Delta'/H'$
and well-ordering of the same weight set is unchanged.
\end{proof}

The real version follows with $\R$ in place of $\C$. Neither needs algebraic
closedness of $\No[i]$, a derivation or a square root. The classification
of \cref{ent:mv:subsec:exceptional} uses ordinary complex directions and is
not asserted over an arbitrary residue field.

Two phenomena are explicit. With $\Gamma=\Q$ and $t=\omega^{-1}$, the series
$\sum_{n\ge0}\omega^{-n^2}(XY)^n$ is entire in that workspace; after
enlarging by the dominant exponent $e_*=\omega$ it is strongly evaluable at
$X=\omega^{\omega}$, $Y=\omega^{-\omega}$, both outside the old field and the
first outside its one-variable extension domain, because the family is
exactly $(\omega^{-n^2})_n$ and every weight is zero. Infinite fibres are
therefore necessary in \cref{ent:mv:thm:domain}. And with pairwise distinct
$c_j\in\C$, the series
$1+\sum_{n\ge1}\omega^{-n^2}\prod_{j\le n}(Y-c_jX)$ admits block evaluation
at $\omega^{\omega}(1,c_m)$, where it is the finite polynomial above at
$T=\omega^\omega$, while at every other ordinary slope and at the direction
at infinity its nonzero blocks have strictly descending leading exponents: a
prescribed countable set of directions crosses a genuinely surreal scale and
every other direction fails.

\begin{proposition}[Whole-class rigidity in both conventions]
\label{ent:mv:prop:allscale}
Let $F\in\No[i][[X_1,\ldots,X_d]]$ be one formal series in finitely many
variables. If its whole monomial family is strongly summable at every
surcomplex tuple, $F$ is a polynomial. The same holds if instead its
sequence of homogeneous blocks is strongly summable at every surcomplex
tuple.
\end{proposition}
\begin{proof}
Localize the coefficients in some $\K$ (\cref{ent:mv:lem:localize}) and choose
a positive surreal $e_*$ above the set $\Gamma$. For a nonpolynomial whole
family at the diagonal point $(\omega^{e_*},\ldots,\omega^{e_*})$, the leading
exponents $v(a_\alpha)-|\alpha|e_*$ along strictly increasing degrees
strictly decrease. For blocks, \cref{ent:mv:thm:exceptional}, which does not
need entireness, gives a complex direction $c$ with nonpolynomial restriction,
and evaluating that restriction at $\omega^{e_*}$ again gives strictly
descending leading exponents.
\end{proof}

The whole-family half is already in this collection, in \cref{ent:cor:all-no}
for one variable and in the analysis report, and is not reclaimed; the block
half shows that cancellation on exceptional lines does not evade global
rigidity. Neither forbids a different analytic category of local germs,
non-power-series functions or another specified summation, and no strong sum
is identified with a limit in the fine surreal topology.

\subsection{What this section does not settle}
\label{ent:mv:subsec:limits}
The following limitations are the fifth source's and are retained without
weakening.
\begin{itemize}[leftmargin=1.4em]
\item Only the scalar-extension clause of \cref{ent:q:several} is answered,
for the whole-family convention. The preparation, interpolation, divisor
and ideal-theoretic clauses are not settled; no several-variable GCD or
B\'ezout classification is asserted, and a positive-dimensional zero set
cannot be replaced by the one-variable finite root divisor of
\cref{ent:sec:products}. \emph{Re-scoped by the seventh source}
(\crefrange{ent:rc:subsec:setting}{ent:rc:subsec:limits}): for configurations
of points that an entire automorphism carries into a coordinate line, the
ideal-theoretic clause is settled by explicit fat-point generators, their
sharp number and non-finite generation for unbounded multiplicities, and the
interpolation of values reduces exactly to the one-variable image theorem
(\cref{ent:rc:thm:fat-main,ent:rc:cor:evaluation}). Preparation, divisors
with positive-dimensional support, GCD and B\'ezout classification, and ideals
of configurations that are not line-rectifiable remain unsettled.
\emph{Re-scoped by source~11 (batch~32)}
(\crefrange{ent:as:subsec:setting}{ent:as:subsec:limits}): for $\kk\in\{\R,\C\}$
and ideals extended from $\kk[X]$, membership, normal forms and contraction are
decided scale by scale and $\kk[X]\to\Ed$ is faithfully flat
(\cref{ent:as:cor:ideal,ent:as:cor:normal-form,ent:as:thm:flat}); the entire
functions vanishing on a set of ordinary points, and for $\kk=\C$ on the
$\K$-points of an algebraic set defined over $\C$, form an extended ideal
(\cref{ent:as:cor:vanishing}, the merge's). Ideals not extended from $\kk[X]$,
preparation, GCD and B\'ezout classification, and divisors with
positive-dimensional support that are not algebraic sets defined over $\kk$
remain unsettled.
\item The full homogeneous-block domain at an arbitrary tuple of new-field
coordinates is not classified. \Cref{ent:mv:thm:rays} covers rays whose
direction is defined over the old field, and new-field direction ratios can
contain new scales, so that the block coefficients leave the old field after
substitution (\cref{ent:mv:q:block-domain}).
\item The number of variables is finite, and this is used; all fields are
full Hahn fields, and nothing is claimed for support-restricted or incomplete
subfields, where an ambient strong sum need not lie in the smaller field.
\item The semilinear test is a finite decision procedure only relative to a
supplied semilinear presentation and exact ordered-group operations on
$\Delta/H$. It does not discover supports, recognize semilinearity, build the
convex-hull quotient from a black-box exponent oracle or decide general
well-ordering, and no finite sampling can certify the general criterion:
every finite support has a well-ordered weight set at every point
(\cref{ent:mv:q:effective}).
\item \Cref{ent:mv:subsec:exceptional} depends on $\kk=\C$ and ordinary
complex directions.
\item The finite-word well-quasi-order theorem behind
\cref{ent:lem:neumann}(b) and the Hilbert basis theorem are imported, and
choice is used for lattice bases, coefficient bounds and algebraically
independent coordinates.
\end{itemize}

% ---------------------------------------------------------------------------
% Seventh source (file prefix 10-entire-rectification), labels ent:rc:.
% Displayed equations in these subsections are numbered (R1), (R2), ... so
% that the numbers of all later equations of the report are unchanged.
% ---------------------------------------------------------------------------
\begingroup
\setcounter{entrcsaveeq}{\value{equation}}%
\setcounter{equation}{0}%
\renewcommand{\theequation}{R\arabic{equation}}%
\renewcommand{\theHequation}{entrc.\arabic{equation}}%

\subsection{Rectification and fat points: setting and main results}
\label{ent:rc:subsec:setting}

This subsection and the six after it are the seventh source's contribution,
the manuscript \emph{Entire Automorphisms at Surreal Scales: an order-unit
dichotomy, omnific rectification, and exact fat-point ideals} (file prefix
\texttt{10-entire-rectification}); their labels carry \texttt{ent:rc:}. They
ask when an infinite configuration of points of $\K^d$ can be carried into a
coordinate line by an invertible entire change of coordinates, and what the
ideals of points with multiplicities look like once it has been. Against this
report they answer part of the ideal-theory and interpolation clauses of
\cref{ent:q:several}, for configurations of points that are line-rectifiable
in the sense below; the extent is recorded after that question. The source
itself described its problem as ``a different several-variable global
problem''; that was accurate as a description of its question and is
superseded here by the placement against \cref{ent:q:several}. Four warnings
govern these subsections.

\begin{enumerate}[label=(\arabic*),leftmargin=2em]
\item \emph{Hypotheses.} Warning (1) at the head of this section does not
apply here; the standing hypotheses are \cref{ent:rc:conv:standing}. In
particular $\Gamma$ is divisible, as in the rest of the report, and the
coefficient field is $\R$ or $\C$. Uncountability of the coefficient field is
used once, in \cref{ent:rc:lem:projection}. Algebraic closedness,
preparation and root counts are not used.
\item \emph{Printed once.} For dependency control the seventh source reproved
the analytic tools it needs: the finite-variable coefficient criterion,
canonical products with multiplicities, division by a linear factor and by an
escaping product, a sufficient value-interpolation condition, and the
countability of radially finite sets. These are
\cref{ent:thm:multi-criterion,ent:thm:product,ent:lem:linear-division,ent:cor:gcd,ent:prop:kernel-direct,ent:thm:image,ent:thm:interpolation,ent:lem:divisor-countable},
and the source credits this report for them. They are printed once, by
citation, in \cref{ent:rc:subsec:calculus}; the one proof kept as a second
route is the source's explicit interpolation series
(\cref{ent:rc:lem:interpolation}).
\item \emph{Words.} ``Entire'' means whole-family entire
(\cref{ent:mv:def:entire}); homogeneous-block evaluation plays no role. An
\emph{ideal} is an algebraic ideal: a generating set permits finite linear
combinations only, and no analytic closure is substituted for ideal
generation.
\item \emph{Letters.} Several of the source's letters collide with this
report's and are renamed; the last table of \cref{ent:subsec:notation} gives
the translation. The three collisions that matter are these. The source's
$\mathcal A_\Gamma$ is a ring of polynomial parts, written $\Kle$ here,
because our $\A$ is $\C[X]((t^\Gamma))$. Its denominators $d_i$ are written
$\varpi_i$, because $d$ is the number of variables. Its $e_j$ denoted both the
basis vectors of $\Gamma_\infty$ and the coordinate vectors of $\K^d$ (the
source flags this); here $e_j$ always denotes the former, and the coordinate
vectors are $\mathbf u_1,\ldots,\mathbf u_d$, which are not the units $u_j$
of \cref{ent:mv:lem:twist}.
\end{enumerate}

\begin{convention}[Standing hypotheses of the seventh source's subsections]
\label{ent:rc:conv:standing}
$\Gamma\ne0$ is a set-sized divisible ordered abelian group. The coefficient
field $\kk$ is $\R$ or $\C$ with the trivial valuation, and
$\K=\kk((t^\Gamma))$; for $\kk=\R$ the one-variable ring $\E$ is the real
form of \cref{ent:prop:real-form}. The number $d\ge1$ of variables is finite,
$\Ed$ is the ring of \cref{ent:mv:def:entire}, and $\E=\mathcal E_{\Gamma,1}$
is written in the variable $Z$. For $p=(p_1,\ldots,p_d)\in\K^d$ put
\[
 v(p)=\min_{1\le j\le d}v(p_j),\qquad v(0)=+\infty,\qquad
 B^d_\gamma=\{p\in\K^d:v(p)\ge\gamma\}\quad(\gamma\in\Gamma),
\]
so that $B^1_\gamma$ is the ball $B_\gamma$ of \cref{ent:sec:zeros}. For
$F=\sum_\alpha a_\alpha X^\alpha\in\Ed$ and $\gamma\in\Gamma$, the weighted
Gauss valuation is
\begin{equation}\label{ent:rc:eq:gauss}
 m_\gamma(F)=\min_{\alpha}\{v(a_\alpha)+|\alpha|\gamma\},
 \qquad m_\gamma(0)=+\infty,
\end{equation}
which for $d=1$ is \eqref{ent:eq:weighted}; the source's $w_\gamma$ is this
$m_\gamma$, with the same sign. ``Tends to $+\infty$'' always means cofinally
in $\Gamma$, not merely increasingly.
\end{convention}

The minimum in \eqref{ent:rc:eq:gauss} exists: choose one nonzero
coefficient; by \cref{ent:thm:multi-criterion}(iii) only finitely many
coefficients have weighted valuation at most its own, and these finitely
many include the minimum. The value bounds $v(F(z))$ from below on all of
$B^d_\gamma$.

\begin{definition}[Radially finite and line-rectifiable configurations]
\label{ent:rc:def:rectifiable}
A subset $\mathcal X\subseteq\K^d$ is \emph{radially finite} if
$\mathcal X\cap B^d_\gamma$ is finite for every $\gamma\in\Gamma$. An
\emph{entire map} $\K^d\to\K^s$ has coordinates in $\Ed$. An \emph{entire
automorphism} of $\K^d$ is an entire map $\Phi$ with an entire map
$\Phi^{-1}$ such that $\Phi^{-1}\circ\Phi$ and $\Phi\circ\Phi^{-1}$ are the
identity. A subset $\mathcal X$ is \emph{line-rectifiable} if some entire
automorphism maps it into the coordinate line $\K\times\{0\}^{d-1}$.
\end{definition}

For $d=1$ the radially finite subsets are the supports of the radially finite
divisors of \cref{ent:def:radial} with all multiplicities one. Radial
finiteness is a condition on the valuation polydisks of the one fixed
workspace $\K^d$; it is not discreteness in the fine surreal topology, which
is weaker and is not a substitute (\cref{ent:rc:subsec:limits}). The order
unit is \cref{ent:def:order-unit}; the source uses the same definition and
notes that it is weaker than asking $\Gamma$ to be Archimedean, since smaller
Archimedean classes may lie below the largest one.

\Needspace{14\baselineskip}
\begin{theorem}[Rectification dichotomy]\label{ent:rc:thm:main}
Suppose $d\ge2$ and $\cf(\Gamma)=\aleph_0$. The following are equivalent:
\begin{enumerate}[label=\textnormal{(\roman*)},leftmargin=2em]
\item $\Gamma$ has an order unit;
\item every radially finite subset of $\K^d$ is line-rectifiable.
\end{enumerate}
When an order unit exists, an affine change of coordinates followed by one
triangular entire shear suffices. When no order unit exists, there is a
radially finite $\mathcal X$ such that, for every entire automorphism $\Phi$,
every sufficiently late block of $d+1$ specified points of
$\Phi(\mathcal X)$ is affinely independent. In particular no
$\Phi(\mathcal X)$ lies in an affine hyperplane.
\end{theorem}

The positive half is \cref{ent:rc:prop:positive}, the negative half
\cref{ent:rc:thm:obstruction}; the proof is assembled after the latter. The
hypothesis $\cf(\Gamma)=\aleph_0$ cannot be dropped: if
$\cf(\Gamma)>\aleph_0$, every radially finite set is finite
(\cref{ent:rc:lem:countable}), hence polynomially line-rectifiable, while no
order unit exists, since an order unit forces countable cofinality
(\cref{ent:rem:one-name}). The three cases are tabulated at the end of
\cref{ent:rc:subsec:obstruction}. There is a substantial positive class on
the negative side of the dichotomy. Put
\begin{equation}\label{ent:rc:eq:Kle}
 \Kle=\{a\in\K:\supp(a)\subseteq\Gamma_{\le0}\},
\end{equation}
the Hahn polynomial-part ring; it is a subring of $\K$, not a field.

\begin{theorem}[Polynomial-part and omnific rectification]
\label{ent:rc:thm:omnific-main}
For every $\Gamma$ as in \cref{ent:rc:conv:standing}, every radially finite
subset of $(\Kle)^d$ is line-rectifiable, again by an affine change of
coordinates followed by one triangular entire shear. For a divisible ordered
subgroup $\Gamma\subseteq\No$ and the embedding \eqref{ent:eq:embedding},
this applies to the omnific integers
$\Z+\R((t^{\Gamma_{<0}}))$ of the real workspace and, after
complexification, to the Gaussian omnific integers
$\Z[i]+\C((t^{\Gamma_{<0}}))$ of the complex one; see
\eqref{ent:rc:eq:omnific-real}--\eqref{ent:rc:eq:omnific-complex}. The
ambient automorphism is not asserted to preserve either ring.
\end{theorem}

Here $\kk((t^{\Gamma_{<0}}))$ denotes the set of Hahn series supported in
$\Gamma_{<0}$, not a field. In particular the negative examples of
\cref{ent:rc:thm:main} need genuine infinitesimal tails: they cannot lie in
$(\Kle)^d$. The contrast is sharper than a statement that some discrete sets
happen to interpolate and others do not.

For $a\in\K^d$ and $F\in\Ed$, $\ord_aF$ is the least total degree of a
nonzero term of the Taylor expansion of $F$ at $a$
(\cref{ent:rc:lem:composition}), with $\ord_a0=+\infty$; multiplicities below
are total multiplicities in all coordinates. For an ideal $I$, $\mu(I)$ is the
least cardinality of an algebraic generating set.

\Needspace{12\baselineskip}
\begin{theorem}[Sharp fat-point generator count]\label{ent:rc:thm:fat-main}
Let $d\ge2$, let $\mathcal X=\{p_i\}$ be a nonempty, radially finite,
line-rectifiable subset of $\K^d$, and prescribe positive integers $m_i$.
The ideal
\[
 I_{\mathbf m}(\mathcal X)=\{F\in\Ed:\ord_{p_i}F\ge m_i\text{ for all }i\}
\]
is finitely generated if and only if the $m_i$ are bounded. If
$\bar m=\max_im_i<\infty$, then
\begin{equation}\label{ent:rc:eq:mu-main}
 \mu\bigl(I_{\mathbf m}(\mathcal X)\bigr)=\binom{\bar m+d-1}{d-1},
\end{equation}
with explicit generators in \cref{ent:rc:thm:fat-generators}. For $m_i=1$
the vanishing ideal is generated by a regular sequence of length $d$. In
dimension one these ideals are principal, even for unbounded multiplicities.
\end{theorem}

\paragraph{Relation to classical tameness.}
Rosay and Rudin studied discrete subsets of ordinary $\C^d$ that a
holomorphic automorphism carries onto $\N\times\{0\}^{d-1}$
\cite{RosayRudin}; related notions, including the extension of injective
self-maps, appear in Winkelmann's treatment \cite{Winkelmann}. The source
uses the distinct term \emph{line-rectifiable}: the image is required to lie
in the line, not to be a prescribed set there, and no equivalence with those
stronger properties is assumed. Over a trivially valued coefficient field
$\N$ itself is valuation-bounded, so it is not the model escaping sequence
here. The projection-and-shear strategy has a classical precedent; the
source's novelty claim concerns the arbitrary-rank dichotomy and its
valuation-scale obstruction, not the strategy in isolation. The classical
comparison suggests questions but is not a transfer principle; there may be
several inequivalent notions of tameness in an arbitrary-rank entire
category, and the dichotomy is proved only for being sent into a coordinate
line.

\subsection{The several-variable calculus, and what is cited}
\label{ent:rc:subsec:calculus}

\paragraph{The coefficient criterion.} The seventh source's all-scale
criterion is \cref{ent:thm:multi-criterion}, proved there by the same
descending-leading-exponent argument, and is printed once. The source adds
two remarks, kept here: the criterion tests the individual monomials, not only
their homogeneous block sums; and a countable family whose valuations tend
cofinally to $+\infty$ is strongly summable, which is
\cref{ent:lem:cofinal-tails} and \cref{ent:mv:lem:tails}(a).

\begin{lemma}[Gauss values and faithful evaluation]\label{ent:rc:lem:gauss}
The ring $\Ed$ is an integral domain, and for $F,G\in\Ed$ and
$\gamma\in\Gamma$
\[
 m_\gamma(FG)=m_\gamma(F)+m_\gamma(G),\qquad
 m_\gamma(F+G)\ge\min\{m_\gamma(F),m_\gamma(G)\}.
\]
A nonzero element of $\Ed$ is nonzero at some point of $\K^d$.
\end{lemma}
\begin{proof}
The inequality for sums is immediate. For products, fix $\gamma$. A term of
the Cauchy product has weighted valuation the sum of the two weighted input
valuations; both input families are bounded below and tend to $+\infty$, so
only finitely many pairs fall below any threshold, all coefficient
regroupings are legitimate, and $FG\in\Ed$ by
\cref{ent:thm:multi-criterion}. Substitute $X=t^\gamma U$ and multiply $F$ by
$t^{-m_\gamma(F)}$: the coefficients acquire nonnegative valuation, and only
finitely many have valuation zero, so the reduction modulo positive valuation
is a nonzero polynomial $\inpol_\gamma(F)\in\kk[U_1,\ldots,U_d]$. The product
of two nonzero such polynomials is nonzero, and it is the reduction for $FG$;
this gives the equality. Since $\kk$ is infinite, choose $c\in\kk^d$ with
$\inpol_\gamma(F)(c)\ne0$; then $F(t^\gamma c)$ has valuation exactly
$m_\gamma(F)$ and is nonzero. The domain property follows from the inclusion
in $\K[[X_1,\ldots,X_d]]$.
\end{proof}

In one variable the multiplicativity is \eqref{ent:eq:gauss-mult} and the
faithfulness is \cref{ent:cor:units} (over $\R$,
\cref{ent:sm:rem:real}); \cref{ent:thm:multi-unit} uses the same regrouping.

\begin{lemma}[All-disk summation]\label{ent:rc:lem:sum}
Let $F_n\in\Ed$. If $m_\gamma(F_n)\to+\infty$ for every $\gamma\in\Gamma$,
then $\sum_nF_n$ is a well-defined element of $\Ed$, it may be evaluated and
differentiated termwise, and its weighted value at $\gamma$ is at least any
common lower bound for the $m_\gamma(F_n)$. More generally, if a formal series
$H$ is the coefficientwise limit of $H_N\in\Ed$ and, for every $\gamma$,
all weighted coefficients of $H-H_N$ at $\gamma$ are bounded below by some
$B_{N,\gamma}$ that tends to $+\infty$ with $N$, then $H\in\Ed$.
\end{lemma}
\begin{proof}
For a fixed coefficient, the contributions of the $F_n$ have valuations
tending to $+\infty$, so their Hahn sum exists (\cref{ent:lem:cofinal-tails}).
Fix $\gamma$ and a threshold $b$. All but finitely many summands have every
weighted coefficient above $b$, and in each of the remaining finitely many
summands only finitely many coefficients fail that inequality; so the sum
satisfies \cref{ent:thm:multi-criterion}(iii), with the asserted lower
bound. The limit assertion follows by choosing one sufficiently accurate
$H_N$. Evaluation obeys the same estimates on every polydisk. A Hasse
derivative of multi-order $\beta$ replaces $a_\alpha$ by
$\binom\alpha\beta a_\alpha$, a nonzero integer multiple of valuation
$v(a_\alpha)$, and lowers the weight by $|\beta|\gamma$; so the criterion and
the uniform estimates survive differentiation, which justifies termwise
differentiation.
\end{proof}

In one variable this is \cref{ent:sm:lem:entire-sum}.

\begin{lemma}[Composition and Taylor expansion]\label{ent:rc:lem:composition}
Entire series are closed under substitution of entire series, including
substitutions with nonzero constant terms, and evaluation commutes with
substitution. Translations and invertible affine maps induce automorphisms of
$\Ed$. The usual Taylor and chain rules hold, with the coefficients in $\K$
held constant.
\end{lemma}
\begin{proof}
On the polydisk of weight $\gamma$ let $\nu$ be a common lower bound for the
weighted values of the inner series $G_1,\ldots,G_s$. For an outer series
$F=\sum a_\alpha U^\alpha$, \cref{ent:rc:lem:gauss} gives
$m_\gamma(a_\alpha G^\alpha)\ge v(a_\alpha)+|\alpha|\nu$, and the right side
tends to $+\infty$. Apply \cref{ent:rc:lem:sum} to the polynomial partial sums
of $F(G)$; all nested coefficient sums can be regrouped because only finitely
many contributions remain below each threshold (\cref{ent:mv:lem:tails}). The
chain rule holds for polynomials, and passes to these limits by the derivative
estimate of \cref{ent:rc:lem:sum}. Translation is a special case, with inverse
the opposite translation; Taylor coefficients are the Hasse derivatives at the
centre, and the same bounds justify the expansion. Affine inversion is
polynomial.
\end{proof}

In one variable this is \cref{ent:prop:operations}.

\begin{remark}[A forbidden shortcut]\label{ent:rc:rem:shortcut}
The source records the trap that \cref{ent:lem:neumann} and
\cref{ent:ex:ranktwo-unit} already isolate: a positive $\lambda\in\Gamma$ need
not have cofinal multiples, so $v(q)>0$ does not imply $q^n\to0$ in the
valuation topology, and Hahn summability of a geometric series is not
convergence of its partial sums. The infinite constructions below are
controlled by cofinal sequences of node radii instead, and no diskwise
expression $1+q$ is assumed to have a cofinally convergent geometric inverse.
\end{remark}

\paragraph{One-variable products and division, cited.}
Let $D$ be a radially finite divisor on $\K$ whose support consists of
nonzero points. Its canonical product $P_D$ is entire with divisor exactly
$D$ and constant coefficient $1$, for arbitrary, possibly unbounded,
multiplicities (\cref{ent:thm:product}); a zero of $f\in\E$ can be divided
out, repeatedly up to its order (\cref{ent:lem:linear-division}); and
$f\in P_D\E$ if and only if $\ord_\rho f\ge D(\rho)$ at every node
(\cref{ent:cor:gcd}, and without preparation \cref{ent:prop:kernel-direct}).
The seventh source proves these three statements for itself. Its division
proof, by finite partial quotients $f/P_N$ and an inverse-tail estimate uniform
in the degree, is the argument of \cref{ent:prop:kernel-direct} and is not
printed twice; its observation that a purely formal quotient by an infinite
product is not automatically entire is the remark after
\cref{ent:thm:factorization}.

\begin{remark}[The real coefficient field]\label{ent:rc:rem:real}
Everything cited in this subsection holds over $\R((t^\Gamma))$ with nodes and
data there. \Cref{ent:thm:multi-criterion} is stated over any trivially valued
field; the constructions of \cref{ent:thm:product,ent:lem:linear-division}
use only the support calculus and units of the Hahn algebra with initial
polynomial~$1$, as the sixth source's use of them over $\R$ already requires;
\cref{ent:prop:kernel-direct} uses neither preparation nor algebraic
closedness; and \cref{ent:thm:image,ent:cor:order-unit-image} hold over $\R$ by
\cref{ent:prop:real-form}. The seventh source proves its own versions over
$\kk\in\{\R,\C\}$ directly. This remark is the merge's bookkeeping.
\end{remark}

\paragraph{Value interpolation at escaping nodes.}
Let $(\rho_i)$ be distinct nonzero points of $\K$ forming a radially finite
set, with
\begin{equation}\label{ent:rc:eq:nodes}
 \beta_i=-v(\rho_i)>0 ,
\end{equation}
enumerated, when infinite, with $\beta_1\le\beta_2\le\cdots$; each value
occurs finitely often and $\beta_i\to+\infty$ (\cref{ent:lem:divisor-countable}).
These are the $\beta_i$ of \cref{ent:prop:kernel-direct}. Let $P$ be the
canonical product of $\sum_i[\rho_i]$, and put
\begin{equation}\label{ent:rc:eq:cofactor}
 \widehat P_i=\prod_{j\ne i}\Bigl(1-\frac Z{\rho_j}\Bigr),\qquad
 \varpi_i=\widehat P_i(\rho_i)\ne0 ,
\end{equation}
the node cofactor and its value, the canonical product of the divisor with
$\rho_i$ removed; for one node per shell, $\widehat P_i$ is the cofactor
$\widehat P_n$ of \cref{ent:sm:subsec:cardinal} and $\varpi_i$ is the class
$\varpi_n$ there.

\begin{lemma}[The exact denominator valuation]\label{ent:rc:lem:denominator}
With the nondecreasing enumeration,
\begin{equation}\label{ent:rc:eq:denominator}
 v(\varpi_i)=\sum_{\substack{j\ne i\\ \beta_j\le\beta_i}}
 \bigl(v(\rho_j-\rho_i)+\beta_j\bigr),
\end{equation}
a finite sum. If every $\rho_i$ lies in $\Kle$ and
$s_i=\#\{j\ne i:\beta_j\le\beta_i\}$, then
\begin{equation}\label{ent:rc:eq:denom-bound}
 -s_i\beta_i\le v(\varpi_i)\le s_i\beta_i .
\end{equation}
\end{lemma}
\begin{proof}
The sum is finite by radial finiteness. Write $\widehat P_i$ as the finite
product over $j\ne i$ with $\beta_j\le\beta_i$ times the tail product $T_i$
over $\beta_j>\beta_i$ (\cref{ent:thm:product}); evaluation at $\rho_i$ is
multiplicative. The coefficient $c_k$ of $Z^k$ in $T_i$ is a strong sum of
signed products of $k$ distinct $1/\rho_j$ with $\beta_j>\beta_i$, so
$v(c_k\rho_i^k)>0$ for $k\ge1$, and $T_i(\rho_i)=1+\sum_{k\ge1}c_k\rho_i^k$
has valuation $0$: a strong sum of terms of positive valuation has positive
valuation. Each remaining factor
contributes $v(1-\rho_i/\rho_j)=v(\rho_j-\rho_i)-v(\rho_j)$, which gives
\eqref{ent:rc:eq:denominator}. If both nodes lie in $\Kle$, their nonzero
difference lies there too and has valuation at most $0$, while the
ultrametric inequality gives $v(\rho_j-\rho_i)\ge-\beta_i$ when
$\beta_j\le\beta_i$. Each term of \eqref{ent:rc:eq:denominator} therefore lies
between $\beta_j-\beta_i\ge-\beta_i$ and $\beta_j\le\beta_i$, which proves
\eqref{ent:rc:eq:denom-bound}.
\end{proof}

\begin{lemma}[Scale-moderate value interpolation]\label{ent:rc:lem:interpolation}
Let $(\rho_i)$ be as in \eqref{ent:rc:eq:nodes} and let $b_i\in\K$. Suppose
that for all but finitely many $i$ there is an integer $c_i\ge0$ with
\begin{equation}\label{ent:rc:eq:moderate}
 v(b_i/\varpi_i)\ge-c_i\beta_i .
\end{equation}
Then some $G\in\E$ satisfies $G(\rho_i)=b_i$ for every $i$.
\end{lemma}
\begin{proof}[Proof, as a special case of \cref{ent:thm:image}; the merge's]
A finite set of nodes is handled by Lagrange interpolation. Otherwise the
divisor $D=\sum_i[\rho_i]$ is infinite and radially finite with empty finite
part, since $v(\rho_i)<0$. Take its shells (\cref{ent:lem:shells}): radii
$\gamma_n$, nodes $\rho_{n,j}=\tau_n\xi_{n,j}$ with $\xi_{n,j}$ units of $\OO$.
For value data the normalized shell Hermite polynomial of
\eqref{ent:eq:normalizedjets} is the Lagrange polynomial
\[
 R_n(T)=\sum_j\frac{b_{n,j}}{\varpi^\circ_{n,j}}\prod_{k\ne j}(T-\xi_{n,k}),
 \qquad \varpi^\circ_{n,j}=\prod_{k\ne j}(\xi_{n,j}-\xi_{n,k}),
\]
whose products have coefficients in $\OO\subseteq V_n$. For the node
$\rho_i=\rho_{n,j}$, the terms of \eqref{ent:rc:eq:denominator} with
$\beta_l<\gamma_n$ are $\beta_l-\gamma_n$, of absolute value at most $\gamma_n$,
and those with $\beta_l=\gamma_n$ are $v(\rho_{n,k}-\rho_{n,j})+\gamma_n=
v(\xi_{n,k}-\xi_{n,j})$. Hence $v(\varpi_i)=v(\varpi^\circ_{n,j})+h_i$ with
$|h_i|\le s_i\gamma_n$, and \eqref{ent:rc:eq:moderate} gives
$v(b_{n,j}/\varpi^\circ_{n,j})\ge-(c_i+s_i)\gamma_n$, that is,
$b_{n,j}/\varpi^\circ_{n,j}\in V_n$. So $R_n\in V_n[T]$ on every shell that contains
no exceptional node, which is all but finitely many shells, and
\cref{ent:thm:image} supplies $G$; over $\R$ use \cref{ent:prop:real-form}.
\end{proof}
\begin{proof}[Second proof: the seventh source's explicit series]
Put $N_i=2c_i+2$ for the nonexceptional $i$ and
\begin{equation}\label{ent:rc:eq:cardinal}
 G(Z)=\sum_{i\text{ exceptional}}\frac{b_i}{\varpi_i}\widehat P_i(Z)
 +\sum_{i\text{ not exceptional}}
 \frac{b_i}{\varpi_i}\,\widehat P_i(Z)\Bigl(\frac Z{\rho_i}\Bigr)^{N_i}.
\end{equation}
Fix $\gamma$. Once $\beta_i+\gamma>0$, the factor $1-Z/\rho_i$ has weighted
value $0$, so $m_\gamma(\widehat P_i)=m_\gamma(P)=:A_\gamma$ by
\cref{ent:rc:lem:gauss}. Once also $\beta_i>2\abs\gamma$, the $i$th summand has
weighted value at least
\[
 A_\gamma-c_i\beta_i+(2c_i+2)(\beta_i+\gamma)
 \ \ge\ A_\gamma-c_i\beta_i+(c_i+1)\beta_i=A_\gamma+\beta_i ,
\]
which tends to $+\infty$. \Cref{ent:rc:lem:sum} makes $G$ entire and allows
termwise evaluation. At $\rho_j$ every summand with $i\ne j$ vanishes and the
$j$th equals $b_j$.
\end{proof}

The exponent $N_i$ may grow without bound with $i$; no fixed Archimedean
class need contain all the losses, only an integer multiple of each node's
own scale. For one node per shell, \eqref{ent:rc:eq:cardinal} is the
cardinal series of \cref{ent:sm:rem:simple-cardinal}; with several nodes on a
shell it corrects node by node, rather than shell by shell as
\cref{ent:sm:subsec:cardinal} does. As the source says, the lemma is a
sufficient value-interpolation condition with an explicit construction, not
the exact image theorem for arbitrary infinite jets; the exact image is
\cref{ent:thm:image}.

\begin{corollary}[Two sufficient conditions]\label{ent:rc:cor:interp-cases}
In \cref{ent:rc:lem:interpolation} the data $b_i$ can be arbitrary if all
sufficiently late $\beta_i$ are order units. Without an order unit, they can
be interpolated whenever every $\rho_i$ lies in $\Kle$ and
$v(b_i)\ge-\beta_i$ for all $i$; one may take $c_i=s_i+1$.
\end{corollary}
\begin{proof}
An order unit bounds the absolute value of each fixed element of $\Gamma$ by
an integer multiple of itself, hence bounds each loss $-v(b_i/\varpi_i)$
separately; the first assertion is also \cref{ent:cor:order-unit-image}. For
the second, \eqref{ent:rc:eq:denom-bound} gives
$v(b_i/\varpi_i)\ge-\beta_i-s_i\beta_i$.
\end{proof}

\subsection{Projections, shears and omnific rectification}
\label{ent:rc:subsec:rectification}

\begin{lemma}[Radial finiteness forces countability]\label{ent:rc:lem:countable}
An infinite radially finite subset of $\K^d$ is countable and forces
$\cf(\Gamma)=\aleph_0$. If $\cf(\Gamma)>\aleph_0$, every radially finite set
is finite and $\Ed=\K[X_1,\ldots,X_d]$.
\end{lemma}
\begin{proof}
This is the argument of \cref{ent:lem:divisor-countable} with $v(p)$ in place
of $v(\rho)$. Choose distinct nonzero $p_n$ in the set. If
$\{-v(p_n)\}$ had an upper bound $B$, the polydisk $B^d_{-B}$ would contain
all of them; so it is a countable cofinal subset of $\Gamma$. The polydisks
$B^d_{v(p_n)}$ cover $\K^d$ and each meets the set finitely, so the set is
countable. The polynomial assertion is the last sentence of
\cref{ent:thm:multi-criterion}.
\end{proof}

\begin{lemma}[Boundedness under entire maps]\label{ent:rc:lem:boundedness}
An entire map sends every valuation polydisk into some valuation polydisk.
An entire automorphism $\Phi$ preserves and reflects radial finiteness:
$\mathcal X$ is radially finite if and only if $\Phi(\mathcal X)$ is.
\end{lemma}
\begin{proof}
For $v(p)\ge\gamma$ each coordinate satisfies $v(F_j(p))\ge m_\gamma(F_j)$,
since $v(a_\alpha p^\alpha)\ge v(a_\alpha)+|\alpha|\gamma$; take the minimum of
these finitely many values. The points of $\mathcal X$ whose images lie in
$B^d_\gamma$ are mapped by $\Phi^{-1}$ into one polydisk, so there are finitely
many when $\mathcal X$ is radially finite and $\Phi$ is injective. Apply the
same argument to $\Phi^{-1}$.
\end{proof}

\begin{lemma}[A simultaneous valuation-preserving projection]
\label{ent:rc:lem:projection}
Let $\mathcal X\subseteq\K^d$ be countable. There is $c\in\kk$ such that the
$\kk$-linear functional $\ell_c(X)=X_1+cX_2+\cdots+c^{d-1}X_d$ satisfies
\begin{equation}\label{ent:rc:eq:project}
 v(\ell_c(p))=v(p)\quad(p\in\mathcal X\setminus\{0\}),\qquad
 v(\ell_c(p)-\ell_c(q))=v(p-q)\quad(p\ne q\text{ in }\mathcal X).
\end{equation}
In particular $\ell_c$ is injective on $\mathcal X$, and it is the first
coordinate of $g_c(X)=(\ell_c(X),X_2,\ldots,X_d)$, which lies in
$\operatorname{GL}_d(\kk)$ with determinant $1$.
\end{lemma}
\begin{proof}
For a nonzero $w\in\K^d$ with $\lambda=v(w)$, the coefficients of $t^\lambda$
in its coordinates form a nonzero vector $\lc(w)\in\kk^d$, and
$v(\ell(w))=\lambda$ for a $\kk$-linear $\ell$ exactly when
$\ell(\lc(w))\ne0$. For $\ell=\ell_c$ this is the nonvanishing
at $c$ of a nonzero polynomial of degree less than $d$, which excludes finitely
many $c$. Apply this to the countably many nonzero points of $\mathcal X$ and
nonzero differences of two of its points; the union of the excluded sets is
countable, and $\kk$ is uncountable.
\end{proof}

This uses only the uncountability of the constant field, not a probabilistic
choice or a compactness argument. Avoiding the leading coefficients of the
\emph{differences} matters when two points share a shell.

\begin{lemma}[All projected nodes beyond one scale]\label{ent:rc:lem:shift}
Let $\mathcal X$ be radially finite, take $g_c$ from
\cref{ent:rc:lem:projection}, and let $\delta\in\Gamma$ be positive. After
$g_c$ and a translation of the first coordinate, each $p\in\mathcal X$ becomes
a point
\[
 (\rho_p,b_{p,2},\ldots,b_{p,d}),\qquad
 \beta_p=-v(\rho_p)\ge\delta,\qquad v(b_{p,j})\ge-\beta_p ,
\]
with the $\rho_p$ distinct and nonzero. The set $\{\rho_p\}$ is radially
finite, so when $\mathcal X$ is infinite it can be enumerated as in
\eqref{ent:rc:eq:nodes}. If $\mathcal X\subseteq(\Kle)^d$, all these
coordinates lie in $\Kle$.
\end{lemma}
\begin{proof}
Replace $\ell_c(p)$ by $\rho_p=\ell_c(p)+c't^{-\delta}$ with $c'\in\kk^\times$.
Only finitely many points of $\mathcal X$ have $v(p)=-\delta$; choose $c'$ so
that none of their leading coefficients cancels. Then, for every $p$ including
$p=0$, $v(\rho_p)=\min\{v(p),-\delta\}$, which is at most $-\delta$; so
$\rho_p\ne0$, and the untouched coordinates $b_{p,j}=p_j$ have valuation at
least $v(p)\ge-\beta_p$. The translation preserves injectivity. Since
$\{p:-\beta_p\ge\gamma\}\subseteq\mathcal X\cap B^d_\gamma$, the nodes are
radially finite. Finally $\Kle$ is closed under $\kk$-linear combinations and
contains $c't^{-\delta}$.
\end{proof}

\begin{proposition}[Rectification by a single shear]\label{ent:rc:prop:positive}
If $\Gamma$ has an order unit, every radially finite subset of $\K^d$ is
line-rectifiable. Without that assumption, every radially finite subset of
$(\Kle)^d$ is line-rectifiable. In both cases an invertible affine map
followed by one triangular entire shear suffices.
\end{proposition}
\begin{proof}
A finite set is separated by a constant projection, and finite Lagrange
polynomials in the first coordinate interpolate its remaining coordinates; the
corresponding polynomial shear straightens it. Assume the set infinite, hence
countable (\cref{ent:rc:lem:countable}), and choose the coordinates of
\cref{ent:rc:lem:shift}. In the order-unit case take $\delta$ to be an order
unit; then every $\beta_p\ge\delta$ is an order unit, and the first assertion
of \cref{ent:rc:cor:interp-cases} supplies, for $j=2,\ldots,d$, entire $G_j$
with $G_j(\rho_p)=b_{p,j}$. In the polynomial-part case use the second
assertion of \cref{ent:rc:cor:interp-cases}, with the bounds of
\cref{ent:rc:lem:shift}. The map
\begin{equation}\label{ent:rc:eq:shear}
 \Psi(X_1,X_2,\ldots,X_d)=\bigl(X_1,\ X_2-G_2(X_1),\ \ldots,\
 X_d-G_d(X_1)\bigr)
\end{equation}
is an entire automorphism, with inverse obtained by changing each minus sign
to a plus sign, and it sends each point to $(\rho_p,0,\ldots,0)$. Compose it
with the affine map.
\end{proof}

\begin{corollary}[Distinct-radius configurations]
\label{ent:rc:cor:distinct-radii}
Suppose $\mathcal X$ is radially finite and all but finitely many of its
points have pairwise distinct values of $v(p)$. Then $\mathcal X$ is
line-rectifiable, even if its coordinates have infinitesimal tails and
$\Gamma$ has no order unit.
\end{corollary}
\begin{proof}
Choose the coordinates of \cref{ent:rc:lem:shift} for some $\delta>0$. Then
$\beta_p=\max\{-v(p),\delta\}$, so the finitely many points with
$v(p)\ge-\delta$ collapse to the radius $\delta$; together with the finitely
many points whose radius is not distinct and the finitely many whose radius
equals one of theirs, they form a finite exceptional set $E$. For
$p\notin E$, every other node $\rho_q$ with $\beta_q\le\beta_p$ has
$\beta_q<\beta_p$, so $v(\rho_q-\rho_p)=-\beta_p$ and the corresponding term of
\eqref{ent:rc:eq:denominator} is $\beta_q-\beta_p<0$. Thus
$v(\varpi_p)\le0$, and $v(b_{p,j}/\varpi_p)\ge v(b_{p,j})\ge-\beta_p$: take
$c_p=1$. The nodes in $E$ are the finitely many exceptions allowed by
\cref{ent:rc:lem:interpolation}, which supplies the $G_j$; the shear
\eqref{ent:rc:eq:shear} applies.
\end{proof}

The source's proof says that after the shift ``the late nodes have distinct
radii''; the collapse of the points with $v(p)\ge-\delta$ onto one radius is
covered by the exceptions allowed in \cref{ent:rc:lem:interpolation}, and the
proof above makes that step explicit. The corollary locates the obstruction:
escaping points alone do not cause the failure without an order unit, but
sufficiently small \emph{clusters on one shell} do. In one variable that is
the lesson of \cref{ent:thm:close-pair}.

\paragraph{Translation to surreal and omnific numbers.}
Let $\Gamma\subseteq\No$ be a divisible ordered subgroup and apply
\eqref{ent:eq:embedding}, $t^\gamma\mapsto\omega^{-\gamma}$: the real workspace
$K_{\R}=\R((t^\Gamma))$ becomes a normal-form subfield of $\No$, the complex
one its complexification in $\No[i]$. This is the standard Hahn reading of
normal forms \cite{Conway,Gonshor,MantovaMatusinski}, not a new construction.
An omnific integer is a normal form whose exponents are nonnegative and whose
constant term is an integer, $\Oz=\Pi\oplus\Z$ with $\Pi$ the purely infinite
normal forms (\texttt{found:eq:omnific}, \texttt{odg:def:rings}). Exponent
$y=-\gamma\ge0$ means $\gamma\le0$, so inside the embedded fields
\begin{align}
 \Oz\cap K_{\R}&=\Z+\R((t^{\Gamma_{<0}})),\label{ent:rc:eq:omnific-real}\\
 (\Oz+i\Oz)\cap K_{\C}&=\Z[i]+\C((t^{\Gamma_{<0}})).
 \label{ent:rc:eq:omnific-complex}
\end{align}
This is the polynomial-part framework of the factorization literature
\cite{LInnocenteMantova}; the ring $\Oz+i\Oz$ is what the omnific
Diophantine-geometry report calls the Gaussian omnific integers, and
$\R\oplus\Pi$ intersected with $K_{\R}$ is the real $\Kle$, the ring that
report writes $\mathcal B_{\R,H}$ for a workspace (\texttt{odg:lem:workspace}).
Both rings in \eqref{ent:rc:eq:omnific-real}--\eqref{ent:rc:eq:omnific-complex}
lie in the corresponding $\Kle$, so \cref{ent:rc:prop:positive} proves
\cref{ent:rc:thm:omnific-main}.

\begin{remark}[Three indispensable restrictions]\label{ent:rc:rem:restrictions}
First, radial finiteness is measured inside the chosen workspace: the ordinary
integers, although omnific, form an infinite valuation-bounded set and are not
radially finite. Second, the affine projection may use real or complex
constants and the shear may have arbitrary Hahn coefficients, so the ambient
automorphism need not preserve the omnific ring. Third, these results concern
coordinate changes of $\K^d$, not ring automorphisms of $\Oz$ or field
automorphisms of $\No$.
\end{remark}

\subsection{Eventual polynomial reduction of entire automorphisms}
\label{ent:rc:subsec:reduction}

The negative half of \cref{ent:rc:thm:main} rests on a restriction that every
fixed entire map eventually satisfies: increasing valuation rank does not give
one map unlimited local freedom at every new scale.

\paragraph{Coarsened valuation rings.}
For a convex subgroup $H\subsetneq\Gamma$ put
\[
 V_H=\{x\in\K:v(x)\ge h\text{ for some }h\in H\},\qquad
 \mathfrak q_H=\{x\in\K:v(x)>h\text{ for every }h\in H\},
\]
with $0$ in both. As in \cref{ent:lem:coarse-residue}, whose proof uses only
the convexity of $H_\gamma$, $V_H$ is the valuation ring of $v$ coarsened by
$H$, with maximal ideal $\mathfrak q_H$, and restriction of the Hahn support
to $H$ is a surjective ring map
\[
 \red_H:V_H\longrightarrow\kappa_H=\kk((t^H))
\]
with kernel $\mathfrak q_H$; for $H=H_\gamma$ these are $V_\gamma$,
$\mathfrak q_\gamma$ and $\kappa_\gamma$ of \cref{ent:def:coarse}. Both $V_H$
and $\mathfrak q_H$ are closed under strong sums: a nonzero strong sum has
valuation at least the least valuation of a term, and that least valuation is
attained. Reduction is applied coefficientwise to series and coordinatewise to
points and matrices.

\begin{theorem}[Eventual polynomial reduction]\label{ent:rc:thm:reduction}
Let $H_1\subseteq H_2\subseteq\cdots$ be proper convex subgroups of $\Gamma$
with $\bigcup_nH_n=\Gamma$, and let $\mathcal F\subseteq\Ed$ be finite. For
all sufficiently large $n$:
\begin{enumerate}[label=\textnormal{(\alph*)},leftmargin=2em]
\item every coefficient of every $F\in\mathcal F$ lies in $V_{H_n}$;
\item all but finitely many of them lie in $\mathfrak q_{H_n}$, so that
coefficientwise reduction gives a polynomial
$\overline F_n\in\kappa_{H_n}[X_1,\ldots,X_d]$;
\item for every $p\in V_{H_n}^d$, $F(p)\in V_{H_n}$ and
\begin{equation}\label{ent:rc:eq:reduce-eval}
 \red_{H_n}\bigl(F(p)\bigr)=\overline F_n\bigl(\red_{H_n}p\bigr);
\end{equation}
\item (a)--(c) hold, for the same $n$, for every Hasse derivative of every
$F\in\mathcal F$, and reduction commutes with Hasse derivatives.
\end{enumerate}
\end{theorem}
\begin{proof}
For nonzero $F$, $m_0(F)=\min_\alpha v(a_\alpha)$ exists; let $L$ be the least of
these values over $\mathcal F$. Once $L\in H_n$, every coefficient lies in
$V_{H_n}$. Fix such an $n$. Since $H_n$ is proper and convex, there is
$\eta\in\Gamma$ with $\eta>H_n$, and by \cref{ent:thm:multi-criterion}(iii)
at $\gamma=0$ only finitely many coefficients have valuation at most $\eta$;
the others lie in $\mathfrak q_{H_n}$ and reduce to $0$. For $p\in
V_{H_n}^d$, split the strongly summable family $(a_\alpha p^\alpha)_\alpha$
into the finitely many terms with $v(a_\alpha)\le\eta$ and the rest
(\cref{ent:mv:lem:tails}(b)). Every term lies in $V_{H_n}$, and every term of
the rest lies in the ideal $\mathfrak q_{H_n}$; by closure under strong sums,
the rest sums into $\mathfrak q_{H_n}$ and $F(p)\in V_{H_n}$. Since $\red_{H_n}$
is a ring map, \eqref{ent:rc:eq:reduce-eval} follows from the finite part. A
Hasse derivative multiplies $a_\alpha$ by a nonzero integer, which does not
change its valuation, and shifts the index; its coefficients therefore obey
the same bounds $L$ and $\eta$, and its reduction is the corresponding Hasse
derivative of $\overline F_n$.
\end{proof}

The degree of $\overline F_n$ may grow with $n$: the theorem asserts a
polynomial reduction at each sufficiently late scale, not one polynomial for
all scales. No coefficient outside $H_n$ is moved into $H_n$; such exponents
vanish under the residue map. The source's proof of (c) also chooses an
auxiliary element $\varrho\in H_n$ below the coordinate valuations of $p$;
closure of $\mathfrak q_{H_n}$ under strong sums makes that step unnecessary,
and it is omitted here.

\paragraph{Antecedents.} The device is not new to this collection. In one
variable, the necessity half of \cref{ent:thm:units} reduces a restricted
series modulo a proper convex subgroup to a polynomial over the coarsened
residue field, and \cref{ent:thm:multi-unit} does the same in finitely many
variables; \cref{ent:lem:coarse-residue} is the residue map. The report
\emph{Holonomic Rigidity for Entire Hahn Functions} proves the one-variable
entire case, \texttt{hol:cf:lem:reduction} \cite{RepoHolonomic}: an entire
series with coefficients in the coarsened ring of a proper convex subgroup
reduces to a polynomial. \Cref{ent:rc:thm:reduction} is the several-variable
form of that lemma, with the evaluation compatibility (c), Hasse derivatives,
and a finite family treated at once; its application below to an automorphism
together with its inverse is the seventh source's. The seventh source cites
none of these antecedents; the credit is the merge's.

\begin{corollary}[Inverse polynomial reductions]
\label{ent:rc:cor:automorphism-reduction}
Let $\Phi$ be an entire automorphism of $\K^d$ and let $(H_n)$ be as in
\cref{ent:rc:thm:reduction}. For all sufficiently large $n$, $\Phi$ and
$\Phi^{-1}$ map $V_{H_n}^d$ into itself and reduce to mutually inverse
polynomial automorphisms of $\kappa_{H_n}^d$, and for every
$p\in V_{H_n}^d$
\begin{equation}\label{ent:rc:eq:J-unit}
 \Jac\Phi(p)\in\operatorname{GL}_d(V_{H_n}),
\end{equation}
where $\Jac\Phi$ is the Jacobian matrix $(\partial\Phi_i/\partial X_j)$.
\end{corollary}
\begin{proof}
Apply \cref{ent:rc:thm:reduction} to the $2d$ coordinates of $\Phi$ and
$\Phi^{-1}$. By (c), both maps preserve $V_{H_n}^d$, and reducing
$\Phi^{-1}(\Phi(p))=p$ and $\Phi(\Phi^{-1}(p))=p$ gives the two composite
polynomial identities at every residue point, since $\red_{H_n}$ is surjective.
The residue field is infinite, so these are identities of polynomials. By (d),
the entries of $\Jac\Phi(p)$ and of $\Jac\Phi^{-1}(\Phi(p))$ lie in
$V_{H_n}$, and by the chain rule (\cref{ent:rc:lem:composition}) the second
matrix is the inverse of the first.
\end{proof}

\begin{lemma}[Integral Taylor remainder]\label{ent:rc:lem:taylor-remainder}
Let $F:\K^d\to\K^d$ be an entire map with all coefficients in $V_H$ for a
proper convex subgroup $H$. For $p,\mathbf w\in V_H^d$ and
$\varepsilon\in\mathfrak q_H$ there is $R\in V_H^d$ with
\begin{equation}\label{ent:rc:eq:taylor}
 F(p+\varepsilon\mathbf w)=F(p)+\varepsilon\,\Jac F(p)\,\mathbf w
 +\varepsilon^2R .
\end{equation}
\end{lemma}
\begin{proof}
By \cref{ent:rc:lem:composition}, $F_i(p+Y)=\sum_\beta c_{i,\beta}Y^\beta$ with
$c_{i,\beta}$ the Hasse derivatives of $F_i$ at $p$. Each $c_{i,\beta}$ is a
strong sum of products of elements of $V_H$, so lies in $V_H$. The family
$(c_{i,\beta}\varepsilon^{|\beta|}\mathbf w^\beta)_\beta$ is strongly summable;
group it by $|\beta|=0$, $|\beta|=1$ and $|\beta|\ge2$
(\cref{ent:mv:lem:tails}(b)). Dividing the last group termwise by the fixed
nonzero $\varepsilon^2$ (the case $\varepsilon=0$ is trivial) shifts all
supports by one element, so it remains strongly summable, and its terms
$c_{i,\beta}\varepsilon^{|\beta|-2}\mathbf w^\beta$ lie in $V_H$; so does their
sum $R_i$.
\end{proof}

\begin{proposition}[Infinitesimal simplices survive good reduction]
\label{ent:rc:prop:simplex}
Let $H\subsetneq\Gamma$ be convex, and let $\Phi$ be an entire automorphism
such that $\Phi$ and $\Phi^{-1}$ have all coefficients in $V_H$. For every
$p\in V_H^d$ and every nonzero $\varepsilon\in\mathfrak q_H$ the $d+1$ points
\[
 \Phi(p),\quad\Phi(p+\varepsilon\mathbf u_1),\quad\ldots,\quad
 \Phi(p+\varepsilon\mathbf u_d)
\]
are affinely independent. More precisely,
\begin{equation}\label{ent:rc:eq:simplexdet}
 \det\bigl[\Phi(p+\varepsilon\mathbf u_j)-\Phi(p)\bigr]_{j=1}^d
 =\varepsilon^du,\qquad u\in V_H^\times .
\end{equation}
\end{proposition}
\begin{proof}
By \cref{ent:rc:lem:taylor-remainder}, the matrix of the differences divided by
$\varepsilon$ is $\Jac\Phi(p)+\varepsilon R$ with $R$ integral, and its
reduction is that of $\Jac\Phi(p)$. As in the proof of
\cref{ent:rc:cor:automorphism-reduction}, which uses only the integrality of
the coefficients of $\Phi$ and $\Phi^{-1}$, $\Jac\Phi(p)$ is invertible over
$V_H$; so the determinant of $\Jac\Phi(p)+\varepsilon R$ is a unit $u$ of
$V_H$. Multiplying each of the $d$ columns by $\varepsilon$ gives
\eqref{ent:rc:eq:simplexdet}, which is nonzero.
\end{proof}

The \emph{inverse} is essential here. Integrality of the forward map alone
would not make its Jacobian a coarsened unit, and bijectivity in some
unrelated function category would not supply the integral inverse series.
The obstruction below uses from \cref{ent:rc:thm:reduction} only the
integrality clause (a) and the Jacobian consequence
\eqref{ent:rc:eq:J-unit}; the polynomiality of the reductions describes what a
fixed automorphism looks like at a late scale and is not needed for the
determinant. That dependency remark is the merge's.

\subsection{The escaping-simplex obstruction and the dichotomy}
\label{ent:rc:subsec:obstruction}

\paragraph{Scales.} Suppose $\cf(\Gamma)=\aleph_0$ and $\Gamma$ has no order
unit. \Cref{ent:lem:escaping-classes} gives a positive cofinal sequence with
$\gamma_{n+1}\ggA\gamma_n$. Put $H_n=H_{\gamma_n}$ (\cref{ent:def:coarse}).
Each $H_n$ is a convex subgroup, proper because $\gamma_n$ is not an order
unit; the $H_n$ increase and exhaust $\Gamma$ because $(\gamma_n)$ is
cofinal; and $\gamma_{n+1}>H_n$, since $\gamma_{n+1}$ exceeds every integer
multiple of $\gamma_n$. The seventh source proves this as its own lemma on
separated cofinal scales; it is printed once, here.

\begin{theorem}[A universal escaping-simplex obstruction]
\label{ent:rc:thm:obstruction}
With these scales let
\begin{equation}\label{ent:rc:eq:Dexplicit}
 p_n=t^{-\gamma_n}\mathbf u_1,\qquad\varepsilon_n=t^{\gamma_{n+1}},\qquad
 \mathcal X=\bigcup_{n\ge1}\{p_n,\ p_n+\varepsilon_n\mathbf u_1,\ \ldots,\
 p_n+\varepsilon_n\mathbf u_d\}.
\end{equation}
Then $\mathcal X$ is radially finite, and for every entire automorphism
$\Phi$ all sufficiently late blocks of $\Phi(\mathcal X)$ are affinely
independent. Neither $\mathcal X$ nor any cofinite subset of it can be carried
into an affine hyperplane by an entire automorphism.
\end{theorem}
\begin{proof}
Every point of the $n$th block has $v=-\gamma_n$: the first coordinate has
leading term $t^{-\gamma_n}$, and each perturbation has positive valuation.
The $\gamma_n$ increase cofinally, so a fixed polydisk meets finitely many
blocks, each finite; distinct blocks are disjoint because their valuations
differ, and the $d+1$ points of a block are distinct. For a fixed $\Phi$,
\cref{ent:rc:thm:reduction}(a) applied to the exhaustion $(H_n)$ gives an
$n_0$ such that $\Phi$ and $\Phi^{-1}$ have coefficients in $V_{H_n}$ for
$n\ge n_0$. The point $p_n$ lies in $V_{H_n}^d$ since
$-\gamma_n\in H_n$, and $\varepsilon_n$ is a nonzero element of
$\mathfrak q_{H_n}$ since $\gamma_{n+1}>H_n$. \Cref{ent:rc:prop:simplex}
applies. An affine hyperplane contains no $d+1$ affinely independent points,
and deleting finitely many points leaves every sufficiently late block intact.
\end{proof}

\begin{proof}[Proof of \cref{ent:rc:thm:main}]
If $\Gamma$ has an order unit, \cref{ent:rc:prop:positive} applies. If not,
\cref{ent:rc:thm:obstruction} gives a radially finite counterexample; for
$d\ge2$ a coordinate line lies in an affine hyperplane, so it is not
line-rectifiable. This proves the equivalence and the stronger negative
assertion.
\end{proof}

\paragraph{An explicit surreal realization.}
In the group $\Gamma_\infty=\bigoplus_{j\ge0}\Q e_j$, $e_j=\omega^j$, of
\eqref{ent:eq:concrete-group}, every $e_{j+1}$ dominates every multiple of
$e_j$, and no element is an order unit. The proof of
\cref{ent:rc:thm:obstruction} needs only $p_n\in V_{H_n}^d$ and
$\varepsilon_n\in\mathfrak q_{H_n}\setminus\{0\}$, so the perturbation
exponent need not be the next centre scale. The seventh source takes the
centre scales $e_{2n}$ and perturbation exponents $\vartheta_n=e_{2n+1}$. Under
$t^\gamma=\omega^{-\gamma}$ its points are
\begin{equation}\label{ent:rc:eq:surreal-D}
 p_n=\omega^{\omega^{2n}}\mathbf u_1,\qquad
 \varepsilon_n=\omega^{-\omega^{2n+1}},\qquad
 p_n,\ p_n+\varepsilon_n\mathbf u_1,\ \ldots,\ p_n+\varepsilon_n\mathbf u_d ,
\end{equation}
where $\mathbf u_j$ are coordinate vectors and the exponents are already
surreal. Each centre is an omnific integer in its nonzero coordinate. The
added $\varepsilon_n$ is a genuine infinitesimal, with a negative normal-form
exponent, so the perturbed points leave the polynomial-part ring. The centres
alone lie on a line; the arbitrarily small full-dimensional clusters cannot all
be made collinear by any entire automorphism of this workspace.

\begin{corollary}[Rectifiability is not stable under cluster refinement]
\label{ent:rc:cor:refinement}
In $K_{\Gamma_\infty}^d$, $d\ge2$, a line-rectifiable sequence of omnific
centres admits finite infinitesimal refinements, all on the same centre
shells, whose union is not line-rectifiable and cannot be sent into an affine
hyperplane by an entire automorphism.
\end{corollary}
\begin{proof}
The centres $p_n$ of \eqref{ent:rc:eq:surreal-D} already lie on the first
coordinate line. The refinement \eqref{ent:rc:eq:surreal-D} satisfies the
hypotheses of the proof of \cref{ent:rc:thm:obstruction} with
$H_n=H_{e_{2n}}$, since $e_{2n+1}>H_n$, and that proof applies verbatim.
\end{proof}

The source states this corollary without a label or proof; the two-line proof
is the merge's. It concerns refinement by \emph{clusters}, not replacement of
each centre by a single perturbed point; that single-node situation can stay
rectifiable by \cref{ent:rc:cor:distinct-radii}.

\begin{remark}[The one-variable shadow of the simplex; the merge's]
\label{ent:rc:rem:close-pair}
With the theorem's own choice $\gamma_n=e_n$ in $\Gamma_\infty$, the first two
points of the $n$th block of \eqref{ent:rc:eq:Dexplicit} lie on the first
coordinate axis at $\rho_n=\omega^{\omega^n}$ and
$\sigma_n=\omega^{\omega^n}+\omega^{-\omega^{n+1}}$, the close pair
\eqref{ent:eq:surreal-roots} behind \cref{ent:thm:no-bezout}. In one variable
the close pair defeats interpolation (\cref{ent:thm:close-pair}); in $d\ge2$
the full simplex on the same shells defeats rectification. The mechanisms are
different --- a valuation bound on values, and a coarsened Jacobian unit ---
and neither result is derived from the other.
\end{remark}

\paragraph{The three cases.} The first row is finite geometry together with
coefficient collapse (\cref{ent:rc:lem:countable}); the distinction the seventh
source adds is between the second and third rows. The same three rows govern
the one-variable trichotomy of \cref{ent:thm:main}.

\begin{center}
\small
\begin{tabular}{@{}>{\raggedright\arraybackslash}p{0.25\textwidth}>{\raggedright\arraybackslash}p{0.25\textwidth}>{\raggedright\arraybackslash}p{0.4\textwidth}@{}}
\toprule
Value group & Radially finite sets & Entire rectification in $d\ge2$\\
\midrule
$\cf(\Gamma)>\aleph_0$ & every such set is finite & polynomial
rectification always exists\\[3pt]
an order unit exists & infinite escaping sets occur & every radially finite
set is rectifiable\\[3pt]
$\cf(\Gamma)=\aleph_0$, no order unit & infinite escaping sets occur &
polynomial-part sets are rectifiable; escaping infinitesimal simplices
obstruct every entire automorphism\\
\bottomrule
\end{tabular}
\end{center}

\subsection{Fat-point ideals on rectifiable configurations}
\label{ent:rc:subsec:ideals}

The rectification theorems turn a geometric configuration into an explicit
algebraic ideal. Let $\mathcal X=\{p_i\}$ be nonempty, radially finite and
line-rectifiable. By \cref{ent:rc:lem:boundedness} its image on the
coordinate line is radially finite, and a translation of the first coordinate
as in \cref{ent:rc:lem:shift} makes the first coordinates nonzero with
positive radii. By \cref{ent:rc:lem:order-invariant} below it therefore
suffices to treat
\begin{equation}\label{ent:rc:eq:line-D}
 p_i=(\rho_i,0,\ldots,0),\qquad\beta_i=-v(\rho_i)>0 ,
\end{equation}
with the nodes as in \eqref{ent:rc:eq:nodes}; finite configurations use the
same notation and finite products. Write $X'=(X_2,\ldots,X_d)$ for the
transverse variables and $X'^\alpha$, $\alpha\in\N^{d-1}$, for their
monomials. Assign positive integers $m_i$, and for $r\in\N$ let
$P_{\langle r\rangle}$ be the canonical product (\cref{ent:thm:product}) of
the divisor
\begin{equation}\label{ent:rc:eq:Pr}
 D_{\langle r\rangle}=\sum_i(m_i-r)_+[\rho_i],\qquad(u)_+=\max\{u,0\},
\end{equation}
with $P_{\langle r\rangle}=1$ when that divisor is zero. It is entire even
for unbounded multiplicities, and
$\ord_{\rho_i}P_{\langle r\rangle}=(m_i-r)_+$.

\begin{lemma}[Transverse coefficient criterion]\label{ent:rc:lem:transverse}
Write $F\in\Ed$ uniquely as
\[
 F=\sum_{\alpha\in\N^{d-1}}f_\alpha(X_1)X'^\alpha .
\]
Each $f_\alpha$ is entire, and the following are equivalent:
\begin{enumerate}[label=\textnormal{(\roman*)},leftmargin=2em]
\item $\ord_{p_i}F\ge m_i$ for all $i$;
\item $\ord_{\rho_i}f_\alpha\ge(m_i-|\alpha|)_+$ for all $\alpha$ and $i$;
\item $f_\alpha\in P_{\langle|\alpha|\rangle}\E$ for every $\alpha$.
\end{enumerate}
\end{lemma}
\begin{proof}
Selecting the coefficients with a fixed transverse index $\alpha$ preserves
\cref{ent:thm:multi-criterion}(iii), after subtracting the constant weight
$|\alpha|\gamma$; so $f_\alpha$ is entire. The Taylor expansion of $F$ at
$(\rho_i,0)$ has coefficients indexed by pairs $(q,\alpha)$, of total degree
$q+|\alpha|$, namely the Hasse jets of $f_\alpha$ at $\rho_i$; distinct pairs
do not interact. Vanishing in all total degrees below $m_i$ is therefore (ii).
The equivalence of (ii) and (iii) is the division statement cited in
\cref{ent:rc:subsec:calculus}, applied to each $f_\alpha$ and the divisor
$D_{\langle|\alpha|\rangle}$.
\end{proof}

\begin{lemma}[Entire changes preserve local order]
\label{ent:rc:lem:order-invariant}
If $\Phi$ is an entire automorphism, then
$\ord_a(F\circ\Phi)=\ord_{\Phi(a)}F$ for all $F\in\Ed$ and $a\in\K^d$.
Consequently pullback by $\Phi$ is a ring automorphism of $\Ed$ that maps the
fat-point ideal of $\Phi(\mathcal X)$ onto that of $\mathcal X$, with the same
multiplicities, and preserves minimal numbers of generators.
\end{lemma}
\begin{proof}
The Taylor expansion of $\Phi(a+U)-\Phi(a)$ has no constant term and has the
invertible linear part $\Jac\Phi(a)$ (chain rule for $\Phi^{-1}\circ\Phi$). If
the first nonzero homogeneous part of $F$ at $\Phi(a)$ has degree $m$, its
composite with that linear part is nonzero of degree $m$, and all other terms
have higher order; so the order cannot drop, and applying the argument to
$\Phi^{-1}$ shows it cannot rise. Pullback is a ring automorphism by
\cref{ent:rc:lem:composition}, so algebraic generating sets correspond.
\end{proof}

\begin{theorem}[The fat-point generators]\label{ent:rc:thm:fat-generators}
Suppose $\bar m=\max_im_i<\infty$. In the coordinates \eqref{ent:rc:eq:line-D},
\begin{equation}\label{ent:rc:eq:ideal-generators}
 I_{\mathbf m}(\mathcal X)=\bigl(P_{\langle|\alpha|\rangle}(X_1)\,X'^\alpha:
 \alpha\in\N^{d-1},\ |\alpha|\le\bar m\bigr).
\end{equation}
This family is minimal in cardinality, and
$\mu(I_{\mathbf m}(\mathcal X))=\binom{\bar m+d-1}{d-1}$. For a configuration
rectified by $\Phi$, pull the generators back through $\Phi$.
\end{theorem}
\begin{proof}
Each displayed generator has order $(m_i-|\alpha|)_++|\alpha|\ge m_i$ at
$p_i$, since orders add under products. Conversely let $F\in
I_{\mathbf m}(\mathcal X)$ and write
\begin{equation}\label{ent:rc:eq:transverse-truncate}
 F=\sum_{|\alpha|<\bar m}f_\alpha(X_1)X'^\alpha
 +\sum_{|\beta|=\bar m}X'^\beta H_\beta,\qquad H_\beta\in\Ed ,
\end{equation}
assigning each monomial of transverse degree at least $\bar m$ to one fixed
degree-$\bar m$ divisor $X'^\beta$ of it by any deterministic rule; after
division by that monomial the selected coefficient family has its weights
shifted by $-\bar m\gamma$, so it is entire, and there are finitely many
$\beta$. Thus \eqref{ent:rc:eq:transverse-truncate} holds in $\Ed$, not only
formally. By \cref{ent:rc:lem:transverse}, each $f_\alpha$ with
$|\alpha|<\bar m$ is divisible by $P_{\langle|\alpha|\rangle}$, and the
high-degree terms use the generators with $|\beta|=\bar m$, for which
$P_{\langle\bar m\rangle}=1$. This proves \eqref{ent:rc:eq:ideal-generators}.

For minimality choose $i$ with $m_i=\bar m$, and for $F$ in the ideal let
$\mathsf h_i(F)$ be the homogeneous degree-$\bar m$ part of its Taylor
expansion at $p_i$, all lower parts being zero. For $H\in\Ed$,
\begin{equation}\label{ent:rc:eq:leading-module}
 \mathsf h_i(HF)=H(p_i)\,\mathsf h_i(F).
\end{equation}
The generator indexed by $\alpha$ has $\mathsf h_i$ a nonzero multiple of
$(X_1-\rho_i)^{\bar m-|\alpha|}X'^\alpha$; as $\alpha$ ranges over
$|\alpha|\le\bar m$ these are all the monomials of degree $\bar m$ in $d$
variables, a basis of a space of dimension $\binom{\bar m+d-1}{d-1}$. By
\eqref{ent:rc:eq:leading-module}, the image of an ideal generated by $q$
elements is spanned by $q$ vectors, so $q$ is at least that dimension, which is
the size of the displayed family.
\end{proof}

\begin{corollary}[Simple points form an explicit complete intersection]
\label{ent:rc:cor:ci}
Let $P=P_{\langle0\rangle}$ for $m_i=1$, that is, $P=\prod_i(1-X_1/\rho_i)$.
Then
\begin{equation}\label{ent:rc:eq:simple-ideal}
 I(\mathcal X)=\{F:F(p_i)=0\ \forall i\}=(P,X_2,\ldots,X_d).
\end{equation}
The sequence $X_2,\ldots,X_d,P$ is regular, its common zero set in $\K^d$ is
exactly $\mathcal X$, and
\begin{equation}\label{ent:rc:eq:quotient}
 \Ed/I(\mathcal X)\cong\E/(P).
\end{equation}
The least number of generators is $d$.
\end{corollary}
\begin{proof}
Take $\bar m=1$ in \cref{ent:rc:thm:fat-generators}. Setting the last
variable to zero identifies $\mathcal E_{\Gamma,k}/(X_k)$ with
$\mathcal E_{\Gamma,k-1}$ (a series vanishing on $X_k=0$ is $X_k$ times an
entire series, by \cref{ent:thm:multi-criterion}), and each of these is a
domain (\cref{ent:rc:lem:gauss}). The quotient before $P$ is imposed is $\E$,
in which the nonzero $P$ is not a zero divisor, and the final ideal is proper,
since evaluation at a node annihilates it. The zero set comes from the exact
divisor of $P$ and the equations $X_j=0$; substitution $X'=0$ gives
\eqref{ent:rc:eq:quotient}; minimality is the case $\bar m=1$.
\end{proof}

The source adds a warning, kept here: the evaluation map from $\E/(P)$ to
$\prod_i\K$ is injective but need not be surjective without an order unit, so
\eqref{ent:rc:eq:quotient} asserts no unrestricted Chinese remainder theorem for
infinitely many nodes. The complete-intersection calculation uses vanishing
and division, not interpolation of arbitrary value sequences. The next
corollary combines \eqref{ent:rc:eq:quotient} with this report's image
theorem.

\begin{corollary}[Evaluation and interpolation on a rectifiable
configuration; the merge's]\label{ent:rc:cor:evaluation}
Let $\mathcal X=\{p_i\}\subseteq\K^d$ be nonempty, radially finite and
line-rectifiable, and let $\Phi$ be an entire automorphism with
$\Phi(p_i)=(\rho_i,0,\ldots,0)$, the $\rho_i$ as in \eqref{ent:rc:eq:nodes}.
Put $D=\sum_i[\rho_i]$.
\begin{enumerate}[label=\textnormal{(\alph*)},leftmargin=2em]
\item $\Ed/I(\mathcal X)\cong\E/(P_D)$, through
$F\mapsto(F\circ\Phi^{-1})(Z,0,\ldots,0)$; for infinite $\mathcal X$ this is
the restricted product $\RD$ of \cref{ent:thm:quotient}.
\item A family $(b_i)$ in $\K$ is $(F(p_i))_i$ for some $F\in\Ed$ if and only
if some $G\in\E$ has $G(\rho_i)=b_i$ for all $i$, that is, if and only if the
value prescription $(b_i)$ at the nodes $\rho_i$ satisfies the shellwise
condition \eqref{ent:eq:maincondition}.
\item If $\Gamma$ has an order unit, arbitrary values on an arbitrary radially
finite subset of $\K^d$ are taken by some entire function.
\end{enumerate}
\end{corollary}
\begin{proof}
(b) For $F\in\Ed$, $f(Z)=(F\circ\Phi^{-1})(Z,0,\ldots,0)$ is entire
(\cref{ent:rc:lem:composition}) and $f(\rho_i)=F(p_i)$; conversely, if $G$ is
given, $F=G\circ\Phi_1$ is entire and $F(p_i)=G(\rho_i)$. The shellwise
description is \cref{ent:thm:image} for values, over $\R$ by
\cref{ent:prop:real-form}. (a) Pullback by $\Phi$ carries
$I(\Phi(\mathcal X))$ onto $I(\mathcal X)$ (\cref{ent:rc:lem:order-invariant}),
and \eqref{ent:rc:eq:quotient} with $P=P_D$ and \cref{ent:thm:quotient} give
the rest. (c) An order unit forces $\cf(\Gamma)=\aleph_0$, so for $d\ge2$,
$\mathcal X$ is line-rectifiable by \cref{ent:rc:thm:main} (a finite
$\mathcal X$ needs only Lagrange interpolation), and after a translation along the line the nodes are
as in \eqref{ent:rc:eq:nodes}; by \cref{ent:cor:order-unit-image}
condition \eqref{ent:eq:maincondition} is empty, and (b) applies. For $d=1$
this is \cref{ent:thm:interpolation}.
\end{proof}

Part (a) is a \emph{restricted} product description, as \cref{ent:thm:quotient}
is; it is not the unrestricted product description that the source disclaims.
The corollary concerns values only: interpolation of higher jets in several
variables is not treated.

\begin{corollary}[Constant multiplicity gives an ordinary ideal power]
\label{ent:rc:cor:powers}
If $m_i=\bar m$ for every node, then
$I_{\mathbf m}(\mathcal X)=I(\mathcal X)^{\bar m}$ and
$\mu(I(\mathcal X)^{\bar m})=\binom{\bar m+d-1}{d-1}$.
\end{corollary}
\begin{proof}
Here $P_{\langle r\rangle}=P^{\bar m-r}$ for $0\le r\le\bar m$, so the
generators \eqref{ent:rc:eq:ideal-generators} are exactly the degree-$\bar m$
products of $P,X_2,\ldots,X_d$; minimality is
\cref{ent:rc:thm:fat-generators}.
\end{proof}

\begin{theorem}[Unbounded multiplicities prevent finite generation]
\label{ent:rc:thm:unbounded}
Suppose $d\ge2$ and the $m_i$ are unbounded. Then $I_{\mathbf m}(\mathcal X)$
is a nonzero ideal that is not finitely generated; a generating set of finite
cardinality $q$ would satisfy
\begin{equation}\label{ent:rc:eq:local-lower}
 q\ge\binom{m_i+d-1}{d-1}\qquad\text{for every }i .
\end{equation}
In dimension one the ideal is $P_{\langle0\rangle}\E$, whatever the
multiplicities.
\end{theorem}
\begin{proof}
The nonzero $P_{\langle0\rangle}$ lies in the ideal, and so does
$P_{\langle|\alpha|\rangle}X'^\alpha$ for every $\alpha$. Fix $i$. The members
with $|\alpha|\le m_i$ have degree-$m_i$ Taylor parts at $p_i$ equal to nonzero
multiples of $(X_1-\rho_i)^{m_i-|\alpha|}X'^\alpha$, so the degree-$m_i$
Taylor image of the ideal at $p_i$ is the whole space of forms of that degree,
on which multiplication acts through evaluation of the multiplier, as in
\eqref{ent:rc:eq:leading-module}. Finitely many generators span at most $q$
dimensions, which gives \eqref{ent:rc:eq:local-lower}; its right side is
unbounded for $d\ge2$. For $d=1$ the vanishing conditions are divisibility by
$P_{\langle0\rangle}$, by the division statement of
\cref{ent:rc:subsec:calculus}.
\end{proof}

No assertion of countable generation is implicit here. Although the functions
$P_{\langle|\alpha|\rangle}X'^\alpha$ are countably many, an infinite convergent
expansion in them is not a finite algebraic combination; the proof establishes
non-finite generation and does not replace it by a stronger classification.

\begin{proposition}[Radical and zero set]\label{ent:rc:prop:radical}
For any multiplicities,
\begin{equation}\label{ent:rc:eq:radical}
 \sqrt{I_{\mathbf m}(\mathcal X)}=
 \bigl\{F\in\Ed:\exists N\ge1\ \forall i,\ N\ord_{p_i}F\ge m_i\bigr\}.
\end{equation}
The common zero set of $I_{\mathbf m}(\mathcal X)$ in $\K^d$ is exactly
$\mathcal X$. If the multiplicities are bounded, the radical is
$I(\mathcal X)$; if they are unbounded, it is strictly smaller than
$I(\mathcal X)$, although the zero sets agree.
\end{proposition}
\begin{proof}
Leading homogeneous parts multiply without vanishing, so
$\ord_a(F^N)=N\ord_aF$, and \eqref{ent:rc:eq:radical} is the definition of the
radical. Every node is a common zero. If $x$ is none of the $\rho_i$, then
$P_{\langle0\rangle}(x)\ne0$ excludes every point above $x$. If the first
coordinate is $\rho_i$ and some transverse coordinate $y_j$ is nonzero, take
an integer $r\ge m_i$: the member $P_{\langle r\rangle}(X_1)X_j^r$ of the ideal
is nonzero there, since $\rho_i$ is not a zero of $P_{\langle r\rangle}$. If
all $m_i\le\bar m$, every $F\in I(\mathcal X)$ has $F^{\bar m}$ in the ideal,
and every element of the radical vanishes on $\mathcal X$. For unbounded
multiplicities, $P=\prod_i(1-X_1/\rho_i)$ lies in $I(\mathcal X)$ with order
exactly one at every node, so no power of it lies in the ideal.
\end{proof}

\begin{proof}[Proof of \cref{ent:rc:thm:fat-main}]
Rectify, apply \cref{ent:rc:thm:fat-generators,ent:rc:thm:unbounded} in the
line coordinates \eqref{ent:rc:eq:line-D}, and transfer back with
\cref{ent:rc:lem:order-invariant}. Simple points are
\cref{ent:rc:cor:ci}; the dimension-one clause is the last sentence of
\cref{ent:rc:thm:unbounded}. Finite configurations use finite products.
\end{proof}

\begin{example}[Multiplicity one or three in the plane]
\label{ent:rc:ex:planar}
Let $d=2$, split the nodes into disjoint sets $N_1$ and $N_3\ne\varnothing$,
assign multiplicity one on $N_1$ and three on $N_3$, and let $P_{N_1}$ and
$P_{N_3}$ be the canonical products of $\sum_{i\in N_1}[\rho_i]$ and
$\sum_{i\in N_3}[\rho_i]$. Then
\begin{equation}\label{ent:rc:eq:example13}
 I_{\mathbf m}=\bigl(P_{N_1}P_{N_3}^3,\ P_{N_3}^2X_2,\ P_{N_3}X_2^2,\
 X_2^3\bigr),
\end{equation}
and exactly four generators are necessary: at a multiplicity-three node their
degree-three Taylor parts are nonzero multiples of $U^3,U^2X_2,UX_2^2,X_2^3$,
where $U=X_1-\rho_i$. This holds without an order unit; arbitrary infinite
value interpolation is not needed.
\end{example}

\begin{example}[Double points in three variables]\label{ent:rc:ex:double}
For constant multiplicity two and $d=3$,
$I_2=(P^2,PX_2,PX_3,X_2^2,X_2X_3,X_3^2)=I(\mathcal X)^2$, and the minimal
number of generators is $\binom42=6$, not the ambient dimension three: the six
quadratic jets at any node are independent.
\end{example}

\begin{example}[Increasing multiplicity on an already straight line]
\label{ent:rc:ex:increasing}
Take any escaping sequence of nonzero nodes, for example
$\rho_i=\omega^{\omega^{2i}}=t^{-e_{2i}}$ in $K_{\Gamma_\infty}$, and
$m_i=i+1$. The product $\prod_{i\ge1}(1-X_1/\rho_i)^{i+1}$ is entire and
nonzero, and in every dimension $d\ge2$ the fat-point ideal of the points
$(\rho_i,0,\ldots,0)$ is not finitely generated. Its zero set is that of the
finitely generated simple-point ideal, but the radicals differ. Unbounded local
complexity, not geometric nonrectifiability, already forces the failure of
finite generation.
\end{example}

\subsection{Workspace dependence, and what these subsections do not settle}
\label{ent:rc:subsec:limits}

\paragraph{Surreal scales, not class-global assertions.}
Every result of \crefrange{ent:rc:subsec:setting}{ent:rc:subsec:ideals}
concerns one fixed $\K=\kk((t^\Gamma))$. The embedding \eqref{ent:eq:embedding}
turns its elements into surreal or surcomplex numbers, and
\eqref{ent:rc:eq:surreal-D} uses such numbers directly. But a set cofinal in a
set-sized exponent workspace is never cofinal in the proper class of all
surreal exponents, and a larger workspace can introduce polydisks containing
all the old nodes: adjoin a positive exponent $\delta$ above every element of
$\Gamma$, and all the old centres $t^{-\gamma_n}$ lie in the one disk
$v(x)\ge-\delta$, so an infinite old configuration is no longer radially
finite there. ``Radially finite'' and ``entire'' keep their workspace
subscripts even when the notation suppresses them. This is consistent with,
and distinct from, the scalar-extension analysis of \cref{ent:sec:extension}.
The omnific Diophantine-geometry report's lemma that every set-sized family of
normal forms lies in one set-sized workspace (\texttt{odg:lem:workspace}) is
consistent with this: it locates a workspace, and radial finiteness is then
measured inside the chosen one.

\paragraph{The coefficient field.}
Constant coefficients carry the trivial valuation, and no infinite summation
of nonzero constants at one Hahn exponent is allowed: an ordinary complex
entire function such as $\exp(X)$, read as a series with constant coefficients,
is not in $\mathcal E_{\Gamma,1}$, because at $X=1$ its terms all have
exponent zero. These maps are not ordinary holomorphic maps tensored with a
Hahn field. Rank-one non-Archimedean function theory supplies the vocabulary
\cite{CherryLectures}; the proofs are written with ordered-group inequalities
so that no real-valued norm is imported into higher rank. Uncountability of
$\kk$ is used in \cref{ent:rc:lem:projection}; the analytic lemmas, the
division statements, the reduction obstruction and the line-coordinate ideal
calculations do not need it in the same way, but the positive rectification
theorems are not inferred for a countable constant field.

\paragraph{What the negative theorem does not forbid.}
The simplex configuration can be collapsed by noninvertible maps, for example
the projection to the first axis, and it satisfies many nonzero entire
equations: a canonical product in $X_1$ over the first coordinates of all its
points (finitely many on each shell) vanishes on it. The theorem forbids only
its inclusion in a hyperplane after an invertible entire change of all
coordinates, which is stronger than finding equations and different from an
interpolation impossibility. The obstruction is eventually uniform over the
blocks for each fixed map, but the threshold depends on the map; no threshold
is asserted that works for all entire automorphisms at once.

\paragraph{Dimension one and finite configurations.}
In dimension one every set already lies in the coordinate line, so
rectification is empty, and the ideal theorem has the principal exception of
\cref{ent:rc:thm:unbounded}. Every finite configuration in every dimension is
polynomially line-rectifiable. This is compatible with
\cref{ent:rc:thm:obstruction}: any finite union of its blocks is straightened
by a polynomial automorphism whose bad reduction occurs at one of those
finitely many scales, but no one entire automorphism can keep doing so along
a cofinal sequence of new scales.

\paragraph{Limits recorded by the seventh source.}
The following are retained without weakening.
\begin{itemize}[leftmargin=1.4em]
\item Every result is relative to one fixed Hahn workspace; there is no
class-global entire theory on $\No$ or $\No[i]$, the fine topology on $\No$ or
$\No[i]$ is not the topology used for the infinite sums, and topological
discreteness is weaker than radial finiteness and not a substitute for it.
\item Ordinary analytic convergence is not extended by summing infinitely many
constant coefficients at one Hahn exponent.
\item The results are stated for $\kk=\R$ or $\C$ only; no countable-field
version is claimed, and no descent of a rectifying automorphism to a
prescribed smaller Hahn workspace is asserted.
\item There is no classification of all ideals of $\Ed$, of all entire
automorphisms, of all nonrectifiable sets, or of all entire functions on the
full surreal or surcomplex class.
\item Line-rectifiability is not claimed equivalent to Rosay--Rudin tameness or
to Winkelmann's extension properties; the classical comparison is not a
transfer principle, and several inequivalent notions of tameness may exist in
this category.
\item The rectifying automorphism is not claimed to preserve $\Oz$ or the
Gaussian omnific integers; nothing concerns ring automorphisms of $\Oz$ or
field automorphisms of $\No$; ordinary integers are not radially finite; it
is not claimed that arbitrary bijections or injections of omnific
configurations extend to entire automorphisms.
\item No omnific factorization, refinement or irreducibility question
\cite{LInnocenteMantova} is settled; the only omnific input is the normal-form
description \eqref{ent:rc:eq:omnific-real}--\eqref{ent:rc:eq:omnific-complex}.
The Berarducci--Mantova derivation is neither used nor strengthened: the
coordinate derivatives here keep the Hahn coefficients fixed
\cite{MantovaMatusinski}.
\item The obstruction detects full-dimensional infinitesimal simplices; it is
not shown to be the only obstruction. No threshold uniform over all entire
automorphisms is asserted.
\item The value-interpolation lemma is a sufficient condition, not an exact
image theorem for arbitrary jets.
\item No countable generation is implied, and infinite analytic sums are never
used as algebraic ideal generation; the generator lower bound assumes no
Noetherianity, coherence or Nakayama lemma for an unproved local analytic ring.
\item The evaluation map $\E/(P)\to\prod_i\K$ is not claimed surjective without
an order unit, and no unrestricted product description of the infinite
evaluation quotient is given without an order unit.
\item The unbounded-multiplicity radical and zero-set distinction is a
consequence in this ring, not a claim that such non-Noetherian phenomena are
new in general.
\item The projection-and-shear strategy is classical and the local jet count
is elementary commutative algebra; neither is claimed as a new method. The
novelty assessment concerns their consequences in Hahn workspaces and the
coarsening obstruction; ``proposed'' means only that no matching result was
located in the inspected sources, after a bounded search that cannot rule out
a statement under different terminology or in an unindexed source.
\item No named published conjecture is claimed settled. The manuscript is an
AI-assisted, unrefereed research draft, not Lean-verified; it supplies no Lean
source and no claim that its results were added to or checked by the
repository's Lean build. Its finite symbolic program is a sanity check only:
it does not verify the infinite Hahn summation arguments, the universal
quantifiers over entire automorphisms, or novelty.
\end{itemize}

\setcounter{equation}{\value{entrcsaveeq}}%
\endgroup

% ---------------------------------------------------------------------------
% Source 11 (batch 32, manuscript 09; file prefix 11-arithmetic-sampling),
% labels ent:as:. Displayed equations in these subsections are numbered
% (S1), (S2), ... so that the numbers of all later equations are unchanged.
% ---------------------------------------------------------------------------
\begingroup
\setcounter{entassaveeq}{\value{equation}}%
\setcounter{equation}{0}%
\renewcommand{\theequation}{S\arabic{equation}}%
\renewcommand{\theHequation}{entas.\arabic{equation}}%

\subsection{Arithmetic sampling at ordinary points: setting and main results}
\label{ent:as:subsec:setting}

This subsection and the ten after it are the contribution of source~11, the
manuscript \emph{Arithmetic Sampling in Surreal Hahn Fields: entire rigidity,
coefficientwise algebraic division, omnific exceptional loci, and a
computability barrier} (batch~32, manuscript~09; file prefix
\texttt{11-arithmetic-sampling}). It is the eighth manuscript merged into this
report and is called source~11 after its file prefix; its labels carry
\texttt{ent:as:}. It asks what happens when \emph{arithmetic values are
prescribed at ordinary coefficient-field points}, rather than along a
sequence escaping through larger and larger valuation scales, and it answers
by one mechanism: at a fixed Hahn exponent, the coefficients of an entire
series across all monomials form an ordinary polynomial. Against this report
it partly answers the ideal-theory clause of \cref{ent:q:several}, for ideals
extended from the coefficient polynomial ring; the extent is recorded after
that question. Four warnings govern these subsections.

\begin{enumerate}[label=(\arabic*),leftmargin=2em]
\item \emph{Hypotheses.} The standing hypotheses are
\cref{ent:as:conv:standing}: $\kk$ is $\R$ or $\C$, and $\Gamma$ is an
arbitrary set-sized ordered abelian group, \emph{not} assumed divisible and
allowed to be $0$. Warning (1) at the head of this section applies, with the
coefficient field restricted to $\R$ or $\C$; the seventh source's hypotheses
(\cref{ent:rc:conv:standing}) do not. No preparation, root count, algebraic
closedness or canonical product is used.
\item \emph{Printed once.} The source proves for itself the locally finite
summation lemma, the finite-variable coefficient criterion, the cofinality
dichotomy and whole-class polynomial rigidity. These are
\cref{ent:mv:lem:tails}(a), \cref{ent:thm:multi-criterion},
\cref{ent:cor:cofinality} with the last clause of
\cref{ent:thm:multi-criterion}, and \cref{ent:cor:all-no,ent:mv:prop:allscale};
the source itself credits this report with the criterion and the cofinality
principle. They are printed once, by citation, in
\cref{ent:as:subsec:scales,ent:as:subsec:realization}. The source says that it
gives ``a proof that allows arbitrary rank and does not assume divisibility''
of the criterion; that adds nothing here, because
\cref{ent:thm:multi-criterion} is already stated and proved under exactly
those hypotheses, over any trivially valued field, by the same
descending-leading-exponent argument.
\item \emph{Two constant terms, and every exponent.} The power-series constant
coefficient $F(0)\in\K$ is not the scale-zero polynomial $\pg_0(F)\in\kk[X]$
(\cref{ent:as:rem:two-constants}). Membership in the omnific ring tests
\emph{every} positive Hahn exponent of a value, not the sign of its leading
valuation (\cref{ent:as:rem:magnitude}).
\item \emph{Letters.} Several of the source's letters collide with this
report's and are renamed; the last table of \cref{ent:subsec:notation} gives
the translation. The collisions that matter are these. Its $D$, meaning $\Z$
or $\Z[i]$, is written $\Zk$, because $D$ is a divisor here. Its omnific ring
$A=D+\Pi$ is written $\OzK$, because $\A$ is the Hahn algebra and $A,B$ name
entire functions in B\'ezout identities. Its ring $B=k+\Pi$ is the seventh
source's polynomial-part ring $\Kle$. Its $\Pi$ is written $\PiK$, apart from
the class $\Pi$ of purely infinite surreals and the shell factors $\Pi_n$. Its
Hahn scale $\gamma$ and weight $\rho$ are written $\eta$ and $\gamma$, as in
\cref{ent:thm:multi-criterion} and \eqref{ent:eq:Halg}, because $\rho$ is a node
here. Its matrices $H$ and its monic generator $H_F$ are renamed, because $H$
is a convex subgroup.
\end{enumerate}

\begin{convention}[Standing hypotheses of source~11's subsections]
\label{ent:as:conv:standing}
$\kk$ is $\R$ or $\C$ with the trivial valuation; $\Gamma$ is a set-sized
ordered abelian group, not assumed divisible, possibly $0$; and
$\K=\kk((t^\Gamma))$ with $v(x)=\min\supp x$ and $v(0)=+\infty$, so positive
exponents are infinitesimal terms. The number $d\ge1$ of variables is finite,
$X=(X_1,\ldots,X_d)$, and $\Ed$ is the ring of \cref{ent:mv:def:entire}; its
elements are written $F=\sum_{\alpha\in\N^d}a_\alpha X^\alpha$. Put
\begin{align}
 \PiK&=\{x\in\K:\supp(x)\subseteq\Gamma_{<0}\},\label{ent:as:eq:pi}\\
 \Kle&=\kk+\PiK=\{x\in\K:\supp(x)\subseteq\Gamma_{\le0}\},\label{ent:as:eq:B}\\
 \OzK&=\Zk+\PiK,\qquad
 \Zk=\begin{cases}\Z,&\kk=\R,\\ \Z[i],&\kk=\C.\end{cases}\label{ent:as:eq:A}
\end{align}
Empty support is permitted, so $0\in\PiK$. The ring $\Kle$ is the one of
\eqref{ent:rc:eq:Kle}, here without divisibility.
\end{convention}

The set $\PiK$ is a $\kk$-vector space, not a field, and it is an ideal of
$\Kle$: a strictly negative exponent plus a nonpositive one is strictly
negative. On $\Kle$ extraction of the coefficient at $0$ is a ring
homomorphism
\[
 \ct:\Kle\longrightarrow\kk,\qquad
 \ct\Bigl(\sum_\eta x_\eta t^\eta\Bigr)=x_0,
\]
with kernel $\PiK$, because two nonpositive exponents sum to zero only when
both are zero; and $\OzK=\ct^{-1}(\Zk)$ is a pullback. For an ordered subgroup
$\Gamma\subseteq\No$ and the embedding \eqref{ent:eq:embedding},
\[
 \OzK=\K\cap\Oz\quad(\kk=\R),\qquad
 \OzK=\K\cap(\Oz+i\Oz)\quad(\kk=\C),
\]
which is \eqref{ent:rc:eq:omnific-real}--\eqref{ent:rc:eq:omnific-complex}; for
$\Gamma=0$ the definitions give $\PiK=0$, $\Kle=\kk$ and $\OzK=\Zk$. The
set-sized-quotients report writes the same real ring $A_\Gamma=\Z\oplus
J_\Gamma$ with $J_\Gamma$ supported in $\Gamma_{>0}$ for the monomials $T^\gamma$
(\texttt{osq:or:cor:workspace}); its exponent sign is the reverse of ours, and
the ring is the same.

\begin{remark}[Two different constant terms]\label{ent:as:rem:two-constants}
For $F=\sum_\alpha a_\alpha X^\alpha$, the power-series constant coefficient is
$a_0=F(0)\in\K$. The scale-zero polynomial $\pg_0(F)=[t^0]F$ of
\cref{ent:as:def:scale} is a polynomial over $\kk$, obtained by extracting the
valuation-zero coefficient from \emph{every} $a_\alpha$. These objects must not
be confused.
\end{remark}

\begin{remark}[Magnitude is not enough]\label{ent:as:rem:magnitude}
The element $t^{-1}+t$ has infinite leading magnitude but lies in neither
$\Kle$ nor $\OzK$: its positive tail $t$ is forbidden. All the sampling and jet
criteria below detect the full positive tail, not just the sign of $v(x)$.
\end{remark}

We do not identify $\K$ with the fraction field of $\OzK$ in general. The
identity holds for the full class, $\operatorname{Frac}(\Oz)=\No$, but its
set-sized analogues have genuine cofinality restrictions; the source cites
\cite[Proposition~2.4.5]{LInnocenteMantova} for that distinction. None of the
proofs below needs it.

\paragraph{The simplest conclusion.} Embed $\R((t))$ in $\No$ by
$t=\omega^{-1}$, and let $F(Z)$ be an ordinary power series strongly summable
at every point of this Hahn field. If $F(n)$ is an omnific integer for
infinitely many distinct ordinary integers $n$, then $F$ is a polynomial
(\cref{ent:as:cor:finite-points}); no growth hypothesis beyond entireness is
used. On the other hand a nonpolynomial entire function can vanish at
infinitely many \emph{infinite} omnific integers
(\cref{ent:as:subsec:boundaries}). The location of the sampling points is
therefore essential.

\paragraph{Main results.} For a matrix $\mathbf M$ over $\kk[X]$, an entire
vector $F$ is $\mathbf MG$ with $G$ entire if and only if every scale
polynomial of $F$ lies in the polynomial image of $\mathbf M$
(\cref{ent:as:thm:matrix}); so $\Ed$ is faithfully flat over $\kk[X]$
(\cref{ent:as:thm:flat}), with exact normal forms modulo every ideal defined
over $\kk$ (\cref{ent:as:cor:normal-form}). If $\mathcal S\subseteq\kk^d$ is
Zariski dense and $F(\mathcal S)\subseteq\Kle$, then $F\in\Kle[X]$
(\cref{ent:as:thm:sampling}), and
\begin{equation}\label{ent:as:eq:int-intro}
 F(\Zk^d)\subseteq\OzK
 \iff F(\OzK^{\,d})\subseteq\OzK
 \iff F\in\Int(\Zk^d,\Zk)+\PiK[X]
\end{equation}
(\cref{ent:as:thm:integer}); the last equivalence is a classification inside
$\Ed$, not an assertion about arbitrary functions on $\OzK$. The ordinary
$\Kle$-valued locus of $F$ is the zero set of the finitely generated ideal
$I_+(F)$ spanned by the positive-scale polynomials, and the $\OzK$-valued locus
adds the ordinary arithmetic condition $\pg_0(F)(\mathbf s)\in\Zk$
(\cref{ent:as:thm:tail}); finitely many actual scales generate $I_+(F)$, the
same scales control every finite jet order (\cref{ent:as:thm:jets}), and every
nonzero ideal occurs for a nonpolynomial entire function when such functions
exist (\cref{ent:as:thm:realization}). Finally, even for explicitly computable
coefficients in $\Q[t]$ with a common quadratic valuation bound and promised
nonpolynomiality, whether $F(1)$ is omnific cannot be decided in general
(\cref{ent:as:thm:undecidable}); that is a statement about one precise
infinite-series presentation, not about every finite symbolic language.

\paragraph{Scope.} Nonpolynomial results concern fixed Hahn workspaces, not an
unrestricted analytic function on the proper class $\No[i]$
(\cref{ent:as:rem:full-class}); strong Hahn summation is not ordinary real or
complex convergence (\cref{ent:as:subsec:boundaries}). The results are proposed
contributions with complete written proofs, and the source claims no
certified historical priority, independent refereeing or proof-assistant
verification (\cref{ent:as:subsec:limits}).

\subsection{The coefficient criterion, affine changes and scale polynomials}
\label{ent:as:subsec:scales}

\paragraph{Printed once.} The source's lemma on locally finite families of
supports --- if for every $b\in\Gamma$ only finitely many $x_j$ satisfy
$v(x_j)\le b$, then $(x_j)$ is strongly summable --- is
\cref{ent:mv:lem:tails}(a). Its entire coefficient criterion is
\cref{ent:thm:multi-criterion}: $F\in\Ed$ if and only if, for all
$\gamma,b\in\Gamma$, only finitely many $\alpha$ with $a_\alpha\ne0$ satisfy
$v(a_\alpha)+|\alpha|\gamma\le b$ \eqref{ent:mv:eq:growth}. Its proof is the
one printed there: diagonal monomial points, a common lower bound for the
leading exponents, a subsequence of strictly increasing total degrees, and a
strictly decreasing sequence of leading exponents at a shifted diagonal point,
with no rank-one or convergence argument. Two remarks of the source are kept.
It suffices to test $\gamma\le0$, since the condition at $\gamma=0$ implies it
at every $\gamma\ge0$. For $\Gamma=0$, which \cref{ent:mv:subsec:setting}
excludes, evaluation at $(1,\ldots,1)$ forces finitely many nonzero terms, so
the criterion says exactly that $F$ is a polynomial, as noted after
\cref{ent:mv:prop:realize-support}. That $\Ed$ is a $\K$-algebra follows from
\cref{ent:mv:lem:tails}(b): for a product use the strongly summable family
indexed by pairs of multi-indices and regroup by their sum; each power-series
coefficient involves finitely many pairs. The phrase ``whole family'' in
\cref{ent:mv:def:entire} rules out evaluating different iterated sums without
proving that their common multivariate family is summable.

\begin{proposition}[Affine changes of coordinates]\label{ent:as:prop:affine}
Let $F\in\Ed$, $\mathbf a\in\K^d$, and let $\mathbf L$ be a $d\times e$ matrix
over $\K$. The binomial expansion of $F(\mathbf a+\mathbf LY)$ in
$Y=(Y_1,\ldots,Y_e)$ defines an element of $\mathcal E_{\Gamma,e}$ that agrees
with pointwise composition. If $e=d$ and $\mathbf L$ is invertible, the
substitution is a $\K$-algebra automorphism of $\Ed$ and preserves
polynomiality in both directions.
\end{proposition}
\begin{proof}
Fix a weight $\gamma$ for the new coordinates. Choose $c\in\Gamma$ no greater
than the valuation of every nonzero entry of $\mathbf a$ and every
$v(\mathbf L_{ij})+\gamma$ with $\mathbf L_{ij}\ne0$. Every expanded summand
coming from $a_\alpha(\mathbf a+\mathbf LY)^\alpha$ has valuation at least
$v(a_\alpha)+|\alpha|c$ at every point with $v(Y_j)\ge\gamma$, and for each
$\alpha$ there are only finitely many summands. By \eqref{ent:mv:eq:growth},
only finitely many $\alpha$ contribute at or below any fixed bound. So the whole
expanded family is strongly summable and can be regrouped
(\cref{ent:mv:lem:tails}), both to form the new coefficients and to evaluate
them; this proves the first two claims. The inverse affine substitution
proves the last; equality of the resulting series also follows from the
uniqueness statement of \cref{ent:as:lem:scale}. Affine substitution of a
polynomial is plainly polynomial.
\end{proof}

For divisible $\Gamma$ this is contained in \cref{ent:rc:lem:composition}, and
its linear part is \cref{ent:mv:prop:linear}, which holds for any $\Gamma$; the
source's direct proof covers translations without divisibility.

\begin{definition}[Scale coefficient polynomials]\label{ent:as:def:scale}
For $F\in\Ed$ and $\eta\in\Gamma$, put
\begin{equation}\label{ent:as:eq:p-eta}
 \pg_\eta(F)(X)=[t^\eta]F(X)
 =\sum_{\alpha}\bigl([t^\eta]a_\alpha\bigr)X^\alpha ,
\end{equation}
abbreviated $\pg_\eta$. For a vector of entire series, $\pg_\eta$ is applied
componentwise.
\end{definition}

In one variable $\pg_\eta(F)$ is the coefficient $P_\eta$ of the expansion
$F=\sum_\eta P_\eta t^\eta$ in the Hahn algebra $\A=\C[X]((t^\Gamma))$ of
\eqref{ent:eq:Halg}, and its lowest nonzero instance is the initial polynomial
there; in $d$ variables $\pg_\eta(F)$ is the sum over $n$ of the homogeneous
polynomials $P_{n,\eta}$ used in the proof of \cref{ent:mv:thm:exceptional}.
(The merge's comparison.)

\begin{lemma}[Scale polynomials]\label{ent:as:lem:scale}
Let $F\in\Ed$. Each $\pg_\eta(F)$ is a polynomial over $\kk$. For each
$b\in\Gamma$ the degrees of all $\pg_\eta(F)$ with $\eta\le b$ are uniformly
bounded. The set of scales $\eta$ with $\pg_\eta(F)\ne0$ is well ordered. For
every ordinary point $\mathbf s\in\kk^d$,
\begin{equation}\label{ent:as:eq:scale-eval}
 [t^\eta]F(\mathbf s)=\pg_\eta(F)(\mathbf s).
\end{equation}
Consequently evaluation on all of $\kk^d$ determines an entire series
uniquely.
\end{lemma}
\begin{proof}
By \eqref{ent:mv:eq:growth} with $\gamma=0$, only finitely many $a_\alpha$ have
$v(a_\alpha)\le b$, and every coefficient at a scale $\eta\le b$ comes from one
of them; this gives polynomiality and the uniform degree bound. Entireness at
$(1,\ldots,1)$ shows that the union of the coefficient supports is well
ordered, and it is exactly the set of nonzero scales. At a point of $\kk^d$,
multiplication by $\mathbf s^\alpha\in\kk$ introduces no new exponent, so
strong summation gives \eqref{ent:as:eq:scale-eval}. A polynomial over the
infinite field $\kk$ vanishing on $\kk^d$ is zero, by induction on $d$;
applying this to every $\pg_\eta$ proves uniqueness.
\end{proof}

\begin{corollary}[The nonpositive part is a polynomial]
\label{ent:as:cor:truncation}
For $F\in\Ed$ the scale truncation
\[
 F_{\le0}=\sum_{\eta\le0}\pg_\eta(F)\,t^\eta
\]
is a polynomial in $\Kle[X]$. The difference $F_+=F-F_{\le0}$ is entire, and
each of its power-series coefficients has support in $\Gamma_{>0}$.
\end{corollary}
\begin{proof}
The degree bound of \cref{ent:as:lem:scale} with $b=0$ makes the truncation a
polynomial, and its coefficients have nonpositive support. Truncating a Hahn
coefficient cannot lower its valuation, so $F_+$ still satisfies
\eqref{ent:mv:eq:growth}.
\end{proof}

\begin{remark}[Two different finite bounds]\label{ent:as:rem:two-bounds}
There can be infinitely many scales $\eta\le0$, because a Hahn coefficient may
have infinite support bounded above. What is finite is the number of
\emph{power-series monomials} involved at those scales. So $F_{\le0}$ is a
polynomial with possibly infinitely supported Hahn coefficients, not
necessarily a finite sum of Hahn monomials.
\end{remark}

In the language of this report, \cref{ent:as:lem:scale} places $\Ed$ inside
$\kk[X]((t^\Gamma))$, the $d$-variable form of the inclusion $\E\subseteq\A$ of
\eqref{ent:eq:chain}, with degrees bounded on every initial segment of scales
(the merge's remark).

\subsection{Coefficientwise division and faithful flatness over
\texorpdfstring{$\kk[X]$}{k[X]}}
\label{ent:as:subsec:division}

In this subsection every polynomial matrix has entries in $\kk[X]$, not in
$\K[X]$. The distinction matters: multiplication by an ordinary polynomial
does not mix valuation scales. For a polynomial vector, the degree is the
largest total degree of its entries, and the zero vector has degree
$-\infty$.

\begin{lemma}[Bounded degree increase preserves entireness]
\label{ent:as:lem:degree-transfer}
Let $F\in\Ed^{\,r}$. For each $\eta\in\Gamma$ choose a polynomial vector
$\mathsf q_\eta\in\kk[X]^s$ with $\mathsf q_\eta=0$ whenever $\pg_\eta(F)=0$,
and suppose that a fixed integer $c\ge0$ satisfies
\begin{equation}\label{ent:as:eq:degree-transfer}
 \deg\mathsf q_\eta\le\deg\pg_\eta(F)+c
 \qquad\text{whenever }\mathsf q_\eta\ne0.
\end{equation}
Then the coefficientwise reconstruction $G=\sum_\eta\mathsf q_\eta t^\eta$
defines an element of $\Ed^{\,s}$ whose scale support is contained in that of
$F$.
\end{lemma}
\begin{proof}
The scale support of $F$ is well ordered, so for a fixed monomial
$X^{\alpha'}$ in a fixed component the corresponding coefficients of the
$\mathsf q_\eta$ define a Hahn series $b_{\alpha'}$. We check
\eqref{ent:mv:eq:growth} for them. Fix $\gamma\le0$ and $b$. Suppose
$b_{\alpha'}\ne0$ and $v(b_{\alpha'})+|\alpha'|\gamma\le b$, and put
$\eta=v(b_{\alpha'})$. The vector $\mathsf q_\eta$ has a nonzero term of degree
$|\alpha'|$, so by \eqref{ent:as:eq:degree-transfer} some component of
$\pg_\eta(F)$ has a monomial $X^\alpha$ with $|\alpha|\ge|\alpha'|-c$. Let
$a_\alpha$ be its original Hahn coefficient in that component. Since
$\eta\in\supp(a_\alpha)$ and $\gamma\le0$,
\[
 v(a_\alpha)+|\alpha|\gamma\ \le\ \eta+(|\alpha'|-c)\gamma\ \le\ b-c\gamma .
\]
Only finitely many component--multi-index pairs satisfy this, by entireness of
$F$; their degrees have a finite maximum, so $|\alpha'|$ is uniformly bounded
and there are only finitely many such $\alpha'$. This is
\eqref{ent:mv:eq:growth} for $G$, which suffices by the first remark of
\cref{ent:as:subsec:scales}. The support assertion is immediate from the
construction.
\end{proof}

\begin{lemma}[Polynomial matrix lifting with bounded degree loss]
\label{ent:as:lem:matrix-degree}
Let $\mathbf M\in\operatorname{Mat}_{r\times s}(\kk[X])$. There is an integer
$c_{\mathbf M}\ge0$ such that every $f$ in the image of
$\mathbf M:\kk[X]^s\to\kk[X]^r$ has a lift $u\in\kk[X]^s$ with
\[
 \mathbf Mu=f,\qquad \deg u\le\deg f+c_{\mathbf M},
\]
the zero vector being lifted to zero.
\end{lemma}
\begin{proof}
Choose a module monomial order that compares total degree first. The leading
monomials of the submodule $\operatorname{im}\mathbf M\subseteq\kk[X]^r$
generate a monomial submodule, which is finitely generated because $\kk[X]$ is
Noetherian. Choose $g_1,\ldots,g_l\in\operatorname{im}\mathbf M$ whose leading
monomials generate it; they form a module Gr\"obner basis. For each $j$ choose
a polynomial lift $u_j$ with $\mathbf Mu_j=g_j$. Division of
$f\in\operatorname{im}\mathbf M$ by this basis has remainder zero, and with a
degree-compatible order no reduction increases total degree, so
$f=\sum_jh_jg_j$ with $\deg h_j\le\deg f$ for every nonzero $h_j$: each
multiplying monomial has degree at most that of the dividend at the stage
where it is used. Then $u=\sum_jh_ju_j$ is a lift with
$\deg u\le\deg f+\max_j\deg u_j$; take that maximum, enlarged to zero if
necessary, as $c_{\mathbf M}$. For $\operatorname{im}\mathbf M=0$ the assertion
is immediate. Existence of a finite Gr\"obner basis uses only the Noetherian
property and the well-founded monomial order, not an effective presentation of
$\kk$; the Noetherian input is the Hilbert basis theorem
\cite[Lemma~10.31.1]{Stacks}.
\end{proof}

\begin{theorem}[Exact coefficientwise image]\label{ent:as:thm:matrix}
For $\mathbf M\in\operatorname{Mat}_{r\times s}(\kk[X])$ and $F\in\Ed^{\,r}$,
\begin{equation}\label{ent:as:eq:matrix-image}
 F\in\mathbf M\,\Ed^{\,s}
 \iff
 \pg_\eta(F)\in\mathbf M\,\kk[X]^s\quad\text{for every }\eta\in\Gamma .
\end{equation}
When these hold, an entire lift $G$ can be chosen with no scales outside those
of $F$.
\end{theorem}
\begin{proof}
If $F=\mathbf MG$, then $\pg_\eta(F)=\mathbf M\,\pg_\eta(G)$, because $\mathbf M$
has ordinary coefficients and finitely many terms. Conversely, lift each
$\pg_\eta(F)$ by \cref{ent:as:lem:matrix-degree}, zero to zero. The uniform
degree bound permits \cref{ent:as:lem:degree-transfer}, which gives
$G=\sum_\eta\mathsf q_\eta t^\eta\in\Ed^{\,s}$, and coefficient extraction
gives $\pg_\eta(\mathbf MG)=\pg_\eta(F)$ at every scale, hence $\mathbf MG=F$.
\end{proof}

\begin{corollary}[Algebraic ideal membership]\label{ent:as:cor:ideal}
For every ideal $I\subseteq\kk[X]$,
\begin{equation}\label{ent:as:eq:ideal-membership}
 F\in I\,\Ed \iff \pg_\eta(F)\in I\quad(\eta\in\Gamma).
\end{equation}
In particular $(I\,\Ed)\cap\kk[X]=I$.
\end{corollary}
\begin{proof}
Choose finitely many generators of $I$ and apply \cref{ent:as:thm:matrix} to
the one-row matrix they form. An ordinary polynomial has only its scale-zero
polynomial, equal to itself, which proves the contraction.
\end{proof}

\begin{theorem}[Faithful flatness over the coefficient polynomial ring]
\label{ent:as:thm:flat}
For every set-sized ordered abelian group $\Gamma$ and $\kk\in\{\R,\C\}$, the
inclusion $\kk[X]\to\Ed$ is faithfully flat.
\end{theorem}
\begin{proof}
Take a relation $\sum_{i=1}^rf_iU_i=0$ with $f_i\in\kk[X]$ and $U_i\in\Ed$. The
syzygy module
\[
 \Syz(f_1,\ldots,f_r)=\Bigl\{(h_1,\ldots,h_r)\in\kk[X]^r:
 \textstyle\sum_if_ih_i=0\Bigr\}
\]
is finitely generated, $\kk[X]$ being Noetherian; let $\mathbf N$ be a matrix
whose columns generate it. Every scale vector $\pg_\eta(U_1,\ldots,U_r)$ lies
in $\Syz(f_1,\ldots,f_r)$, so by \cref{ent:as:thm:matrix} the entire vector $U$
is $\mathbf NV$ for an entire vector $V$. Thus the relation is an entire linear
combination of polynomial syzygies, which is the equational criterion of
flatness \cite[Lemma~10.39.11]{StacksFlat}. For every proper ideal $I$ of
$\kk[X]$, \cref{ent:as:cor:ideal} gives $1\notin I\,\Ed$, so $\Ed/I\,\Ed\ne0$;
with flatness this is faithful flatness \cite[Lemma~10.39.15]{StacksFlat}.
Explicitly, a nonzero cyclic submodule $\kk[X]/I$ of a nonzero module stays
nonzero after tensoring, and its inclusion stays injective by flatness.
\end{proof}

\begin{corollary}[Canonical entire normal forms]\label{ent:as:cor:normal-form}
Fix a reduced Gr\"obner basis of an ideal $I\subseteq\kk[X]$ for a
degree-compatible order, with normal-form map $\NF_I$. Every coset of $I\,\Ed$
has a unique representative $N_F\in\Ed$ all of whose scale polynomials are
linear combinations of standard monomials for that basis, namely
\[
 N_F=\sum_\eta\NF_I\bigl(\pg_\eta(F)\bigr)t^\eta .
\]
It introduces no new valuation scales.
\end{corollary}
\begin{proof}
Polynomial normal-form reduction does not increase total degree, so
\cref{ent:as:lem:degree-transfer} makes $N_F$ entire. Every
$\pg_\eta(F-N_F)$ lies in $I$, so \cref{ent:as:cor:ideal} gives
$F-N_F\in I\,\Ed$. If two such representatives differ by an element of
$I\,\Ed$, each scale polynomial of the difference lies in $I$ and in the span of
the standard monomials, whose intersection is zero by uniqueness of polynomial
normal forms.
\end{proof}

\begin{corollary}[Ordinary fibres]\label{ent:as:cor:fiber}
For $\mathbf s\in\kk^d$, evaluation induces
$\Ed/(X_1-s_1,\ldots,X_d-s_d)\Ed\cong\K$.
\end{corollary}
\begin{proof}
Evaluation is a ring map, and its image contains $\K$ as constant functions.
Its kernel consists of the $F$ with $\pg_\eta(F)(\mathbf s)=0$ for every $\eta$
(\cref{ent:as:lem:scale}); these polynomials lie in the ideal of the ordinary
point, so \cref{ent:as:cor:ideal} identifies the kernel.
\end{proof}

For divisible $\Gamma$ this is also the one-point case of
\cref{ent:rc:cor:ci}, whose proof divides by linear factors; the proof above
needs no divisibility (the merge's comparison).

\begin{remark}[What flatness does not say]\label{ent:as:rem:flatness}
The base is $\kk[X]$. Nothing is claimed for polynomial matrices with
arbitrary Hahn coefficients (\cref{ent:as:q:matrices}), and $\Ed$ is not made
Noetherian, a principal ideal domain or a B\'ezout domain; this is compatible
with the arbitrary-rank distinctions of this report. Nor is the omnific ring
$\OzK$ identified with an integer polynomial ring. The explicit content is that
finite polynomial systems over the ordinary coefficient field acquire no new
linear relations in $\Ed$ other than combinations of their existing syzygies.
\end{remark}

\subsection{Arithmetic sampling and restriction to ordinary algebraic sets}
\label{ent:as:subsec:sampling}

A subset $\mathcal S\subseteq\kk^d$ is \emph{Zariski dense} if the only
polynomial in $\kk[X]$ vanishing on $\mathcal S$ is zero; its vanishing ideal is
$I_\kk(\mathcal S)=\{g\in\kk[X]:g|_{\mathcal S}=0\}$, an ideal of $\kk[X]$, not
of $\Ed$.

\begin{theorem}[Dense arithmetic sampling forces polynomiality]
\label{ent:as:thm:sampling}
Let $F\in\Ed$ and let $\mathcal S\subseteq\kk^d$ be Zariski dense. If
$F(\mathbf s)\in\Kle$ for all $\mathbf s\in\mathcal S$, then $F\in\Kle[X]$. In
particular the same holds under the stronger hypothesis
$F(\mathcal S)\subseteq\OzK$.
\end{theorem}
\begin{proof}
For $\eta>0$, \eqref{ent:as:eq:scale-eval} gives
$\pg_\eta(F)(\mathbf s)=[t^\eta]F(\mathbf s)=0$ on $\mathcal S$, so
$\pg_\eta(F)=0$ by density. Thus every Hahn coefficient $a_\alpha$ has support
in $\Gamma_{\le0}$, and every nonzero one has $v(a_\alpha)\le0$;
\eqref{ent:mv:eq:growth} with $\gamma=b=0$ permits only finitely many. So $F$ is
a polynomial with coefficients in $\Kle$.
\end{proof}

\begin{corollary}[One-variable finiteness]\label{ent:as:cor:finite-points}
If $F\in\mathcal E_{\Gamma,1}$ is not a polynomial, the sets
$\{s\in\kk:F(s)\in\Kle\}$ and $\{s\in\kk:F(s)\in\OzK\}$ are finite.
Conversely, infinitely many distinct ordinary points with values in $\Kle$
force polynomiality.
\end{corollary}
\begin{proof}
Every infinite subset of the infinite field $\kk$ is Zariski dense in one
variable, since a nonzero polynomial has at most its degree many roots. Apply
\cref{ent:as:thm:sampling}.
\end{proof}

\begin{corollary}[Product grids and affine arithmetic lattices]
\label{ent:as:cor:affine-sampling}
If each $\mathcal S_j\subseteq\kk$ is infinite and
$F(\mathcal S_1\times\cdots\times\mathcal S_d)\subseteq\Kle$, then
$F\in\Kle[X]$. More generally, let $\mathcal S\subseteq\kk^d$ be Zariski dense,
$\mathbf a\in\K^d$ and $\mathbf L\in\mathrm{GL}_d(\K)$. If
$F(\mathbf a+\mathbf L\mathbf s)\in\Kle$ for all $\mathbf s\in\mathcal S$, then
$F$ is a polynomial, and its expression in the normalized coordinates lies in
$\Kle[X]$.
\end{corollary}
\begin{proof}
A polynomial vanishing on a product of infinite sets is zero: treat it as a
polynomial in the last variable with infinitely many roots and induct on $d$
for the coefficients. This proves density of the product. For the affine
statement apply \cref{ent:as:thm:sampling} to $G(X)=F(\mathbf a+\mathbf LX)$
and then \cref{ent:as:prop:affine}. The coefficients of $F$ in the original
coordinates need not lie in $\Kle$, because $\mathbf L^{-1}$ can have
positive-valuation entries.
\end{proof}

This includes infinitesimally spaced sampling grids: for any nonzero
$h\in\K$, omnific values at infinitely many points $a+hn$ with ordinary
integers $n$ force a one-variable entire function to be a polynomial. No
metric separation condition is imposed.

Let $\Ed^{+}$ be the set of entire series all of whose power-series
coefficients are supported in $\Gamma_{>0}$. It is an additive subgroup and a
$\kk[X]$-module, but generally not an ideal of $\Ed$: multiplication by a
constant of negative valuation can shift its scales to zero or below.

\begin{theorem}[Arithmetic restriction theorem]\label{ent:as:thm:restriction}
For any $\mathcal S\subseteq\kk^d$ and $F\in\Ed$ the following are
equivalent:
\begin{enumerate}[label=\textnormal{(\roman*)},leftmargin=2em]
\item $F(\mathcal S)\subseteq\Kle$;
\item $\pg_\eta(F)\in I_\kk(\mathcal S)$ for every $\eta>0$;
\item there are $P\in\Kle[X]$, finitely many $g_j\in I_\kk(\mathcal S)$ and
$G_j\in\Ed^{+}$ with
\begin{equation}\label{ent:as:eq:restriction}
 F=P+\sum_jg_jG_j .
\end{equation}
\end{enumerate}
In (iii), necessarily $P=F_{\le0}$; in particular
$F|_{\mathcal S}=F_{\le0}|_{\mathcal S}$. Moreover $F(\mathcal S)\subseteq\OzK$
exactly when these conditions hold and $\pg_0(F)(\mathcal S)\subseteq\Zk$.
\end{theorem}
\begin{proof}
(i) and (ii) are equivalent by \eqref{ent:as:eq:scale-eval} and the support
description of $\Kle$. Assume (ii). Put $F_+=F-F_{\le0}\in\Ed^{+}$ and choose
generators $g_1,\ldots,g_l$ of the Noetherian ideal $I_\kk(\mathcal S)$. Every
scale polynomial of $F_+$ lies in that ideal, so \cref{ent:as:thm:matrix} for
the row $(g_1,\ldots,g_l)$ gives a support-preserving lift, with no
nonpositive scales: each $G_j\in\Ed^{+}$. This is (iii) with $P=F_{\le0}$.
Conversely, if \eqref{ent:as:eq:restriction} holds, evaluation on $\mathcal S$
kills every $g_jG_j$ and leaves the $\Kle$-valued polynomial $P$; taking
nonpositive truncations in \eqref{ent:as:eq:restriction} shows $P=F_{\le0}$.
Membership in $\OzK$ adds exactly the scale-zero condition in $\Zk$.
\end{proof}

\begin{example}[A positive-dimensional exception]\label{ent:as:ex:parabola}
In $\K=\R((t))$ put
\[
 F(X,Y)=t^{-1}+(Y-X^2)\sum_{n\ge0}t^{(n+1)^2}X^n .
\]
It is entire by \eqref{ent:mv:eq:growth} and nonpolynomial. At every ordinary
point of the parabola $Y=X^2$ its value is $t^{-1}\in\OzK$; away from the
parabola its first positive scale is nonzero. So ordinary arithmetic values
can persist on a whole algebraic curve, but the restriction of the function
there is the polynomial $t^{-1}$.
\end{example}

\begin{corollary}[Vanishing ideals of ordinary and algebraic point sets; the
merge's]\label{ent:as:cor:vanishing}
\leavevmode
\begin{enumerate}[label=\textnormal{(\alph*)},leftmargin=2em]
\item For every $\mathcal S\subseteq\kk^d$, the entire functions vanishing on
$\mathcal S$ form the finitely generated ideal $I_\kk(\mathcal S)\,\Ed$.
\item Let $\kk=\C$, let $I\subseteq\C[X]$ be an ideal, and let
$\mathbb V_{\K}(I)=\{\mathbf x\in\K^d:g(\mathbf x)=0\text{ for all }g\in I\}$. The
entire functions vanishing on $\mathbb V_{\K}(I)$ form the ideal
$\sqrt I\,\Ed$.
\end{enumerate}
\end{corollary}
\begin{proof}
(a) A Hahn series is zero exactly when all its coefficients are, so by
\eqref{ent:as:eq:scale-eval} $F$ vanishes on $\mathcal S$ exactly when every
$\pg_\eta(F)$ lies in $I_\kk(\mathcal S)$; apply \cref{ent:as:cor:ideal}.
(b) The ordinary points $\mathbb V_\C(I)=\mathbb V_{\K}(I)\cap\C^d$ are among
those of $\mathbb V_{\K}(I)$, and $I_\C(\mathbb V_\C(I))=\sqrt I$ by Hilbert's
Nullstellensatz; so a function vanishing on $\mathbb V_{\K}(I)$ lies in
$\sqrt I\,\Ed$ by (a). Conversely, if $g^n\in I$ then $g(\mathbf x)^n=0$ at every
$\mathbf x\in\mathbb V_{\K}(I)$, so each $g\in\sqrt I$ vanishes there, and so does
every $\Ed$-combination of such $g$, evaluation being a ring homomorphism.
\end{proof}

Part (b) is a first positive-dimensional statement in several variables, for
zero sets that are algebraic and defined over the coefficient field; it says
nothing about zero sets of entire functions that are not of this form, and
nothing for $\kk=\R$, where the real Nullstellensatz would be needed. The
source states (a) only implicitly, through \cref{ent:as:cor:ideal} and
\cref{ent:as:lem:scale}; both parts are the merge's.

\subsection{Entire maps preserving the omnific integers of a workspace}
\label{ent:as:subsec:integer}

For $\Zk$ with fraction field $\Q$ or $\Q(i)$ put
\[
 \Int(\Zk^d,\Zk)=\{P\in\operatorname{Frac}(\Zk)[X]:P(\Zk^d)\subseteq\Zk\}.
\]
Here the analytic hypothesis does the first step: once polynomiality is
known, what remains is the classical integer-valued-polynomial mechanism and
the constant-term pullback already used in the collection's omnific reports.

\begin{lemma}[Rationality of ordinary coefficients]\label{ent:as:lem:rational}
If $P\in\kk[X]$ takes values in $\Zk$ on $\Zk^d$, then
$P\in\operatorname{Frac}(\Zk)[X]$.
\end{lemma}
\begin{proof}
Choose a bound $m_j$ for the degree in $X_j$ and interpolate on the grid
$\prod_j\{0,1,\ldots,m_j\}$. The tensor-product Lagrange formula expresses
$P$ as a finite linear combination of basis polynomials with rational
coefficients, with coefficients the values of $P$, which lie in $\Zk$.
\end{proof}

\begin{theorem}[Exact entire integer-part stabilizer]\label{ent:as:thm:integer}
For $F\in\Ed$ the three conditions of \eqref{ent:as:eq:int-intro} are
equivalent. More explicitly,
\begin{equation}\label{ent:as:eq:integer-classification}
 \{F\in\Ed:F(\OzK^{\,d})\subseteq\OzK\}
 =\Int(\Zk^d,\Zk)\oplus\PiK[X]
\end{equation}
as additive groups; the decomposition is unique, with ordinary part
$\pg_0(F)$.
\end{theorem}
\begin{proof}
If $F(\OzK^{\,d})\subseteq\OzK$ then certainly $F(\Zk^d)\subseteq\OzK$.
Conversely assume the latter. The grid $\Zk^d$ is Zariski dense
(\cref{ent:as:cor:affine-sampling}), so \cref{ent:as:thm:sampling} gives
$F\in\Kle[X]$, and coefficientwise constant extraction gives a unique
decomposition $F=P+Q$ with $P=\pg_0(F)\in\kk[X]$ and $Q\in\PiK[X]$. For
$\mathbf z\in\Zk^d$, $P(\mathbf z)=\ct F(\mathbf z)\in\Zk$, so
$P\in\Int(\Zk^d,\Zk)$ by \cref{ent:as:lem:rational}. Now start from such a
decomposition and let $\mathbf x=\mathbf z+\mathbf u\in\OzK^{\,d}$ with
$\mathbf z\in\Zk^d$ and $\mathbf u\in\PiK^{\,d}$. The finite polynomial
expansion gives $P(\mathbf z+\mathbf u)-P(\mathbf z)\in\PiK$, since every term of
the difference contains a factor from $\PiK$, a $\kk$-vector-space ideal of
$\Kle$; likewise $Q(\mathbf x)\in\PiK$. Therefore
$F(\mathbf x)\in P(\mathbf z)+\PiK\subseteq\OzK$. Finally
$\kk[X]\cap\PiK[X]=0$, so the sum is direct and the decomposition unique.
\end{proof}

\paragraph{Printed once: the polynomial half.} Once $F$ is known to be a
polynomial, \eqref{ent:as:eq:integer-classification} is the omnific
numerical-polynomial theorem of the set-sized-quotients report,
$\operatorname{Num}_r(\Oz)=\Pi[\mathbf X]\oplus\Int(\Z^r)$ and its Gaussian form
(\texttt{osq:pm:thm:numreal}, \texttt{osq:pm:thm:numgaussian}), which that report
transfers to the fixed workspace rings $A_\Gamma=\Z\oplus J_\Gamma$ in
\texttt{osq:or:cor:workspace}; there $\Int(\Z^r)$ is our $\Int(\Zk^d,\Zk)$ and
$J_\Gamma$, with the reversed exponent sign, is our $\PiK$. The source credits
that mechanism to the collection's omnific reports and proposes as new only
the analytic step, the collapse of an entire series to a polynomial. That
report states the limit this step removes: ``entire Hahn series and
expressions of nonstandard degree are outside the theorems'' (the paragraph
after \texttt{osq:or:cor:workspace}, from its source 18). The equivalence with
$F(\OzK^{\,d})\subseteq\OzK$ for polynomials is also there, and the source
reproves it directly. (The merge's comparison; the source names
those reports but not these labels.)

\paragraph{The real binomial normal form.} Write
$\binom{X}{\beta}=\prod_{j=1}^d\binom{X_j}{\beta_j}$, with
$\binom{X}{m}=X(X-1)\cdots(X-m+1)/m!$ and $\binom X0=1$, let
$\nabla_jF(X)=F(X+\mathbf u_j)-F(X)$ be the forward difference in the $j$th
coordinate, $\mathbf u_j$ the $j$th coordinate vector, and
$\nabla^\beta=\prod_j\nabla_j^{\beta_j}$. Here $\nabla$ is a forward
difference, not a gradient.

\begin{corollary}[Binomial normal form over the real integer part]
\label{ent:as:cor:binomial}
Let $\kk=\R$. An entire function preserves $\OzK$ if and only if it has a
finite expansion
\begin{equation}\label{ent:as:eq:binomial}
 F(X)=\sum_\beta b_\beta\binom X\beta,\qquad b_\beta\in\OzK .
\end{equation}
The coefficients are unique, and $b_\beta=(\nabla^\beta F)(0)$.
\end{corollary}
\begin{proof}
In one variable $\nabla\binom Xm=\binom X{m-1}$ for $m\ge1$ and $\nabla1=0$, so
evaluation at zero shows that the finite Newton expansion of every polynomial
has coefficients $\nabla^mF(0)$; applying this in each variable gives the
formula. If $F$ preserves $\OzK$, it is a polynomial by
\cref{ent:as:thm:integer}, and every finite difference at zero is an integral
combination of values on $\N^d$, hence lies in $\OzK$. Conversely every
$\binom X\beta$ lies in $\Int(\Z^d,\Z)$ and so preserves $\OzK$ by
\cref{ent:as:thm:integer}; multiplying by $b_\beta\in\OzK$ and adding finitely
many terms preserves $\OzK$. Uniqueness holds because the binomial polynomials
form a triangular basis in each variable.
\end{proof}

The coefficients $b_\beta$ may be infinite and may have infinitely many
negative Hahn exponents; the finiteness in \eqref{ent:as:eq:binomial} concerns
the number of binomial basis elements, not the length of their coefficients.
For polynomials this is the finite-grid test \texttt{osq:pm:thm:grid}(iii) with
\texttt{osq:pm:lem:binomial}, read in the workspace through
\texttt{osq:or:cor:workspace}; what is added is again the analytic collapse.

\begin{example}[The Gaussian distinction]\label{ent:as:ex:gaussian}
$\binom X2$ is integer valued on $\Z$, but $\binom i2=(-1-i)/2\notin\Z[i]$. So
the ordinary binomial basis must not be used as a basis of
$\Int(\Z[i],\Z[i])$ without further work; the Gaussian classification in
\cref{ent:as:thm:integer} is exact as stated with that ring. Testing a
complex-coefficient function only on $\Z^d$ proves polynomiality but does not
by itself establish preservation of the full Gaussian omnific ring.
\end{example}

\begin{remark}[An infinite subset is not the whole arithmetic grid]
\label{ent:as:rem:subset}
$X/2$ is integer valued on the even ordinary integers but not on all of them.
So infinitely many omnific values prove polynomiality
(\cref{ent:as:cor:finite-points}), not membership in $\Int(\Z,\Z)$; the
full-grid hypothesis is required in \cref{ent:as:thm:integer}.
\end{remark}

\subsection{The positive-tail ideal and the exceptional locus}
\label{ent:as:subsec:tail}

\begin{definition}[Positive-tail ideal]\label{ent:as:def:tail}
For $F\in\Ed$, its \emph{positive-tail ideal} is
\begin{equation}\label{ent:as:eq:tail-ideal}
 I_+(F)=\bigl(\pg_\eta(F):\eta>0\bigr)\subseteq\kk[X].
\end{equation}
The name refers to the forbidden positive valuation exponents, not to large
powers of the ordinary variables. For an ideal $I\subseteq\kk[X]$ write
$\mathbb V_\kk(I)=\{\mathbf s\in\kk^d:g(\mathbf s)=0\text{ for all }g\in I\}$.
\end{definition}

\begin{theorem}[Finite algebraic classification of ordinary exceptions]
\label{ent:as:thm:tail}
For $F\in\Ed$,
\begin{align}
 \{\mathbf s\in\kk^d:F(\mathbf s)\in\Kle\}
 &=\mathbb V_\kk\bigl(I_+(F)\bigr),\label{ent:as:eq:B-locus}\\
 \{\mathbf s\in\kk^d:F(\mathbf s)\in\OzK\}
 &=\bigl\{\mathbf s\in\mathbb V_\kk\bigl(I_+(F)\bigr):
 \pg_0(F)(\mathbf s)\in\Zk\bigr\}.\label{ent:as:eq:A-locus}
\end{align}
Finitely many positive scales $\eta_1,\ldots,\eta_l$ have actual scale
polynomials that generate $I_+(F)$. If $F$ is not a polynomial, then
$I_+(F)\ne0$ and its zero set is a proper algebraic subset of $\kk^d$.
\end{theorem}
\begin{proof}
By \eqref{ent:as:eq:scale-eval}, $F(\mathbf s)$ has no positive-valuation term
exactly when $\pg_\eta(F)(\mathbf s)=0$ for every $\eta>0$, which is
\eqref{ent:as:eq:B-locus}; the scale-zero condition gives
\eqref{ent:as:eq:A-locus}. The ring $\kk[X]$ is Noetherian: choose finitely
many generators of $I_+(F)$ and express each as a finite polynomial
combination of the original $\pg_\eta$; only finitely many scales occur, and
together they generate the same ideal. If the ideal is zero, every positive
scale vanishes and \cref{ent:as:cor:truncation} makes $F$ a polynomial.
Finally a nonzero polynomial over the infinite field $\kk$ cannot vanish on
all of $\kk^d$.
\end{proof}

The $\Kle$-locus is algebraic; the $\OzK$-locus need not be, because the
additional condition $\pg_0(\mathbf s)\in\Zk$ is arithmetic. For example
$F(X)=X$ has $\Kle$-locus $\kk$ and $\OzK$-locus $\Zk$. This distinction
prevents replacing the exact formula by a blanket algebraicity assertion for
all omnific-valued points.

\begin{corollary}[Several output functions]\label{ent:as:cor:several}
For $F_1,\ldots,F_m\in\Ed$ the simultaneous $\Kle$-valued ordinary locus is
$\mathbb V_\kk\bigl(I_+(F_1)+\cdots+I_+(F_m)\bigr)$; the simultaneous
$\OzK$-valued locus adds $\pg_0(F_j)(\mathbf s)\in\Zk$ for every $j$. Finitely
many actual component--scale polynomials suffice.
\end{corollary}
\begin{proof}
Intersect \eqref{ent:as:eq:B-locus} and \eqref{ent:as:eq:A-locus} over the
finitely many functions; intersections of zero sets correspond to sums of
ideals.
\end{proof}

\paragraph{A sharp univariate bound.} In one variable write $\gcd_+(F)$ for the
monic generator of $I_+(F)$ when that ideal is nonzero, and $\gcd_+(F)=0$ when
it is zero. So $\gcd_+(F)=1$ means that the ordinary $\Kle$-locus is empty, and
$\gcd_+(F)=0$ that it is all of $\kk$.

\begin{theorem}[Finite scale certificate in one variable]
\label{ent:as:thm:sharp}
Let $d=1$, suppose $I_+(F)\ne0$, and let $\pg_{\eta_0}(F)$ be any nonzero
positive-scale polynomial, of degree $n_0$. Then
\[
 \deg\gcd_+(F)\le n_0,\qquad
 \#\{s\in\kk:F(s)\in\OzK\}\le\#\{s\in\kk:F(s)\in\Kle\}\le\deg\gcd_+(F).
\]
At most $n_0+1$ actual positive-scale polynomials, including
$\pg_{\eta_0}(F)$, generate $I_+(F)$, and the bound $n_0+1$ is sharp in general.
\end{theorem}
\begin{proof}
The ideal is generated by the greatest common divisor of its scale
polynomials, so $\gcd_+(F)$ divides $\pg_{\eta_0}(F)$; the degree and root
bounds follow. Start with the monic associate of $\pg_{\eta_0}(F)$. While the
current gcd is not the gcd of the whole family, some further member strictly
reduces it, and each strict reduction lowers the degree by at least one; so at
most $n_0$ further reductions occur. This proves existence of the bound; it
gives no algorithm for recognizing that no further reduction will occur.

For sharpness choose distinct $c_0,\ldots,c_{n_0}\in\kk$ and put
$f_j(X)=\prod_{l\ne j}(X-c_l)$, $0\le j\le n_0$. All $f_j$ have degree $n_0$ and
their gcd is $1$, while omitting $f_j$ leaves every remaining polynomial
divisible by $X-c_j$. For any $\delta>0$ the polynomial
$F=\sum_{j=0}^{n_0}t^{(j+1)\delta}f_j$ has exactly these nonzero scale
polynomials, and no proper subcollection generates its positive-tail ideal.
For $n_0=0$ a single nonzero constant is the sharp one-scale example. The
sharpness concerns selecting actual scales, not arbitrary precomputed
polynomial combinations of them.
\end{proof}

\paragraph{A parallel in this report.} \Cref{ent:as:thm:tail} and the
exceptional-direction theorem \cref{ent:mv:thm:exceptional} use the same
finiteness mechanism --- the Hilbert basis theorem applied to the polynomial
coefficients of an entire series at fixed Hahn exponents --- for different
questions: there the homogeneous components $P_{n,\eta}$ decide on which
complex lines the series restricts to a polynomial, here their sums $\pg_\eta$
decide at which ordinary points the value has no positive tail. Neither theorem
implies the other, and the realization theorems
\cref{ent:mv:thm:realization,ent:as:thm:realization} use the same damping idea.
The source does not cite the fifth source's theorems; the comparison is the
merge's.

\subsection{Jets and multiplicities}
\label{ent:as:subsec:jets}

For a multi-index $\beta$ let $\partial^{[\beta]}$ be the Hasse derivative,
$\partial^{[\beta]}X^\alpha=\binom\alpha\beta X^{\alpha-\beta}$, zero unless
$\alpha\ge\beta$ coordinatewise; in characteristic zero it is the usual partial
derivative divided by $\beta_1!\cdots\beta_d!$. In one variable
$\partial^{[k]}F(b)$ is the Hasse jet coefficient $F^{[k]}(b)$ of
\cref{ent:subsec:notation}, and these are the Hasse derivatives of
\cref{ent:rc:lem:sum}.

\begin{lemma}[Hasse derivatives and scales]\label{ent:as:lem:hasse}
Hasse derivatives preserve $\Ed$ and commute with every scale map:
$\pg_\eta(\partial^{[\beta]}F)=\partial^{[\beta]}\pg_\eta(F)$.
\end{lemma}
\begin{proof}
Each new coefficient is an old one multiplied by an ordinary binomial integer,
so its valuation is unchanged or becomes $+\infty$, while its total degree is
lower by $|\beta|$. For a fixed weight $\gamma$ the growth expression is shifted
by the constant $-|\beta|\gamma$, so \eqref{ent:mv:eq:growth} persists.
Coefficient extraction commutes with this finite scalar multiplication term by
term.
\end{proof}

For divisible $\Gamma$ the first assertion is in the proof of
\cref{ent:rc:lem:sum}; the argument is the same.

\begin{definition}[Positive-tail jet ideals]\label{ent:as:def:jet-ideal}
For an ordinary integer $m\ge1$ put
\[
 I_{+,m}(F)=\bigl(\partial^{[\beta]}\pg_\eta(F):\eta>0,\ |\beta|<m\bigr)
 \subseteq\kk[X],
\]
so that $I_{+,1}(F)=I_+(F)$.
\end{definition}

\begin{theorem}[One finite set of scales controls every jet order]
\label{ent:as:thm:jets}
The ordinary points at which every Hasse derivative of $F$ of total order less
than $m$ is $\Kle$-valued form $\mathbb V_\kk(I_{+,m}(F))$; for $\OzK$-valued
jets add the conditions $(\partial^{[\beta]}\pg_0(F))(\mathbf s)\in\Zk$ for
$|\beta|<m$. If $\pg_{\eta_1}(F),\ldots,\pg_{\eta_l}(F)$ generate $I_+(F)$, then
for \emph{every} $m\ge1$
\begin{equation}\label{ent:as:eq:finite-jets}
 I_{+,m}(F)=\bigl(\partial^{[\beta]}\pg_{\eta_j}(F):
 1\le j\le l,\ |\beta|<m\bigr).
\end{equation}
Raising the jet order needs more derivatives, but no new valuation scales.
\end{theorem}
\begin{proof}
The pointwise assertions follow from \cref{ent:as:thm:tail} applied to the
finitely many derivatives, with \cref{ent:as:lem:hasse}. For
\eqref{ent:as:eq:finite-jets} write a positive-scale polynomial as
$\pg_\eta=\sum_ja_j\pg_{\eta_j}$ with $a_j\in\kk[X]$. The Hasse product rule
gives
\[
 \partial^{[\beta]}\pg_\eta
 =\sum_j\sum_{\beta'\le\beta}
 \bigl(\partial^{[\beta']}a_j\bigr)\bigl(\partial^{[\beta-\beta']}\pg_{\eta_j}\bigr),
\]
and if $|\beta|<m$ every derivative of a selected scale on the right has total
order less than $m$. That is one inclusion; the other holds because the
selected scales belong to the original family.
\end{proof}

\begin{corollary}[The univariate contact divisor]\label{ent:as:cor:contact}
Let $d=1$ and $I_+(F)\ne0$. At $s\in\kk$ all Hasse derivatives of orders
$0,\ldots,m-1$ are $\Kle$-valued if and only if $\ord_s\gcd_+(F)\ge m$. In
particular
\[
 \sum_{s\in\kk:\ \gcd_+(F)(s)=0}\ord_s\gcd_+(F)\ \le\ \deg\gcd_+(F)\ \le\ n_0
\]
for any nonzero positive scale of degree $n_0$.
\end{corollary}
\begin{proof}
Vanishing of the first $m$ Hasse derivatives of a polynomial at $s$ is
divisibility by $(X-s)^m$, by its Taylor expansion in $X-s$. Requiring this
for every positive-scale polynomial is requiring it for their gcd. The
multiplicity sum is bounded by the degree; over $\C$ the sum over all roots
equals the degree, while over $\R$ nonreal roots do not contribute.
\end{proof}

\begin{remark}[The residue condition remains independent]
\label{ent:as:rem:residue}
For $F(X)=\binom X2$ there is no positive tail at any derivative order, but
$F'(0)=-\tfrac12\notin\Z$. So a $\Kle$-valued jet need not be an $\OzK$-valued
jet, and the extra scale-zero conditions of \cref{ent:as:thm:jets} are needed
even when $\gcd_+(F)=0$.
\end{remark}

\subsection{Realizing every algebraic obstruction}
\label{ent:as:subsec:realization}

Positive-tail ideals are not confined to a small family: apart from the
cofinality obstruction to nonpolynomial entireness itself, there is no
algebraic restriction.

\paragraph{Printed once: the cofinality dichotomy.} The source's
``background cofinality dichotomy'' --- for $\Gamma\ne0$ a nonpolynomial member
of $\Ed$ exists if and only if $\cf(\Gamma)=\aleph_0$, and for $\Gamma=0$ all
entire series are polynomials --- is the last clause of
\cref{ent:thm:multi-criterion} with \cref{ent:mv:prop:realize-support}, and in
one variable \cref{ent:cor:cofinality}; the source credits it to this report.
Its witness $\sum_{n\ge1}t^{nu_n}X_1^n$, for a strictly increasing positive
cofinal sequence $(u_n)$, is the series of \cref{ent:cor:cofinality}, and its
proof --- $n(u_n+\gamma)>b$ eventually, because $u_n+\gamma$ is eventually
positive and larger than $b$, and multiplying by a positive integer cannot
decrease it --- needs no divisibility.

\begin{theorem}[Universal ideal realization]\label{ent:as:thm:realization}
Suppose $\Gamma\ne0$ and let $I\subseteq\kk[X]$ be an ideal. There is a
polynomial $F\in\Ed$ with $I_+(F)=I$ and $\pg_0(F)=0$. If $I\ne0$ and
$\cf(\Gamma)=\aleph_0$, then $F$ can be chosen nonpolynomial. Consequently
every proper algebraic subset of $\kk^d$ is the exact ordinary $\OzK$-valued
locus of a nonpolynomial entire function when $\cf(\Gamma)=\aleph_0$; the empty
locus is allowed, and the whole space cannot be the locus of a nonpolynomial
entire function.
\end{theorem}
\begin{proof}
Choose generators $g_1,\ldots,g_l$ of $I$ and $\delta>0$. Then
$F_0=\sum_{j=1}^lt^{j\delta}g_j$ has the required ideal and zero scale-zero
polynomial; for $I=0$ take $F_0=0$. Now let $I\ne0$, with $g_1\ne0$, choose a
strictly increasing positive cofinal sequence $(u_n)$ with $u_1>l\delta$, and
put
\begin{equation}\label{ent:as:eq:realization}
 F(X)=F_0(X)+\sum_{n\ge1}t^{nu_n}X_1^n\,g_1(X).
\end{equation}
The scales $nu_n$ strictly increase and lie above every $j\delta$; every new
scale polynomial $X_1^ng_1$ lies in $I$, while the initial scales already
generate $I$, so $I_+(F)=I$, and there is no scale-zero term. For entireness,
expand $g_1$ into its finitely many monomials: each summand of the $n$th block
has total degree $n+e$ for one of finitely many $e\ge0$ and valuation $nu_n$,
and its weighted valuation $n(u_n+\gamma)+e\gamma$ exceeds every prescribed
bound for all large $n$, uniformly in these $e$. This is
\eqref{ent:mv:eq:growth}; only finitely many blocks affect any fixed power
coefficient. The scale polynomials $X_1^ng_1$ have unbounded degrees, so $F$ is
not a polynomial. Finally \cref{ent:as:thm:tail} identifies the $\OzK$-valued
locus with $\mathbb V_\kk(I)$, because $\pg_0=0\in\Zk$. A proper algebraic set
has nonzero vanishing ideal over the infinite field $\kk$; the empty locus is
realized by $I=\kk[X]$; the whole space is excluded by
\cref{ent:as:thm:sampling}.
\end{proof}

\begin{remark}[Nonreduced information is retained]\label{ent:as:rem:nonreduced}
The realization concerns the ideal itself, not only its radical or its zero
set. For example $(X^m)$ is realized with its prescribed order of contact at
zero, and the finite jet ideals of \cref{ent:as:thm:jets} then recover the
multiplicity. This is why $I_+(F)$ is kept rather than replaced at once by its
zero set.
\end{remark}

\begin{remark}[The full surreal class has a different boundary; printed once]
\label{ent:as:rem:full-class}
The source's corollary that an ordinary power series with coefficients in
$\No$ or $\No[i]$, strongly Hahn-summable at every point of the full class, is
a polynomial is \cref{ent:cor:all-no} in one variable and the whole-family half
of \cref{ent:mv:prop:allscale} in several, which credit the collection's
all-scale theme (\cref{ent:rem:rigidity-credit}). Its short route is the same
descending-leading-exponent argument: the necessity half of the coefficient
criterion uses only set-sized families, so it applies to the normal-form
valuation with value group the surreal class; infinitely many nonzero
coefficients would then have valuations cofinal in the whole class, but they
form a set, and every set of surreals has a surreal upper bound by the cut
construction. So the nonpolynomial realization theorem is a theorem about
set-sized Hahn workspaces; it constructs no globally entire nonpolynomial
function on $\No[i]$.
\end{remark}

\subsection{Existence of finite certificates versus computability}
\label{ent:as:subsec:computability}

Noetherianity proves that finitely many scales control all ordinary
exceptions. It gives no way to recognize, while reading an infinite
coefficient stream, that the last necessary scale has already appeared. The
next construction shows that the distinction is unavoidable, even with
completely explicit growth bounds.

\paragraph{An explicit presentation.} Work in $\K=\R((t))$, with value group
$\Z$ (not divisible, which \cref{ent:as:conv:standing} allows), and restrict
the coefficients to $\Q[t]$; this is a real surreal workspace under
$t=\omega^{-1}$, and the same example works in $\C((t))$ with the Gaussian
omnific ring. Put
\begin{equation}\label{ent:as:eq:base-machine}
 G_{\mathrm{sq}}(Z)=Z(Z-1)\sum_{n\ge0}t^{(n+1)^2}Z^n .
\end{equation}
For a Turing machine $\mathtt M$ on its specified input let
$h_{\mathtt M}(N)=1$ if $\mathtt M$ halts for the first time exactly at step
$N$, and $h_{\mathtt M}(N)=0$ otherwise, a computable function of
$(\mathtt M,N)$ by bounded simulation, and define
\begin{equation}\label{ent:as:eq:machine-family}
 F_{\mathtt M}(Z)=G_{\mathrm{sq}}(Z)
 +\sum_{N\ge0}h_{\mathtt M}(N)\,t^{(N+2)^2+1}Z^{N+1}.
\end{equation}
At most one summand of the second series is nonzero, but no prior decision
about its existence is needed to compute any power coefficient.

\begin{lemma}[The machine presentation]\label{ent:as:lem:machine}
The coefficients $a_m$ of $F_{\mathtt M}$ are uniformly computable
polynomials in $t$ with integer coefficients:
\begin{equation}\label{ent:as:eq:machine-coefficients}
 a_0=0,\qquad a_1=-t+h_{\mathtt M}(0)\,t^5,\qquad
 a_m=t^{(m-1)^2}-t^{m^2}+h_{\mathtt M}(m-1)\,t^{(m+1)^2+1}\quad(m\ge2).
\end{equation}
For every $m\ge1$,
\begin{equation}\label{ent:as:eq:quadratic-bound}
 v(a_m)\ge(m-1)^2 .
\end{equation}
Each $F_{\mathtt M}$ is nonpolynomial and entire, and $\pg_0(F_{\mathtt M})=0$.
\end{lemma}
\begin{proof}
Expanding \eqref{ent:as:eq:base-machine} gives the displayed coefficients, and
the extra summand in degree $m$ is found by simulating $\mathtt M$ through step
$m-1$. The smallest exponent of $a_m$ cannot cancel: for $m\ge2$ it is
$(m-1)^2$, and the exceptional exponent exceeds $m^2$; for $m=1$ it is $1$. So
every $a_m$ with $m\ge1$ is nonzero and satisfies
\eqref{ent:as:eq:quadratic-bound}. Quadratic growth gives
\eqref{ent:mv:eq:growth} at every integer weight, so $F_{\mathtt M}$ is entire.
All exponents are positive, so $\pg_0=0$.
\end{proof}

\begin{theorem}[An undecidable ordinary omnific-value test]
\label{ent:as:thm:undecidable}
The family \eqref{ent:as:eq:machine-family} satisfies
\begin{equation}\label{ent:as:eq:machine-ideal}
 I_+(F_{\mathtt M})=
 \begin{cases}(Z(Z-1)),&\mathtt M\text{ never halts},\\
 (Z),&\mathtt M\text{ halts},\end{cases}
\end{equation}
so its ordinary $\Kle$-valued and $\OzK$-valued loci coincide and are
\begin{equation}\label{ent:as:eq:machine-locus}
 \{0,1\}\ \text{if }\mathtt M\text{ never halts},\qquad
 \{0\}\ \text{if }\mathtt M\text{ halts}.
\end{equation}
No algorithm decides $F(1)\in\OzK$ for all computably presented entire series
of this form; nor can any algorithm always output their complete exceptional
locus, their positive-tail ideal, or a certified complete finite collection of
coefficient scales. The index set $\{\mathtt M:F_{\mathtt M}(1)\in\OzK\}$ is
$\Pi^0_1$-complete.
\end{theorem}
\begin{proof}
At the square scale $(n+1)^2$ the scale polynomial of $G_{\mathrm{sq}}$ is
$Z^{n+1}(Z-1)$; their gcd is $Z(Z-1)$, already supplied by $n=0$. If
$\mathtt M$ halts at $N$, one further scale appears, with polynomial $Z^{N+1}$;
its scale $(N+2)^2+1$ lies strictly between $(N+2)^2$ and $(N+3)^2$, so it is
not a square and cannot collide with a base scale. The gcd becomes
$\gcd(Z(Z-1),Z^{N+1})=Z$. This proves \eqref{ent:as:eq:machine-ideal}, and
since $\pg_0=0$, \cref{ent:as:thm:tail} gives \eqref{ent:as:eq:machine-locus}.
Directly, $F_{\mathtt M}(1)=0$ if $\mathtt M$ never halts, and
$F_{\mathtt M}(1)=t^{(N+2)^2+1}\notin\OzK$ if it halts at $N$. A deciding
algorithm would decide nonhalting, hence halting, which is impossible by the
usual diagonal argument; computing any of the listed complete finite outputs
would decide the same test, by a gcd or by evaluating finitely many
polynomials at $1$. Finally nonhalting is a $\Pi^0_1$ condition, and every
condition $\forall n\,R(e,n)$ with computable $R$ reduces to it through the
machine that searches for a counterexample and halts on finding one; so
\eqref{ent:as:eq:machine-locus} gives $\Pi^0_1$-completeness of this explicit
index set.
\end{proof}

\paragraph{What remains effectively available.} The obstruction is not
hidden in a convergence oracle: the same quadratic bound holds for every input
of the reduction, and for any positive integer exponent $e$ only the
coefficients with $m\le1+\lfloor\sqrt e\rfloor$ contribute to
$[t^e]F_{\mathtt M}(1)$, each by a finite calculation. Nonmembership is
eventually witnessed by one nonzero scale coefficient, whereas membership may
require knowing that no later witness will ever appear. For a \emph{finite}
Hahn-polynomial input over an effective field such as $\Q$ the scale
polynomials form a finite list, so ordinary Gr\"obner-basis computations
produce $I_+(F)$, and in one variable a finite gcd suffices;
\cref{ent:as:cor:normal-form} also gives effective coefficientwise reduction
when an algebraic ideal is supplied in advance. Neither is an algorithm for
the remaining integer conditions $\pg_0(\mathbf s)\in\Zk$ in arbitrary
dimension. For an infinite presentation a finite list of scales is only a
\emph{candidate} certificate until all later scales are known to lie in its
ideal; \cref{ent:as:thm:undecidable} says that no universal terminating
procedure supplies that last assurance in this presentation. It does not rule
out such procedures for restricted recurrence, differential, rational or
finite-expression languages.

Undecidability of other problems about non-Archimedean entire-function rings
already has a literature; for example Garcia-Fritz and Pasten treat Hilbert's
tenth problem through positive existential interpretations
\cite{GarciaFritzPasten}, a paper whose abstract only the source consulted.
\Cref{ent:as:thm:undecidable} is not a new negative solution to Hilbert's
tenth problem; it is an explicit one-point positive-tail certificate
obstruction with a uniform coefficient-growth bound. It also does not answer
\cref{ent:mv:q:effective} or the decision clause of
\cref{ent:q:image-remaining}, which ask about other tests (well-ordering of
weight sets, membership in the rings $V_n$) on other presentations (the
merge's comparison).

\subsection{Sharp boundaries and counterexamples}
\label{ent:as:subsec:boundaries}

\paragraph{Strong evaluation at every ordinary point is not entireness.} In
$\K=\kk((t^\Q))$ consider
\begin{equation}\label{ent:as:eq:local-counterexample}
 F_{\mathrm{loc}}(Z)=\sum_{n\ge1}t^{-1/n}Z^n .
\end{equation}
For every ordinary $c\ne0$ the family is strongly summable, with the
increasing well-ordered support $\{-1,-\tfrac12,-\tfrac13,\ldots\}$, and for
$c=0$ all terms of positive degree vanish; so $F_{\mathrm{loc}}(c)\in\PiK\subseteq
\OzK$ for every ordinary $c$, although $F_{\mathrm{loc}}$ is not a polynomial.
At $Z=t^{-1}$ the term exponents $-n-1/n$ strictly decrease, so the family is
not summable there. The global growth step, not coefficient extraction on
$\kk$ alone, is what \cref{ent:as:thm:sampling} needs.

\paragraph{Infinitely many infinite omnific zeros are possible.} In
$\K=\kk((t))$ form
\begin{equation}\label{ent:as:eq:qproduct}
 P_{\mathrm{esc}}(Z)=\prod_{m\ge1}\bigl(1-t^mZ\bigr).
\end{equation}
Its coefficient of $Z^n$ is the strong sum over the $n$-element subsets of the
positive integers, with sign $(-1)^n$ and exponent the sum of the subset; only
finitely many subsets have a given sum, and the least exponent,
$n(n+1)/2$, is attained only by $\{1,\ldots,n\}$. So every coefficient is
nonzero with valuation exactly $n(n+1)/2$, which proves nonpolynomiality and
entireness. Removing the $m$th factor gives an entire product with the same
quadratic bound, and $P_{\mathrm{esc}}=(1-t^mZ)\cdot(\text{that product})$ gives
$P_{\mathrm{esc}}(t^{-m})=0$; every $t^{-m}$ is an infinite element of $\OzK$.
So $P_{\mathrm{esc}}$ has infinitely many omnific zeros, none at ordinary
points: its first positive scale polynomial is $-Z$, its scale-zero polynomial
is $1$, and its ordinary $\OzK$-valued locus is exactly $\{0\}$. The product
mechanism is classical and is this report's canonical-product setting
(\cref{ent:thm:product}, stated there for divisible groups; the value group
here is $\Z$, and the source's direct argument above is what applies); here it
marks the boundary of the sampling theorem.

\paragraph{Rational maps, classical entire functions and finite tests.} The
rational function $t^{-1}/(1+Z^2)$ takes purely infinite values at every
ordinary real integer; it is neither a polynomial nor a Hahn-entire power
series. Over a real Hahn field it has no poles at field-valued points, so
absence of real poles is no substitute for the entire power-series hypothesis;
over the complex field it has the expected poles at $\pm i$. The ordinary
function $\sin(\pi Z)$ has infinitely many integer zeros; this does not
contradict the theorem, because its Taylor coefficients all have valuation
zero, so its nonzero terms do not form a strongly summable family at an
ordinary nonzero point. Ordinary complex convergence and strong Hahn summation
are different notions. Finally, no fixed finite sample can replace a
Zariski-dense one: for finite $\Theta\subset\Z$ the entire polynomial
$t\prod_{n\in\Theta}(Z-n)$ vanishes on $\Theta$ and has a forbidden positive
term at every ordinary integer outside $\Theta$; this is a simple polynomial
obstruction, not a new general theorem about finite tests.

\paragraph{Why an ordinary coefficient field must be named.} The scale
polynomials use the chosen coefficient field and monomial section. The proofs
are invariant under isomorphisms that transport these data together with the
positive-tail cut, and they do not show that arbitrary pure-field automorphisms
preserve the construction. This matters in the surcomplex setting, where a
pure field structure does not by itself specify the real form, conjugation or
the valuation; the omnific-preserving-automorphisms report studies these
distinctions \cite{RepoOPA}, and they must not be removed silently from the
hypotheses of an arithmetic sampling theorem.

\subsection{Source~11's audit, and what these subsections do not claim}
\label{ent:as:subsec:limits}

\paragraph{The proposed addition.} The analytic-to-algebraic sequence is:
entire growth, then polynomial scale coefficients
(\cref{ent:as:lem:scale}), then degree-controlled lifting
(\cref{ent:as:lem:degree-transfer,ent:as:lem:matrix-degree}), then arithmetic
and ideal classification. The bounded-degree transfer is the key analytic step
of the module theorem: finite generation alone would not justify lifting
infinitely many scale equations to entire functions, because arbitrary
polynomial lifts could have uncontrolled degree, and the Gr\"obner degree bound
supplies the missing uniform control. Likewise finite generation of $I_+(F)$
gives an exact finite certificate but no effective stopping rule; the machine
family separates the two while keeping a uniform quadratic bound,
nonpolynomiality, zero ordinary residue and a two-point candidate exceptional
set. The realization theorem shows that the ideal classification has no
hidden algebraic restriction when nonpolynomial entire functions exist.

\paragraph{Proposed contributions and established ingredients.} The source
proposes as new: the uniform bounded-degree transfer at arbitrary rank; the
coefficientwise lifting over $\kk[X]$, faithful flatness and exact entire
normal forms; dense ordinary sampling, including invertible affine Hahn images
of dense sets, and the restriction theorem; the classification
\eqref{ent:as:eq:int-intro}, where the new step is the analytic collapse; the
finite positive-tail ideal, the sharp one-variable bound and the finite-scale
jet certificate; realization of every nonzero ideal exactly in the regime that
permits nonpolynomial series; and the explicit quadratic-growth family with
locus $\{0\}$ or $\{0,1\}$ by halting. It lists as established, not claimed:
Hahn normal forms and the omnific support cut; strong summability, support
closure and regrouping; the entire coefficient criterion and the cofinality
obstruction, already in this report; the Hilbert basis theorem, module
Gr\"obner bases, polynomial division and the equational criterion of flatness;
integer-valued polynomials and finite differences; the constant-term pullback
already developed in the collection; the halting obstruction and the basic
infinite-product example; and whole-class polynomial rigidity. Against this
report, the printed-once items are listed in warning (2) of
\cref{ent:as:subsec:setting}; everything else in these subsections is the
source's new material, except the passages marked as the merge's.

\paragraph{Critical checkpoints recorded by the source.} Summability comes
before regrouping: the criterion is applied to the full multivariate family,
affine substitution is bounded before regrouping, and matrix lifts are shown
entire before their reconstruction is used in an equality. Degree control is
uniform: $c_{\mathbf M}$ depends on the fixed matrix, not on the scale, which is
what lets \cref{ent:as:lem:degree-transfer} hold at every weight, and no rank-one
or Archimedean hypothesis is hidden. The arithmetic cut uses all coefficients:
a negative leading valuation does not imply membership in $\OzK$ or $\Kle$, and
the scale-zero polynomial is kept apart from the power-series constant and from
$I_+(F)$. Set and class scopes are not interchanged: nonpolynomial realization
uses a countable cofinal sequence in a set-sized group, the full surreal class
has no set cofinal subset, and $\operatorname{Frac}(\OzK)=\K$ is not assumed.
Noetherian is not computable: the finite certificate is an existence theorem,
its sharp size bound is not a stopping bound, and the halting construction
isolates the difference without noncomputable real constants.

\paragraph{The source's comparison with this repository.} It pinned the
repository at \texttt{343dc2c}, read it through a connector (a local clone
failed for lack of name resolution, so no local build followed), and inspected
the root and documentation guides, the initial portion of this report's
README, a passage of this article beginning at its line~330, the initial
portion of the set-sized-quotients README, the question and comparison portion
of the omnific-preserving-automorphisms README and that report's source audits
12 and 13. Several connector responses were truncated, and the source does not
treat requested ranges as proof that all their contents were returned. This
report's directory is unchanged between that pin and the base of this merge,
so its account of this report is current and was checked: the coefficient
criterion, cofinality principle, scalar-extension distinctions and escaping-node
interpolation are here. The omnific Diophantine-geometry report was credited
only through the guides and not read. Its non-claim ``Polynomial lifting does
not license evaluating infinite power series'' is the boundary that
\cref{ent:as:thm:integer} and \cref{ent:as:subsec:boundaries} address for
entire series in a fixed workspace; this is the merge's observation, recorded
in \cref{ent:sec:collection}.

\paragraph{Limits recorded by source~11.} The following are retained without
weakening.
\begin{enumerate}[label=(\arabic*),leftmargin=2em]
\item Strong Hahn summation is not ordinary real or complex convergence;
$\sin(\pi Z)$ and $t^{-1}/(1+Z^2)$ mark that boundary.
\item Nonpolynomial results concern fixed set-sized Hahn workspaces, not an
unrestricted analytic function on the proper class $\No[i]$; an ordinary power
series strongly entire on the whole class is a polynomial
(\cref{ent:as:rem:full-class}).
\item The exceptional-locus theorem is algebraic for $\Kle$-valued points only;
the $\OzK$-valued locus carries the residue condition $\pg_0(\mathbf s)\in\Zk$
and need not be Zariski closed.
\item The leading valuation is not an omnific test; every positive exponent is
tested.
\item Flatness and the exact image theorem are over $\kk[X]$, not $\K[X]$;
nothing is claimed for polynomial matrices with Hahn coefficients
(\cref{ent:as:q:matrices}).
\item $\Ed$ is not claimed to be Noetherian, a principal ideal domain or a
B\'ezout domain.
\item $\OzK$ is not identified with an integer polynomial ring, and
$\operatorname{Frac}(\OzK)=\K$ is not assumed for a set-sized workspace.
\item The ordinary binomial basis is not a basis of $\Int(\Z[i]^d,\Z[i])$;
testing a complex-coefficient function on $\Z^d$ proves polynomiality but not
preservation of the Gaussian omnific ring.
\item An infinite sample gives polynomiality, not membership in
$\Int(\Zk^d,\Zk)$; the classification needs the full grid.
\item \eqref{ent:as:eq:int-intro} is a classification inside $\Ed$, not an
assertion about arbitrary functions on $\OzK$.
\item Sharpness of $n_0+1$ concerns selecting actual scales, not arbitrary
polynomial combinations of them.
\item Finite certificates exist but are not computable in general. The
undecidability concerns the explicit presentation
\eqref{ent:as:eq:machine-family}; restricted recurrence, differential, rational
or finite-expression classes are not ruled out, and the integer conditions
$\pg_0(\mathbf s)\in\Zk$ in arbitrary dimension are not decided.
\item \Cref{ent:as:thm:undecidable} is not a new negative solution to
Hilbert's tenth problem.
\item The scale polynomials depend on the coefficient field and the monomial
section; invariance is claimed only under isomorphisms transporting them with
the positive-tail cut, and the unrestricted Gaussian automorphism problems of
the collection's companion reports are not settled.
\item The finite checks (\cref{ent:app:checks}, Suite H) are not proofs: they
do not establish infinite summability, arbitrary-rank lifting, faithful
flatness or undecidability, and the machine-family checks use specified
bounded examples and decide no halting problem.
\item There is no Lean or other proof-assistant certificate, no independent
referee report, and no claim that a repository proof build was run.
Historical priority is not certified, no named published conjecture is claimed
settled, and the ten questions are proposed directions, not a list of
previously published open problems.
\item The repository comparison was targeted: not every manuscript of the
collection was read, the omnific Diophantine-geometry article was not audited,
and a failed search is not evidence that a result has never appeared. The
Wikipedia overview the source cites \cite{WikiSurreal} is orientation only.
\end{enumerate}

\setcounter{equation}{\value{entassaveeq}}%
\endgroup

\section{Prime spectra above a one-node-per-shell product}
\label{ent:ps:sec:spectrum}

This section is the sixth source's contribution. It describes every prime
ideal of $\E$ that contains a canonical product with one simple node per
shell, not only the maximal ones of \cref{ent:sm:thm:maximal-all}; it decides
when the quotient by such a product is zero-dimensional; it computes the
contraction of these primes under a cofinal enlargement of the value group,
with an explicit surreal example in which the prime spectrum refines while
the function and its zeros do not change; and it locates the multiplicity
threshold beyond which the reduced description is incomplete. It answers the
``prime ideals'' clause of \cref{ent:q:spectrum} for one simple node per
shell, and only for that case; the non-claim about prime ideals recorded at
the end of \cref{ent:sm:subsec:all-maximal} is kept there as written, now
with a pointer to this section. Three warnings govern the section.

\begin{enumerate}[label=(\arabic*),leftmargin=2em]
\item \emph{What is inherited.} The analytic input is the quotient theorem
\cref{ent:thm:quotient} in its one-node form \cref{ent:cor:single-node}, over
$\R((t^\Gamma))$ by \cref{ent:prop:real-form}. The sixth source reproved the
simple-node case so as not to depend on this report's preparation theory; its
argument is the first route of \cref{ent:thm:image} specialized to one simple
node per shell, and it is printed once, by citation
(\cref{ent:ps:rem:route}). The maximal ideals, their residue fields and the
maximal spectrum are \cref{ent:thm:maximal,ent:sm:thm:maximal-all,ent:sm:thm:beta};
the sixth source recovers them from the full spectrum and does not claim them.
\item \emph{What is classical.} The idempotent and ultraproduct mechanism for
the primes of a product ring, including products of Pr\"ufer domains described
by ultrafilters, ultraproduct valuations and cuts in value groups, is due to
Finocchiaro, Frisch and Windisch \cite{FFW}, a direct precedent for
\cref{ent:ps:lem:ultrafilter,ent:ps:prop:fibre}. The correspondence between
the primes of a valuation domain and the convex subgroups of its value group
is standard \cite{StacksValuation}. The Hahn support calculus is Neumann's
\cite{Neumann}, and the normal forms used in \cref{ent:ps:subsec:splitting}
are Conway's and Gonshor's \cite{Conway,Gonshor}. The sixth source proposes
as new only the exact analytic realization of these mechanisms, the
order-unit dichotomy for the reduced divisor and the cofinal splitting
construction, and it certifies none of them against the whole literature
(\cref{ent:ps:subsec:limits}).
\item \emph{Words and letters.} \emph{Spectrum} means prime spectrum
$\Spec$ throughout this section; the maximal spectrum is written
$\operatorname{Max}$, as in \cref{ent:sm:thm:beta}. The sixth source's
letters are translated by the sixth manuscript's table in
\cref{ent:subsec:notation}. In
particular its sequence ring $\mathcal A$ is written $\RD$ here, because our
$\A$ is the Hahn algebra of \cref{ent:sec:growth}; its convex subgroups $C$
and $D$ are written $\mathcal C$ and $\mathcal C'$, because our $C_n$ are
corrections and our $D$ is a divisor; and its coefficient field $k$ is
written $\kk$.
\end{enumerate}

\subsection{Setting and the inherited quotient}
\label{ent:ps:subsec:setting}

\begin{convention}[Standing hypotheses of this section]
\label{ent:ps:conv:standing}
Unless a second group is introduced explicitly, $\Gamma$ is a nonzero
set-sized divisible ordered abelian group with $\cf(\Gamma)=\aleph_0$. The
coefficient field $\kk$ is $\R$ or $\C$, with the trivial valuation; in this
section $K_\Gamma$ denotes $\kk((t^\Gamma))$ and $\mathcal E_\Gamma$ its
ring of strongly entire series (for $\kk=\C$ these are $\K$ and $\E$, for
$\kk=\R$ the real forms of \cref{ent:prop:real-form}). In
\cref{ent:sec:multivariable} the letter $\kk$ denoted an arbitrary trivially
valued field; here it is $\R$ or $\C$, although the sixth source records that
most of its proofs work over any trivially valued field. Fix a strictly
increasing positive sequence $\gamma_1<\gamma_2<\cdots$ cofinal in $\Gamma$
and put
\begin{equation}\label{ent:ps:eq:product}
 \rho_n=t^{-\gamma_n},\qquad
 D=\sum_{n\ge1}[\rho_n],\qquad
 P_D(Z)=\prod_{n\ge1}\bigl(1-t^{\gamma_n}Z\bigr)
       =\prod_{n\ge1}\Bigl(1-\frac Z{\rho_n}\Bigr).
\end{equation}
The letters $H_n=H_{\gamma_n}$, $V_n=V_{\gamma_n}$ and
$\mathfrak q_{\gamma_n}$ are those of \cref{ent:def:coarse} and
\cref{ent:lem:coarse-residue}, and $\kappa_n=\kk((t^{H_n}))$. Ultrafilters
are ultrafilters on $\N_{>0}$; nonprincipal ones exist by the axiom of choice,
which is used, and no particular one is computed. All rings, products and
ultraproducts below are sets.
\end{convention}

The node $\rho_n$ is the shell normalizer $\tau_n$ of
\eqref{ent:eq:shellpoly}, and $P_D$ is the canonical product of
\cref{ent:thm:product}. The product in \eqref{ent:ps:eq:product} is not the
product $P$ of \eqref{ent:eq:infinite-product}, which runs over $j\ge0$; for
$\Gamma=\Gamma_\infty$ and $\gamma_n=e_n$ it is the function $f$ of
\eqref{ent:eq:f-explicit}. By \cref{ent:thm:quotient,ent:cor:single-node}
(and \cref{ent:prop:real-form} over $\R$), evaluation at the nodes induces a
$K_\Gamma$-algebra isomorphism
\begin{equation}\label{ent:ps:eq:quotient}
 \begin{gathered}
 \mathcal E_\Gamma/(P_D)\ \xrightarrow{\ \sim\ }\ \RD,
 \qquad F\longmapsto\bigl(F(\rho_n)\bigr)_n ,\\
 \RD=\bigl\{(b_n)_{n\ge1}\in K_\Gamma^{\N_{>0}}:
      b_n\in V_n\text{ for all sufficiently large }n\bigr\},
 \end{gathered}
\end{equation}
the ring displayed in \cref{ent:subsec:invisible}; operations
are coordinatewise. We write $\RD^\Gamma$ when the group varies. From now on
$\mathcal E_\Gamma/(P_D)$ is identified with $\RD$. A prime of $\RD$ is the
same as a prime of $\mathcal E_\Gamma$ containing $P_D$, through its inverse
image; in particular a \emph{minimal} prime of $\RD$ is minimal \emph{above}
$(P_D)$, not a minimal prime of the domain $\mathcal E_\Gamma$.

\begin{remark}[The sixth source's analytic route, printed once]
\label{ent:ps:rem:route}
The sixth source proves, for the monomial nodes \eqref{ent:ps:eq:product}, a
coefficient criterion, Gauss multiplicativity, the canonical product with a
weighted bound on the formal inverse of its tails, division by a linear
factor, translation and additivity of orders, the principal kernel and the
exact value image. These are \cref{ent:thm:criterion},
\cref{ent:prop:operations} with \eqref{ent:eq:gauss-mult},
\cref{ent:thm:product} with \eqref{ent:eq:productcoef},
\cref{ent:lem:linear-division}, \cref{ent:prop:operations} with
\cref{ent:cor:evaluation}, \cref{ent:prop:kernel-direct} and
\cref{ent:cor:single-node}; its weighted value
$\min_j\{v(a_j)-jr\}$ is $m_{-r}$ of \eqref{ent:eq:weighted}. Its sufficiency
construction fits finitely many values by a polynomial with coefficients in
some $V_L$ and then adds
\[
 C_n(Z)=y_n\prod_{j<n}\frac{Z-\rho_j}{\rho_n-\rho_j}\Bigl(\frac
 Z{\rho_n}\Bigr)^{d_n},\qquad d_n\ge\max\{n,2c_n+2\},
\]
where $y_n$ is the residual value and $v(y_n)\ge-c_n\gamma_n$. This is the
correction \eqref{ent:eq:correction} of the first route for one simple node
per shell: with $\tau_n=\rho_n$ the normalized node is $1$, the shell algebra
is $K_\Gamma$, $\Psi_{n,d}$ is the constant $y_n/Q_n(1)$, and
$(Z/\tau_n)^dQ_n(Z/\tau_n)/Q_n(1)$ is the displayed Lagrange factor times
$(Z/\rho_n)^d$. The initialization over $V_L$ and the damping bound
\eqref{ent:eq:damping-target} are the same. Its kernel proof, like
\cref{ent:prop:kernel-direct}, divides by the finite product $P_{\le N}$ and
multiplies by the formal inverse of the tail; it controls the quotients
through the stabilized Gauss value
$m_{-r}(F/P_{\le N})=m_{-r}(F)-\sum_{n\le N}\min\{0,\gamma_n-r\}$ rather than
through a coefficient bound uniform in $N$. Both are therefore printed once,
here by citation; the kernel argument differs from
\cref{ent:prop:kernel-direct} only in that bookkeeping and is not kept as a
separate route. Like the first route, neither uses preparation or algebraic
closedness.
\end{remark}

\begin{lemma}[Indicator idempotents]\label{ent:ps:lem:idempotents}
For every $S\subseteq\N_{>0}$ the indicator sequence
$e_S=(\mathbf 1_S(n))_n$ lies in $\RD$, and
$e_Se_T=e_{S\cap T}$, $e_S+e_{\N_{>0}\setminus S}=1$. Some $F\in\mathcal
E_\Gamma$ interpolates it. The ring $\RD$ is reduced, so $(P_D)$ is a radical
ideal of $\mathcal E_\Gamma$.
\end{lemma}
\begin{proof}
The values $0$ and $1$ lie in every $V_n$, so \cref{ent:cor:single-node}
supplies an interpolant; it need not be idempotent in $\mathcal E_\Gamma$ and
becomes idempotent only modulo $(P_D)$. The identities are coordinatewise.
The ring $\RD$ is a subring of a product of fields, hence reduced.
\end{proof}

The subscript of $e_S$ is a set; an integer subscript, as in $e_j$, denotes a
basis vector of $\Gamma_\infty$, as elsewhere in this report. The same
idempotents were used in the proof of \cref{ent:sm:thm:maximal-all}.

\begin{remark}[Other nodes on the same radii]\label{ent:ps:rem:nodes}
The sixth source takes the monomial nodes $\rho_n=t^{-\gamma_n}$. The ring
$\RD$ of \eqref{ent:ps:eq:quotient} depends only on the radii: by
\cref{ent:thm:quotient,ent:cor:single-node} it is also the quotient of
$\mathcal E_\Gamma$ by the canonical product of any divisor with one simple
node $\rho_n'$ of valuation $-\gamma_n$ on each shell and empty finite part,
through $F\mapsto(F(\rho_n'))_n$. Consequently
\cref{ent:ps:subsec:fibres,ent:ps:subsec:classification,ent:ps:subsec:topology,ent:ps:subsec:dichotomy,ent:ps:subsec:finite-algebra},
which are statements about $\RD$ and about inverse images under that
evaluation map, hold verbatim for such nodes. This observation is the
merge's. \Cref{ent:ps:subsec:contraction,ent:ps:subsec:splitting,ent:ps:subsec:multiplicity}
are stated only in the source's setting.
\end{remark}

\subsection{The ultrafilter fibres}
\label{ent:ps:subsec:fibres}

For an ultrafilter $\Ucal$ put
\begin{equation}\label{ent:ps:eq:nU}
 \mathfrak n_\Ucal=\{a\in\RD:\{n:a_n=0\}\in\Ucal\}.
\end{equation}
For the principal ultrafilter at $n$ this is the coordinate ideal
\[
 \mathfrak e_n=\{a\in\RD:a_n=0\},\qquad \RD/\mathfrak e_n\cong K_\Gamma,
\]
whose inverse image in $\mathcal E_\Gamma$ is the evaluation ideal
$(Z-\rho_n)\mathcal E_\Gamma$ of \cref{ent:cor:evaluation}.

\begin{lemma}[Idempotents select exactly one ultrafilter]
\label{ent:ps:lem:ultrafilter}
For a prime ideal $\mathfrak p$ of $\RD$ the set
$\Ucal_{\mathfrak p}=\{S\subseteq\N_{>0}:e_S\notin\mathfrak p\}$ is an
ultrafilter, and $\mathfrak p\supseteq\mathfrak n_{\Ucal_{\mathfrak p}}$.
Primes with different selected ultrafilters are incomparable.
\end{lemma}
\begin{proof}
In the domain $\RD/\mathfrak p$ every idempotent is $0$ or $1$, and the
identities of \cref{ent:ps:lem:idempotents} give the ultrafilter axioms. If
$a$ vanishes on some $S\in\Ucal_{\mathfrak p}$, then $ae_S=0$ with
$e_S\notin\mathfrak p$, so $a\in\mathfrak p$. If
$\mathfrak p\subseteq\mathfrak p'$ and $e_S\notin\mathfrak p'$, then
$e_S\notin\mathfrak p$; thus $\Ucal_{\mathfrak p'}\subseteq\Ucal_{\mathfrak
p}$, and two ultrafilters one of which contains the other are equal.
\end{proof}

\begin{proposition}[The fibre ring]\label{ent:ps:prop:fibre}
Let $\Ucal$ be nonprincipal. Replacing the finitely many coordinates of
$a\in\RD$ that lie outside $V_n$ by zero and passing to the class modulo
$\Ucal$ defines a surjective ring homomorphism
\[
 \red_\Ucal:\RD\longrightarrow\mathcal V_\Ucal=\prod_\Ucal V_n
\]
with kernel $\mathfrak n_\Ucal$, so $\RD/\mathfrak n_\Ucal\cong\mathcal
V_\Ucal$. Here an ultraproduct is a quotient by agreement on a member of
$\Ucal$, not a space of convergent sequences. The ring $\mathcal V_\Ucal$ is
a valuation domain with fraction field $\mathcal K_\Ucal=\prod_\Ucal K_\Gamma$
and value group
\[
 \Gamma_\Ucal=\prod_\Ucal(\Gamma/H_n),
\]
the valuation of a nonzero class being
\begin{equation}\label{ent:ps:eq:wU}
 w_\Ucal\bigl([a_n]_\Ucal\bigr)=[v(a_n)+H_n]_\Ucal ,
\end{equation}
where zero coordinates on a set outside $\Ucal$ may be changed arbitrarily.
\end{proposition}
\begin{proof}
Changing finitely many coordinates does not change a class at a nonprincipal
ultrafilter, which gives well-definedness and compatibility with the ring
operations. Every sequence in $\prod_nV_n$ already lies in $\RD$, so
$\red_\Ucal$ is onto, and its kernel is $\mathfrak n_\Ucal$.

The ultraproduct $\mathcal K_\Ucal$ is a field: a nonzero class has a
representative that is nonzero on a member of $\Ucal$, and reciprocals there
invert it. Every class in $\mathcal K_\Ucal^\times$ lies in $\mathcal V_\Ucal$
or has its inverse there, because the alternative holds in each coordinate
and $\Ucal$ contains one of the two index sets. So $\mathcal V_\Ucal$ is a
valuation domain with fraction field $\mathcal K_\Ucal$. Formula
\eqref{ent:ps:eq:wU} respects agreement modulo $\Ucal$, multiplication and
the ultrametric inequality coordinatewise. Its range is the whole ordered
group ultraproduct: choose an exponent in each coordinate coset and take the
corresponding monomials. The classes of nonnegative value are exactly the
nonzero classes in $\mathcal V_\Ucal$. No model-theoretic transfer theorem is
used.
\end{proof}

Composing $\red_\Ucal$ with the coordinatewise coarsened residue maps of
\cref{ent:lem:coarse-residue} gives $\Phi_\Ucal$ of \eqref{ent:eq:Phi} on
$\RD$: the kernel of $[a_n]_\Ucal\mapsto[\res_n(a_n)]_\Ucal$ on $\mathcal
V_\Ucal$ consists of the classes with $a_n\in\mathfrak q_{\gamma_n}$ on a
member of $\Ucal$, which is the maximal ideal $\{w_\Ucal>0\}\cup\{0\}$, and
the map is onto $L_\Ucal=\prod_\Ucal\kappa_n$. This remark is the merge's
bridge between the two notations.

\begin{remark}[Relation to the product-ring literature]
\label{ent:ps:rem:ffw}
\Cref{ent:ps:lem:ultrafilter,ent:ps:prop:fibre} are the familiar idempotent
and ultraproduct mechanisms of product-ring theory; compare \cite{FFW}. What
the analytic theory supplies is the particular restricted sequence ring
$\RD$, with coarsenings that change with $n$. This is why a statement about
product rings becomes a statement about the values of entire functions at
the nodes.
\end{remark}

\begin{lemma}[Primes of a valuation domain]
\label{ent:ps:lem:valuation-primes}
Let $\mathcal V$ be a valuation domain with surjective valuation
$w:\operatorname{Frac}(\mathcal V)^\times\to\mathcal G$. For a convex subgroup
$\mathcal C\le\mathcal G$ put
\[
 \mathfrak p_{\mathcal C}(\mathcal V)=\{0\}\cup
 \{x\in\mathcal V\setminus\{0\}:w(x)>\mathcal C\},
\]
where $g>\mathcal C$ means that $g$ exceeds every element of $\mathcal C$.
Every prime of $\mathcal V$ is $\mathfrak p_{\mathcal C}(\mathcal V)$ for
exactly one $\mathcal C$, and $\mathcal C\subseteq\mathcal C_1$ if and only
if $\mathfrak p_{\mathcal C_1}(\mathcal V)\subseteq\mathfrak p_{\mathcal
C}(\mathcal V)$. The endpoints are $\mathfrak p_{\mathcal G}(\mathcal V)=0$ and
$\mathfrak p_0(\mathcal V)$, the maximal ideal.
\end{lemma}
\begin{proof}
For $g\ge0$, lying outside a convex subgroup is the same as exceeding all of
it. Hence $\mathfrak p_{\mathcal C}(\mathcal V)$ is zero together with the
elements of positive value for the coarsening $w+\mathcal C$; it is an ideal
by the valuation inequalities, and prime because two nonnegative values in
$\mathcal C$ have their sum in $\mathcal C$.

Conversely let $\mathfrak p$ be prime. The values of the nonzero elements of
$\mathcal V\setminus\mathfrak p$ form an additive submonoid $S$ of
$\mathcal G_{\ge0}$ that is closed downwards: if $0\le g\le w(s)$ with
$s\notin\mathfrak p$, choose $x$ of value $g$; then $s=xu$ with
$u\in\mathcal V$, so $x\notin\mathfrak p$. The set $S\cup(-S)$ is a convex
subgroup whose nonnegative part is $S$, and $\mathfrak p=\mathfrak
p_{S\cup(-S)}(\mathcal V)$. Surjectivity of $w$ separates distinct convex
subgroups and gives the inclusion rule. This is the standard correspondence
\cite{StacksValuation}.
\end{proof}

\subsection{The complete classification}
\label{ent:ps:subsec:classification}

\begin{theorem}[All primes above the reduced product]
\label{ent:ps:thm:spectrum}
The prime ideals of $\RD\cong\mathcal E_\Gamma/(P_D)$ are exactly the
following.
\begin{enumerate}[label=\textnormal{(\alph*)},leftmargin=2em]
\item For each $n\ge1$, the maximal ideal $\mathfrak e_n$, with inverse image
$(Z-\rho_n)\mathcal E_\Gamma$.
\item For each nonprincipal ultrafilter $\Ucal$ and each convex subgroup
$\mathcal C\le\Gamma_\Ucal$, the prime
\begin{equation}\label{ent:ps:eq:pUC}
 \mathfrak p_{\Ucal,\mathcal C}
 =\red_\Ucal^{-1}\bigl(\mathfrak p_{\mathcal C}(\mathcal V_\Ucal)\bigr).
\end{equation}
\end{enumerate}
The parameters are unique. Primes over different ultrafilters are
incomparable; within one fibre, inclusion of primes reverses inclusion of
convex subgroups. The fibre over $\Ucal$ has the unique minimal prime
$\mathfrak n_\Ucal=\mathfrak p_{\Ucal,\Gamma_\Ucal}$ and the unique maximal
prime
\[
 \overline{\mathfrak M}_\Ucal=\mathfrak p_{\Ucal,0},
\]
whose inverse image in $\mathcal E_\Gamma$ is the ideal $\mathfrak M_\Ucal$
of \cref{ent:thm:maximal}. For a principal ultrafilter at $n$ both endpoints
are $\mathfrak e_n$.
\end{theorem}
\begin{proof}
Let $\mathfrak p$ be prime and $\Ucal=\Ucal_{\mathfrak p}$. If $\Ucal$ is
principal at $n$, then $1-e_{\{n\}}\in\mathfrak p$, and evaluation
identifies $\RD/(1-e_{\{n\}})\RD$ with the field $K_\Gamma$; hence
$\mathfrak p=\mathfrak e_n$. If $\Ucal$ is nonprincipal,
\cref{ent:ps:lem:ultrafilter} makes $\mathfrak p$ the inverse image of a prime
of $\RD/\mathfrak n_\Ucal\cong\mathcal V_\Ucal$
(\cref{ent:ps:prop:fibre}), and \cref{ent:ps:lem:valuation-primes} applies.
The converse and uniqueness follow from the same correspondence and from the
idempotents, which recover $\Ucal$. Convex subgroups of an ordered group are
linearly ordered by inclusion, which gives the uniqueness of the endpoints.
The kernel of $\Phi_\Ucal$ on $\RD$ is $\mathfrak p_{\Ucal,0}$ by the
remark after \cref{ent:ps:prop:fibre}, and its inverse image in
$\mathcal E_\Gamma$ is $\mathfrak M_\Ucal$ by \eqref{ent:eq:Mexplicit}.
\end{proof}

The maximal ideals among these are the $\mathfrak e_n$ and the
$\overline{\mathfrak M}_\Ucal$. For simple nodes this is
\cref{ent:sm:thm:maximal-all}, which the sixth source recovers and does not
claim; the new content is the nonmaximal part of each fibre.

\begin{corollary}[Membership test for entire functions]
\label{ent:ps:cor:membership}
The inverse image of $\mathfrak p_{\Ucal,\mathcal C}$ in $\mathcal E_\Gamma$
consists of the $F$ with $\{n:F(\rho_n)=0\}\in\Ucal$ and, in the
complementary case, of the $F$ whose ultraproduct value satisfies
\[
 \bigl[v(F(\rho_n))+H_n\bigr]_\Ucal>\mathcal C .
\]
Thus an entire function can lie in a non-evaluation prime even when none of
its node values vanishes. The condition measures an asymptotic scale of the
values, not an additional zero.
\end{corollary}

\begin{theorem}[Local rings and residue towers]\label{ent:ps:thm:local}
For nonprincipal $\Ucal$ and convex $\mathcal C\le\Gamma_\Ucal$ there are
canonical isomorphisms
\[
 (\RD)_{\mathfrak p_{\Ucal,\mathcal C}}
 \cong(\mathcal V_\Ucal)_{\mathfrak p_{\mathcal C}(\mathcal V_\Ucal)}
 =\{0\}\cup\{x\in\mathcal K_\Ucal^\times:w_\Ucal(x)+\mathcal C\ge0\}.
\]
This is a valuation domain with value group $\Gamma_\Ucal/\mathcal C$. Its
residue field carries a residual valuation with value group $\mathcal C$,
whose residue field is
\begin{equation}\label{ent:ps:eq:residue}
 L_\Ucal=\prod_\Ucal\kk\bigl((t^{H_n})\bigr).
\end{equation}
At a principal prime $\mathfrak e_n$ the local ring and the residue field are
$K_\Gamma$.
\end{theorem}
\begin{proof}
The quotient map to $\mathcal V_\Ucal$ localizes to a surjection, and its
kernel dies after localization: if $a\in\mathfrak n_\Ucal$ vanishes on
$S\in\Ucal$, then $ae_S=0$ and $e_S\notin\mathfrak p_{\Ucal,\mathcal C}$.

For a valuation domain $\mathcal V$ with value group $\mathcal G$, localizing
at $\mathfrak p_{\mathcal C}(\mathcal V)$ inverts exactly the elements whose
nonnegative values lie in $\mathcal C$, so the localization is the ring of
the coarsening $w+\mathcal C$. In the one case needing an argument, if
$w(x)\in\mathcal C$ is negative, choose $h\in\mathcal C_{\ge0}$ with
$h\ge-w(x)$ and $s\in\mathcal V$ of value $h$; then
$s\notin\mathfrak p_{\mathcal C}(\mathcal V)$ and $sx\in\mathcal V$. This
gives the description and the value group. A nonzero residue of a unit for
$w+\mathcal C$ has a well-defined residual value in $\mathcal C$, every
element of $\mathcal C$ occurs by surjectivity of $w$, and taking residues
once more gives the residue field of $\mathcal V$. For $\mathcal V_\Ucal$ the
coordinatewise residue maps and \cref{ent:lem:coarse-residue} identify that
field with \eqref{ent:ps:eq:residue}. The principal case follows by inverting
$e_{\{n\}}$.
\end{proof}

\begin{remark}[What is and is not new about the residue fields]
The fields \eqref{ent:ps:eq:residue} are the residue fields $L_\Ucal$ of
\cref{ent:thm:maximal,ent:sm:thm:maximal-all}, algebraically closed for
$\kk=\C$ and real closed for $\kk=\R$ by \cref{ent:sm:cor:boundary-fields}.
At intermediate primes the residual valuation is additional structure. No
ultraproduct field is asserted to be canonically a Hahn field on its
ultraproduct value group; no such identification is needed.
\end{remark}

\subsection{Topology: a Stone space of valuation fibres}
\label{ent:ps:subsec:topology}

Give $\Spec\RD$ the Zariski topology, with basic open sets
$D(a)=\{\mathfrak p:a\notin\mathfrak p\}$, and give $\beta\N_{>0}$ the basic
clopen sets $\widehat S=\{\Ucal:S\in\Ucal\}$, as in \cref{ent:sm:thm:beta}.

\begin{theorem}[The spectrum over the ultrafilter space]
\label{ent:ps:thm:topology}
The map $\pi:\Spec\RD\to\beta\N_{>0}$, $\mathfrak p\mapsto\Ucal_{\mathfrak
p}$, is continuous and onto, with $\pi^{-1}(\widehat S)=D(e_S)$. Its fibre at
a nonprincipal $\Ucal$ is the closed irreducible subspace
$V(\mathfrak n_\Ucal)\cong\Spec\mathcal V_\Ucal$, with generic point
$\mathfrak n_\Ucal$ and unique closed point $\overline{\mathfrak M}_\Ucal$.
Principal fibres are open points. Both endpoint maps are homeomorphisms
\[
 \operatorname{Min}\RD\ \cong\ \beta\N_{>0},\qquad
 \operatorname{Max}\RD\ \cong\ \beta\N_{>0}
\]
for the subspace topologies, and the evaluation primes $\mathfrak e_n$ form a
discrete dense subset of the \emph{whole} prime spectrum.
\end{theorem}
\begin{proof}
The formula for $\pi^{-1}(\widehat S)$ is the definition of
$\Ucal_{\mathfrak p}$, and $\pi$ is onto by \cref{ent:ps:thm:spectrum}. A
prime contains $\mathfrak n_\Ucal$ exactly when it selects $\Ucal$: if
$S\notin\Ucal$ then $e_S\in\mathfrak n_\Ucal$, and if $S\in\Ucal$ the
complementary idempotent lies there. So the fibre is $V(\mathfrak n_\Ucal)$,
homeomorphic to $\Spec\mathcal V_\Ucal$; the spectrum of a domain is
irreducible with generic point $0$, and a valuation domain has one maximal
ideal. Principal fibres are $D(e_{\{n\}})=\{\mathfrak e_n\}$.

For the minimal primes put $S_a=\{n:a_n\ne0\}$; then $a\notin\mathfrak
n_\Ucal$ exactly when $S_a\in\Ucal$, so $D(a)$ meets the minimal primes in
$\widehat{S_a}$, and the idempotents give every $\widehat S$. For the maximal
primes put $T_a=\{n:a_n\in V_n^\times\}\subseteq S_a$. At a nonprincipal
$\Ucal$, $a\notin\overline{\mathfrak M}_\Ucal$ exactly when $T_a\in\Ucal$; at
the principal point $n$, exactly when $n\in S_a$. So $D(a)$ meets the maximal
primes in
$\widehat{T_a}\cup\{\text{principal point at }n:n\in S_a\setminus T_a\}$,
which is open, and the idempotents again give continuity of the inverse. This
recovers \cref{ent:sm:thm:beta} for simple nodes. Finally a nonempty $D(a)$
has $a\ne0$, hence some $a_n\ne0$ and $\mathfrak e_n\in D(a)$; together with
the open points this gives density and discreteness.
\end{proof}

\begin{remark}[Not a disjoint union]
The fibres partition the primes but are not open in general, so the spectrum
is not a topological disjoint union of valuation spectra. A nonprincipal
fibre is closed, and its position relative to the principal points is
governed by all indicator idempotents at once.
\end{remark}

\begin{corollary}[Semiprimitive, and Jacobson only in the degenerate case]
\label{ent:ps:cor:semiprimitive}
The Jacobson radical of $\RD$ is zero. If some nonprincipal fibre has more
than one point, then $\RD$ is not a Jacobson ring.
\end{corollary}
\begin{proof}
The intersection of the $\mathfrak e_n$ is zero, hence so is the intersection
of all maximal ideals. A nonmaximal prime of a nonprincipal fibre lies in
exactly one maximal ideal, by \cref{ent:ps:thm:spectrum}, so it is not the
intersection of the maximal ideals containing it; that is the failure of the
Jacobson property.
\end{proof}

\subsection{The order-unit dichotomy and continuum prime chains}
\label{ent:ps:subsec:dichotomy}

\emph{Order unit} is the notion of \cref{ent:def:order-unit}; countable
cofinality does not imply that one exists (\cref{ent:rem:one-name}). In this
subsection $\theta$ denotes a real number in $(0,1)$, and $\lfloor\cdot\rfloor$
the integer part.

\begin{lemma}[Uniform Archimedean gaps in an ultraproduct]
\label{ent:ps:lem:gaps}
Let $\mathcal G_n$ be nonzero ordered abelian groups, choose
$\bar\eta_n>0$ in $\mathcal G_n$, and let $\Ucal$ be nonprincipal. For
$0<\theta<1$ put
\[
 \zeta_\theta=\bigl[\lfloor n^\theta\rfloor\,\bar\eta_n\bigr]_\Ucal
 \in\prod_\Ucal\mathcal G_n,
 \qquad
 \mathcal C_\theta=\conv(\Z\zeta_\theta),
\]
the convex subgroup generated by $\zeta_\theta$. If $\theta<\theta'$ then
$m\zeta_\theta<\zeta_{\theta'}$ for every positive integer $m$, and hence
$\mathcal C_\theta\subsetneq\mathcal C_{\theta'}$.
\end{lemma}
\begin{proof}
For fixed $\theta<\theta'$ and $m$, the integer inequality
$m\lfloor n^\theta\rfloor<\lfloor n^{\theta'}\rfloor$ holds for all large
$n$; multiplying by $\bar\eta_n>0$ preserves it, and a nonprincipal
ultrafilter contains every cofinite set. So $\zeta_{\theta'}$ exceeds every
ordinary multiple of $\zeta_\theta$. The convex subgroup generated by a
positive $\zeta$ consists of the elements bounded in absolute value by an
ordinary finite multiple of $\zeta$, which gives the strict inclusion.
\end{proof}

\begin{theorem}[Reduced-divisor dichotomy]\label{ent:ps:thm:dichotomy}
Under \cref{ent:ps:conv:standing} the following are equivalent.
\begin{enumerate}[label=\textnormal{(\alph*)},leftmargin=2em]
\item $\Gamma$ has an order unit.
\item $V_n=K_\Gamma$ for all sufficiently large $n$, and
$\RD=K_\Gamma^{\N_{>0}}$.
\item $\RD$ is von Neumann regular: for every $a\in\RD$ there is $b\in\RD$
with $a^2b=a$.
\item Every prime of $\RD$ is maximal; equivalently $\dim\RD=0$.
\item $\RD$ is a Jacobson ring.
\end{enumerate}
If they fail, then for every nonprincipal $\Ucal$ the fibre over $\Ucal$
contains a strict chain of primes indexed by $(0,1)$ with its order reversed.
In particular $\mathcal E_\Gamma/(P_D)$ has infinite Krull dimension,
although every zero of $P_D$ is simple.
\end{theorem}
\begin{proof}
If $h$ is an order unit, cofinality gives $\gamma_n\ge h$ eventually, so
$H_n=\Gamma$ and $V_n=K_\Gamma$. Conversely $V_n=K_\Gamma$ forces
$H_n=\Gamma$, by testing monomials of arbitrarily negative value, so
$\gamma_n$ is an order unit. This is (a)$\Leftrightarrow$(b), which is
\cref{ent:cor:order-unit-image} together with \cref{ent:sm:cor:no-order-unit}.
In a product of fields put $b_n=a_n^{-1}$ where $a_n\ne0$ and $b_n=0$
elsewhere; this gives (c). Under (b) every $\Gamma/H_n$ is zero, so every
$\Gamma_\Ucal$ is zero and \cref{ent:ps:thm:spectrum} gives (d); and (d)
implies (e).

Suppose there is no order unit. Every $H_n$ is proper; choose $\eta_n>H_n$
and put $\bar\eta_n=\eta_n+H_n>0$ in $\Gamma/H_n$.
\Cref{ent:ps:lem:gaps} gives strictly increasing convex subgroups
$\mathcal C_\theta\le\Gamma_\Ucal$, and the primes $\mathfrak
p_{\Ucal,\mathcal C_\theta}$ strictly decrease by \cref{ent:ps:thm:spectrum},
so (d) fails. Each nonprincipal fibre then has a nonmaximal prime, so (e)
fails by \cref{ent:ps:cor:semiprimitive}. For (c) take $a_n=t^{\eta_n}$,
which lies in $\RD$ and has no zero coordinate: $a^2b=a$ would force
$b_n=t^{-\eta_n}$, which lies in no $V_n$, so $b\notin\RD$. A real interval
contains finite strict chains of every length, which gives the assertion
about Krull dimension.
\end{proof}

The assertion concerns the quotient $\mathcal E_\Gamma/(P_D)$. It does
\emph{not} say that $\mathcal E_\Gamma$ itself is zero-dimensional in the
order-unit case.

\begin{proposition}[Entire witnesses for every strict inclusion]
\label{ent:ps:prop:witnesses}
Without an order unit, choose $\eta_n>H_n$ as above. For each $0<\theta<1$
there is $F_\theta\in\mathcal E_\Gamma$ with
$F_\theta(\rho_n)=t^{\lfloor n^\theta\rfloor\eta_n}$ for all $n$. For the
subgroups $\mathcal C_{\theta'}$ of the proof of
\cref{ent:ps:thm:dichotomy}, the class of $F_\theta$ lies in
$\mathfrak p_{\Ucal,\mathcal C_{\theta'}}$ when $\theta>\theta'$ and not when
$\theta\le\theta'$. None of these node values is zero.
\end{proposition}
\begin{proof}
All prescribed values lie in $V_n$, so \cref{ent:cor:single-node} supplies
$F_\theta$, whose ultraproduct value is $\zeta_\theta$. If $\theta>\theta'$
it exceeds every element of $\mathcal C_{\theta'}$ by
\cref{ent:ps:lem:gaps}; if $\theta\le\theta'$ it lies in
$\mathcal C_{\theta'}$. Apply \cref{ent:ps:cor:membership}.
\end{proof}

\begin{remark}[Cardinality and order type]
The theorem asserts a chain with continuum many distinct members. It assigns
no cardinal-valued Krull dimension, does not assert that the chain is well
ordered, and does not bound the number of convex subgroups of an
ultraproduct; those are different set-theoretic questions.
\end{remark}

\begin{remark}[Rank is not the dividing invariant]\label{ent:ps:rem:rank}
The sixth source's example is \cref{ent:subsec:finite-rank} together with
\cref{ent:subsec:gamma-infinity}, printed once there. Every nonzero
Archimedean divisible group, and $\Q^d$ in lexicographic order, has an order
unit, so for them $\RD=K_\Gamma^{\N_{>0}}$ and every prime above $P_D$ is
maximal. The group $\Gamma_\infty$ has countable cofinality and no order
unit, so every nonprincipal fibre carries a continuum chain. What divides the
two regimes is one cofinal cyclic scale, not rank greater than one; this is
the same dividing line as case (ii) versus case (iii) of \cref{ent:thm:main}.
\end{remark}

\subsection{Finite algebra in an infinite-dimensional spectrum}
\label{ent:ps:subsec:finite-algebra}

The quotient can have a very deep spectrum and still admit simple finite
linear algebra: every coordinate ring $V_n$ is a valuation ring, and every
subset of the index set has an indicator idempotent.

\begin{theorem}[Finitely generated ideals and annihilators]
\label{ent:ps:thm:bezout}
Every finitely generated ideal of $\RD$ is principal. For $a\in\RD$,
\[
 \operatorname{Ann}(a)=e_{\{n:a_n=0\}}\RD,\qquad
 a\RD\cong e_{\{n:a_n\ne0\}}\RD .
\]
Hence every finitely generated ideal is projective and $\RD$ is
semihereditary. It is a B\'ezout ring with zero divisors, not a B\'ezout
domain.
\end{theorem}
\begin{proof}
Let $a^{(1)},\ldots,a^{(s)}\in\RD$. Outside a finite exceptional set all their
coordinates lie in $V_n$. At such a coordinate with some nonzero value choose
the first index $j$ whose coarsened valuation $v+H_n$ is least among the
nonzero values; at an exceptional coordinate choose any index with a nonzero
value, if there is one. This partitions the coordinates with some nonzero
value into sets $S_1,\ldots,S_s$; put $g=\sum_je_{S_j}a^{(j)}$, an element of
the ideal. Define $b^{(j)}_n=a^{(j)}_n/g_n$ where $g_n\ne0$ and $0$ elsewhere.
The minimality of the chosen valuation puts these ratios in $V_n$ at every
nonexceptional coordinate, so $b^{(j)}\in\RD$ and $a^{(j)}=gb^{(j)}$; the
ideal is $g\RD$. The annihilator of $a$ consists of the sequences supported
on its zero set, and multiplication by $a$ maps $e_{\{a_n\ne0\}}\RD$
isomorphically onto $a\RD$, being onto by definition and injective because
each surviving coordinate is a nonzero element of a field. A direct summand
of the free module $\RD$ is projective.
\end{proof}

\begin{remark}[Translation to $\mathcal E_\Gamma$]\label{ent:ps:rem:fg-ideals}
Every finitely generated ideal $I$ of $\mathcal E_\Gamma$ that contains
$P_D$ is $(P_D,G)$ for a single $G\in\mathcal E_\Gamma$: the image of $I$ in
$\RD$ is finitely generated, hence $g\RD$, and any lift $G$ of $g$ works.
Membership of an arbitrary element in $(P_D,G)$ is decided by
\cref{ent:sm:cor:membership}. This translation is the merge's; it bears on
\cref{ent:q:ideals} only for ideals containing this product.
\end{remark}

\begin{theorem}[Smith forms for finite matrices]\label{ent:ps:thm:smith}
For every $r\times s$ matrix $M$ over $\RD$ there are
$U\in\operatorname{GL}_r(\RD)$ and $V\in\operatorname{GL}_s(\RD)$ with
\[
 UMV=\operatorname{diag}(d_1,\ldots,d_\ell,0,\ldots),\qquad
 \ell=\min(r,s),\qquad d_i\mid d_{i+1},
\]
the diagonal taken in the rectangular sense and zero entries allowed. Thus
$\RD$ is an elementary divisor ring, and every finitely presented
$\RD$-module is a direct sum of finitely many cyclic modules and a finite
free module.
\end{theorem}
\begin{proof}
Choose a finite exceptional set outside which every entry of $M$ lies in
$V_n$. Over a valuation domain a finite matrix has the usual pivot
reduction: a nonzero entry of least valuation divides every entry; swaps move
it to the corner, elementary operations clear its row and column, the
remaining entries stay divisible by it, and one inducts on the smaller block.
A zero matrix takes identity transformations. Perform this reduction over
$V_n$ at each nonexceptional coordinate and over $K_\Gamma$ at the
exceptional ones. The row and column operations and their inverses have
entries in $V_n$ eventually, so their coordinate sequences define invertible
matrices $U,V$ over $\RD$. The diagonal entries lie in $\RD$, and the
coordinatewise quotients $d_{i+1}/d_i$ lie in $V_n$ eventually, so
$d_i\mid d_{i+1}$ in $\RD$; a zero pivot is followed only by zero pivots at
that coordinate, and the quotient there is taken to be $0$. Applying this to
a presentation matrix gives the module decomposition. No uniqueness of the
diagonal representatives is claimed.
\end{proof}

In this subsection $U$ and $V$ are invertible matrices, not an ultrafilter or
a coarsened valuation ring.

\begin{proposition}[Total quotient ring and integral closedness]
\label{ent:ps:prop:total-quotient}
An element of $\RD$ is a non-zero-divisor exactly when none of its
coordinates is zero, and the total quotient ring is
$Q(\RD)=K_\Gamma^{\N_{>0}}$. The ring $\RD$ is integrally closed in it, and
it is not Noetherian.
\end{proposition}
\begin{proof}
A zero coordinate gives a nonzero annihilating indicator, and a sequence with
no zero coordinate acts injectively; so the total quotient ring embeds in
$K_\Gamma^{\N_{>0}}$. Given $b\in K_\Gamma^{\N_{>0}}$ put $s_n=1$ if
$b_n\in V_n$ and $s_n=b_n^{-1}$ otherwise; then $s$ is a non-zero-divisor of
$\RD$, $sb\in\RD$ and $b=(sb)/s$. If $b$ satisfies a monic equation over
$\RD$, all its coefficients lie in $V_n$ for large $n$, and a valuation ring
is integrally closed: an element of negative valuation cannot satisfy a monic
equation, since its leading power would have strictly smaller valuation than
every other term. So $b_n\in V_n$ eventually and $b\in\RD$. The ideals
generated by $e_{\{1,\ldots,N\}}$ increase strictly with $N$.
\end{proof}

\begin{remark}[Two B\'ezout statements that do not conflict]
Without an order unit the full ring $\mathcal E_\Gamma$ is not B\'ezout
(\cref{ent:thm:main}(iii)); \cref{ent:ps:thm:bezout} says that its quotient by
this particular simple one-node-per-shell product \emph{is} a B\'ezout ring.
There is no contradiction: the quotient acquires indicator idempotents,
which permit coordinatewise choices of generator.
\end{remark}

\subsection{Cofinal scalar extension and contraction of primes}
\label{ent:ps:subsec:contraction}

Let $\Gamma\subseteq\Delta$ be an ordered-group inclusion of divisible
ordered abelian groups with $\Gamma$ cofinal in $\Delta$, and keep the radii
$\gamma_n\in\Gamma$ and the nodes $\rho_n=t^{-\gamma_n}$. Superscripts record
the group:
\[
 H_n^\Gamma=H_{\gamma_n}\subseteq\Gamma,\qquad
 H_n^\Delta=\{\delta\in\Delta:\abs\delta\le m\gamma_n
   \text{ for some }m\in\N\},
\]
and similarly $V_n^\Gamma$, $V_n^\Delta$, $\RD^\Gamma$, $\RD^\Delta$ and
$\Gamma_\Ucal=\prod_\Ucal(\Gamma/H_n^\Gamma)$,
$\Delta_\Ucal=\prod_\Ucal(\Delta/H_n^\Delta)$. The notation $\Delta$ is that
of \cref{ent:sec:extension}; the convex hull there is all of $\Delta$,
because $\Gamma$ is cofinal.

\begin{lemma}[The comparison maps]\label{ent:ps:lem:comparison}
Coefficientwise inclusion gives
$K_\Gamma\subseteq K_\Delta$, $\mathcal E_\Gamma\subseteq\mathcal E_\Delta$
and $\RD^\Gamma\subseteq\RD^\Delta$, the last induced on the quotients by the
same series $P_D$. Moreover
\[
 H_n^\Delta\cap\Gamma=H_n^\Gamma,\qquad
 V_n^\Delta\cap K_\Gamma=V_n^\Gamma ,
\]
and the zero set of $P_D$ in either field is $\{\rho_n:n\ge1\}$, every zero
being simple.
\end{lemma}
\begin{proof}
An ordered-group inclusion preserves well-ordered supports and the Hahn
operations. Coefficient slopes cofinal in $\Gamma$ are cofinal in $\Delta$,
so \cref{ent:thm:criterion} gives the second inclusion; for $\kk=\C$ this is
\cref{ent:cor:entire-extension}. The equality of convex subgroups is the
common inequality $\abs g\le m\gamma_n$ for $g\in\Gamma$, and the equality of
valuation rings follows by testing the leading exponent. The evaluation maps
commute with the inclusions, and both kernels are $(P_D)$ by
\cref{ent:prop:kernel-direct}, so the induced map of quotients is injective
and is the displayed inclusion of sequence rings. The radii remain cofinal in
$\Delta$. At $z\ne0$ choose $N$ with $\gamma_{N+1}+v(z)>0$: the tail
$\prod_{n>N}(1-t^{\gamma_n}z)$ is $1$ plus an element of positive valuation,
hence nonzero, so only the listed nodes are zeros; removing one factor
leaves a cofactor that does not vanish at its node, so each zero is simple.
\end{proof}

\begin{remark}[What scalar extension means here]
Rings are compared by coefficientwise inclusion and the induced map of
quotients. It is \emph{not} asserted that
$\mathcal E_\Gamma\otimes_{K_\Gamma}K_\Delta=\mathcal E_\Delta$, or that the
corresponding tensor product of quotients is the larger quotient; these
surjectivity assertions are neither needed nor proved. Cofinality cannot
simply be dropped: without it a nonpolynomial series entire over $K_\Gamma$
is not entire over $K_\Delta$ (\cref{ent:thm:extension-criterion}).
\end{remark}

For nonprincipal $\Ucal$ the coordinate maps
$\Gamma/H_n^\Gamma\to\Delta/H_n^\Delta$ reflect order and equality, so they
induce an order embedding $\Gamma_\Ucal\hookrightarrow\Delta_\Ucal$,
by which $\Gamma_\Ucal$ is identified with a subgroup of $\Delta_\Ucal$.

\begin{lemma}[All convex extensions of a convex subgroup]
\label{ent:ps:lem:convex-interval}
Let $\mathcal G\subseteq\mathcal G'$ be an inclusion of ordered abelian
groups with $\mathcal G$ divisible, and let $\mathcal C$ be a convex subgroup
of $\mathcal G$. Put
\begin{align}
 \mathcal C_{\min}&=\{x\in\mathcal G':\abs x\le c\text{ for some }
   c\in\mathcal C_{\ge0}\},\label{ent:ps:eq:Cmin}\\
 \mathcal C_{\max}&=\{x\in\mathcal G':\abs x<g\text{ for every }
   g\in\mathcal G_{>0}\setminus\mathcal C\},\label{ent:ps:eq:Cmax}
\end{align}
with $\mathcal C_{\max}=\mathcal G'$ when the condition is empty. Both are
convex subgroups of $\mathcal G'$ with trace $\mathcal C$ on $\mathcal G$, and
the convex subgroups $\mathcal C'\le\mathcal G'$ with
$\mathcal C'\cap\mathcal G=\mathcal C$ are exactly those with
\begin{equation}\label{ent:ps:eq:interval}
 \mathcal C_{\min}\subseteq\mathcal C'\subseteq\mathcal C_{\max}.
\end{equation}
\end{lemma}
\begin{proof}
$\mathcal C_{\min}$ is the convex subgroup of $\mathcal G'$ generated by
$\mathcal C$, and its trace is $\mathcal C$ by convexity of $\mathcal C$ in
$\mathcal G$. The set $\mathcal C_{\max}$ is symmetric and convex. If
$\mathcal C=\mathcal G$ there is nothing more to prove; otherwise, for $g>0$
outside $\mathcal C$ also $g/2$ is outside $\mathcal C$, by divisibility, and
$\abs{x+y}<g/2+g/2=g$ for $x,y\in\mathcal C_{\max}$, so $\mathcal C_{\max}$ is
a subgroup. A positive element of $\mathcal G$ outside $\mathcal C$ exceeds
all of $\mathcal C$, so $\mathcal C\subseteq\mathcal C_{\max}$, while no
$g\in\mathcal G\setminus\mathcal C$ lies in $\mathcal C_{\max}$; its trace is
$\mathcal C$, and it contains $\mathcal C_{\min}$. If $\mathcal C'$ is convex
with trace $\mathcal C$, it contains $\mathcal C_{\min}$; and for
$x\in\mathcal C'$ and $g\in\mathcal G_{>0}\setminus\mathcal C$, the
inequality $g\le\abs x$ would put $g$ in $\mathcal C'\cap\mathcal G=\mathcal
C$, so $\mathcal C'\subseteq\mathcal C_{\max}$. Conversely every convex
subgroup in the interval has trace $\mathcal C$.
\end{proof}

\begin{theorem}[Exact contraction and all extensions of a prime]
\label{ent:ps:thm:contraction}
Contraction of primes along $\RD^\Gamma\subseteq\RD^\Delta$ preserves the
selected ultrafilter. At a principal point it sends $\mathfrak e_n^\Delta$ to
$\mathfrak e_n^\Gamma$. At a nonprincipal $\Ucal$,
\begin{equation}\label{ent:ps:eq:contraction}
 \mathfrak p^\Delta_{\Ucal,\mathcal C'}\cap\RD^\Gamma
 =\mathfrak p^\Gamma_{\Ucal,\mathcal C'\cap\Gamma_\Ucal}.
\end{equation}
Every prime of $\RD^\Gamma$ extends to $\RD^\Delta$. The primes contracting
to $\mathfrak p^\Gamma_{\Ucal,\mathcal C}$ are exactly the
$\mathfrak p^\Delta_{\Ucal,\mathcal C'}$ with $\mathcal C'$ in the interval
\eqref{ent:ps:eq:interval} for $\mathcal G=\Gamma_\Ucal$ and
$\mathcal G'=\Delta_\Ucal$, so each fibre of the contraction map is
anti-isomorphic, as an ordered set, to that interval.
\end{theorem}
\begin{proof}
Each indicator $e_S$ maps to the same indicator, so contraction preserves
the zero-or-one test and hence the ultrafilter; the principal case is
coordinate evaluation. For nonprincipal $\Ucal$ the maps of ultraproduct
valuation rings are injective and compatible with
$\Gamma_\Ucal\hookrightarrow\Delta_\Ucal$. A nonzero element of
$\mathcal V_\Ucal^\Gamma$ of value $g\ge0$ maps into the prime indexed by
$\mathcal C'$ exactly when $g\notin\mathcal C'$, that is, when
$g\notin\mathcal C'\cap\Gamma_\Ucal$; zero is preserved. This is
\eqref{ent:ps:eq:contraction}. The group $\Gamma_\Ucal$ is divisible, since
division by a positive integer is performed coordinatewise, so
\cref{ent:ps:lem:convex-interval} supplies at least one convex extension and
describes all of them, and \cref{ent:ps:thm:spectrum} translates this into
primes with the order reversed.
\end{proof}

\begin{corollary}[Uniqueness of extension]\label{ent:ps:cor:unique-extension}
The prime $\mathfrak p^\Gamma_{\Ucal,\mathcal C}$ has a unique extension
exactly when $\mathcal C_{\min}=\mathcal C_{\max}$ in $\Delta_\Ucal$. If
$\Gamma_\Ucal=\Delta_\Ucal$, every prime of that fibre extends uniquely. If
$\Gamma$ has an order unit, the spectral map
$\Spec\RD^\Delta\to\Spec\RD^\Gamma$ is the identity of the common
ultrafilter space $\beta\N_{>0}$.
\end{corollary}
\begin{proof}
Both endpoints of the interval are admissible, so uniqueness means that they
coincide; if the groups coincide, both endpoints are $\mathcal C$. An order
unit of $\Gamma$ remains one in $\Delta$: first bound a given element by an
element of $\Gamma$, then by a multiple of the order unit. Both quotients are
then products of fields (\cref{ent:ps:thm:dichotomy}), whose spectra are the
ultrafilter spaces, and indicators are preserved.
\end{proof}

\begin{corollary}[Stable endpoint spectra]\label{ent:ps:cor:endpoints}
For every cofinal extension of this subsection, contraction restricts to
homeomorphisms of the minimal and of the maximal spectra, both identified
with the identity of $\beta\N_{>0}$. Every minimal prime of $\RD^\Gamma$ has a
unique prime extension to $\RD^\Delta$. Every maximal ideal has a unique
\emph{maximal} extension, but it can have further nonmaximal extensions.
\end{corollary}
\begin{proof}
The embedding $\Gamma_\Ucal\subseteq\Delta_\Ucal$ is cofinal: dominate a
representative of each new coordinate by an old exponent and pass to the
ultraproduct. A convex subgroup of $\Delta_\Ucal$ containing all of
$\Gamma_\Ucal$ is therefore all of $\Delta_\Ucal$, which by
\cref{ent:ps:thm:contraction} gives unique extension at the minimal endpoint
$\mathcal C=\Gamma_\Ucal$. At the maximal endpoint $\mathcal C'=0$ contracts
to $\mathcal C=0$, and \cref{ent:ps:thm:spectrum} gives exactly one maximal
prime over the selected ultrafilter. Principal points behave in the same way
by coordinate evaluation, and \cref{ent:ps:thm:topology} makes both endpoint
maps the identity of $\beta\N_{>0}$.
\end{proof}

\subsection{An explicit continuum splitting with no new zeros}
\label{ent:ps:subsec:splitting}

The sixth source's main construction separates new cuts in the value group
from new zeros, inside the surreal numbers. It enlarges the group
$\Gamma_\infty$ of \cref{ent:subsec:gamma-infinity} by rational indices:
\begin{equation}\label{ent:ps:eq:groups}
 \Gamma_\infty=\bigoplus_{j\in\Z_{\ge0}}\Q e_j
 \quad\subseteq\quad
 \Delta_\Q=\bigoplus_{q\in\Q_{\ge0}}\Q e_q ,
\end{equation}
finitely supported rational combinations, a nonzero one being positive
exactly when its coefficient at the largest occupied index is positive. The
integer-indexed basis vectors of $\Delta_\Q$ are those of $\Gamma_\infty$; the
sixth source writes $f_q$ for the basis of $\Delta_\Q$ and $f_j=e_j$, and the
single letter $e_q$ is used here because our $f$ is the function of
\eqref{ent:eq:f-explicit}. Take $\gamma_n=e_n$ for $n\ge1$. Then $P_D$ is that
function $f$, and $\rho_n=t^{-e_n}$. In this subsection $q$ and $u$ denote
rational numbers, not the node count $q_n$ or a unit.

\begin{lemma}[The two groups]\label{ent:ps:lem:groups}
Both groups in \eqref{ent:ps:eq:groups} are divisible and of countable
cofinality, neither has an order unit, and $\Gamma_\infty$ is cofinal in
$\Delta_\Q$. Their radius subgroups are
\[
 H_n^{\Gamma_\infty}=\bigoplus_{0\le j\le n}\Q e_j,
 \qquad
 H_n^{\Delta_\Q}=\bigoplus_{q\in\Q\cap[0,n]}\Q e_q ,
\]
and the quotients are identified with the finite combinations supported on
indices strictly greater than $n$.
\end{lemma}
\begin{proof}
Every nonzero finite support has a largest index, and a basis vector with a
larger integer index dominates every ordinary multiple of the given vector;
this gives the absence of an order unit, countable cofinality and cofinality
of $\Gamma_\infty$ in $\Delta_\Q$. A vector supported on indices at most $n$
is bounded by a multiple of $e_n$, and a vector with a larger leading index
is not; this gives the radius subgroups, the first of which is
\cref{ent:cor:residue-fields}. Discarding the coordinates of index at most
$n$ gives the quotient identifications.
\end{proof}

Fix a nonprincipal $\Ucal$ and write $\Gamma_{\infty,\Ucal}$ and
$\Delta_{\Q,\Ucal}$ for the two ultraproduct value groups. Define the convex
subgroup
\begin{equation}\label{ent:ps:eq:Cint}
 \mathcal C_{\mathrm{int}}=\prod_\Ucal(\Q e_{n+1})\ \subseteq\
 \Gamma_{\infty,\Ucal},
\end{equation}
the internal ultraproduct of the coordinate subgroups $\Q e_{n+1}$ under the
identifications of \cref{ent:ps:lem:groups}, and for real $0<\theta<1$
\begin{equation}\label{ent:ps:eq:Ctheta}
 \mathcal C'_\theta=\prod_\Ucal\Bigl(\bigoplus_{\substack{q\in\Q\\
   n<q\le n+1+\theta}}\Q e_q\Bigr)\ \subseteq\ \Delta_{\Q,\Ucal}.
\end{equation}
These are convex subgroups: each coordinate subgroup is convex, and an order
interval bounded by an element of one of them stays inside it on a member
of $\Ucal$.

\begin{theorem}[Cofinal prime splitting without new zeros]
\label{ent:ps:thm:splitting}
For the inclusion \eqref{ent:ps:eq:groups}, the function
$P_D(Z)=\prod_{n\ge1}(1-t^{e_n}Z)$ is entire over both $K_{\Gamma_\infty}$
and $K_{\Delta_\Q}$ and has exactly the same simple zeros $t^{-e_n}$ in both
fields. Nevertheless, for every nonprincipal $\Ucal$ the primes
$\mathfrak p^{\Delta_\Q}_{\Ucal,\mathcal C'_\theta}$, $0<\theta<1$, form a
strictly decreasing chain with continuum many members, all of which contract
to the single intermediate prime
$\mathfrak p^{\Gamma_\infty}_{\Ucal,\mathcal C_{\mathrm{int}}}$.
Intermediate means strictly between the unique minimal and the unique maximal
prime of its fibre.
\end{theorem}
\begin{proof}
Entireness and the zeros are \cref{ent:ps:lem:comparison}. At coordinate $n$
the only integer index in $(n,n+1+\theta]$ is $n+1$, so
$\mathcal C'_\theta\cap\Gamma_{\infty,\Ucal}=\mathcal C_{\mathrm{int}}$; the
equality holds for ultraproduct classes, because membership on a member of
$\Ucal$ can be represented coordinatewise there and changes off it are
irrelevant. If $\theta<\theta'$, choose a rational $u$ with
$\theta<u<\theta'$; then $b_u=[e_{n+1+u}+H_n^{\Delta_\Q}]_\Ucal$ lies in
$\mathcal C'_{\theta'}$ but not in $\mathcal C'_\theta$, so
$\mathcal C'_\theta\subsetneq\mathcal C'_{\theta'}$. The subgroup
$\mathcal C_{\mathrm{int}}$ is nonzero, since it contains
$[e_{n+1}]_\Ucal$, and proper, since it omits $[e_{n+2}]_\Ucal$, so its prime
is intermediate. The inclusion rule of \cref{ent:ps:thm:spectrum} and
\eqref{ent:ps:eq:contraction} finish the proof.
\end{proof}

\begin{proposition}[Entire functions separating the new primes]
\label{ent:ps:prop:splitting-witness}
For each rational $u\in(0,1)$ there is $G_u\in\mathcal E_{\Delta_\Q}$ with
$G_u(t^{-e_n})=t^{e_{n+1+u}}$ for all $n\ge1$. Whenever $\theta<u<\theta'$,
its class lies in $\mathfrak p^{\Delta_\Q}_{\Ucal,\mathcal C'_\theta}$ and not
in $\mathfrak p^{\Delta_\Q}_{\Ucal,\mathcal C'_{\theta'}}$. So the
distinctions of \cref{ent:ps:thm:splitting} are witnessed by entire functions
over the enlarged workspace, not only by elements of an ordered group.
\end{proposition}
\begin{proof}
The prescribed values have positive coarsened valuation for every $n$, so
\cref{ent:cor:single-node} over $K_{\Delta_\Q}$ realizes them. Their
ultraproduct value is $b_u$, whose leading index $n+1+u$ exceeds every index
allowed in the coordinate subgroups of $\mathcal C'_\theta$, so
$b_u>\mathcal C'_\theta$, while $b_u\in\mathcal C'_{\theta'}$. Apply
\cref{ent:ps:cor:membership}.
\end{proof}

\begin{corollary}[A maximal ideal with continuum many prime extensions]
\label{ent:ps:cor:maximal-splitting}
In the same example, for every nonprincipal $\Ucal$ the old maximal ideal
$\overline{\mathfrak M}{}^{\Gamma_\infty}_\Ucal$ is the contraction of a
strictly decreasing chain of continuum many nonmaximal primes of
$\RD^{\Delta_\Q}$, all strictly below
$\overline{\mathfrak M}{}^{\Delta_\Q}_\Ucal$. Unchanged zero sets and
homeomorphic maximal spectra therefore do not make the full spectral
contraction map injective.
\end{corollary}
\begin{proof}
Replace \eqref{ent:ps:eq:Ctheta} by
$\mathcal C'^{\,0}_\theta=\prod_\Ucal\bigl(\bigoplus_{q\in\Q\cap(n,n+\theta]}
\Q e_q\bigr)$. There is no integer strictly between $n$ and $n+1$, so these
convex subgroups meet $\Gamma_{\infty,\Ucal}$ in $0$; they are nonzero, and
they increase strictly with $\theta$ by a rational witness between two
cutoffs. Their primes are nonmaximal, strictly decreasing, and contract to
$\mathfrak p^{\Gamma_\infty}_{\Ucal,0}=\overline{\mathfrak
M}{}^{\Gamma_\infty}_\Ucal$. An analytic witness at a rational cutoff $u$
interpolates the values $t^{e_{n+u}}$ at the nodes.
\end{proof}

\begin{remark}[Why the internal subgroup is needed]
The subgroup $\mathcal C_{\mathrm{int}}$ is not the convex subgroup generated
by the single class $[e_{n+1}]_\Ucal$: it contains $[q_ne_{n+1}]_\Ucal$ for
arbitrary rational sequences $(q_n)$, unbounded ones included, and the
intersection in the proof of \cref{ent:ps:thm:splitting} keeps every such
sequence. Replacing $\mathcal C_{\mathrm{int}}$ by that smaller convex hull
would give a false contraction statement.
\end{remark}

\paragraph{What the splitting does and does not mean.}
Every prime of the chain lies over the same ultrafilter of ordinary nodes,
and $P_D$ vanishes only at the same node coordinates. The new distinctions
come from scales of the enlarged value group lying between consecutive old
integer-indexed levels. They are invisible to the values of old-field
functions after contraction and are detected by the $G_u$. The theorem is a
failure of the inference ``no new zeros under cofinal extension, hence no new
primes above the divisor''. It is not a failure of injectivity of the
coefficient map, nor a claim that a polynomial acquires more roots than its
degree: the divisor is infinite, the primes are algebraic objects of its
quotient, and the new points are non-evaluation points. The cofinal
extension of \cref{ent:subsec:two-extensions}, which adjoins $\omega^{1/2}$
to $\Gamma_\infty$, lies between the two groups of \eqref{ent:ps:eq:groups};
nothing is asserted here about its own prime spectrum.

\paragraph{The surreal realization.}
Under the normal-form embedding \eqref{ent:eq:embedding},
$t^g\mapsto\omega^{-g}$, take $e_q=\omega^q$ for $q\in\Q_{\ge0}$ as elements
of the additive group of $\No$; finite rational combinations of distinct
such monomials are ordered by their largest exponent, as
\eqref{ent:ps:eq:groups} requires, and $e_j=\omega^j$ is the choice of
\eqref{ent:eq:concrete-group}. The real workspaces
$\R((\omega^{-\Gamma_\infty}))\subset\R((\omega^{-\Delta_\Q}))\subset\No$ and
their complexifications in $\No[i]$ are set-sized subfields, and
\[
 P_D(Z)=\prod_{n\ge1}\bigl(1-\omega^{-\omega^n}Z\bigr),\qquad
 \rho_n=\omega^{\omega^n},\qquad
 G_u(\omega^{\omega^n})=\omega^{-\omega^{n+1+u}}\quad(u\in\Q\cap(0,1)),
\]
the first being \eqref{ent:eq:f-explicit} and the nodes being
\eqref{ent:eq:surreal-roots}. These formulas specify values of genuine entire
power series over the larger fixed Hahn workspace.

\begin{corollary}[Surreal-workspace prime splitting]
\label{ent:ps:cor:surreal-splitting}
There are explicit set-sized subfields
$\R((\omega^{-\Gamma_\infty}))\subset\R((\omega^{-\Delta_\Q}))\subset\No$, and
their complexifications in $\No[i]$, such that the entire divisor
$\prod_{n\ge1}(1-\omega^{-\omega^n}Z)$ has the same simple surreal zeros
$\omega^{\omega^n}$ in both workspaces, while a single prime of the smaller
divisor quotient is the contraction of a chain of continuum many primes of
the larger one. This holds over real and over complex coefficients.
\end{corollary}
\begin{proof}
Apply \cref{ent:ps:thm:splitting} through \eqref{ent:eq:embedding}. The
interpolation, classification and contraction arguments do not use
algebraic closedness and hold for either coefficient field.
\end{proof}

\begin{remark}[Workspace-relative, not class-wide]
``The same zeros'' means the zeros in the two specified fields. No strongly
summable value of this product is asserted at every point of $\No[i]$: the
ambient class contains scales beyond any set-sized cofinal sequence, and at
such a scale strong summability fails by \cref{ent:thm:extension-criterion}. The ultraproduct residue
fields are not identified canonically with subfields of the surreal field.
The construction shows a dependence on the chosen set-sized universe of
scales: an analytic object keeps its formula and every zero under an
enlargement while its prime spectrum refines, the refinement being encoded
by convex subgroups of an ultraproduct of quotient value groups. All
carriers are sets and all embeddings are explicit; no paradox is involved.
The Boolean layer is stable --- every stage has the same indicators and the
same ultrafilter space --- and what changes inside a fibre is the ordered
hierarchy of value cuts. Ordinary zero coordinates, ultrafilter choices of
infinitely many shells, and convex subgroups of their asymptotic value group
are three separate pieces of information.
\end{remark}

\subsection{Multiplicity: the limit of the reduced description}
\label{ent:ps:subsec:multiplicity}

The results so far concern one \emph{simple} node per shell. Multiplicity is
another, logically separate, source of nonmaximal primes. For positive
integers $m_n$ put
\[
 D_{\mathbf m}=\sum_{n\ge1}m_n[\rho_n],\qquad
 P_{D_{\mathbf m}}(Z)=\prod_{n\ge1}\Bigl(1-\frac Z{\rho_n}\Bigr)^{m_n},
\]
the canonical product of \cref{ent:thm:product} and the setting of
\cref{ent:sm:subsec:all-maximal}; the sixth source calls it $Q$, a letter
reserved here for \eqref{ent:eq:Qn}. Write $\ord_{\rho_n}(F)$ for the order of
a nonzero $F$ at $\rho_n$, zero when $F(\rho_n)\ne0$.

\begin{lemma}[Divisibility with arbitrary finite multiplicities]
\label{ent:ps:lem:multiplicity-division}
$P_{D_{\mathbf m}}$ is entire, with zeros exactly the $\rho_n$ of
multiplicities $m_n$, and $F\in\mathcal E_\Gamma$ is divisible by
$P_{D_{\mathbf m}}$ if and only if $\ord_{\rho_n}(F)\ge m_n$ for every $n$.
\end{lemma}
\begin{proof}
The first assertion is \cref{ent:thm:product}; over $\R$, a zero in
$\R((t^\Gamma))$ is a zero of the same order in $\C((t^\Gamma))$. The
divisibility criterion is \cref{ent:prop:kernel-direct}, since
$\ord_{\rho_n}(F)\ge m_n$ says that the prescribed zero jets vanish; over
$\R$ it holds by \cref{ent:prop:real-form}. The sixth source reproves this
lemma along the lines of \cref{ent:ps:rem:route}, with each index repeated
$m_n$ times.
\end{proof}

\begin{theorem}[Radicals and the bounded-multiplicity threshold]
\label{ent:ps:thm:multiplicity}
For every $F\in\mathcal E_\Gamma$,
\begin{equation}\label{ent:ps:eq:radical}
 F\in\sqrt{P_{D_{\mathbf m}}\mathcal E_\Gamma}
 \iff
 \text{some integer }d\ge1\text{ has }
 d\ord_{\rho_n}(F)\ge m_n\text{ for every }n .
\end{equation}
In particular $\sqrt{P_{D_{\mathbf m}}\mathcal E_\Gamma}=P_D\mathcal E_\Gamma$
if and only if $\sup_nm_n<\infty$. If $m_n\le M$ for all $n$, the natural map
$\mathcal E_\Gamma/(P_{D_{\mathbf m}})\to\mathcal E_\Gamma/(P_D)$ has nilpotent
kernel of exponent at most $M$ and induces a homeomorphism of prime spectra,
so the prime and dimension statements of
\cref{ent:ps:thm:spectrum,ent:ps:thm:dichotomy} persist for uniformly bounded
multiplicities. This transfers the prime spaces and dimensions, not the
reduced-ring conclusions such as von Neumann regularity or semiheredity. If
$(m_n)$ is unbounded, some prime ideal of $\mathcal E_\Gamma$ contains
$P_{D_{\mathbf m}}$ but not $P_D$, so the reduced description cannot account
for all primes above such a divisor.
\end{theorem}
\begin{proof}
$F\in\sqrt{P_{D_{\mathbf m}}\mathcal E_\Gamma}$ means $F^d\in
P_{D_{\mathbf m}}\mathcal E_\Gamma$ for some $d\ge1$; orders multiply by $d$,
and \cref{ent:ps:lem:multiplicity-division} gives \eqref{ent:ps:eq:radical}.
If $m_n\le M$, then $P_{D_{\mathbf m}}$ divides $P_D^M$ and $P_D$ divides
$P_{D_{\mathbf m}}$, so
$P_D^M\mathcal E_\Gamma\subseteq P_{D_{\mathbf m}}\mathcal E_\Gamma\subseteq
P_D\mathcal E_\Gamma$. Since $P_D\mathcal E_\Gamma$ is radical
(\cref{ent:ps:lem:idempotents}), the radicals agree. The kernel
$P_D\mathcal E_\Gamma/P_{D_{\mathbf m}}\mathcal E_\Gamma$ has vanishing $M$th
power, every prime contains a nilpotent ideal, and the quotient
correspondence gives the homeomorphism. Conversely, if $P_D\in
\sqrt{P_{D_{\mathbf m}}\mathcal E_\Gamma}$, some $P_D^M$ is divisible by
$P_{D_{\mathbf m}}$, and its order $M$ at each $\rho_n$ forces $m_n\le M$.
When $(m_n)$ is unbounded, the class of $P_D$ in
$\mathcal E_\Gamma/(P_{D_{\mathbf m}})$ is therefore not nilpotent;
localizing at its powers gives a nonzero ring, a maximal ideal of which
contracts to a prime avoiding that class, and its inverse image in
$\mathcal E_\Gamma$ is the required prime.
\end{proof}

\begin{example}[Pointwise nilpotence is not uniform nilpotence]
Take $m_n=n$. In the finite jet algebra at each $\rho_n$ the class of $P_D$
is nilpotent, since $P_D$ has a simple zero there. Its class in
$\mathcal E_\Gamma/(P_{D_{\mathbf m}})$ is not nilpotent, because no single
power has order at least $n$ at every node. Unbounded multiplicity cannot be
removed by observing that every jet variable is nilpotent; the exponent must
be uniform over the whole sequence.
\end{example}

\begin{remark}[What the threshold does not do]
\Cref{ent:ps:thm:multiplicity} does not classify the primes that avoid $P_D$
when the multiplicities are unbounded. It identifies the exact threshold at
which the reduced spectrum suffices and proves that beyond it there are
additional primes; no assertion about the complete spectrum in that regime is
hidden in the other results. By \cref{ent:sm:thm:maximal-all}, every maximal
ideal containing $P_{D_{\mathbf m}}$ is an evaluation ideal or some
$\mathfrak M_\Ucal$, and each of these contains $P_D$; so the additional
primes are not maximal. That consequence is the merge's.
\end{remark}

\subsection{Sensitive steps and what this section does not establish}
\label{ent:ps:subsec:limits}

\paragraph{Logically sensitive steps.}
The sixth source lists the following as parts of its proofs, not as
qualifications. The coarsened ultrafilter values are read in the quotient
groups $\Gamma/H_n$; the scalar field is not embedded as a constant sequence
into the residue fields $\kappa_n$, and a scalar may have support in no single
$H_n$ although its leading exponent lies in $H_n$ eventually (compare
\cref{ent:rem:coeff-action}). The convex-subgroup parameters of
\cref{ent:ps:thm:spectrum} need not be internal ultraproduct subgroups; every
convex subgroup of $\Gamma_\Ucal$ is allowed, and only the subgroups of the
splitting example are internal by construction. Every localization of the
reduced quotient is a domain although the quotient has many zero divisors;
the annihilating indicator kills the kernel of the fibre map after
localization. Prime-chain length is not inferred from the number of zeros:
the witnesses of \cref{ent:ps:prop:witnesses,ent:ps:prop:splitting-witness}
have no zero node values and distinguish only asymptotic valuation cuts. The
analytic steps it lists --- a second evaluation radius in the coefficient
criterion, a coefficient bound for a formal inverse that is not itself
entire, and a finite initialization before damping --- are those of
\cref{ent:thm:criterion,ent:prop:kernel-direct,ent:thm:image}. Finally, the
quantifiers differ in statements that look alike: $x_n\in V_n$ means
$v(x_n)\ge-m\gamma_n$ for some $m$ depending on $n$; a class of
$\Gamma_\Ucal$ lies in the convex subgroup generated by $\zeta$ when
$\abs g\le m\zeta$ for one ordinary integer $m$ outside the ultraproduct,
whereas an internal subgroup of sequences of rational multiples is in general
larger; $g>\mathcal C$ quantifies over every element of $\mathcal C$; a
family of locally nilpotent elements need not be globally nilpotent in the
product; and a chain with continuum many members is not thereby well
ordered.

\paragraph{Non-claims of the sixth source.}
The following are retained without weakening.
\begin{itemize}[leftmargin=1.4em]
\item ``Entire'' is strong summability at every point of one fixed set-sized
Hahn field, not entireness on $\No[i]$; nothing is claimed for the proper
class.
\item The exact value image and quotient, the canonical product and
divisibility, the nonprincipal maximal ideals with their residue fields, the
identification of the maximal spectrum with $\beta\N_{>0}$, and the
preservation of entireness and zeros under cofinal extension are inherited
from this report and are not new; the specialized simple-node proof is
supplied only for dependency control (\cref{ent:ps:rem:route}).
\item The ultraproduct and valuation-domain mechanism for primes is classical
\cite{FFW,StacksValuation}; the Hahn support machinery is classical
\cite{Neumann}; the normal forms are classical \cite{Conway,Gonshor}; the
valuation-ring elimination behind the Smith forms is standard. No new abstract
theory of product rings is claimed.
\item The proposed contributions --- the analytic prime fibres, the
reduced-divisor dichotomy and the cofinal splitting construction --- were
not found in the sources it compared. That is not proof of their absence
from the literature; the search was targeted, not exhaustive, and no
priority is certified.
\item No named published conjecture is solved. The manuscript has not been
refereed; no Lean file is supplied and no machine-checked proof is claimed;
its suggested Lean dependency order (Hahn-series, analytic,
commutative-algebra, ordered-group and ultraproduct interfaces) is a
suggestion, not code or a claim about present Mathlib coverage. It does not
claim to have verified this repository completely: its comparison was
concentrated on this report at its pin, and it states that the formalization
status of those reports is not inferred from the presence of their LaTeX
statements.
\item The complete classification is for one simple node per occupied shell,
extended to bounded multiplicities only for prime spaces and dimension. Shells
carrying several nodes, unbounded multiplicities and the proper class
$\No[i]$ are not classified.
\item In the order-unit case the quotient is zero-dimensional; the full ring
$\mathcal E_\Gamma$ is not claimed to be. The bounded-multiplicity transfer
covers prime spaces and dimension, not von Neumann regularity or
semiheredity.
\item No assertion is made that zeros change under cofinal extension; the
point of \cref{ent:ps:thm:splitting} is that they do not.
\item The finite checks (Suite~F of \cref{ent:app:checks}) do not establish
strong summability, construct ultrafilters, prove continuum chains or certify
novelty.
\end{itemize}

\section{Assembly of the classifications}
\label{ent:sec:classification}

\begin{proof}[Proof of \cref{ent:thm:main}]
The ring is an integral domain because it is a subring of $\K[[Z]]$. Its
units are exactly the nonzero constants by \cref{ent:cor:units}: a unit is
zero-free, and nonzero constants are invertible. GCDs exist by
\cref{ent:cor:gcd}.

If $\Gamma$ has no countable cofinal subset, \cref{ent:cor:cofinality} gives
$\E=\K[Z]$, which is a PID over a field. If $\Gamma$ has an order unit, its
positive integer multiples form a countable cofinal subset, so nonpolynomial
functions exist, and \cref{ent:thm:bezout} supplies the B\'ezout property. If
$\Gamma$ has countable cofinality and no order unit,
\cref{ent:thm:no-bezout} supplies a proper ideal generated by two coprime
functions. Such an ideal cannot be principal, since any principal generator
would be a common divisor and hence a unit. Non-Noetherianity in both
nonpolynomial cases is \cref{ent:prop:nonnoetherian}.

These cases are mutually exclusive and exhaustive. For $\Gamma=0$, strong
summability at $1$ says that only finitely many coefficient supports may
contain $0$, hence only finitely many coefficients are nonzero.
\end{proof}

\begin{proof}[Proof of \cref{ent:thm:main-extension}]
Part (a) is \cref{ent:thm:extension} together with
\cref{ent:lem:coarsened-ring}; part (b) is
\cref{ent:cor:entire-extension}, whose coarse form is
\cref{ent:thm:extension-criterion}; part (c) is
\cref{ent:thm:zero-conservation}; part (d) is
\cref{ent:thm:no-continuation}; part (e) is
\cref{ent:cor:intrinsic-boundary}. The residue field statement is
\cref{ent:lem:coarsened-ring}.
\end{proof}

\begin{center}
\small
\begin{tabular}{@{}>{\raggedright\arraybackslash}p{0.34\textwidth}>{\raggedright\arraybackslash}p{0.23\textwidth}>{\raggedright\arraybackslash}p{0.34\textwidth}@{}}
\toprule
Value-group condition & Entire series & Ideal structure\\
\midrule
No countable cofinal subset & Only polynomials & PID\\[3pt]
An order unit exists & Nonpolynomial series exist &
B\'ezout, non-Noetherian\\[3pt]
Countable cofinality, no order unit & Nonpolynomial series exist &
GCD, not B\'ezout; two-generated proper ideal with no common zero\\
\bottomrule
\end{tabular}
\end{center}

The same three rows govern rectification in $d\ge2$ variables: only finite,
hence polynomially rectifiable, radially finite sets in the first; every
radially finite set line-rectifiable in the second; and in the third,
polynomial-part configurations rectifiable but escaping infinitesimal
simplices not (\cref{ent:rc:thm:main,ent:rc:thm:omnific-main}, and the table
at the end of \cref{ent:rc:subsec:obstruction}).
The arithmetic-sampling results of source~11 do not depend on the row, and
need no divisibility: coefficientwise division and faithful flatness over
$\kk[X]$ hold for every $\Gamma$, dense arithmetic sampling forces a polynomial
in every row, and nonpolynomial realization of positive-tail ideals needs only
countable cofinality
(\cref{ent:as:thm:flat,ent:as:thm:sampling,ent:as:thm:realization}).

\section{Explicit higher-rank and surreal examples}
\label{ent:sec:examples}

\subsection{An explicit rank-infinite subgroup of the surreal field}
\label{ent:subsec:gamma-infinity}
Let
\begin{equation}\label{ent:eq:concrete-group}
 \Gamma_\infty=\bigoplus_{j\ge0}\Q e_j,
 \qquad e_j=\omega^j\in\No,
\end{equation}
with the order inherited from $\No$: every element has finite support, and
the largest index with a nonzero coefficient determines its sign. Thus
$e_{j+1}\ggA e_j$ and $(e_j)$ is cofinal, so $\cf(\Gamma_\infty)=\aleph_0$.
No element is an order unit: the largest index in its finite support can
always be exceeded. This is case (iii) of \cref{ent:thm:main}, and by
\cref{ent:cor:unit-structure} also
$\mathcal T_{\Gamma_\infty}^\times=K_{\Gamma_\infty}^\times$.

\subsubsection*{The B\'ezout counterexample realized}
Take $\gamma_n=e_n$ for $n\ge1$, a positive cofinal sequence with
$\gamma_{n+1}\ggA\gamma_n$, as \cref{ent:lem:escaping-classes} demands.
Using $t^\gamma=\omega^{-\gamma}$, \cref{ent:thm:no-bezout} gives roots
\begin{equation}\label{ent:eq:surreal-roots}
 \rho_n=\omega^{\omega^n},\qquad
 \sigma_n=\omega^{\omega^n}+\omega^{-\omega^{n+1}}
 \qquad(n\ge1),
\end{equation}
and the functions
\begin{align}
 f(Z)&=\prod_{n\ge1}\bigl(1-\omega^{-\omega^n}Z\bigr)
      =\prod_{n\ge1}\bigl(1-t^{e_n}Z\bigr),
      \label{ent:eq:f-explicit}\\
 g(Z)&=\prod_{n\ge1}\left(
 1-\frac{Z}{\omega^{\omega^n}+\omega^{-\omega^{n+1}}}\right)
\end{align}
are entire on the fixed surcomplex Hahn field $K_{\Gamma_\infty}$, not on all
of $\No[i]$. They have real Hahn coefficients and all the displayed roots are
real surreal numbers; by \cref{ent:cor:conjugation} this is visible from the
conjugation-invariance of their divisors. The failure of a B\'ezout identity
remains true even when complex Hahn coefficients are allowed in $A,B$, and
consequently also when $A,B$ are restricted to real Hahn coefficients. By
\cref{ent:sm:cor:boundary-pair}, the maximal ideals of
$\mathcal E_{\Gamma_\infty}$ containing $f$ and $g$ are exactly the
$\mathfrak M_\Ucal$ for nonprincipal $\Ucal$, with the residue fields of
\cref{ent:cor:residue-fields}. The fourth source records the same instance
in its group $\Gamma_*=\bigoplus_{j\ge1}\Q e_j$, which omits $e_0$; there
the residue fields are built on $\bigoplus_{1\le j\le n}\Q e_j$, of rank $n$.
The values $t^{-n!e_n}=\omega^{n!\,\omega^n}$ at the nodes $\omega^{\omega^n}$
are interpolated by the explicit series of
\cref{ent:sm:rem:simple-cardinal}, and the values
$t^{-e_{n+1}}=\omega^{\omega^{n+1}}$ are not interpolable.

The valuation formula \eqref{ent:eq:lambda} becomes
\[
 v(g(\rho_n))=e_{n+1}+(2-n)e_n+\sum_{1\le j<n}e_j,
\]
whose first values are
\begin{align*}
 \lambda_1&=e_2+e_1, & \lambda_2&=e_3+e_1,\\
 \lambda_3&=e_4-e_3+e_2+e_1,
 &\lambda_4&=e_5-2e_4+e_3+e_2+e_1.
\end{align*}
The negative lower-level coefficients do not affect domination: the
coefficient of the highest basis element $e_{n+1}$ is $1$.

For additional coefficient information, the $k$th coefficient of $f$ is
\[
 [Z^k]f=(-1)^k\sum_{1\le j_1<\cdots<j_k}
                         t^{e_{j_1}+\cdots+e_{j_k}},
 \qquad
 v\bigl([Z^k]f\bigr)=e_1+\cdots+e_k,
\]
the minimum coming from the unique subset $\{1,\ldots,k\}$, whose coefficient
is $(-1)^k$ and therefore cannot cancel. Division by the finite integer $k$
still leaves a cofinal sequence of Archimedean classes, which independently
exhibits the coefficient growth required by \cref{ent:thm:criterion}. The
reciprocal $1/\sigma_j$ has the explicit support-certified form
\[
 \frac1{\sigma_j}
   =t^{e_j}\sum_{r\ge0}(-1)^r t^{r(e_j+e_{j+1})},
\]
a valid positive-support geometric identity by
\cref{ent:lem:neumann}(c). Its support need not itself be cofinal in the
rank-infinite group; that is irrelevant to the legitimacy of this single
coefficient. Entirety of the product is supplied by the cofinal escape of the
factor indices $j$, not by a false claim about $r(e_j+e_{j+1})$.

\subsubsection*{The same workspace seen from its canonical product}
The natural product over the \emph{whole} basis,
\begin{equation}\label{ent:eq:infinite-product}
 P(Z)=\prod_{j\ge0}(1-t^{e_j}Z),
\end{equation}
differs from \eqref{ent:eq:f-explicit} by the single linear factor of index
zero:
\begin{equation}\label{ent:eq:P-vs-f}
 P(Z)=(1-t^{e_0}Z)f(Z)=(1-\omega^{-1}Z)f(Z).
\end{equation}
It is therefore an equally explicit member of the same ring, with one extra
zero at $\omega$.

\begin{theorem}[A genuine infinite-rank surcomplex entire function]
\label{ent:thm:infinite-example}
The product \eqref{ent:eq:infinite-product} belongs to
$\mathcal E_{\Gamma_\infty}$ and is nonpolynomial. Its zeros are exactly
\begin{equation}\label{ent:eq:surreal-zeros}
 r_j=t^{-e_j}=\omega^{\omega^j}\qquad(j\ge0),
\end{equation}
all simple and positive surreal. Its coefficient of $Z^k$ is
\begin{equation}\label{ent:eq:elementary-coeff}
 c_k=(-1)^k\sum_{j_1<\cdots<j_k}t^{e_{j_1}+\cdots+e_{j_k}},
 \qquad
 v(c_k)=e_0+e_1+\cdots+e_{k-1}
\end{equation}
for $k\ge1$, with $c_0=1$. At a scale $s\in\Gamma_\infty$ put
\[
 N_s=\#\{j:e_j+s<0\},\qquad M_s=\#\{j:e_j+s=0\}.
\]
These are finite, $M_s\in\{0,1\}$, and
\begin{align}
 m_s(P)&=\sum_{e_j+s<0}(e_j+s),\label{ent:eq:product-profile}\\
 I_s(P)(X)&=(-1)^{N_s}X^{N_s}(1-X)^{M_s}.\label{ent:eq:product-initial}
\end{align}
Finally
\begin{equation}\label{ent:eq:product-derivative}
 v\bigl(P'(r_j)\bigr)=e_j+\sum_{k<j}(e_k-e_j).
\end{equation}
\end{theorem}
\begin{proof}
Since $e_j\to+\infty$ cofinally, the proposed zero divisor is radially
finite, and \cref{ent:thm:product} gives the product's existence, entireness
and exact simple zero set. Expanding over finite subsets gives
\eqref{ent:eq:elementary-coeff}; among the $k$-element subsets, the sum of
the first $k$ basis elements is uniquely least in the largest-index order and
its coefficient is $(-1)^k$, so it cannot cancel and every $c_k$ is nonzero.

The slope sequence $(e_0+\cdots+e_{k-1})/k$ is cofinal: once $k-1$ exceeds
the largest index occurring in a prescribed group element, its positive
$e_{k-1}$ coordinate dominates that element. This is a second, direct
verification of entireness via \cref{ent:thm:criterion}.

At scale $s$ only finitely many factors have $e_j+s\le0$. A factor of
negative value contributes that value to $m_s$ and contributes $-X$ to the
normalized initial polynomial; a factor of value zero contributes $1-X$; all
positive factors contribute $1$. This proves
\eqref{ent:eq:product-profile}--\eqref{ent:eq:product-initial}.

Differentiating the factor of index $j$ and evaluating at $r_j$ kills all
other product-rule terms, giving
\[
 P'(r_j)=-t^{e_j}\prod_{k\ne j}(1-t^{e_k-e_j}).
\]
The factors with $k>j$ form a positive-support unit, and the finitely many
factors with $k<j$ have valuations $e_k-e_j$; adding these proves
\eqref{ent:eq:product-derivative}.
\end{proof}

This example is confined to no single Archimedean scale, yet its entire
theory is nontrivial and its zero divisor completely explicit. Together with
\cref{ent:rem:two-existence} it shows that nonpolynomial entire functions and
nonconstant restricted units are genuinely different existence questions.

\subsection{Two series with the same leading valuations and different
boundary behaviour}
\label{ent:subsec:hidden-tails}
Let $\Gamma=\Q\epsilon\oplus\Q\eta$ with $0<\epsilon\ll\eta$ and define
\begin{align}
 A(Z)&=\sum_{n\ge1}t^{n^2\epsilon}Z^n,\label{ent:eq:hidden-A}\\
 B(Z)&=\sum_{n\ge1}
       \left(t^{n^2\epsilon}+t^{\eta-n^2\epsilon}\right)Z^n.
       \label{ent:eq:hidden-B}
\end{align}
Their nonzero coefficients have exactly the same valuations $n^2\epsilon$.
Nevertheless their \emph{monomial} strong evaluation domains differ. For
$s=a\eta+b\epsilon$ with $a,b\in\Q$,
\begin{equation}\label{ent:eq:hidden-profile}
 \begin{array}{c|ccc}
 &a<0&a=0&a>0\\\hline
 A(t^s)&\text{not strong}&\text{strong}&\text{strong}\\
 B(t^s)&\text{not strong}&\text{not strong}&\text{strong}
 \end{array}
\end{equation}
For $a>0$ the leading valuations have $\eta$ coordinate $na\to+\infty$, and
\cref{ent:lem:cofinal-tails} proves both positive assertions. For $a<0$ those
leading valuations strictly decrease. For $a=0$ the exponents of $A(t^s)$ are
$(n^2+nb)\epsilon$, which after finitely many terms strictly increase with
only finite repetition at any exponent, so they form a well-ordered support.
The extra exponents of $B(t^s)$ are $\eta+(nb-n^2)\epsilon$, which strictly
decrease for all sufficiently large $n$; they destroy strong summability
without affecting a single leading coefficient valuation.

Neither series is entire over the whole rank-two field: their coefficient
slopes $n\epsilon$ remain below $\eta$, so \cref{ent:thm:criterion} correctly
excludes both. There is no tension between this example and the
valuation-only criterion. Leading valuations can fail to decide strong
summability at a \emph{boundary} scale, while global summability at
\emph{every} scale removes the ambiguity. Formula
\eqref{ent:eq:hidden-profile} is stated for monomial arguments only; it does
not silently assert a classification for arbitrary nonmonomial arguments of
$B$.

\subsection{One extension preserves the function; another destroys entireness}
\label{ent:subsec:two-extensions}
Adjoin $\omega^{1/2}$ to $\Gamma_\infty$ and take the divisible additive span.
This is an ordered-group extension in which $\Gamma_\infty$ remains cofinal,
since a sufficiently large integer power of $\omega$ dominates every finite
combination of the added and old generators. The series $P$, and likewise the
pair $f,g$, remain entire by
\cref{ent:cor:entire-extension,ent:cor:no-repair}; in particular the B\'ezout
failure persists.

Instead adjoin the new generator $e_*=\omega^\omega>\Gamma_\infty$ and set
$\Delta=\Gamma_\infty\oplus\Q e_*$. At the point $Z=t^{-e_*}$ the nonzero
summands of $P$ have leading exponents
$(e_0+\cdots+e_{k-1})-ke_*$, whose consecutive differences are
$e_k-e_*<0$; the defining family is therefore not strongly summable, exactly
as \cref{ent:thm:extension-criterion} predicts.
\Cref{ent:thm:no-continuation} rules out any alternative entire power series
over $K_\Delta$ agreeing with $P$ on the old field, or even on $\C$. On the
still-defined domain $\mathcal D_{\Gamma_\infty,\Delta}$, however, $P$ is
unchanged and its zeros remain exactly \eqref{ent:eq:surreal-zeros}, by
\cref{ent:thm:extension,ent:thm:zero-conservation}. This concrete pair is the
clearest illustration of \cref{ent:rem:reconcile}.

\subsection{Why finite rank does not give the counterexample}
\label{ent:subsec:finite-rank}
For $\Gamma=\Q^d$ in lexicographic order let $h$ lie in the largest
Archimedean class. Integer multiples of $h$ are cofinal, so $h$ is an order
unit and the entire ring is B\'ezout by \cref{ent:thm:bezout}, even when
$d>1$. The same holds for $\Q e_1\oplus\Q e_2$ with $e_2\ggA e_1$: take
$h=e_2$. It is \emph{absence of a greatest Archimedean class}, together with
countable cofinality, that forces the negative case --- not high rank.

\subsection{An uncountable hierarchy gives polynomial rigidity}
\label{ent:subsec:omega-one}
Let
\begin{equation}\label{ent:eq:Gamma-omegaone}
 \Gamma_{\omega_1}=\bigoplus_{\alpha<\omega_1}\Q e_\alpha,
 \qquad e_\alpha=\omega^\alpha,
\end{equation}
again ordered by the largest nonzero index. Its cofinality is $\aleph_1$: the
basis is cofinal, whereas any countable family of finite supports has its
indices bounded below $\omega_1$ and a larger basis element bounds the whole
family. Therefore
\begin{equation}\label{ent:eq:uncountable-rigidity}
 \mathcal E_{\Gamma_{\omega_1}}
 =\mathcal T_{\Gamma_{\omega_1}}
 =K_{\Gamma_{\omega_1}}[Z],
\end{equation}
the middle equality because a nonpolynomial restricted series would itself
supply a countable sequence of coefficient valuations tending cofinally to
infinity. This example shows that infinite rank by itself does not even
guarantee nonpolynomial entire functions: rank is not the driver, the order
type above the scales is.

\subsection{The separation table}
\label{ent:subsec:table}
\begin{center}\small
\begin{tabularx}{\textwidth}{@{}p{0.24\textwidth}cYYY@{}}
\toprule
Value group & $\cf$ & Has an order unit & Nonpolynomial entire series &
Nonconstant restricted units\\\midrule
$\Q$ or $\R$ & $\aleph_0$ & Every positive element & Yes & Yes\\[0.45em]
$\Q\epsilon\oplus\Q\eta$, $\epsilon\ll\eta$ & $\aleph_0$ &
Yes; not every positive element & Yes & Yes, with the sharper gap
test\\[0.45em]
$\Gamma_\infty$ & $\aleph_0$ & No & Yes & No\\[0.45em]
$\Gamma_{\omega_1}$ & $\aleph_1$ & No & No & No\\
\bottomrule
\end{tabularx}
\end{center}

The second and third rows are the two nonpolynomial cases of
\cref{ent:thm:main}; the third is the non-B\'ezout one. Note that $\R$ is
uncountable of rank one with $\cf=\aleph_0$, while $\Gamma_\infty$ has
infinite rank with $\cf=\aleph_0$: neither cardinality nor rank is the
invariant at work.

\section{Relation to other reports in this collection}
\label{ent:sec:collection}

This report is a merge of eight manuscripts, each prepared against a pinned
snapshot of the same repository, and each states explicitly which of its
neighbours' results it does \emph{not} claim. Those statements are preserved
here, and the cross-references are reciprocal.

\paragraph{Rank-one Berkovich geometry --- this report extends it.}
The report \emph{From Surcomplex Hahn Fields to Berkovich Geometry}
\cite{RepoRankOne} fixes $K=\C((t^\R))$: the value group is $\R$, the field
is spherically complete, and $|a|_v=\exp(-v(a))$ is a genuine real-valued
non-Archimedean absolute value. It defines intrinsic convergence using that
real-valued absolute value. Its entire-function section proves
zero-free entire rigidity and a convergent canonical product \emph{there},
and none of the source manuscripts merged here claims to originate either
point.
The present report is the arbitrary-rank generalization. Three specific
relations should be recorded.
\begin{itemize}[leftmargin=1.4em]
\item \emph{Where rank one sits in the classification.} For $\Gamma=\R$ we
have $\cf(\R)=\aleph_0$ and $1$ is an order unit, so rank one falls squarely
in case (ii) of \cref{ent:thm:main}: nonpolynomial entire functions exist,
$\mathcal E_\R$ is a non-Noetherian B\'ezout domain, and arbitrary radially
discrete Hermite interpolation is available. The pathology of case (iii) is
invisible at rank one.
\item \emph{Why rank one is not general.} That report's
\texttt{prop:rankobstruction} proves that $\Q+\Q\omega$ admits no
order-preserving \emph{injective} additive map to $\R$. That is exactly why
\cref{ent:lem:coarsening} constructs a \emph{coarsening} with a possibly
nonzero kernel $H$ rather than an embedding, and why it is available only
when an order unit exists. The two statements are complementary: one forbids
faithful scalarization in general, the other supplies the best available
non-faithful substitute in the order-unit case.
\item \emph{The ring chains resemble each other and are not equal.} That
report's strict chain $K[Z]\subsetneq\mathcal T\subsetneq\mathcal P
\subsetneq\mathcal A\subsetneq\mathcal F\subsetneq\mathcal E^-$ locates its
polynomial-coefficient and formal algebras. Our chain
\eqref{ent:eq:chain} resembles it but is not the same chain, and the symbols
$\mathcal T$, $\mathcal A$, $\mathcal E$ do not denote the same rings in the
two reports. \Cref{ent:conv:ring} forbids transferring a statement across
that boundary.
\end{itemize}
Conversely, the rank-one report's own entire-function section should be read
with the knowledge that the arbitrary-rank classification now exists and that
its rank-one instance lands in the B\'ezout case.

\paragraph{Analysis and Foundations --- this report is published beside them.}
\Cref{ent:cor:all-no} re-derives a conclusion the collection already contains
twice: the \emph{Analysis} report proves all-scale polynomial rigidity
directly (its Theorem \texttt{f:thm-rigidity}), and the \emph{Foundations}
report records the exact boundary of the phenomenon with a pair of series
over $\C((t^\Q))$ that differ only in the sign of an exponent, noting that
$\sum t^{+n^2}z^n$ is strongly summable at every point of that field and is
\emph{nonpolynomial} while remaining coherent at every radius, and that this
does not refute the all-scale theorem but shows that the all-scale theorem
does not transfer to a \emph{fixed} value group. That example is precisely
\cref{ent:rem:criterion-strength} of this report, and the trichotomy of
\cref{ent:thm:main} explains why it works over $\Q$ --- $\cf(\Q)=\aleph_0$
and $1$ is an order unit, so $\Q$ is a case (ii) group --- and what it would
take to break it, namely replacing $\Q$ by a group of uncountable cofinality.
Conversely, our reduction in \cref{ent:cor:all-no} consumes the foundations
report's \emph{localization principle}: every set of surcomplex numbers lies
in one divisible set-sized workspace. That principle is what makes the
reduction legitimate, and it is cited here, not reproved.

\paragraph{Global divisors and analytic geometry --- different rings, and
two restricted products that restrict different things.}
The global-divisor report \cite{RepoDivisors} works in
$\mathcal O(U)((t^\Gamma))$, with ordinary holomorphic coefficient functions
on a common complex domain; its divisor, interpolation and Picard
obstructions live in a different function ring from the strongly entire ring
studied here, and nothing in this report contradicts or restates them. The
analytic-geometry report states the ring discipline that
\cref{ent:conv:ring} invokes by name, for its own four coefficient rings.
Neither the objects nor the theorems transfer.

That report and \cref{ent:thm:main-image} nevertheless have the same moral,
and the resemblance is close enough that the difference has to be stated
rather than left to be inferred. Its \texttt{global:thm:interpolation} and
\texttt{global:cor:CRT} prove a \emph{support-restricted} Chinese remainder
theorem: the correct target of global interpolation is a restricted
subalgebra $\mathcal B$ of the product of the local finite algebras, not the
full product, and its \texttt{global:ex:selector} shows the failure already
for prescribed values in $\{0,1\}$ --- a $0/1$ selector on clusters of
growing size whose intrinsic quotient coordinates acquire an infinite
strictly decreasing sequence of exponents. \Cref{ent:thm:quotient} has the
same shape, and \cref{ent:thm:close-pair} is likewise a $0/1$ prescription
that is not interpolable. The two restricted products are not the same
construction and neither implies the other:
\begin{itemize}[leftmargin=1.4em]
\item \emph{What is restricted.} There, the condition is that the
\emph{union of the coefficient supports} of the local remainders be
\emph{well ordered}. Here, the condition is \emph{eventual membership} of the
shell remainder's coefficients in a \emph{coarsened valuation ring} $V_n$.
Well-ordering of a support is not a valuation bound, and a valuation bound is
not a well-ordering condition.
\item \emph{What the index set is.} There, the index set is a family of
clusters of \emph{infinitesimally close moving points}, each cluster bounded
by a polynomial $P_c$, and the obstruction comes from negative powers of the
separations \emph{inside} one cluster. Here, the index set is the sequence of
\emph{valuation shells escaping through the Archimedean hierarchy}, and the
obstruction comes from how far out in that hierarchy the data reaches. A
single shell may itself contain close nodes, and then
\cref{ent:thm:close-pair} shows the same divided-difference mechanism acting
inside the shell --- but the shellwise condition is a tail condition on
$n$, which has no counterpart there.
\item \emph{Which ring.} That theorem is stated in
$\mathcal O(U)((t^\Gamma))$, which \cref{ent:conv:ring} and
\cref{ent:subsec:scope-limits} explicitly disclaim for every statement in
this report, and which the third source manuscript disclaims for its own. Its
value group may be $\Q$, with an order unit, where
\cref{ent:cor:order-unit-image} says our restriction disappears entirely;
its obstruction survives there, because it is not our obstruction.
\end{itemize}
Both reports therefore say ``the image of an infinite Chinese remainder map
is a proper restricted subproduct, and the restriction is forced, not
chosen''. They compute different restrictions, for different reasons, in
different rings. The fourth source's restricted product
$\prod_n'(A_n,B_n)$ is this report's $\RD$, and the distinction applies to
it verbatim; that source, too, disclaims every ring other than the strongly
entire one.

A related collision of words concerns the analytic-geometry report. There an
ideal is called \emph{free} when its elements have no common geometric zero,
and \texttt{analytic:thm:fixed-structure} produces, in $\mathcal
O(D)((t^\R))$ on a fixed disk, maximal ideals containing no evaluation kernel
and two-generated ideals with empty zero set; ``boundary'' there refers to a
fixed disk. The fourth source calls a nonprincipal ultrafilter \emph{free}
and calls the non-evaluation maximal ideals of
\cref{ent:sm:subsec:all-maximal} \emph{boundary ideals}. This report writes
``nonprincipal'' and ``non-evaluation'' throughout. The objects live in
different rings and no statement transfers, although the moral --- a proper
finitely generated ideal with no common zero --- is shared, and both reports
credit Henriksen \cite{Henriksen}.

The same report also has a prime spectrum with continuum chains, and the
resemblance to \cref{ent:ps:sec:spectrum} is again close enough to state the
difference. Its section on the prime spectrum of the fixed-polydisk ring
$\mathcal O(D)((t^\R))$ (\texttt{analytic:sec:fixed-spectrum},
\texttt{analytic:thm:fixed-gauge-prime}, \texttt{analytic:thm:fixed-fibre})
builds, from a nonprincipal ultrafilter and the gauges $w(n)=n^r$, a
continuum-sized chain of primes in one fibre over a classical maximal
ideal; \cref{ent:ps:lem:gaps} uses the same power gauges $\lfloor
n^\theta\rfloor$. The rings are different --- ordinary holomorphic
coefficients on a fixed disk and value group $\R$ there
\cite{RepoAnalyticGeometry}, strongly entire
series over a fixed Hahn field here --- and so are the reasons: its chain
comes from the growth of Hahn valuations of values at a fixed simple divisor
inside one classical fibre, ours from the quotient value groups
$\Gamma/H_n$, which are zero when $\Gamma$ has an order unit, as $\R$ does
(\cref{ent:ps:thm:dichotomy}). No statement transfers. That report credits
Henriksen's work on prime ideals of the ring of classical entire functions,
which attaches primes to the growth of zero multiplicities
\cite{Henriksen53}; the sixth source does not cite it, and we record it as
context for \cref{ent:ps:thm:multiplicity}, not as a checked comparison.

\paragraph{Holonomic rigidity and several variables --- one device, two
uses.}
The report \emph{Holonomic Rigidity for Entire Hahn Functions}
\cite{RepoHolonomic} works with the same whole-family entire functions of
several variables. Its Lemma \texttt{hol:lem:generic} chooses a complex point
avoiding countably many nonzero polynomials over $\K$ by picking coordinates
algebraically independent over a countable field; the covering step of
\cref{ent:mv:thm:exceptional} is the same device, applied to complex
polynomials and refined to open disks. Its several-variable theorem uses the
whole-family convention of \cref{ent:mv:def:entire} and groups a family into
total-degree blocks only as a step inside a proof, which is consistent with
the separation of conventions in \cref{ent:mv:subsec:blocks}. No theorem is
shared: that report proves D-finite rigidity, not an extension domain. The
locus $\Pdir(F)$ is written so to avoid the nonabelian-support report's
$\operatorname{Pol}$, which denotes a normalized polar factor.
Its coarsening-descent subsection does share one step with the seventh
source: \texttt{hol:cf:lem:reduction} reduces a one-variable entire series with
coefficients in the coarsened ring of a proper convex subgroup to a
polynomial, and \cref{ent:rc:thm:reduction} is the several-variable form of
that lemma, with evaluation compatibility and Hasse derivatives, applied to
an automorphism together with its inverse. The seventh source does not cite
it; the credit is the merge's. That report itself describes its polynomial
half as a coefficient-field form of this report's obstruction for noncofinal
extensions, and both descend from the coarsened reduction in the proof of
\cref{ent:thm:units}.

\paragraph{Omnific integers --- a shared normal form, no shared theorem.}
At its pin the seventh source could not compare with the collection's
omnific-integer reports, which were placed later, both in
\texttt{docs/surreal/}: \emph{Omnific Integers and Omnific--Diophantine
Geometry} and the report on set-sized quotients of the omnific integers. Its
only omnific input, $\Oz=\Pi\oplus\Z$ with $\Pi$ the purely infinite normal
forms, read in a Hahn workspace as
\eqref{ent:rc:eq:omnific-real}--\eqref{ent:rc:eq:omnific-complex}, agrees with
the foundations report (\texttt{found:eq:omnific}) and with the former's
rings (\texttt{odg:def:rings}), and its workspace-relative radial finiteness
is consistent with the former's set-sized workspaces
(\texttt{odg:lem:workspace}); \cref{ent:rc:subsec:limits} explains why radial
finiteness does not transfer to a larger workspace. No theorem is shared: the
rectifying maps are coordinate changes of $\K^d$ and are not asserted to
preserve $\Oz$, while those reports study the arithmetic of $\Oz$ itself.

\paragraph{Omnific integers --- source~11 shares a theorem.}
Unlike the seventh source, source~11 shares mathematics with the omnific
reports. Its classification \eqref{ent:as:eq:int-intro} has a polynomial half
and an analytic half. The polynomial half ---
$\operatorname{Num}_r(\Oz)=\Pi[\mathbf X]\oplus\Int(\Z^r)$ with its Gaussian form,
the finite-grid test and binomial closure --- is \texttt{osq:pm:thm:numreal},
\texttt{osq:pm:thm:numgaussian}, \texttt{osq:pm:thm:grid} and
\texttt{osq:pm:lem:binomial} of the set-sized-quotients report
\cite{RepoQuotients}, which transfers them to fixed Hahn workspaces
$A_\Gamma=\Z\oplus J_\Gamma$ (\texttt{osq:or:cor:workspace}; its $J_\Gamma$ is
our $\PiK$ with the exponent sign reversed, and its $A_\Gamma$ our $\OzK$); its
matrix and jet reductions (\texttt{osq:or:thm:matrix}) are the polynomial
counterpart of the Hasse-jet conditions of \cref{ent:as:thm:jets}. The
analytic half, the collapse of an entire series to a polynomial
(\cref{ent:as:thm:sampling}), is what that report leaves out: after
\texttt{osq:or:cor:workspace} it says that ``entire Hahn series and expressions
of nonstandard degree are outside the theorems'', and it repeats the exclusion
in its lists of limits. Likewise the omnific Diophantine-geometry report
\cite{RepoODG} lists among its non-claims that ``Polynomial lifting does not
license evaluating infinite power series''; source~11 does not contradict that
sentence, it shows what does hold for entire series in one fixed workspace.
Source~11 names both reports as antecedents of the constant-term pullback but
read neither article in full; the label-level comparison is the merge's, and
neither report is changed by it. The omnific-preserving-automorphisms report
\cite{RepoOPA} studies the automorphism distinctions on which the scale
polynomials depend (\cref{ent:as:subsec:boundaries}); no theorem of it is
used.

\section{What the results do and do not establish}
\label{ent:sec:scope}

\subsection{Three independent distinctions}
\label{ent:subsec:three-distinctions}
The first distinction is between strong Hahn summation and intrinsic
valuation convergence. Positive well-ordered support suffices for the former;
cofinal progress through every group threshold is needed for the latter. On a
single disk they differ. For a series strongly evaluable at \emph{every}
point, \cref{ent:thm:criterion} shows the distinction disappears within its
original field.

The second is between countable cofinality and an order unit. Countable
cofinality supplies a sequence of ever larger scales, possibly moving through
infinitely many Archimedean classes; an order unit supplies all those bounds
by the integer multiples of one element. Entire functions need only the
first; a nonconstant restricted inverse requires the second. The group
$\Gamma_\infty$ separates them, and that separation is the whole content of
the gap between cases (ii) and (iii) of \cref{ent:thm:main}.

The third is between an intrinsic topology in one workspace and the fine
topology tested against every surreal scale. The same countable sequence may
converge intrinsically in $\K$ and fail to converge after a noncofinal
enlargement. None of our proofs identifies these topologies. See
\cref{ent:subsec:three-meanings}.

\subsection{Established, and no more}
The results isolate three logically different phenomena. Coefficient slopes
determine whether an ordinary countable power series can be evaluated at
every scale in a fixed field. Support-controlled preparation and canonical
products completely describe zeros and divisibility within that category. The
order structure above those scales determines whether arbitrary values and
jets can be prescribed and whether coprime functions generate the unit ideal.

In particular the factorization monoid alone cannot decide the B\'ezout
property. In both nonpolynomial cases a function up to a scalar is exactly a
radially finite zero divisor, and every effective subdivisor is again
principal; nevertheless only one of those cases admits unrestricted discrete
interpolation. The additive structure of the ring, not merely its
multiplicative divisibility relation, is essential.

The non-B\'ezout example is not a delicate choice of a transcendental
ordinary complex function or of an uncomputable coefficient. Its nodes,
displacements and factor coefficients are explicit Hahn expressions over
rational data. What cannot be replaced by a finite calculation is the
quantification over all coefficient indices and all Archimedean scales.

\subsection{Scope that must not be silently enlarged}
\label{ent:subsec:scope-limits}
The following restrictions are stated by the source manuscripts and are
retained without weakening.

\paragraph{Hypotheses that cannot be relaxed.}
All supports and index sets are sets; there is no proper-class summation.
The field is a \emph{full} Hahn field with complex residue coefficients, and
$\Gamma$ is divisible. Algebraic closedness is used essentially, to split the
prepared finite polynomial; it is the sole import of
\cref{ent:thm:zero-conservation}. A nondivisible group, or a
support-restricted Hahn subfield, may lose the roots needed to split the
prepared polynomial, so the factorization theorem cannot simply be copied
there. The Hahn--Neumann support lemmas
(\cref{ent:lem:neumann}) and algebraic closedness \cite{Poonen} are imported
inputs, not proved here. \Cref{ent:sec:multivariable} alone works without
divisibility and over an arbitrary trivially valued coefficient field, under
its own statement of hypotheses, and none of the divisible-group results of
the other sections is transferred to it. Its subsections
\crefrange{ent:rc:subsec:setting}{ent:rc:subsec:limits}, the seventh source's,
are the exception: they keep divisibility, take $\kk\in\{\R,\C\}$, use the
uncountability of $\kk$ once, and use no algebraic closedness. Its last
subsections \crefrange{ent:as:subsec:setting}{ent:as:subsec:limits},
source~11's, need no divisibility, allow $\Gamma=0$, and take
$\kk\in\{\R,\C\}$.

\paragraph{Categories excluded.}
Outside \cref{ent:sec:multivariable} the ring is one-variable and
strong-summation entire. It is \emph{not} the ring of
ordinary-holomorphic-coefficient Hahn families on one common complex domain,
not a ring of arbitrary fine-local surcomplex germs, and not a
several-variable analytic algebra; \cref{ent:sec:multivariable} treats
finitely many variables only for coefficient and unit criteria, extension
domains and exceptional directions, with no preparation, divisor or ideal
theory. The seventh source's subsections of that section add,
for configurations of points only, fat-point ideals and value interpolation on
line-rectifiable configurations; they give no preparation, no
positive-dimensional divisor theory and no GCD theory. Source~11's
subsections add the ideals extended from $\kk[X]$, with faithful flatness, and
the vanishing ideals of ordinary point sets and, for $\kk=\C$, of algebraic
sets defined over $\kk$; they give no preparation and no GCD theory, and
nothing for ideals not extended from $\kk[X]$. The arbitrary-rank
preparation theorem is
proved in $\C[X]((t^\Gamma))$, not asserted in a larger power-series
coefficient ring. There is no contour integral here, no global surcomplex
exponential at unrestricted imaginary arguments, and nothing about the
Berarducci--Mantova derivation; the statement about primitives in
\cref{ent:prop:operations} concerns termwise $Z$-integration with
coefficients fixed and says nothing about differential equations in the
coefficient field. All entire functions here are represented by one ordinary
formal power series at the origin: translation preserves this class, but does
not identify it with every locally analytic or surreal-differentiable
function. Piecewise functions on fine-topological domains, Gonshor's
exponential, locally represented classes and other support or summation
conventions lie outside the hypotheses, so
\cref{ent:cor:all-no} is \emph{not} a general no-go theorem for surreal
analysis.

\paragraph{Limits of the coarsening and of the extension theorem.}
The rank-one coarsening of \cref{ent:lem:coarsening} changes residue
information even though it preserves the topology and the entire ring;
residue-direction data is not identified across the two valuations. The
scalar-extension results keep the coefficient field $\C$ fixed and use
canonical ordered-group inclusions of full Hahn fields only: they are not
theorems about arbitrary extensions of incomplete subfields, about changing
residue fields, or about analytic sheaves. The hidden-tail table
\eqref{ent:eq:hidden-profile} is stated for monomial arguments and does not
classify arbitrary nonmonomial arguments of $B$. The single-variable
uniform-domain statement is \emph{false} in several variables --- a series
independent of one coordinate imposes no constraint in that coordinate. Its
exact replacement for whole-family summation is now
\cref{ent:mv:thm:domain}, whose limits are recorded in
\cref{ent:mv:subsec:limits}; what remains open is recorded in
\cref{ent:q:several,ent:mv:q:block-domain,ent:mv:q:effective}.
No density of $\C$ in any fine surreal topology is asserted, and strong Hahn
summability, intrinsic valuation convergence and the fine surreal topology
are deliberately never identified.

\paragraph{Limits of the exact image theorem.}
\Cref{ent:thm:main-image} concerns ordinary power series in \emph{one}
coordinate and radially finite divisors in a full Hahn field. It is not a
theorem about $\mathcal O(U)((t^\Gamma))$, about common-domain holomorphic
coefficient germs, or about functions continuous for the fine surreal
topology; the third source manuscript disclaims that ring by name, as
\cref{ent:conv:ring} does. It does not extend automatically to several
variables, where a shell need not carry one finite monic quotient controlling
all the prescribed data, and it supplies no interpolation criterion for an
infinite node set confined to a single valuation ball, where
\cref{ent:lem:shells} gives no shell decomposition at all. The phrase
``restricted product'' in \cref{ent:def:restricted} names an algebraic union
and asserts no locally compact topology, no adelic structure and no Banach
norm. \Cref{ent:thm:maximal} constructs a specified family of maximal ideals
with a field presentation for each; it does \emph{not} classify all maximal
ideals of $\E$ and does not assert that every non-evaluation maximal ideal
arises this way, and its nonprincipal ultrafilter is an ordinary
set-theoretic choice, not an algorithmically computed object. The fourth
source's \cref{ent:sm:thm:maximal-all} shows that, for one node per shell,
these and the evaluation ideals are all the maximal ideals containing $P_D$;
it too does not classify the maximal ideals of $\E$ that contain no such
product, does not treat shells with several nodes, and does not classify
prime ideals. The sixth source's \cref{ent:ps:thm:spectrum} classifies the
prime ideals containing such a product when every node is simple, and
\cref{ent:ps:thm:multiplicity} extends that to uniformly bounded
multiplicities for prime spaces and dimension; primes above products with
unbounded multiplicities or with several nodes on a shell, and primes
containing no such product, remain unclassified.

\paragraph{Limits recorded by the fourth source.}
Its unit-ideal criteria are not a corona theorem: no bounded function algebra
and no norm lower bound is specified. Its ultrafilter characters are ring
homomorphisms and are not asserted to preserve strong Hahn sums, so
$\Phi_\Ucal(P_D)$ must not be computed termwise
(\cref{ent:sm:rem:no-substitution}); they are not evaluation maps on a larger
surreal domain. Its word ``entire'' has fixed-field scope, and its
nonpolynomial interpolants make no claim against all-scale polynomial
rigidity. Its restricted product is algebraic, with no adelic or locally
compact topology. Its jets are formal derivatives in $Z$, not the
Berarducci--Mantova derivation. Its construction is a finite-step procedure
relative to exact field operations and to bounds $c_n$, cofactor classes and
remainders supplied as certificates; it does not make arbitrary Hahn
coefficients computable and does not decide membership in $H_n$ from an
unspecified representation of $\Gamma$. It claims nothing about surreal
exponential automorphisms, birthdays, integration or the proper-class
analytic theory. It reproved the support, product and division ingredients
it needed so that its interpolation theorem would not depend on this
report's unrefereed preparation and factorization arguments;
\cref{ent:app:audit} records that the second route indeed uses neither.

The construction behind \cref{ent:thm:image} is constructive only in a
limited and stated sense. Given explicit finite representations of a shell's
nodes and jets, \eqref{ent:eq:psi} is a finite computation in a quotient of
$\K[T]$, and a certificate $v([T^j]\Theta_n)\ge-c\gamma_n$ yields the explicit
damping exponent $d_n\ge\max\{n,2c+2\}$; each coefficient of the assembled
$F$ receives only finitely many corrections. This is not a decision procedure
for arbitrary input Hahn series: deciding their equality, computing a
valuation, or deciding membership in $V_n$ may require information absent
from a finite encoding. The theorem is an existence theorem with a
finite-step construction \emph{relative to} exact field operations and
valuation-bound certificates, and the maximal ideals of
\cref{ent:thm:maximal} are not constructive even in that sense.

Finally, \cref{ent:prop:real-form} states exactly which results survive over
$\R((t^\Gamma))$, and no more. Algebraic closedness remains essential to
\cref{ent:thm:roots,ent:thm:factorization} and hence to
\cref{ent:cor:gcd,ent:thm:zero-conservation}, as stated above; the real form
is available for the image, quotient and B\'ezout statements only because
their proofs split the relevant polynomials at the given nodes by hypothesis
and because \cref{ent:prop:kernel-direct} supplies the kernel without
preparation. Nothing in \cref{ent:prop:real-form} should be read as weakening
the algebraic-closedness hypothesis of the factorization theory.

\paragraph{Limits recorded by the sixth source.}
They are listed in full in \cref{ent:ps:subsec:limits}. In summary: its
``entire'' has fixed-field scope and nothing is claimed for $\No[i]$; the
quotient, the maximal ideals, their residue fields, the maximal spectrum and
the cofinal-extension results on entireness and zeros are inherited, not new;
the ultraproduct and valuation-domain mechanism, the Hahn support calculus,
the normal forms and the elimination behind the Smith forms are classical;
its classification covers one simple node per shell, and bounded
multiplicities only for prime spaces and dimension; in the order-unit case
only the quotient, not $\E$, is claimed zero-dimensional; and no change of
zeros under extension is asserted.

\paragraph{Limits recorded by the seventh source.}
They are listed in full in \cref{ent:rc:subsec:limits}. In summary: every
result is relative to one fixed workspace, and radial finiteness is not fine
discreteness and does not survive enlargement of the workspace; $\kk$ is $\R$
or $\C$, with no countable-field or descent statement; nothing classifies all
ideals, automorphisms, nonrectifiable sets or entire functions on the class;
line-rectifiability is not claimed equivalent to classical tameness; the
rectifying maps are not claimed to preserve $\Oz$, and no omnific
factorization question is settled; the obstruction is not shown to be the
only one; ideal generation is algebraic, and no unrestricted product
description of the evaluation quotient is given without an order unit.

\paragraph{Limits recorded by source~11.}
They are listed in full, numbered, in \cref{ent:as:subsec:limits}. In summary:
strong Hahn summation is not complex convergence; nonpolynomial results live
in fixed set-sized workspaces; the omnific locus carries a residue condition
and need not be algebraic, and the leading valuation is not an omnific test;
flatness and the image theorem are over $\kk[X]$, not $\K[X]$, and $\Ed$ is
not claimed Noetherian, principal or B\'ezout; $\OzK$ is not an integer
polynomial ring and $\operatorname{Frac}(\OzK)=\K$ is not assumed; the ordinary
binomial basis is not a Gaussian basis, and an infinite sample gives
polynomiality, not integer-valuedness; certificates exist but are not computable
in general, the undecidability concerns one presentation and is not a new
negative solution to Hilbert's tenth problem; and the scale polynomials depend
on the coefficient field and monomial section.

\paragraph{What is and is not claimed about priority and verification.}
The results of this report are self-declared proposed original contributions,
not a solution of any named published conjecture. The first two source
manuscripts state that priority is \emph{not} certified and that the literature checks
were targeted rather than exhaustive; in one of them several general web
searches returned irrelevant results and were explicitly not treated as
evidence of absence, and in the other a connector keyword search returning
nothing was likewise not treated as proof. The rank-one factorization
specialization is stated to be classical; standard Hahn algebra, algebraic
closedness and support lemmas are not claimed; Cherry himself calls his
factorization results well known; the rank-two distinction between positive
valuation and topological nilpotence is established
\cite[Section~6.2]{Conrad}; and the all-surcomplex polynomial corollary is
credited to the repository, as recorded in \cref{ent:rem:rigidity-credit}.
One source's audit further notes that Lazard was checked bibliographically
only and that Conway and Gonshor were not re-audited in full; its repository
review was a targeted read of four files, not a line-by-line proof audit or
an independent build. None of the eight manuscripts was independently
refereed and none was machine-checked: one records
\texttt{formal\_verification false},
\texttt{independent\_peer\_review false} and \texttt{priority\_certified
false} in its build record, the second states flatly that no Lean
verification is claimed, and the third states that it has not been
independently refereed or verified in Lean, that its literature comparison
was targeted and not a certification of priority, and that its repository
review was a targeted read of repository files at a pinned commit rather than
a successful clone, a line-by-line audit of every report, or a Lean build.
The third also states that its main image theorem is proved independently of
this repository's unrefereed general factorization and preparation claims,
and \cref{ent:app:audit} records where that independence does and does not
hold in the merged text. The fourth states that its proofs have not been
independently refereed or checked in a proof assistant, that its targeted
literature search is not an exhaustive priority determination, that
irrelevant search results were not used as evidence, that its repository
review was a targeted read through a connector at its pinned commit and not
an audit of every theorem, that no changes were made to the repository, and
that it credits the general ultrafilter method and the finite algebra as
classical; its package manifest records \texttt{independent\_peer\_review}
and \texttt{proof\_assistant\_verification} as false. The fifth records
\texttt{priority\_certified}, \texttt{formal\_verification} and
\texttt{independent\_peer\_review} as false; states that its comparison was
targeted, that no search returning nothing was treated as proof of novelty,
and that binary repository archives, retired historical manuscripts and the
Lean modules were not audited; and states that a successful build shows only
that its source typesets. The sixth states that no named published conjecture
is declared solved, that it has not been refereed or verified in Lean and
supplies no Lean file, that its novelty language means only ``not found in
the compared sources'' after a targeted comparison, not a priority
certification, that its review concentrated on this report and the
catalogue at its pin and is not a complete verification of the repository;
its build record gives the verification boundary as written proofs and
finite sanity checks, with no peer review or Lean verification claimed. The
seventh describes itself as an AI-assisted, unrefereed research draft with
explicit proofs and finite symbolic checks, not Lean-verified; it supplies no
Lean source and makes no claim that its results were added to or checked by
the repository's Lean build; it claims to settle no named published
conjecture; its ``proposed'' means only that no matching result was located in
the inspected sources, after searches whose coverage it calls necessarily
incomplete and unable to rule out a statement under different terminology; it
calls the projection-and-shear strategy classical and the jet counting
elementary; and its build record sets formal mathematical verification and
independent peer review to false.
Source~11 describes its results as proposed contributions with complete
written proofs, and claims no certified historical priority, no independent
refereeing and no proof-assistant verification; it settles no named published
conjecture; its repository comparison was targeted, through a connector whose
long responses were sometimes truncated, after a local clone failed, and did
not read every manuscript; its literature search was limited, and a failed
search is not evidence of originality; and its artifact manifest sets
\texttt{proof\_assistant\_verification} and
\texttt{independent\_referee\_review} to false.
No theorem below depends for its validity on the
correctness of an unrefereed repository report. The repository's Lean modules
on finite algebra and Hahn supports are plausible \emph{prerequisites} for a
future formalization, not a certificate that these proofs have been
formalized; missing implementation is not treated here as an open
mathematical problem. No repository files under \texttt{code/} or
\texttt{data/} have been modified.

\paragraph{What the finite checks do not do.}
See \cref{ent:app:checks} for the full statement. In summary: the first six
delivered suites verify $831$, $664$, $1958$, $1876$, $1072$ and $10889$
finite exact assertions respectively, the seventh runs five groups of exact
symbolic checks without an aggregate count, the eighth verifies $8774$ exact
assertions in ten groups, and they prove \emph{no} infinite
theorem. The first
three articles and their READMEs state that the finite truncations of the
counterexample pair \emph{are} coprime polynomials and \emph{do} satisfy
polynomial B\'ezout identities, so no finite computation can witness
\cref{ent:thm:no-bezout}. The preparation checks run modulo $t^9$ in an
integer exponent subgroup, not an arbitrary-rank group, and the finite
dominance samples test only $k\in\{1,2,10,100,10000\}$, not all integer
multiples. The checks cannot verify arbitrary well-ordered supports, the
cofinality dichotomy, the necessity direction of \cref{ent:thm:units}, or any
infinite product identity. The third suite additionally cannot verify
infinite strong summability, the existence or nonexistence of an infinite
entire interpolant, the restricted-product quotient, or the ultrafilter
construction; its rational polynomial specializations verify algebraic
identities only and are not offered as examples of positive fine valuation in
$\C$. The fourth suite likewise uses rational proxies for the finite shell
algebras and finite rational vectors for the valuation inequalities; it
verifies neither infinite summability nor cofinality nor the theorems, and
is not evidence of priority. The fifth verifies finite identities and
example certificates; it cannot prove that an infinite support is well
ordered, that a series is entire, that a general quotient fibre has finite
incidence, or that every countable union of algebraic sets is realized, and
it does not test the irrational example of \cref{ent:mv:ex:dense} through a
rational approximation. The sixth checks finite identities, ordered-group
inequalities in finite-support rational vectors, finite multiplicity tests
and Smith reductions over $\Z_{(2)}$; it does not establish strong
summability, construct a nonprincipal ultrafilter, enumerate convex
subgroups, prove the continuum-chain assertions, treat a proper-class
construction, certify entireness from finite data or certify novelty. The
seventh checks polynomial and rational instances only; it does not verify the
infinite Hahn summation arguments, the universal quantifiers over entire
automorphisms, or novelty. The
eighth checks finite Hahn polynomials with integer exponents, finite
Gr\"obner, syzygy and gcd instances and a bounded machine family; it does not
establish infinite summability, arbitrary-rank lifting, faithful flatness or
undecidability, and decides no halting problem.

\subsection{Questions opened by the classification}
\label{ent:subsec:questions}
The following are research questions generated by this work, not claims that
any particular publication lists them as open.

\begin{question}[Exact interpolation image]\label{ent:q:image}
When $\Gamma$ has countable cofinality but no order unit, characterize the
image of
\[
 \E\longrightarrow\prod_{x\in\supp D}\K[U]/(U^{D(x)})
\]
for a prescribed radially finite divisor $D$.
\end{question}
\emph{Answered, by \cref{ent:thm:image}.} The question is retained in the
form in which the two earlier manuscripts posed it, together with the note
they attached to it, because the note turned out to be the right guess and
locating the answer against it is the clearest statement of what was added.
That note read: the growth barrier gives explicit necessary inequalities and
proves that the image can be proper, but sufficiency for arbitrary nodes and
jets is not established, and a satisfactory answer should encode the
compatibility of divided differences across convex subgroups rather than only
the size of individual values. The shell Hermite polynomial $R_n$ of
\eqref{ent:eq:normalizedjets} is exactly that record of divided differences,
and $V_n$ is exactly the convex-subgroup datum;
\cref{ent:thm:close-pair} is the instance in which individual values are
harmless and the divided difference is not. What remains open after
\cref{ent:thm:image} is narrower and is recorded in
\cref{ent:q:image-remaining}.

\begin{question}[What the image theorem does not decide]
\label{ent:q:image-remaining}
\Cref{ent:thm:image} concerns radially finite divisors in a full Hahn field
and one ordinary variable. Is there a criterion for interpolation on an
infinite node set confined to a single valuation ball, where no shell
decomposition exists? Is there a several-variable replacement, where a shell
need not carry a single finite monic quotient controlling all the prescribed
data? Does the membership condition \eqref{ent:eq:maincondition} admit a
decision procedure for any natural class of finitely encoded inputs?
\end{question}
The first two are genuinely open here; \cref{ent:lem:shells} is what fails in
each case. The several-variable results of \cref{ent:sec:multivariable}
concern extension domains, not interpolation, and do not bear on the second.
\emph{Re-scoped by the seventh source:} its subsections of
\cref{ent:sec:multivariable} now bear on the second for values. On a
configuration of points that an entire automorphism carries into a coordinate
line, the values taken by entire functions are exactly those allowed by
\cref{ent:thm:image} at the image nodes (\cref{ent:rc:cor:evaluation}, the
merge's corollary of the seventh source's results); jets in several
variables and configurations that are not line-rectifiable remain open.
The third is discussed in \cref{ent:subsec:scope-limits}: the
construction behind \cref{ent:thm:image} is finite-step relative to exact
field operations and valuation-bound certificates, and is not by itself an
algorithm for arbitrary Hahn input.
\emph{Related (batch~32), not an answer.} Source~11's
\cref{ent:as:thm:undecidable} proves a different decision problem undecidable:
whether an entire series given by a computable coefficient stream with a
uniform quadratic bound takes an omnific value at one ordinary point. It says
nothing about the membership condition \eqref{ent:eq:maincondition}.

\begin{question}[Finitely generated ideals beyond their common divisor]
\label{ent:q:ideals}
Can one describe the ideal $(f,g)$ by its common zero divisor together with
an explicit obstruction in an interpolation quotient?
\end{question}
\emph{Answered for the case in which one generator is a canonical product,
by \cref{ent:thm:bezout-criterion}.} The exact reduction the earlier
manuscripts recorded --- after dividing out the GCD, the B\'ezout problem is
whether the inverse jets of one function along the other's zero divisor lie
in the interpolation image --- is still the right reduction, and
\cref{ent:thm:image} now supplies the missing intrinsic description of that
image, while \cref{ent:lem:shellunits,ent:lem:restrictedunits} turn it into
the pointwise test \eqref{ent:eq:valueupper}. \Cref{ent:thm:no-bezout} still
proves that a zero-divisor description alone is insufficient; the criterion
says what must be added, namely an upper valuation bound on the values of the
second function along the shells. Two parts of the question were left
unsettled by the criterion: it decides membership of $1$ in $(P_D,F)$, not
the structure of $(f,g)$ as an ideal for arbitrary $f$; and
\cref{ent:thm:maximal} constructs many maximal ideals containing such a
proper ideal without classifying the maximal spectrum of $\E$, or the
relations between restricted products attached to different divisors.

\emph{Partly answered further by the fourth source.}
\Cref{ent:sm:thm:finite-generators} decides membership of $1$ in
$(P_D,G_1,\ldots,G_s)$ for any finite number of generators, and
\cref{ent:sm:cor:membership} decides membership of an \emph{arbitrary}
element in $(P_D,G)$ when $D$ has one simple node per shell; since every
nonzero entire function is a scalar multiple of the canonical product of its
divisor (\cref{ent:thm:factorization}), so that $(f,g)=(P_{\Div f},g)$, this
decides membership in $(f,g)$ whenever the zeros of $f$ are simple, of
negative valuation, and one on each shell.
\Cref{ent:sm:thm:maximal-all} lists every maximal ideal containing such an
ideal. Still open: membership in $(P_D,G)$ when shells carry several nodes or
multiplicities, and a description of $(f,g)$ as an ideal beyond membership
tests. See also \cref{ent:q:spectrum}.

\emph{A related structural fact from the sixth source.} For one simple node
per shell the quotient $\E/(P_D)$ is a B\'ezout ring
(\cref{ent:ps:thm:bezout}), so every finitely generated ideal of $\E$
containing $P_D$ equals $(P_D,G)$ for a single $G$
(\cref{ent:ps:rem:fg-ideals}), and \cref{ent:ps:cor:membership} lists every
prime ideal containing it. This does not describe $(f,g)$ for arbitrary $f$,
and the open parts listed above remain open.

\begin{question}[The maximal spectrum]\label{ent:q:spectrum}
Do all non-evaluation maximal ideals of $\E$ arise as $\mathfrak M_\Ucal$ of
\cref{ent:thm:maximal} for some divisor and some nonprincipal ultrafilter,
and how do the restricted products attached to different divisors relate?
\end{question}
\emph{Partly answered by the fourth source.} Above a canonical product with
one node per shell, with arbitrary multiplicities and empty finite part, the
answer is yes, exactly: the maximal ideals containing $P_D$ are the
evaluation ideals and the $\mathfrak M_\Ucal$
(\cref{ent:sm:thm:maximal-all}), and they form a copy of $\beta\N_{>0}$
(\cref{ent:sm:thm:beta}). Still open, and not a hidden part of
\cref{ent:thm:maximal} or \cref{ent:sm:thm:maximal-all}: maximal ideals
containing only canonical products whose shells carry several nodes, where the
finite residue algebras have several maximal ideals; whether every
non-evaluation maximal ideal of $\E$ contains a product of the treated kind;
the prime ideals; and the relations between the restricted products of
different divisors.

\emph{The ``prime ideals'' clause is partly answered by the sixth source.}
Above a product with one \emph{simple} node per shell and empty finite part,
every prime ideal is an evaluation ideal or one of the primes
$\mathfrak p_{\Ucal,\mathcal C}$, indexed by a nonprincipal ultrafilter and a
convex subgroup of $\prod_\Ucal(\Gamma/H_n)$
(\cref{ent:ps:thm:spectrum}), with local rings, topology and Krull dimension
described in \cref{ent:ps:thm:local,ent:ps:thm:topology,ent:ps:thm:dichotomy}
and contraction under cofinal extension in \cref{ent:ps:thm:contraction}.
For uniformly bounded multiplicities the prime spectrum is the same
(\cref{ent:ps:thm:multiplicity}). Still open after the sixth source: the
primes above products whose shells carry several nodes
(\cref{ent:ps:q:several-nodes}); the primes above products with unbounded
multiplicities, where only the existence of primes not containing the reduced
product is known (\cref{ent:ps:q:unbounded}); primes and maximal ideals of
$\E$ containing no product of the treated kind; and the relations between the
restricted products of different divisors. The first clause of the question
is not affected.

\begin{question}[Several variables]\label{ent:q:several}
Which parts of the cofinality criterion and of the B\'ezout obstruction
extend to several ordinary variables, and which local preparation statements
retain explicit Hahn support certificates? In particular, since the
single-variable uniform extension domain is false in several variables, what
is the correct replacement --- presumably one depending on the unbounded
degree directions of the support, including the boundary where higher Hahn
tails matter?
\end{question}
One must separate the cofinal growth of total degree from the geometry of
positive-dimensional zero sets. Cherry's several-variable work is an
important rank-one antecedent \cite{Cherry}; no arbitrary-rank multivariable
GCD theorem is established here.
\Cref{ent:app:multivariate} proves only the finite-variable convergence and
unit criteria, and it does not extend the one-variable divisor
classification to higher-dimensional zero sets.

\emph{The ``In particular'' clause is answered by the fifth source,}
for whole-family summation: \cref{ent:mv:thm:domain} is the exact
replacement. It depends only on the ordinary support, through its weight
image in $\Delta/H$, and for semilinear supports only the period vectors,
the unbounded degree directions, enter (\cref{ent:mv:thm:semilinear}), in
line with the question's guess. The guess about the boundary was wrong in one
respect: within the entire category, higher Hahn tails do
\emph{not} matter (\cref{ent:mv:cor:tails}); they matter only for non-entire
series (\cref{ent:mv:ex:hidden}, and in rank two
\cref{ent:subsec:hidden-tails}). The other clauses remain open: whether the
B\'ezout obstruction and the ideal theory extend, and which local
preparation statements retain support certificates, in several variables.
The cofinality criterion does extend (\cref{ent:thm:multi-criterion}).

\emph{The ideal-theory clause, and interpolation of values, are partly
answered by the seventh source,} for configurations of points, not
positive-dimensional zero sets, and only for configurations that an entire
automorphism carries into a coordinate line. On such a configuration with
bounded multiplicities, the fat-point ideal has the explicit generators of
\cref{ent:rc:thm:fat-generators} and exactly $\binom{\bar m+d-1}{d-1}$ minimal
generators; for simple points it is a complete intersection with
$\Ed/I\cong\E/(P_D)$ (\cref{ent:rc:cor:ci}); for unbounded multiplicities and
$d\ge2$ it is not finitely generated (\cref{ent:rc:thm:unbounded}); and the
values taken there are described exactly by \cref{ent:thm:image}
(\cref{ent:rc:cor:evaluation}). Which configurations qualify is decided by the
order unit: all radially finite ones when $\Gamma$ has an order unit, and
without one a large class (\cref{ent:rc:thm:omnific-main,ent:rc:cor:distinct-radii})
but not all (\cref{ent:rc:thm:obstruction}). Still open after the seventh
source: whether the B\'ezout obstruction extends and what replaces the GCD
theory in several variables; which local preparation statements retain
support certificates; divisors with positive-dimensional support; ideals of
configurations that are not line-rectifiable (\cref{ent:rc:q:ideals}); and
jet interpolation in several variables.

\emph{Status (batch~32): the ideal-theory clause is partly answered further by
source~11,} for ideals extended from the coefficient polynomial ring, with
$\kk\in\{\R,\C\}$ and any value group. Membership in $I\,\Ed$ for an ideal
$I\subseteq\kk[X]$ is decided scale by scale, with exact normal forms and
$(I\,\Ed)\cap\kk[X]=I$
(\cref{ent:as:cor:ideal,ent:as:cor:normal-form}); $\kk[X]\to\Ed$ is faithfully
flat (\cref{ent:as:thm:flat}), so finite polynomial systems over $\kk$ acquire
no new entire relations; and the entire functions vanishing on a set of
ordinary points, or, for $\kk=\C$, on the $\K$-points of an algebraic set
defined over $\C$, form an extended ideal (\cref{ent:as:cor:vanishing}, the
merge's corollary), a first positive-dimensional case.
Still open after source~11: ideals not extended from $\kk[X]$, in particular
those extended from $\K[X]$ (\cref{ent:as:q:matrices}); whether the B\'ezout
obstruction extends and what replaces GCD theory; local preparation with
support certificates; zero sets that are not algebraic over $\kk$; ideals of
configurations that are not line-rectifiable; and jet interpolation. Source~11
does not claim $\Ed$ to be Noetherian, principal or B\'ezout. The text of the
question and of the earlier notes is kept.

The fifth source adds two questions, suggested by its results; like the
others, they are not claims that any publication lists them as open.

\begin{question}[The full block domain]\label{ent:mv:q:block-domain}
Can the homogeneous-block domain $\mathcal D^{\mathrm{hom}}(F)$ at arbitrary
tuples of new-field coordinates be described by a finite hierarchy of
quotient valuations and initial homogeneous forms, with a precise condition
for indefinite cancellation? Such a description should be invariant under
linear changes of coordinates and recover \cref{ent:mv:thm:rays} on old-field
rays.
\end{question}

\begin{question}[Effective support classes]\label{ent:mv:q:effective}
For which finitely encoded supports does the well-ordering test of
\cref{ent:mv:thm:domain} admit an effective exact procedure over specified
finite-rank ordered groups? Polynomially parametrized supports already give
nonclosed domains (\cref{ent:mv:thm:nonclosed}), and \cref{ent:mv:ex:dense}
shows that lower bounds on the weights are not enough for arbitrary
supports.
\end{question}
\emph{Related (batch~32), not an answer.} Source~11 shows that complete
finite certificates for its positive-tail ideals cannot be computed from
unrestricted computable coefficient streams (\cref{ent:as:thm:undecidable}),
and asks for effective subclasses (\cref{ent:as:q:effective}); the
well-ordering test of \cref{ent:mv:thm:domain} is a different test, and this
question is unchanged.

\begin{question}[Restricted and incomplete subfields]\label{ent:q:subfields}
How much of the divisor theory survives for incomplete surreal subfields or
for support-restricted Hahn subfields? Our proofs isolate two requirements:
closure under the positive-support preparation construction, and enough
algebraic roots to split the finite prepared polynomials. A classification of
natural surreal subfields meeting both would turn these full-Hahn results
into a more computationally selective theory.
\end{question}

The sixth source adds three questions, suggested by
\cref{ent:ps:sec:spectrum}; like the others, they are not claims that any
publication lists them as open.

\begin{question}[Several nodes on a shell]\label{ent:ps:q:several-nodes}
\Cref{ent:thm:image} uses finite shell algebras when several nodes share a
radius. Which additional prime fibres appear when distinct normalized nodes
collide after coarsening? The one-node valuation-domain description of
\cref{ent:ps:thm:spectrum} does not apply directly to those finite algebras.
\end{question}

\begin{question}[Unbounded multiplicities]\label{ent:ps:q:unbounded}
For $P_{D_{\mathbf m}}=\prod_n(1-Z/\rho_n)^{m_n}$ with $(m_n)$ unbounded, can
the primes not containing $P_D$ be described directly from asymptotic zero
orders and coarsened coefficient values? \Cref{ent:ps:thm:multiplicity} gives
an exact radical test but not that classification.
\end{question}

\begin{question}[When is the prime spectrum invariant?]
\label{ent:ps:q:invariance}
\Cref{ent:ps:thm:contraction} reduces unique extension of a prime to the
equality $\mathcal C_{\min}=\mathcal C_{\max}$. Which convenient conditions on
a cofinal inclusion $\Gamma\subseteq\Delta$ make these equalities hold for
every nonprincipal ultrafilter and every convex subgroup? The order-unit case
is one sufficient condition (\cref{ent:ps:cor:unique-extension}), but it does
not exhaust the conceivable cases.
\end{question}

The seventh source adds five questions, suggested by
\crefrange{ent:rc:subsec:setting}{ent:rc:subsec:limits}; it calls them natural
continuation problems, not assertions of known published open status.

\begin{question}[Rectifiability without an order unit]\label{ent:rc:q:criterion}
When $\cf(\Gamma)=\aleph_0$ and $\Gamma$ has no order unit, which radially
finite subsets of $\K^d$, $d\ge2$, are line-rectifiable? The polynomial-part
and distinct-radius classes are rectifiable
(\cref{ent:rc:thm:omnific-main,ent:rc:cor:distinct-radii}), and
full-dimensional infinitesimal simplices are an obstruction
(\cref{ent:rc:thm:obstruction}), which is not shown to be the only one. A
criterion might combine the shellwise geometry of finite residue
configurations with the compatibility of coordinate interpolation.
\end{question}

\begin{question}[Lower-dimensional targets]\label{ent:rc:q:targets}
For $1<r<d$, when is a radially finite set carried into an $r$-dimensional
affine subspace by an entire automorphism? The simplex example obstructs every
proper affine subspace; a hierarchy of sharp lower-dimensional cluster
criteria remains to be developed.
\end{question}

\begin{question}[Ideals of nonrectifiable configurations]\label{ent:rc:q:ideals}
Is every radially finite simple-point ideal of $\Ed$ finitely generated? Does
the configuration of \cref{ent:rc:thm:obstruction} have a similarly explicit
minimal generating family? \Cref{ent:rc:thm:fat-generators} is restricted to
line-rectifiable configurations.
\end{question}

\begin{question}[Extension of prescribed maps]\label{ent:rc:q:extension}
Which bijections or injections between radially finite configurations, for
instance omnific ones, extend to entire automorphisms? The shear of
\cref{ent:rc:prop:positive} straightens a set while keeping the
first-coordinate labels of a projection; that is where analogues of the
stronger classical notions of tameness should be investigated, rather than
assumed from line-rectifiability.
\end{question}

\begin{question}[Countable coefficient fields and descent]
\label{ent:rc:q:descent}
Can the uncountability of $\kk$ used in \cref{ent:rc:lem:projection} be
removed, or replaced by an exact field-of-definition criterion? Can a
rectifying automorphism be chosen over a prescribed smaller Hahn workspace?
\end{question}

Source~11 (batch~32) adds ten questions, suggested by
\crefrange{ent:as:subsec:setting}{ent:as:subsec:limits}. It calls them proposed
directions, not a list of previously published open problems, and asks that
nearby literature, which may contain special cases, be compared before any
priority claim.

\begin{question}[Polynomial matrices with Hahn coefficients]
\label{ent:as:q:matrices}
Does a degree- and support-controlled image theorem hold for matrices over
$\K[X]$ in arbitrary rank? In particular, is $\K[X]\to\Ed$ faithfully flat in
several variables without an order unit? The proof of
\cref{ent:as:thm:matrix} uses that ordinary polynomial coefficients do not mix
valuation scales; a positive extension needs a replacement for that
separation, not the same Gr\"obner argument.
\end{question}
In one variable the question has a positive answer (the merge's remark).
The ring $\mathcal E_{\Gamma,1}$ is a torsion-free module over the principal
ideal domain $\K[Z]$, hence flat. It is faithfully flat because no polynomial
$p$ of positive degree $n$ is a unit of $\mathcal E_{\Gamma,1}$: at a weight
$\gamma$ so negative that $v(a_n)+n\gamma<v(a_j)+j\gamma$ for all $j<n$, the
initial polynomial of the rescaling $p(t^\gamma X)$ has degree $n$, and initial
polynomials multiply \eqref{ent:eq:gauss-mult}, an identity in
$\kk[X]((t^\Gamma))$ that needs no divisibility. The question concerns
$d\ge2$.

\begin{question}[Meromorphic arithmetic rigidity]\label{ent:as:q:meromorphic}
For a quotient of entire Hahn functions, which hypotheses on an infinite
ordinary arithmetic sample force rationality? Polynomiality is too strong, as
$t^{-1}/(1+Z^2)$ shows, and a satisfactory statement must distinguish poles,
cancellation of common factors and the growth of denominators. A useful start
is to prescribe a fixed algebraic denominator and classify the numerators.
\end{question}

\begin{question}[Effective certificate classes]\label{ent:as:q:effective}
Find natural effectively presented subclasses of entire functions for which
$I_+(F)$ can be computed with a proved stopping bound: recurrences for scale
coefficients, finite differential systems and $q$-difference equations are
candidates. The target is a bound on when the generated ideal stabilizes, not
an algorithm that keeps producing larger partial ideals;
\cref{ent:as:thm:undecidable} rules out a solution for unrestricted computable
coefficient streams. (Compare \cref{ent:mv:q:effective}, which asks about a
different test.)
\end{question}

\begin{question}[Sampling sets beyond affine ordinary grids]
\label{ent:as:q:sampling-sets}
Characterize the subsets $\mathcal S\subseteq\K^d$ for which
$F(\mathcal S)\subseteq\OzK$ forces an entire $F$ to be a polynomial.
Invertible affine images of Zariski-dense ordinary sets have this property,
while the escaping zeros of \eqref{ent:as:eq:qproduct} do not. Intermediate
cases include moving nodes with colliding ordinary residues, several nested
infinitesimal scales, and mixed grids in which some coordinates escape and
others remain ordinary.
\end{question}

\begin{question}[Intrinsic arithmetic exceptional ideals]
\label{ent:as:q:intrinsic}
Can a useful substitute for $I_+(F)$ be defined without choosing the full
monomial section, while keeping a finite certificate theorem, for instance
through an associated graded object or compatible truncation operators? A
surcomplex version should state whether conjugation and the natural valuation
are part of the structure; the unrestricted Gaussian automorphism problems of
the collection's companion reports are not settled by source~11.
\end{question}

\begin{question}[Function theory on algebraic varieties]
\label{ent:as:q:varieties}
Develop the quotient $\Ed/I\,\Ed$ as a function algebra on a fixed ordinary
affine variety. \Cref{ent:as:cor:normal-form} gives exact coefficientwise
representatives and \cref{ent:as:thm:restriction} an arithmetic restriction
principle. Is the algebra independent of the chosen affine presentation, how
does it glue across charts, and how do nonreduced schemes and their jet modules
behave?
\end{question}

\begin{question}[Composition and exceptional ideals]
\label{ent:as:q:composition}
Give formulas or bounds for $I_+(F\circ G)$ under entire composition, with a
proof of the relevant composition domain. For affine substitutions over $\kk$
ordinary polynomial pullback controls the scale polynomials directly;
Hahn-coefficient substitutions can mix scales and cancel, and the right
invariant may need more than an ideal and a scale-zero polynomial.
\end{question}

\begin{question}[Other integer parts and coefficient rings]
\label{ent:as:q:integer-parts}
Which integer parts of a real closed Hahn field admit a sampling theorem of
this form? The proofs use an exact support cut and a coefficient-field
retraction, which an arbitrary integer part need not supply. On the arithmetic
side, for which subrings $\mathfrak o\subseteq\kk$ in place of $\Zk$ does
dense-grid sampling give an equally explicit classification by
$\Int(\mathfrak o^d,\mathfrak o)$?
\end{question}

\begin{question}[Characteristic and cardinal extensions]
\label{ent:as:q:characteristic}
The lifting and Hasse-jet arguments make sense over arbitrary fields. Over an
infinite coefficient field of positive characteristic, which arithmetic
versions remain meaningful? Over a finite field ordinary points cannot detect
all polynomials, so finite extensions or scheme-valued points are needed.
Infinitely many variables create another obstruction: polynomial ideals need
not be finitely generated, and the finite-scale conclusion may fail.
\end{question}

\begin{question}[Formalization and proof certificates]
\label{ent:as:q:formalization}
Can the development be formalized along the route below, and what is the right
certificate format? The source suggests proceeding through the growth
criterion, the scale-polynomial construction, the degree-transfer lemma,
finite polynomial matrix lifting, and then sampling and ideal membership, with
a separate arithmetic layer for constant extraction and finite differences.
An effective certificate should record a finite set of scales together with a
genuine proof that all remaining scale polynomials lie in their ideal; a
finite list without that proof is not a complete certificate, as the halting
example shows.
\end{question}

\subsection{A proof-assistant route}
The research statement does not require a formalization of the entire proper
class $\No$. A first formal development can work over a set-sized divisible
linearly ordered abelian group and a Hahn field. The most useful first target
is the global criterion \cref{ent:thm:criterion}, whose nontrivial
implication is a short ordered-group argument plus the definition of strong
summability; the next is the coarsened polynomial argument for
\cref{ent:thm:units}. Support preparation should be formalized only after the
finite-word Neumann lemma and arbitrary regrouping are available. The
explicit obstruction then requires only ordered-group inequalities, finite
product valuations and evaluation homomorphisms. This is a dependency plan,
not a formal-verification claim. The current repository ledger
\texttt{docs/FORMALIZATION.md} distinguishes checked Hahn support,
summability and algebraic-closure prerequisites from the still-pending
statements of this report. No Lean code is included in this report package.
For \cref{ent:ps:sec:spectrum} the sixth source suggests five interfaces in
dependency order: Hahn series (strong summability, the coefficient
criterion, weighted Gauss multiplication, linear division); the analytic
layer (canonical product, divisibility, exact evaluation quotient);
commutative algebra of the eventual-integrality sequence ring (idempotents,
annihilators, finite generators, localization); ordered groups (the
convex-subgroup prime correspondence and the interval
\eqref{ent:ps:eq:Cmin}--\eqref{ent:ps:eq:Cmax}); and ultraproducts, combining
these and specializing the explicit example. It is a suggestion, not code or
a claim about Mathlib coverage.
For \crefrange{ent:rc:subsec:setting}{ent:rc:subsec:limits} the seventh
source suggests a coefficient-controlled entire-series structure over a valued
field with the Hahn summation property, weighted Gauss multiplicativity,
finite-variable Taylor expansion and coarsening by a convex subgroup; the
simplex obstruction then reduces to matrices over a valuation ring, the
generator lower bound to finite-dimensional homogeneous polynomials and the
leading-order multiplication law, and only the omnific corollary needs the
normal-form bridge to actual surreal numbers. Its finite regression program
cannot replace a proof of the support conditions or of the normal-form
embedding. This too is a suggestion, not code.
For \crefrange{ent:as:subsec:setting}{ent:as:subsec:limits} source~11
suggests the route recorded in \cref{ent:as:q:formalization}: the growth
criterion, scale polynomials, the degree-transfer lemma and finite polynomial
matrix lifting, then sampling and ideal membership, with a separate arithmetic
layer for constant extraction and finite differences, and certificates that
carry a proof that all remaining scale polynomials lie in their ideal. It
supplies no Lean code.

\section{Conclusion}
\label{ent:sec:conclusion}
The threshold for nonpolynomial entire behaviour in a Hahn--surcomplex
workspace is countable cofinality, not rank one. Among workspaces admitting
nonpolynomial entire series, the universal B\'ezout property and full
radially finite Hermite interpolation require the stronger condition that an
order unit exist. This is also the threshold for existence of a nonconstant
restricted unit; for a particular restricted series, its least nonconstant
coefficient gap relative to the constant term must itself be a positive
order unit. Individual coprime pairs may admit a B\'ezout identity even
when the universal property fails. The two existence conditions
separate in an explicit infinite-rank subgroup of the surreal numbers, and
the separation is witnessed by two entire functions with disjoint simple zero
sets whose ideal is proper.

After enlarging the value group, every nonpolynomial entire series has the
same exact strong domain: the valuation ring of the convex hull of its
original scales. Cofinality is precisely what preserves entireness, while
support preparation preserves every existing zero and prevents new ones
throughout the surviving domain. The canonical product theorem then shows
that coefficients, zeros and the extension boundary encode the same cofinal
scale structure. Together this gives one explanation of both nonpolynomial
fixed-workspace functions and polynomial rigidity on the whole surcomplex
class, without confusing Hahn identities with convergence in the fine surreal
topology.

Between these thresholds the obstruction is measured exactly. An infinite
Hermite prescription is realizable precisely when its shell data are
eventually integral for the coarsening that kills the Archimedean scale of
each shell, and two independent constructions, sequential and cardinal,
realize it. Above a product with one node per shell, the maximal ideals are
the evaluations and one ideal for each nonprincipal ultrafilter on the shells,
with algebraically closed residue fields assembled from lower-rank Hahn
fields. In finitely many variables the uniform extension domain gives way to
a well-ordering condition on the quotient weights of the ordinary support,
while for nonpolynomial series homogeneous-block summation along complex
directions sees exactly the proper countable unions of projective algebraic
sets, all of which occur when the value group has countable cofinality.

Below the maximal ideals, a product with one simple node per shell carries
two kinds of boundary data. Its idempotents record which infinite
collections of shells an ultrafilter selects, and produce the Stone space
$\beta\N_{>0}$ common to its minimal and maximal primes; its ultraproduct
value groups record how small the other functions become along those
shells, and produce the prime chains inside each fibre. With an order unit
the second layer disappears from the quotient; without one it contains a
continuum chain in every nonprincipal fibre. Under a cofinal enlargement the
nodes and the Stone base stay fixed while the second layer can refine, and
the rational-indexed enlargement of $\Gamma_\infty$ realizes this inside
surreal and surcomplex Hahn workspaces: preserving zeros is not the same as
preserving the prime spectrum of an entire divisor.

In several variables the order unit returns as the threshold for a geometric
property. With countable cofinality, it decides whether every radially finite
configuration can be straightened by an entire automorphism. Without it, the
eventual polynomial reductions of any one automorphism preserve sufficiently
fine full-dimensional simplices at escaping scales, yet polynomial-part
coordinates, and so the omnific configurations of a workspace, exclude the
denominator losses behind that obstruction and admit an explicit
rectification. Once a configuration is straightened, total multiplicity
conditions have an exact global description: bounded multiplicities give a
finite minimal generating family indexed by a simplex of transverse monomials,
and unbounded ones force non-finite generation in dimension at least two. All
of this stays inside one fixed Hahn workspace, not the full surreal or
surcomplex class.

At ordinary points of $\kk^d$ the picture is algebraic, for every value group.
At each Hahn exponent an entire series is an ordinary polynomial, and
bounded-degree lifting makes the entire ring faithfully flat over the
coefficient polynomial ring, with exact coefficientwise membership in every
ideal defined over the coefficient field. So arithmetic values on a
Zariski-dense set of ordinary points force a polynomial, the entire maps
preserving the omnific integers of a workspace are integer-valued polynomials
plus polynomials with purely infinite coefficients, and the ordinary
exceptional locus is cut out by finitely many scale polynomials, which also
control every jet order; every nonzero ideal occurs once the value group has
countable cofinality. The finite certificate exists but cannot in general be
found: an explicit family with quadratic growth encodes nonhalting in the
omnific value at one ordinary point.

\appendix
\section{Proof-dependency and hypothesis audit}
\label{ent:app:audit}

\subsection{The dependency chain}
The main proof can be read along the following chain.
\[
 \begin{gathered}
 \text{Hahn support calculus}\\
 \Downarrow\\[-2pt]
 \text{cofinal coefficient criterion and the algebra }\A\\
 \Downarrow\\[-2pt]
 \text{support preparation and finite root counts}\\
 \Downarrow\\[-2pt]
 \text{canonical products, divisibility, and GCDs}\\
 \Downarrow\\[-2pt]
 \begin{array}{c@{\qquad\qquad}c}
 \text{new Archimedean scales}&\text{one order unit}\\
 \Downarrow&\Downarrow\\
 \text{growth barrier and no B\'ezout}
 &\text{Hermite interpolation and B\'ezout}
 \end{array}
 \end{gathered}
\]
The arrow from preparation to root counting additionally uses algebraic
closedness of the Hahn field. The arrow to interpolation uses completeness of
the real-valued coarsening, proved directly in \cref{ent:lem:coarsening}. No
proof of the negative result depends on the positive interpolation theorem.

The scalar-extension results of \cref{ent:sec:extension} sit on a deliberately
separate branch. The restricted-unit theorem
(\cref{ent:thm:units}) uses Gauss multiplicativity
\eqref{ent:eq:gauss-mult} and a convex coarsening only; it uses neither zero
counting nor canonical products. The universal-domain theorem
(\cref{ent:thm:extension}) uses the global criterion and integral evaluation;
it does not use factorization. \emph{Only} zero conservation
(\cref{ent:thm:zero-conservation}) imports both algebraic closedness and
preparation. The divisor classification is downstream of the local tools and
the elementary entire algebra. Thus the principal classifications are not
consequences of an unproved global factorization assertion, and conversely
the factorization proof never assumes that an arbitrary zero-free restricted
series is restricted-invertible --- by
\cref{ent:ex:ranktwo-unit} that is exactly the assertion that can fail in
higher rank.

The fourth source's cardinal route to the image theorem
(\cref{ent:sm:subsec:cardinal}) uses the canonical product
(\cref{ent:thm:product}), Gauss multiplicativity including its equality
case, the integral inverse of the cofactor class and the summation lemma
\cref{ent:sm:lem:entire-sum}. Like the first route it uses neither
preparation nor algebraic closedness, and it does not use the first route.
The fourth source's ideal-theoretic results
(\cref{ent:sm:thm:finite-generators,ent:sm:cor:membership,ent:sm:thm:paired}
and \cref{ent:sm:subsec:all-maximal}) use the quotient isomorphism, and hence
either route together with \cref{ent:prop:kernel-direct}; algebraic
closedness enters only the statement about the type of the residue fields,
\cref{ent:sm:cor:boundary-fields}, through the Hahn-field theorem for
$\C((t^{H_n}))$.

\Cref{ent:sec:multivariable} sits on a third separate branch. It uses the
Hahn--Neumann support calculus, the finite lattice estimate, unit twisting,
the Hilbert basis theorem and countable-field avoidance. It uses no
divisibility, preparation, root count, factorization or algebraic closedness.
Of the other sections it uses only \cref{ent:lem:neumann}, the definitions
\eqref{ent:eq:convex-hull} and \eqref{ent:eq:D} and the normal-form map
\eqref{ent:eq:embedding}, none of which needs them;
\cref{ent:thm:multi-unit} alone keeps the report's standing hypotheses and
cites \cref{ent:thm:units}.

\Cref{ent:ps:sec:spectrum} sits downstream of the quotient isomorphism. Of
the analytic theory it uses \cref{ent:thm:quotient} in the form
\cref{ent:cor:single-node}, hence either route to the image together with
\cref{ent:prop:kernel-direct}, and, for the canonical product with
multiplicities, \cref{ent:thm:product}; for the comparison of two workspaces
it uses \cref{ent:thm:criterion} and \cref{ent:prop:kernel-direct} over the
larger group. It uses no preparation, root count, factorization or algebraic
closedness, and it does not use zero conservation: the zeros of its product
in the larger field are located directly (\cref{ent:ps:lem:comparison}).
The rest is commutative algebra of the sequence ring, ultraproducts and the
valuation-domain prime correspondence. Algebraic closedness enters only
through the type of the residue fields $L_\Ucal$, quoted from
\cref{ent:sm:cor:boundary-fields}.

The seventh source's subsections
\crefrange{ent:rc:subsec:setting}{ent:rc:subsec:limits} form a fourth branch,
inside \cref{ent:sec:multivariable} but under the standing divisibility and
with $\kk\in\{\R,\C\}$; the sentence above that the section uses no
divisibility refers to its first eight subsections. The source records three
chains. The positive one is the coefficient criterion, all-disk summation,
cardinal interpolation with the denominator bound, and a projection followed
by a shear; through \cref{ent:rc:lem:interpolation} it uses
\cref{ent:thm:image} (first proof) or \cref{ent:thm:product} and
\cref{ent:rc:lem:sum} (second proof). The negative one is the coefficient
criterion, eventual reduction of a map and its inverse, an integral Taylor
remainder and a nonzero simplex determinant; it uses neither the image theorem
nor any classification of ideals or automorphisms. The fat-point chain is
division by escaping products (\cref{ent:prop:kernel-direct}), transverse
Taylor expansion, finite transverse-degree truncation and the dimension of a
space of homogeneous jets; its lower bound assumes no Noetherianity, coherence
or Nakayama lemma. None of the three uses preparation, root counts or
algebraic closedness. The source lists the failure modes its proofs guard
against: increasing valuations confused with cofinal ones; an infinite product
inverted diskwise rather than through quotients by finite products; a generic
projection that cancels leading terms of points or of differences; an integral
reduction assumed polynomial without a coefficient tail beyond the convex
subgroup; a Jacobian called a coarsened unit without the inverse map in the
finite family; infinite analytic sums used as algebraic generation; and a
generator lower bound not tied to an explicit finite-dimensional jet quotient
at a node of maximal multiplicity.

Source~11's subsections
\crefrange{ent:as:subsec:setting}{ent:as:subsec:limits} form a fifth branch,
inside \cref{ent:sec:multivariable}, under \cref{ent:as:conv:standing}. Its
chain is recorded in \cref{ent:as:subsec:limits}: the coefficient criterion
\cref{ent:thm:multi-criterion}, scale polynomials, the degree-transfer lemma
and the Gr\"obner degree bound, then the image theorem, flatness, and the
arithmetic and ideal classifications; the jet results add the Hasse product
rule, the realization theorem the damping construction, and the computability
theorem the halting problem. It uses no divisibility, preparation, root count,
canonical product or algebraic closedness; it imports the Hilbert basis
theorem and the equational criterion of flatness; of the other sections it uses
only \cref{ent:mv:def:entire,ent:mv:lem:tails,ent:thm:multi-criterion}; and the
merge's \cref{ent:as:cor:vanishing}(b) adds Hilbert's Nullstellensatz. The
failure modes its proofs guard against are listed there: regrouping before
summability, a degree bound that depends on the scale, a leading-valuation
test in place of the full positive tail, confusion of set and class scopes,
and a Noetherian finiteness read as a stopping rule.

\subsection{Where an attained minimum is load-bearing}
\paragraph{Attained minima.}
The proof of \cref{ent:thm:criterion} uses a minimum attained at an ordinary
\emph{finite} coefficient index, not a termwise inequality. The proof of
\cref{ent:lem:growth-barrier} uses an attained minimum in the same way:
termwise strict inequalities alone would not justify a strict inequality for
an arbitrary limiting process. These are the two places in the article where
the distinction is load-bearing, and neither can be replaced by an estimate
valid only term by term.

\paragraph{Two independent infinities.}
The finite index $k$ in a power series is not a Hahn exponent. Integer
multiples of a positive exponent need not be cofinal. The support lemma
\cref{ent:lem:neumann}(b) handles geometric-series expansions without that
property; the positive interpolation proof uses it only after constructing a
legitimate real-valued coarsening, which is available only under the
order-unit hypothesis.

\paragraph{Divisibility is not assumed in its own proof.}
The quotient $F/G$ in \cref{ent:thm:factorization} is first formed, after
removing the shared zero at the origin, in $\K[[Z]]$. Only preparation at
every scale proves it entire. This avoids assuming the global divisibility
theorem in its own proof.

\paragraph{The counterexample is irreducibly infinite.}
No finite truncation of the pair \eqref{ent:eq:paired-products} can witness
non-B\'ezout behaviour: the finite truncations are polynomials with disjoint
root sets and therefore generate the unit ideal. The contradiction requires a
single candidate coefficient lower bound to fail at cofinally many
successively dominating scales.

\paragraph{Why preparation must be a support recursion.}
\Cref{ent:cor:prep-extension} --- the fact that the prepared factorization is
unchanged by enlarging the value group --- is available precisely because
preparation was proved by well-founded recursion over $S^*$. An iteration
whose convergence had been asserted in $\K$ would have to have that
convergence re-established in $\KD$, and \cref{ent:ex:ranktwo-unit} shows it
can genuinely fail.

\section{Reproducible finite checks}
\label{ent:app:checks}
Eight independent exact-arithmetic suites accompany this report, one from
each source manuscript, preserved unmodified in \texttt{code/} and
\texttt{data/} under the file names given below. The first six use only the
Python standard library, with exact integer and \texttt{fractions.Fraction}
arithmetic and no floating point; the seventh and eighth use exact rational
and symbolic arithmetic in SymPy. The first four represent group elements as
tuples of rational coordinates ordered by the largest nonzero coordinate;
the fifth needs no ordered-group tuples beyond integer relation vectors; the
sixth uses finitely supported rational vectors with rational indices, ordered
in the same way.

\paragraph{Suite A (831 checks).}
The source program \texttt{verify\_examples.py} is shipped as
\begin{quote}\small
\path{code/06-cofinality-bezout-trichotomy-verify-examples.py}
\end{quote}
It checks finite instances of the
reverse-lexicographic exponent inequalities, the valuation formula
\eqref{ent:eq:lambda} in the first nine coordinates of $\Gamma_\infty$,
leading exponents of finite canonical products, a finite-truncation
implementation of the preparation recursion
\eqref{ent:eq:prep-rec1}--\eqref{ent:eq:prep-rec2}, and the local inverse-jet
identity underlying \eqref{ent:eq:interpolation-r}. The delivered run passed
all \textbf{831} checks, with category counts $72$ paired-root checks, $40$
finite dominance samples, $88$ product-coefficient checks, $456$ preparation
checks and $175$ inverse-jet checks; the record is the delivered
verification JSON. The script takes a \emph{required} output path argument
and exits nonzero without one; a fresh output path preserves the historical
record. A review run on a temporary copy passed $831$ of $831$ again.

\paragraph{Suite B (664 checks).}
The source program \texttt{verify.py} is shipped as
\begin{quote}\small
\path{code/03-cofinality-and-scalar-extension-verify.py}
\end{quote}
It checks finite truncations of
\eqref{ent:eq:infinite-product} for $N=1,\ldots,8$ --- each coefficient's
leading exponent, its binomial subset count and its signs --- exactness and
simplicity of each root $t^{-e_j}$ and the simple-root derivative
\eqref{ent:eq:product-derivative}, closed- and open-ball Newton counts at
sampled scales including half-shifts by $e_0$, descent of the valuations
after adjoining the dominant scale $e_9$, the rank-two hidden-tail comparison
of \cref{ent:subsec:hidden-tails} for $n=1,\ldots,100$ and
$\lambda\in\{-7,-\tfrac12,0,\tfrac32,10\}$, and the preparation recursion for
\eqref{ent:eq:cubic} through perturbation order eight. The delivered run
reports \textbf{664} passing checks. A maintained review on a temporary copy
passed all $664$ checks, leaving historical records unchanged. The delivered
results are recorded in
\texttt{data/03-cofinality-and-scalar-extension-verification.json} and
\texttt{.txt}. The script takes no arguments and writes
\texttt{verification.json} and \texttt{verification.txt} into the
\texttt{data/} directory next to its own folder, so it should be run on a
copy of the directory.

\paragraph{Suite C (1958 checks).}
\texttt{code/01-exact-jet-image-verify.py} accompanies the third source and
the material of \cref{ent:subsec:shells,ent:subsec:jet-image,ent:subsec:image-proof}.
It checks, over $24$ finite polynomial cases, $5$ damping exponents and rank
vectors of dimension $12$, the finite Hermite jet identities and degree
bounds, preservation of earlier jets under a shell correction, the shell
correction congruence and its low-degree vanishing, the finite quotient
inverse and the multiplication-determinant formula, the paired-root
valuation formula and its coarse maximal ideal, and the close-pair corrected
value and forbidden coefficient. The delivered record
\texttt{data/01-exact-jet-image-checks.json} reports \textbf{1958} passing
checks in $17$ categories, with random seed $20260922$. The script takes a
required output path and refuses to overwrite an existing file. Its own
record lists what it does not verify: infinite strong summability or
cofinality, the existence or nonexistence of an infinite entire interpolant,
the infinite quotient and ultrafilter theorems, and independent refereeing or
Lean verification.

\paragraph{Suite D (1876 checks).}
\texttt{code/07-scale-moderate-interpolation-verify.py} accompanies the fourth
source, chiefly \cref{ent:sm:subsec:cardinal} and the paired threshold. Over
$18$ random cases with three shells each, at the rational scales $2$, $17$
and $103$ and with near-coincident rational roots $1$ and $\tfrac{102}{101}$
of multiplicity one or two, it checks exact division of the product by each
shell factor ($54$), the cofactor inverse and the inverse of $T$ in the finite
algebra ($54$ each), the cardinal shell jets --- the prescribed class on its
own shell and zero on the others ($162$) --- the multiplication-determinant
formula \eqref{ent:eq:detnorm} ($54$), and simultaneous interpolation by the
sum of the cardinal corrections ($54$), together with the two values of the
close-pair polynomial ($2$). In nine-dimensional rational vectors ordered by
the largest coordinate it checks the scale condition $\gamma\ge-\gamma_n/2$ and
the suppression inequality behind \eqref{ent:sm:eq:d-choice} ($700$ each),
that the paired value \eqref{ent:sm:eq:paired-value} lies in $H_n$ exactly
when the displacement does ($21$) and is positive when it does not ($14$),
and that the gain $\varkappa_n$ is nonnegative ($7$). Its random seeds are
$22092026$ and $20260922$. The delivered record
\texttt{data/07-scale-moderate-interpolation-verification.json} reports
\textbf{1876} passing assertions; it lists as not verified infinite Hahn
summability, the full theorems, novelty and Lean formalization. The file
\texttt{data/07-scale-moderate-interpolation-package\_manifest.json} records
the source package's pin, page count, assertion count, the two verification
flags, and SHA-256 hashes of the six files of that package; of those, the
script, the record and the source audit (shipped as
\texttt{07-scale-moderate-interpolation-SOURCE\_AUDIT.md}) are in this
directory and match their hashes, while the manuscript's own
\texttt{article.tex}, \texttt{article.pdf} and \texttt{README.md} are not
shipped. The \texttt{--output} argument is optional; without it the script
only prints. Its documented command writes to \texttt{data/verification.json},
so a re-run should name a scratch path instead.

\paragraph{Suite E (1072 checks).}
\texttt{code/08-multivariate-extension-verify\_examples.py} accompanies the
fifth source and \cref{ent:sec:multivariable}. With bivariate polynomials over
$\Q$ it checks the homogeneous products of the explicit direction example and
their value at infinity ($20$), the prescribed polynomial directions and the
nonvanishing products off them ($210$), binomial block cancellation in
$(X-Y)^n$ ($80$), the quadratic weight-difference identity and eventual-sign
certificates of \cref{ent:mv:thm:nonclosed} ($100$ and $25$), samples
approaching its nonclosed boundary ($60$), semilinear period-weight identities
($225$), a two-generator relation lattice and its residual bound ($243$),
coordinate-face realization from minimal excluded faces ($57$), and
degree-separated projective block zero loci ($52$). The delivered record
\texttt{data/08-multivariate-extension-verification.json} reports
\textbf{1072} passing assertions and sets \texttt{formal\_verification},
\texttt{independent\_peer\_review} and \texttt{priority\_certified} to false.
\texttt{data/08-multivariate-extension-research-audit.json} records the pin,
the blob of this article at that pin, the answered clause, the items not
claimed solved and the audit scope, and
\texttt{data/08-multivariate-extension-build.json} records the source's own
build, including SHA-256 hashes of its \texttt{article.tex} and
\texttt{article.pdf}, which are not shipped here. The \texttt{--output}
argument is optional; the documented command replaces
\texttt{data/verification.json}, so a re-run should name a scratch path.

For this merge both new suites were re-run with \texttt{--output} pointing
outside the repository; each reproduced its delivered record exactly, up to
platform line endings.

\paragraph{Suite F (10889 checks).}
\texttt{code/09-prime-spectra-verify\_finite.py} accompanies the sixth source
and \cref{ent:ps:sec:spectrum}. It checks the node values of the
simple-node Lagrange corrections of \cref{ent:ps:rem:route} at the rational
nodes $2,4,\ldots,2^n$ ($153$; the source calls them cardinal corrections,
which are not the entire cardinal corrections of
\cref{ent:sm:subsec:cardinal}) and their damping inequalities in rational
vectors ($390$); the
intersection of old integer scales with new rational cutoffs, at an
intermediate subgroup and at the maximal endpoint ($3600$ each), with $192$
strict-containment witnesses, in finite-support rational vectors; $30$
representative power-gap inequalities $mn^{\theta}<n^{\theta'}$ at the single
value $n=2^{60}$, with $\theta,\theta'\in\{\tfrac15,\ldots,\tfrac45\}$, behind
\cref{ent:ps:lem:gaps}; $200$ Smith diagonalizations of random matrices of
sizes up to $5\times5$ over the valuation ring $\Z_{(2)}$, checking the
diagonal shape ($200$), the identity $UMV=\operatorname{diag}$ ($200$), the
integrality of the transformations and their inverses ($800$), the pivot
ratios ($1190$) and the divisibility chain ($240$); and the finite
multiplicity tests behind \cref{ent:ps:thm:multiplicity} ($196$ uniform-power
and $98$ radical-order checks). Its random seeds are $23092026$, $10050$ and
$31$. The delivered record \texttt{data/09-prime-spectra-verification\_results.json}
reports \textbf{10889} passing assertions in $13$ categories and states that
the checks are finite sanity checks, not proofs of infinite or ultrafilter
theorems; \texttt{data/09-prime-spectra-BUILD\_REPORT.json} records the
source's own build of its $28$-page article, which is not shipped. The
\texttt{--json} argument is optional; without it the script only prints, and
the documented command rewrites the delivered record in place, so a re-run
should name a scratch path. For this merge it was re-run with \texttt{--json}
pointing outside the repository and reproduced the delivered record exactly,
up to platform line endings.

\paragraph{Suite G (five groups of symbolic checks).}
\texttt{code/10-entire-rectification-verify.py} accompanies the seventh source
and \crefrange{ent:rc:subsec:setting}{ent:rc:subsec:limits}; it needs Python
3.10 or later and SymPy, pinned in
\texttt{data/10-entire-rectification-requirements.txt} to version $1.14.0$.
It checks the padded cardinal series \eqref{ent:rc:eq:cardinal} at four
rational nodes with padding degrees $2,4,6,8$ ($4$ node values); the
coefficient formula for division by one linear factor; for two triangular
polynomial automorphisms, in dimensions two and three, that the Jacobian
determinant is $1$ and that the normalized simplex determinant of
\eqref{ent:rc:eq:simplexdet} is $1$ plus a polynomial multiple of
$\varepsilon$; the monomial count behind
\cref{ent:rc:thm:fat-generators} in $40$ cases ($d=2,\ldots,6$,
$\bar m=1,\ldots,8$); and two finite polynomial fat-point ideals over $\Q$, in
$d=2$ with nodes $1,2,4$ and multiplicities $1,3,2$ and in $d=3$ with nodes
$1,2$ and multiplicities $1,2$, where every proposed generator satisfies every
prescribed jet ($40$ and $30$ equalities), the finite jet map has the expected
rank ($10$ and $5$), and a Gr\"obner basis of the generators reduces every
kernel basis vector to zero ($35$ and $51$). The record
\texttt{data/10-entire-rectification-verification-results.json} reports every
group passed and states its scope as finite symbolic sanity checks, not formal
verification. \texttt{data/10-entire-rectification-BUILD\_REPORT.json} records
the source's own $25$-page build, its pin and the flags for formal
verification and peer review, both false; that article is not shipped. Without
\texttt{--output}, the script writes \texttt{verification-results.json} next
to itself; the source's documented command therefore rewrites its delivered
record in place, and in this directory it would add a new file to
\texttt{code/}. A re-run should name a scratch path, on a copy. For this merge
it was run on a copy with \texttt{--output} outside the repository, on Python
$3.14.4$ with SymPy $1.14.0$; its output is identical to the delivered record
except the recorded Python version ($3.13.5$ there).

\paragraph{Suite H (8774 checks).}
\path{code/11-arithmetic-sampling-verify.py} accompanies source~11 and
\crefrange{ent:as:subsec:setting}{ent:as:subsec:limits}. It needs Python 3.9
or later and SymPy, pinned in \path{data/11-arithmetic-sampling-requirements.txt}
to version $1.14.0$, and uses exact rational and symbolic polynomial
arithmetic with random seed $20260923$. In ten groups it checks: scale
extraction commuting with evaluation for random finite Hahn polynomials in two
variables at five rational points, with the full positive-tail and residue
tests ($450$); the constant-term pullback, closure of the finite model of
$\OzK$ under sums, products and binomial coefficients, $t^{-1}+t\notin\OzK$ and
$\binom i2=(-1-i)/2$ ($1193$); multivariate Newton coefficients recovered by
finite differences on a $5\times5$ grid ($1000$); graded Gr\"obner division by
$(Y-X^2,X^3-X)$ with degree control, idempotent normal forms and scalewise
reduction ($720$); an explicit polynomial syzygy of a $2\times3$ matrix, scale
by scale ($521$); the sharp actual-scale gcd certificates for $n_0\le6$ and a
greedy-reduction count ($421$); Hasse-derivative gcds and contact
multiplicities for jet orders below $7$, and the Hasse product rule ($260$);
the realization ideal with scales $n(n+4)$ and its growth ($87$); the machine
family with its first $49$ complete base blocks, in the nonhalting case and for
every halting step below $32$, with the quadratic modulus ($3818$); and the
boundary examples: the descending and ascending exponent sequences of
\eqref{ent:as:eq:local-counterexample}, finite truncations of
\eqref{ent:as:eq:qproduct} with their zeros $t^{-j}$, and finite samples
($304$). The delivered record \path{data/11-arithmetic-sampling-verification.json}
reports \textbf{8774} passing assertions and states its scope as exact finite
algebraic checks, not a proof-assistant certificate or a proof of infinite
theorems. The \texttt{--output} argument is optional; without it the script
only prints, and the source's documented command
(\texttt{--output verification.json}) rewrites the delivered record in place,
so a re-run should name a scratch path, on a copy.
\path{code/11-arithmetic-sampling-build.py} builds the source's own article,
which is not shipped, in a temporary directory, copies the PDF back and writes
\texttt{build\_report.json} next to it; here it stops, the article being absent.
\path{data/11-arithmetic-sampling-build_report.json} records that build ($24$
pages, three pdfLaTeX passes, no diagnostics) with SHA-256 hashes of the
unshipped source and PDF, and
\path{data/11-arithmetic-sampling-ARTIFACT_MANIFEST.json} records the pin,
the page count, the assertion count, the flags for proof-assistant
verification and referee review (both false) and SHA-256 hashes of nine files:
the six shipped here (the source audit, both scripts, the build report, the
requirements and the verification record) match, and the manuscript, its PDF
and its delivery README are not shipped. For this merge the script was run on
a copy with \texttt{--output} outside the repository, on Python $3.14.4$ with
SymPy $1.14.0$; its output is identical to the delivered record except the
recorded Python version ($3.13.5$ there) and platform line endings.

\paragraph{What the checks cannot do.}
All eight suites print or record their own scope disclaimers, and all eight
are correct to do so.
The counts are counts of finite assertions, not confidence scores for any
infinite theorem. No suite can prove cofinality statements, Hahn support
well-ordering for every permitted group, any infinite product identity, the
necessity direction of \cref{ent:thm:units}, or the nonexistence of a global
B\'ezout solution. In the preparation checks the $t$-adic truncation is
explicitly finite, so higher coefficients are discarded only after the
claimed congruence has been specified, and the exponent subgroup used is an
integer one rather than an arbitrary-rank group. The finite dominance samples
test only $k\in\{1,2,10,100,10000\}$. Most importantly, the finite
truncations of \eqref{ent:eq:paired-products} \emph{are} coprime polynomials
and \emph{do} satisfy polynomial B\'ezout identities, so no finite
computation can witness \cref{ent:thm:no-bezout}. What the checks can detect
is an indexing error, a sign error in a paired root valuation, a reversed
order convention, a wrong root-count endpoint or a wrong finite preparation
coefficient. Everything else rests on the proofs in the body.

\section{Source and novelty ledger}
\label{ent:app:sources}
The first two of this report's eight source manuscripts are both dated
21 September 2026, and each pinned the repository at its own snapshot: one at commit
\texttt{4896a2808ce30e01b1c8\allowbreak 6ae3ba2295a64762d246} and the other at
\texttt{aa846271b4dcae2c055b\allowbreak 216126a87210292ec19b}. The first re-read four
files at its pinned commit --- the root README, the documentation catalogue,
and the READMEs of the rank-one and global-divisor reports --- after its
working branch changed during preparation; the second cited the documentation
catalogue and the rank-one article. Both reviews were targeted reads, not
line-by-line proof audits or independent builds. Neither reinterprets missing
Lean implementation as an unsolved mathematical problem.

The primary literature checked across the two manuscripts comprises Poonen's
Hahn-field construction and algebraic-closedness corollary, Neumann's ordered
division-ring paper for the support calculus, Cherry's complete real-valued
non-Archimedean setting together with his one-variable factorization
statement and his lectures, Conrad's higher-rank valuation examples, Lazard's
bibliographic record as the historical reference for one-variable zeros, van
der Hoeven's generalized-series operator framework, and the Conway--Gonshor
surreal normal-form background with cross-checks against
\cite{ADH,BournezGuilmant}. Lazard was checked bibliographically only, and
Conway and Gonshor were not re-audited in full. General web searches were
attempted; several returned irrelevant results and were not treated as
evidence that a theorem is absent from the literature.

A third manuscript, dated 22 September 2026, was written later against the
merged report itself, at the snapshot
\texttt{a3124af79f66b8b9c196\allowbreak d76b4cbc5ac3938907c4}. Its decisive comparison
was that snapshot's subsection recording the missing image of the infinite
jet map, now answered by \cref{ent:thm:main-image}. It was a targeted source
inspection through repository file reads, not a local clone, a line-by-line
audit of every report, or a Lean build. Its literature comparison covered
Lazard, Cherry and van der Hoeven and located no matching statement of the
shellwise restricted-product theorem; that is a reason to propose it as new,
not a proof that it occurs nowhere. The ultrafilter construction is standard
algebra; what is proposed there is the concrete analytic surjection and the
explicitly identified residue fields. It credits the coefficient criterion,
the canonical product, the order-unit interpolation theorem and the explicit
non-B\'ezout pair to the first two sources and reproves them only for
dependency control.

A fourth manuscript, \emph{Scale-Moderate Interpolation and Boundary Ideals
in Surcomplex Hahn Fields}, dated 22 September 2026, pinned the repository
at \texttt{0097304c7ae9d5de46d2\allowbreak ea342e2494f8fc126303}. At that snapshot the
merged report still contained the subsection ``The missing image of the
infinite jet map'' and recorded \cref{ent:q:image} as open, so the
manuscript's statement that it resolves ``the repository's explicitly
unidentified infinite jet image'' was accurate for its pin. It is stale for
this tree: the third source, pinned earlier at \texttt{a3124af}, was merged
at commit \texttt{de84d8b} about an hour after \texttt{0097304} and had
already answered the question. The fourth source is therefore recorded as an
independent, near-simultaneous second proof of \cref{ent:thm:image},
\cref{ent:thm:quotient} and \cref{ent:thm:bezout-criterion}, under the same
hypotheses ($\Gamma$ divisible and set-sized, coefficients $\C$ or $\R$), and
not as their source. Its review read the documentation index, the
formalization ledger and the relevant parts of this report through a
connector at the pin; it was targeted, not an audit of every theorem. Its
literature comparison covered Poonen, both Cherry references, Helmer
(metadata only, a historical antecedent and not a proof input), Bruno's
dissertation (whose value group is explicitly a subgroup of $\R$), Henriksen
and Gonshor. Where the two texts prove the same statement it is printed once:
its coefficient criterion and Gauss estimates, canonical product, kernel
theorem, exact sequence, one-node criterion, close-pair example, factorial
example, order-unit check and unit test for shell algebras are
\cref{ent:thm:criterion}, \cref{ent:prop:operations} with
\eqref{ent:eq:gauss-mult}, \cref{ent:thm:product} with
\eqref{ent:eq:productcoef}, \cref{ent:prop:kernel-direct},
\cref{ent:thm:quotient}, \cref{ent:cor:single-node},
\cref{ent:thm:close-pair}, \cref{ent:cor:no-interpolation},
\cref{ent:cor:order-unit-image} and \cref{ent:lem:shellunits}. Its
sufficiency construction is genuinely different and is kept as the second
route of \cref{ent:sm:subsec:cardinal}; its finite-prefix corollary is the
completion step there. Its new material is
\cref{ent:sm:cor:no-order-unit,ent:sm:lem:entire-sum,ent:sm:lem:inverse-shell,ent:sm:lem:cardinal,ent:sm:cor:error},
\cref{ent:sm:rem:simple-cardinal},
\cref{ent:sm:thm:finite-generators,ent:sm:cor:membership,ent:sm:thm:paired},
all of \cref{ent:sm:subsec:all-maximal}, and \cref{ent:sm:rem:real}. Where
the merge had to choose: its letters were translated into this report's
(\cref{ent:subsec:notation}); its surreal instance was restated in
$\Gamma_\infty$, which contains $e_0$, instead of its $\Gamma_*$; and its
finite-generator and membership results are stated, as in the source, for
divisors with empty finite part.

A fifth manuscript, \emph{Multivariate Entire Hahn Series across New Surreal
Scales}, dated 22 September 2026, pinned the repository at
\texttt{4cf691c7d951e037739d\allowbreak 32d9f5c387dcce724f3c}, after the third source
had been merged, and recorded the blob \texttt{f034bdb} of this article at
that pin; the body it read is unchanged in the present base, which differs
only in the ledger and checks appendices. It read this report's
introduction, extension section, questions, scope discussion and
finite-variable appendix, together with the documentation inventory and some
neighbouring summaries; binary archives, retired manuscripts and the Lean
modules were not audited. Its claims about this report --- that
\cref{ent:q:several} asks for the several-variable replacement, that a
series independent of one variable cannot constrain it, and that the appendix
already proved the finite-variable coefficient criterion and cofinality
dichotomy --- were checked and hold; that appendix is now
\cref{ent:app:multivariate}. Its literature comparison covered van der
Hoeven's operators, Achinger's notes on rank-one non-Archimedean geometry
\cite{Achinger}, Joswig and Smith on convergent Hahn series and higher-rank
tropical geometry \cite{JoswigSmith}, Ma and Neelon, the Stacks Project for
the Hilbert basis theorem, and Berarducci for Conway normal forms. Its growth
criterion replaces the earlier appendix's ratio-form proof as the statement
of \cref{ent:thm:multi-criterion}, under weaker hypotheses; the earlier
\cref{ent:thm:multi-unit} is kept with its original hypotheses; everything
else in \cref{ent:sec:multivariable} is its new material. Where the merge had
to choose: its homogeneous block, admissible set, support, quotient, scale
and polynomial-direction symbols were renamed (\cref{ent:subsec:notation});
its standalone localization proposition is printed as
\cref{ent:mv:lem:localize} with credit to the foundations report; and none
of this report's divisible-group or algebraically closed results was
attached to its weaker hypotheses.

A sixth manuscript, \emph{Prime Spectra at Surreal Infinity: Reduced entire
divisors, valuation fibres, and splitting under cofinal scalar extension},
dated September 2026, pinned the repository at
\texttt{465a54b479a1ee842cbf\allowbreak 7db1689a7d2f6bfe25e1}, after the fifth
source had been merged; this report's directory is identical at that pin and
at the base of the sixth merge, and the seven labels its audit cites
(\texttt{ent:thm:image}, \texttt{ent:thm:quotient},
\texttt{ent:cor:single-node}, \texttt{ent:thm:product},
\texttt{ent:thm:maximal}, \texttt{ent:sm:thm:maximal-all},
\texttt{ent:sm:thm:beta}) exist with the meaning it gives them. Its
statements about this report --- that it identifies the quotient by a
one-simple-node product with the eventually integral sequences, describes
all maximal ideals of that quotient with ultraproduct residue fields,
identifies the maximal spectrum with $\beta\N_{>0}$, proves that the full
entire ring fails to be B\'ezout without an order unit, and uses finite shell
algebras for several nodes on a shell --- were checked and hold; no stale
statement was found. Its review concentrated on this report, its README and
the documentation catalogue; it compared the literature of Finocchiaro,
Frisch and Windisch \cite{FFW}, the Stacks Project on valuation rings
\cite{StacksValuation}, Neumann, and Conway and Gonshor (the last two not
reviewed page by page). It is \cref{ent:ps:sec:spectrum}, all of which is
its material except where stated: its analytic section is printed once by
citation (\cref{ent:ps:rem:route}); its recovered maximal ideals, residue
fields and maximal spectrum are
\cref{ent:thm:maximal,ent:sm:thm:maximal-all,ent:sm:thm:beta}; its example
``rank is not the dividing invariant'' is \cref{ent:subsec:finite-rank} with
\cref{ent:subsec:gamma-infinity}; and its divisibility lemma for arbitrary
multiplicities is \cref{ent:thm:product} with \cref{ent:prop:kernel-direct}.
Its three questions are \cref{ent:ps:q:several-nodes,ent:ps:q:unbounded,ent:ps:q:invariance},
and its suggested Lean dependency order is summarized in
\cref{ent:subsec:questions}. Where the merge had to choose: its letters were
renamed (\cref{ent:subsec:notation}), including $e_q$ for its $f_q$; its
kernel argument was not kept as a separate route, being a bookkeeping
variant of \cref{ent:prop:kernel-direct}; and four statements are the
merge's, not the source's, and are marked where they occur --- that the
ring-theoretic results hold for any nodes on the same radii
(\cref{ent:ps:rem:nodes}), that $\Phi_\Ucal$ is $\red_\Ucal$ followed by
residues (after \cref{ent:ps:prop:fibre}), that finitely generated ideals containing the
product are $(P_D,G)$ (\cref{ent:ps:rem:fg-ideals}), and that the additional
primes for unbounded multiplicities are not maximal (the remark after
\cref{ent:ps:thm:multiplicity}). Its finite checks are Suite~F of
\cref{ent:app:checks}.

A seventh manuscript, \emph{Entire Automorphisms at Surreal Scales: An
order-unit dichotomy, omnific rectification, and exact fat-point ideals},
dated 23 September 2026, pinned the repository at
\texttt{71e9606d297c9b98c732\allowbreak 070dc462682cc6948981}, after the sixth
source had been merged; this report's directory and that of the
holonomic-rigidity report are identical at that pin and at the base of the
seventh merge. It read the root README, the report inventory, and this
report's README and relevant source passages. Its statements about this
report --- that it contains the all-scale coefficient criterion, canonical
products and division, the order-unit interpolation case, finer
restricted-product interpolation results, and several-variable coefficient and
scalar-extension results --- were checked and hold
(\cref{ent:thm:multi-criterion,ent:thm:product,ent:thm:factorization,ent:thm:interpolation,ent:thm:image}
and \cref{ent:sec:multivariable}). Its statement that a repository search for
``rectification'' returned no match was true at its pin; the word now occurs
only in its own placed files, and ``rectifiable'' elsewhere only for curves.
Its description of its problem as different from this report's
several-variable question is corrected by the placement: it partly answers
\cref{ent:q:several}. At its pin the collection had no omnific-integer
report; it now has two, recorded in \cref{ent:sec:collection}. Its literature
comparison covered Rosay and Rudin, Winkelmann, Cherry's lectures, Mantova and
Matusinski, and L'Innocente and Mantova; it searched combinations of
``Hahn'', ``entire automorphisms'', ``tame discrete sets'' and ``fat point'',
found the classical tameness literature and standard Hahn and omnific
background but not the precise classification, and calls its coverage
necessarily incomplete. For this merge the journal data of L'Innocente and
Mantova and the Winkelmann reference were not re-verified.

Its own contribution table, kept here in prose, reads as follows. The
coefficient criterion and escaping products with division are existing
mathematics of this report, reproved with explicit arbitrary-rank bounds, with
no novelty claim. Projection and triangular shear are a classical strategy;
the simultaneous valuation-preserving projection and its application are
proved there. The order-unit dichotomy is proposed as a new classification
for Hahn and surcomplex workspaces, its positive half using moderate
interpolation and its negative half coarsened polynomial reduction. The
eventual polynomial reductions of an automorphism and its inverse are a
proposed structural result, with the finite residue-polynomial argument and
the Jacobian-unit consequence proved directly. The escaping-simplex
obstruction is a proposed explicit construction. Polynomial-part and omnific
rectification without an order unit is a proposed application of the
denominator bound, concerning ambient coordinate changes, not arithmetic
automorphisms of omnific integers. The fat-point generators and sharp count
are a proposed arbitrary-rank global formulation whose local lower bound is
elementary commutative algebra, not a new method. The unbounded-multiplicity
radical distinction is a consequence in this ring, not a claim that such
non-Noetherian phenomena are new.

Where the two texts prove the same statement it is printed once: its
criterion, canonical products, linear division, escaping division, separated
scales and countability lemma are
\cref{ent:thm:multi-criterion,ent:thm:product,ent:lem:linear-division,ent:prop:kernel-direct,ent:lem:escaping-classes,ent:lem:divisor-countable},
and its interpolation theorem is a special case of \cref{ent:thm:image} and,
with order units, of \cref{ent:thm:interpolation}; its explicit series is kept
as the second proof of \cref{ent:rc:lem:interpolation}. Where the merge had
to choose: the material was placed as the then closing subsections of
\cref{ent:sec:multivariable} rather than as a new section, so that no existing
section, theorem, question or equation number changed, and its displayed
equations are numbered (R1), (R2), \ldots; its letters were renamed
(\cref{ent:subsec:notation}), above all its polynomial-part ring, its
denominators and its coordinate vectors; an auxiliary step in its proof of
\cref{ent:rc:thm:reduction}(c) was replaced by closure under strong sums; the
proof of \cref{ent:rc:cor:distinct-radii} makes explicit the collapse of the
points of one disk onto the shift radius, which the source's proof leaves
implicit; and its unnumbered corollary on cluster refinement is
\cref{ent:rc:cor:refinement}. The following are the merge's and are marked
where they occur: the first proof of \cref{ent:rc:lem:interpolation},
\cref{ent:rc:rem:real}, the credit of the antecedents of
\cref{ent:rc:thm:reduction}, the dependency remark after
\cref{ent:rc:prop:simplex}, the proof of \cref{ent:rc:cor:refinement},
\cref{ent:rc:rem:close-pair} and \cref{ent:rc:cor:evaluation}. Its finite
checks are Suite~G of \cref{ent:app:checks}.

An eighth manuscript, source~11 (batch~32, manuscript~09),
\emph{Arithmetic Sampling in Surreal Hahn Fields: Entire rigidity,
coefficientwise algebraic division, omnific exceptional loci, and a
computability barrier}, dated 23 September 2026, $24$ pages, pinned the
repository at \texttt{343dc2c471212bb9b53f\allowbreak f4623bace2e1943f255b},
after the seventh source had been merged; this report's directory is identical
at that pin and at the base of the eighth merge. It read the collection
through a connector (a local clone failed, so nothing was built), inspecting
the root and documentation guides, the initial portion of this report's README
and an article passage from its line~330, the initial portion of the
set-sized-quotients README, the later portion of the
omnific-preserving-automorphisms README and that report's source audits 12 and
13; several responses were truncated. Its statements about this report ---
that it contains the coefficient criterion, the cofinality obstruction, the
scalar-extension distinctions and escaping-node interpolation --- were checked
and hold
(\cref{ent:thm:multi-criterion,ent:cor:cofinality,ent:thm:extension-criterion,ent:thm:image}).
It noted that older merged guides can carry superseded open-question wording
and relied on the later companion audits rather than on such wording. Its
literature comparison covered L'Innocente and Mantova, for the distinction
between the class and set-sized fraction fields, Kuhlmann and Serra, for strong
summability, the Stacks Project, for the Hilbert basis theorem and the flatness
criteria, and the abstract of Garcia-Fritz and Pasten, with a Wikipedia
overview as orientation only; it calls the search limited. Where the two texts
prove the same statement it is printed once: its locally finite families
lemma, coefficient criterion, cofinality dichotomy and whole-class corollary
are \cref{ent:mv:lem:tails}(a), \cref{ent:thm:multi-criterion},
\cref{ent:cor:cofinality} with \cref{ent:mv:prop:realize-support}, and
\cref{ent:cor:all-no,ent:mv:prop:allscale}. Everything else in
\crefrange{ent:as:subsec:setting}{ent:as:subsec:limits} is its material. Where
the merge had to choose: the material was placed as the last subsections of
\cref{ent:sec:multivariable}, after the seventh source's, so that no existing
section, theorem, question or equation number changed; its displayed equations
are numbered (S1), (S2), \ldots; its letters were renamed
(\cref{ent:subsec:notation}), above all $D$, $A$, $B$, $\Pi$, $\rho$, $\gamma$
and $H$; its ten questions were moved to \cref{ent:subsec:questions}; and its
source-comparison section is summarized in \cref{ent:as:subsec:limits}. The
following are the merge's and are marked where they occur: the identification
of $\pg_\eta$ with the coefficients of \eqref{ent:eq:Halg} and of
\cref{ent:mv:thm:exceptional}, the remark after \cref{ent:as:cor:truncation},
the comparison after \cref{ent:as:cor:fiber}, \cref{ent:as:cor:vanishing}, the
comparisons with the set-sized-quotients report after \cref{ent:as:thm:integer}
and \cref{ent:as:cor:binomial}, the parallel with the fifth source at the end of
\cref{ent:as:subsec:tail}, the last paragraph of
\cref{ent:as:subsec:computability}, the remark after \cref{ent:as:q:matrices},
and the notes after \cref{ent:q:several,ent:q:image-remaining,ent:mv:q:effective}.
Its finite checks are Suite~H of \cref{ent:app:checks}.

The strongest proposed original statements of the first two sources are
\cref{ent:thm:no-bezout} and its use in \cref{ent:thm:main}, together with
\cref{ent:thm:extension,ent:thm:zero-conservation,ent:thm:no-continuation}
and their assembly in \cref{ent:thm:main-extension}. The sharp workspace
criterion and the persistence result are
\cref{ent:thm:extension-criterion,ent:cor:no-repair}; the restricted-unit
criterion \cref{ent:thm:units} is closely related to established topological
nilpotence distinctions, and it is its precise necessity-and-sufficiency
formulation for \eqref{ent:eq:Talg}, including groups with no order unit,
that is proposed here. The rank-one factorization formula is explicitly
classical. The all-surcomplex polynomial corollary is credited to the
pre-existing all-scale theme in the repository
(\cref{ent:rem:rigidity-credit}). A later discovery of an antecedent would
change the priority discussion, not the content or the hypotheses of the
displayed proofs.

\clearpage
\phantomsection\addcontentsline{toc}{section}{References}
\begingroup\raggedright
\begin{thebibliography}{99}
\bibitem{Conway}
J. H. Conway, \emph{On Numbers and Games}, London Mathematical Society
Monographs, vol.~6, Academic Press, 1976.

\bibitem{Gonshor}
H. Gonshor, \emph{An Introduction to the Theory of Surreal Numbers},
London Mathematical Society Lecture Note Series, vol.~110,
Cambridge University Press, 1986.

\bibitem{Neumann}
B. H. Neumann, On ordered division rings, \emph{Transactions of the American
Mathematical Society} \textbf{66} (1949), 202--252.
\href{https://doi.org/10.1090/S0002-9947-1949-0032593-5}{doi:10.1090/S0002-9947-1949-0032593-5}.

\bibitem{Poonen}
B. Poonen, Maximally complete fields, \emph{L'Enseignement
Math\'ematique} (2) \textbf{39} (1993), 87--106.
Section~3 and Corollary~4 are the principal inputs used here.
\url{https://math.mit.edu/~poonen/papers/amsval.pdf}.

\bibitem{Lazard}
M. Lazard, Les z\'eros d'une fonction analytique d'une variable sur un
corps valu\'e complet, \emph{Publications Math\'ematiques de l'IH\'ES}
\textbf{14} (1962), 47--75.
\href{https://doi.org/10.1007/BF02684326}{doi:10.1007/BF02684326}.
\url{https://www.numdam.org/item/PMIHES_1962__14__47_0/}.

\bibitem{Cherry}
W. Cherry, Existence of GCD's and factorization in rings of
non-Archimedean entire functions, \emph{Contemporary Mathematics}
\textbf{551} (2011), 57--69.
\url{https://arxiv.org/abs/1007.0984}.

\bibitem{CherryLectures}
W. Cherry, \emph{Lectures on Non-Archimedean Function Theory}, 2009,
arXiv:0909.4509. In particular the preparation and product results in
Lectures~2--3. \url{https://arxiv.org/abs/0909.4509}.

\bibitem{Conrad}
B. Conrad, \emph{Lecture 6: More constructions with Huber rings}, Stanford
Number Theory Learning Seminar on perfectoid spaces, 31 October 2014,
informal lecture notes. In particular Section~6.2 and Example~6.2.1.
\url{https://virtualmath1.stanford.edu/~conrad/Perfseminar/Notes/L6.pdf}.

\bibitem{vdH}
J. van der Hoeven, Operators on generalized power series,
\emph{Illinois Journal of Mathematics} \textbf{45} (2001), no.~4,
1161--1190.
\href{https://doi.org/10.1215/ijm/1258138061}{doi:10.1215/ijm/1258138061}.

\bibitem{ADH}
M. Aschenbrenner, L. van den Dries, and J. van der Hoeven,
On numbers, germs, and transseries, arXiv:1711.06936, revised version.
\url{https://arxiv.org/abs/1711.06936}.
Used for bibliographic cross-checks of the surreal background.

\bibitem{BournezGuilmant}
O. Bournez and Q. Guilmant, \emph{Surreal fields stable under exponential and
logarithmic functions}, 2022, arXiv:2201.08199.
\url{https://arxiv.org/abs/2201.08199}.

\bibitem{Helmer}
O. Helmer, Divisibility properties of integral functions, \emph{Duke
Mathematical Journal} \textbf{6} (1940), no.~2, 345--356.
\href{https://doi.org/10.1215/S0012-7094-40-00626-3}{doi:10.1215/S0012-7094-40-00626-3}.
Cited by the fourth source as a historical antecedent, not a proof input.

\bibitem{Henriksen}
M. Henriksen, On the ideal structure of the ring of entire functions,
\emph{Pacific Journal of Mathematics} \textbf{2} (1952), no.~2, 179--184.
\href{https://doi.org/10.2140/pjm.1952.2.179}{doi:10.2140/pjm.1952.2.179}.
Free maximal ideals, their residue fields, and the transcendental coordinate
of Theorem~5.

\bibitem{Henriksen53}
M. Henriksen, On the prime ideals of the ring of entire functions,
\emph{Pacific Journal of Mathematics} \textbf{3} (1953), no.~4, 711--720.
\url{https://msp.org/pjm/1953/3-4/pjm-v3-n4-p04-p.pdf}.
Cited as recorded in the analytic-geometry report of this collection; not
cited by the sixth source and not compared here.

\bibitem{FFW}
C. A. Finocchiaro, S. Frisch and D. Windisch, \emph{Prime ideals in
infinite products of commutative rings}, arXiv:2009.03069v2, revised
24 August 2023. \url{https://arxiv.org/abs/2009.03069v2}.
Credited by the sixth source as the direct precedent for the ultrafilter and
valuation-domain mechanism of \cref{ent:ps:sec:spectrum}.

\bibitem{Bruno}
N. Bruno, \emph{Ideal Structure of Rings of Analytic Functions with
non-Archimedean Metrics}, Ph.D. dissertation, The Ohio State University,
2021. Definition~2.3 assumes a value group inside $\R$.

\bibitem{MaNeelon}
D. Ma and T. S. Neelon, On convergence sets of formal power series,
\emph{Complex Analysis and its Synergies} \textbf{1} (2015), article~4.
\href{https://doi.org/10.1186/s40627-015-0004-4}{doi:10.1186/s40627-015-0004-4}.
In particular Corollary~4.19.

\bibitem{Stacks}
The Stacks Project Authors, \emph{The Stacks Project}, Section~10.31,
Noetherian rings, accessed 22 September 2026.
\url{https://stacks.math.columbia.edu/tag/00FM}.

\bibitem{StacksValuation}
The Stacks Project Authors, \emph{The Stacks Project}, Section~10.50,
Valuation rings, Tag~00I8. \url{https://stacks.math.columbia.edu/tag/00I8}.
Cited by the sixth source for the correspondence between primes of a
valuation ring and convex subgroups of its value group.

\bibitem{Berarducci}
A. Berarducci, \emph{Surreal numbers, exponentiation and derivations}, 2020,
arXiv:2008.06878. In particular Theorem~10.2 and Definition~10.3.
\url{https://arxiv.org/abs/2008.06878}.

\bibitem{Achinger}
P. Achinger, \emph{Introduction to non-Archimedean Geometry}, course notes,
version of 28 January 2021. \url{https://achinger.impan.pl/rigid/notes.pdf}.

\bibitem{JoswigSmith}
M. Joswig and B. Smith, Convergent Hahn series and tropical geometry of
higher rank, \emph{Journal of the London Mathematical Society} (2)
\textbf{107} (2023), no.~4, 1450--1481.
\href{https://doi.org/10.1112/jlms.12716}{doi:10.1112/jlms.12716}.

\bibitem{RosayRudin}
J.-P. Rosay and W. Rudin, Holomorphic maps from $\C^n$ to $\C^n$,
\emph{Transactions of the American Mathematical Society} \textbf{310} (1988),
47--86. \url{https://www.jstor.org/stable/2001110}.

\bibitem{Winkelmann}
J. Winkelmann, Rosay--Rudin-spaces, tame sets and Danielewski surfaces,
arXiv:2309.07678v1, 2023. \url{https://arxiv.org/abs/2309.07678}.

\bibitem{MantovaMatusinski}
V. Mantova and M. Matusinski, Surreal numbers with derivation, Hardy fields
and transseries: a survey, arXiv:1608.03413v2.
\url{https://arxiv.org/abs/1608.03413}.

\bibitem{LInnocenteMantova}
S. L'Innocente and V. Mantova, A factorisation theory for generalised power
series and omnific integers, \emph{Advances in Mathematics} (2024), article
109513, \href{https://doi.org/10.1016/j.aim.2024.109513}{doi:10.1016/j.aim.2024.109513};
author version arXiv:1710.07304v5.

\bibitem{RepoRoot}
V. Reshetnikov, \emph{Surreal}: repository README and reported
formalization scope, inspected 21 September 2026.
\url{https://github.com/VladimirReshetnikov/Surreal}.
AI-assisted research collection; the reports are unrefereed drafts, not
independent peer-reviewed validation.

\bibitem{RepoCatalogue}
V. Reshetnikov, \emph{The surreal and surcomplex reports}, documentation
catalogue \texttt{docs/README.md} in the \emph{Surreal} repository,
inspected 21 September 2026.

\bibitem{RepoRankOne}
\emph{From Surcomplex Hahn Fields to Berkovich Geometry: rank-one
convergence, Tate algebras, analytic annuli, and an entire-function case
study}, research report in the \emph{Surreal} repository,
\texttt{docs/surcomplex/rank-one-berkovich/}, inspected 21 September 2026.
Already contains rank-one zero-free rigidity and a canonical product theorem;
its \texttt{prop:rankobstruction} is the rank-one boundary discussed in
\cref{ent:sec:collection}.

\bibitem{RepoDivisors}
\emph{Global Support Obstructions: Moving Divisors, Mittag--Leffler, and a
Picard Dichotomy}, research report in the \emph{Surreal} repository,
\texttt{docs/surcomplex/global-divisors/}, inspected 21 September 2026.
Stated over $\mathcal O(U)((t^\Gamma))$, a different coefficient ring.

\bibitem{RepoAnalyticGeometry}
\emph{Analytic Geometry over Surcomplex Hahn Fields}, research report in the
\emph{Surreal} repository, \texttt{docs/surcomplex/analytic-geometry/}.
Source of the ring-discipline convention invoked in \cref{ent:conv:ring}.

\bibitem{RepoAnalysis}
\emph{Surcomplex Analysis}, research report in the \emph{Surreal}
repository, \texttt{docs/surcomplex/analysis/}. Contains the all-scale
polynomial rigidity theorem discussed in \cref{ent:rem:rigidity-credit}.

\bibitem{RepoFoundations}
\emph{Foundations}, research report in the \emph{Surreal} repository,
\texttt{docs/foundations-and-computation/foundations/}. Source of the
localization principle and of the sign-flipped pair
$\sum t^{\pm n^2}z^n$ discussed in
\cref{ent:rem:criterion-strength,ent:sec:collection}.

\bibitem{RepoHolonomic}
\emph{Holonomic Rigidity for Entire Hahn Functions}, research report in the
\emph{Surreal} repository,
\texttt{docs/surcomplex/holonomic-rigidity-for-entire-hahn-functions/}. Its
Lemma \texttt{hol:lem:generic} is the countable-field device used in
\cref{ent:mv:thm:exceptional}, and its Lemma \texttt{hol:cf:lem:reduction} is
the one-variable antecedent of \cref{ent:rc:thm:reduction}.

\bibitem{KuhlmannSerra}
S. Kuhlmann and M. Serra, The automorphism group of a valued field of
generalised formal power series, arXiv:2107.03362v3, 2022. Cited by source~11
for the strong summability conditions (Section~4).

\bibitem{StacksFlat}
The Stacks Project Authors, \emph{The Stacks Project}, Section~10.39, Flat
modules and flat ring maps, Tag~00H9, Lemmas~10.39.11 and 10.39.15.
\url{https://stacks.math.columbia.edu/tag/00H9}. Cited by source~11 for the
equational and faithful-flatness criteria.

\bibitem{GarciaFritzPasten}
N. Garcia-Fritz and H. Pasten, Uniform positive existential interpretation of
the integers in rings of entire functions of positive characteristic,
arXiv:1411.7109, 2014. Source~11 consulted the abstract only.

\bibitem{WikiSurreal}
Wikipedia contributors, \emph{Surreal number},
\url{https://en.wikipedia.org/wiki/Surreal_number}, accessed 23 September 2026
by source~11 as orientation only.

\bibitem{RepoQuotients}
\emph{Set-Sized Quotients of the Omnific Integers}, research report in the
\emph{Surreal} repository,
\texttt{docs/surreal/set-sized-quotients-of-omnific-integers/}. Its
numerical-polynomial theorems and their fixed-workspace transfer are the
polynomial half of \cref{ent:as:thm:integer}.

\bibitem{RepoODG}
\emph{Omnific Integers and Omnific--Diophantine Geometry}, research report in
the \emph{Surreal} repository,
\texttt{docs/surreal/omnific-diophantine-geometry/}.

\bibitem{RepoOPA}
\emph{Omnific-Preserving Automorphisms}, research report in the
\emph{Surreal} repository,
\texttt{docs/surreal/omnific-preserving-automorphisms/}.
\end{thebibliography}
\endgroup
\end{document}
